\input{preamble}

% OK, start here.
%
\begin{document}

\title{More on Algebra}


\maketitle

\phantomsection
\label{section-phantom}

\tableofcontents



\section{Introduction}
\label{section-introduction}

\noindent
In this chapter we prove some results in commutative algebra which
are less elementary than those in the first chapter on commutative
algebra, see
Algebra, Section \ref{algebra-section-introduction}.
A reference is \cite{MatCA}.






\section{Advice for the reader}
\label{section-advice}

\noindent
More than in the chapter on commutative algebra, each of the sections in
this chapter stands on its own. Starting with
Section \ref{section-derived-modules}
we freely use the (unbounded) derived category of modules over rings and all
the machinery that comes with it.




\section{Stably free modules}
\label{section-stably-free}

\noindent
Here is what seems to be the generally accepted definition.

\begin{definition}
\label{definition-stably-free}
Let $R$ be a ring. 
\begin{enumerate}
\item Two modules $M$, $N$ over $R$ are said to be
{\it stably isomorphic} if there exist $n, m \geq 0$ such
that $M \oplus R^{\oplus m} \cong N \oplus R^{\oplus n}$
as $R$-modules.
\item A module $M$ is {\it stably free} if it is stably isomorphic
to a free module.
\end{enumerate}
\end{definition}

\noindent
Observe that a stably free module is projective.

\begin{lemma}
\label{lemma-exact-category-stably-free}
Let $R$ be a ring. Let $0 \to P' \to P \to P'' \to 0$ be a short
exact sequence of finite projective $R$-modules. If $2$ out of $3$
of these modules are stably free, then so is the third.
\end{lemma}

\begin{proof}
Since the modules are projective, the sequence is split. Thus we can
choose an isomorphism $P = P' \oplus P''$. If $P' \oplus R^{\oplus n}$
and $P'' \oplus R^{\oplus m}$ are free, then we see that
$P \oplus R^{\oplus n + m}$ is free. Suppose that $P'$ and $P$ are
stably free, say $P \oplus R^{\oplus n}$ is free and $P' \oplus R^{\oplus m}$
is free. Then
$$
P'' \oplus (P' \oplus R^{\oplus m}) \oplus R^{\oplus n} =
(P'' \oplus P') \oplus R^{\oplus m} \oplus R^{\oplus n} =
(P \oplus R^{\oplus n}) \oplus R^{\oplus m}
$$
is free. Thus $P''$ is stably free. By symmetry we get the last of the
three cases.
\end{proof}

\begin{lemma}
\label{lemma-lift-stably-free}
Let $R$ be a ring. Let $I \subset R$ be an ideal. Assume that
every element of $1 + I$ is a unit (in other words $I$ is contained
in the Jacobson radical of $R$). For every finite stably free $R/I$-module $E$
there exists a finite stably free $R$-module $M$ such that $M/IM \cong E$.
\end{lemma}

\begin{proof}
Choose a $n$ and $m$ and an isomorphism
$E \oplus (R/I)^{\oplus n} \cong (R/I)^{\oplus m}$.
Choose $R$-linear maps $\varphi : R^{\oplus m} \to R^{\oplus n}$
and $\psi : R^{\oplus n} \to R^{\oplus m}$ lifting the
projection $(R/I)^{\oplus m} \to (R/I)^{\oplus n}$
and injection $(R/I)^{\oplus n} \to (R/I)^{\oplus m}$.
Then $\varphi \circ \psi : R^{\oplus n} \to R^{\oplus n}$
reduces to the identity modulo $I$. Thus the determinant of
this map is invertible by our assumption on $I$. Hence
$P = \Ker(\varphi)$ is stably free and lifts $E$.
\end{proof}

\begin{lemma}
\label{lemma-lift-projective}
Let $R$ be a ring. Let $I \subset R$ be an ideal.
Assume that every element of $1 + I$ is a unit
(in other words $I$ is contained in the Jacobson radical of $R$).
Let $M$ be a finite flat $R$-module such that
$M/IM$ is a projective $R/I$-module.
Then $M$ is a finite projective $R$-module.
\end{lemma}

\begin{proof}
By Algebra, Lemma \ref{algebra-lemma-finite-flat-local}
we see that $M_\mathfrak p$ is finite free for all prime ideals
$\mathfrak p \subset R$.
By
Algebra, Lemma \ref{algebra-lemma-finite-projective}
it suffices to show that the function $\rho_M : \mathfrak p \mapsto
\dim_{\kappa(\mathfrak p)} M \otimes_R \kappa(\mathfrak p)$
is locally constant on $\Spec(R)$. Because $M/IM$ is finite projective, this
is true on $V(I) \subset \Spec(R)$. Since every closed point
of $\Spec(R)$ is in $V(I)$ and since
$\rho_M(\mathfrak p) = \rho_M(\mathfrak q)$
whenever $\mathfrak p \subset \mathfrak q \subset R$
are prime ideals, we conclude by
an elementary argument on topological spaces which we omit.
\end{proof}

\noindent
The lift of Lemma \ref{lemma-lift-stably-free}
is unique up to isomorphism by the following lemma.

\begin{lemma}
\label{lemma-isomorphic-finite-projective-lifts}
Let $R$ be a ring. Let $I \subset R$ be an ideal. Assume that
every element of $1 + I$ is a unit (in other words $I$ is contained
in the Jacobson radical of $R$). If $P$ and $P'$ are finite
projective $R$-modules, then
\begin{enumerate}
\item if $\varphi : P \to P'$ is an $R$-module map inducing an
isomorphism $\overline{\varphi} : P/IP \to P'/IP'$, then $\varphi$
is an isomorphism,
\item if $P/IP \cong P'/IP'$, then $P \cong P'$.
\end{enumerate}
\end{lemma}

\begin{proof}
Proof of (1). As $P'$ is projective as an $R$-module we may
choose a lift $\psi : P' \to P$ of the map
$P' \to P'/IP' \xrightarrow{\overline{\varphi}^{-1}} P/IP$.
By Nakayama's lemma (Algebra, Lemma \ref{algebra-lemma-NAK})
$\psi \circ \varphi$ and $\varphi \circ \psi$ are surjective.
Hence these maps are isomorphisms (Algebra, Lemma \ref{algebra-lemma-fun}).
Thus $\varphi$ is an isomorphism.

\medskip\noindent
Proof of (2). Choose an isomorphism $P/IP \cong P'/IP'$.
Since $P$ is projective we can choose a lift $\varphi : P \to P'$ of the map
$P \to P/IP \to P'/IP'$. Then $\varphi$ is an isomorphism by (1).
\end{proof}




\section{A comment on the Artin-Rees property}
\label{section-artin-rees}

\noindent
Some of this material is taken from \cite{conrad-dejong}. A general
discussion with additional references can be found in
\cite[Section 1]{Eis}.

\medskip\noindent
Let $A$ be a Noetherian ring and let $I \subset A$ be an ideal. Given a
homomorphism $f : M \to N$ of finite $A$-modules there exists a $c \geq 0$
such that
$$
f(M) \cap I^nN \subset f(I^{n - c}M)
$$
for all $n \geq c$, see Algebra, Lemma \ref{algebra-lemma-map-AR}. In this
situation we will say {\it $c$ works for $f$ in the Artin-Rees lemma}.

\begin{lemma}
\label{lemma-approximate-complex}
Let $A$ be a Noetherian ring. Let $I \subset A$ be an ideal contained in
the Jacobson radical of $A$. Let
$$
S : L \xrightarrow{f} M \xrightarrow{g} N
\quad\text{and}\quad
S' : L \xrightarrow{f'} M \xrightarrow{g'} N
$$
be two complexes of finite $A$-modules as shown. Assume that
\begin{enumerate}
\item $c$ works in the Artin-Rees lemma for $f$ and $g$,
\item the complex $S$ is exact, and
\item $f' = f \bmod I^{c + 1}M$ and $g' = g \bmod I^{c + 1}N$.
\end{enumerate}
Then $c$ works in the Artin-Rees lemma for $g'$ and the
complex $S'$ is exact.
\end{lemma}

\begin{proof}
We first show that $g'(M) \cap I^nN \subset g'(I^{n - c}M)$ for $n \geq c$.
Let $a$ be an element of $M$ such that $g'(a) \in I^nN$. We want to
adjust $a$ by an element of $f'(L)$, i.e, without changing $g'(a)$, so
that $a \in I^{n-c}M$. Assume that $a \in I^rM$, where $r < n - c$.
Then
$$
g(a) = g'(a) + (g - g')(a) \in
I^n N + I^{r + c + 1}N = I^{r + c + 1}N.
$$
By Artin-Rees for $g$ we have $g(a) \in g(I^{r + 1}M)$. Say $g(a) = g(a_1)$
with $a_1 \in I^{r + 1}M$. Since the sequence $S$ is exact, $a - a_1 \in f(L)$.
Accordingly, we write $a = f(b) + a_1$ for some $b \in L$.
Then $f(b) = a - a_1 \in I^rM$. Artin-Rees for $f$ shows that
if $r \geq c$, we may replace $b$ by an element of $I^{r - c}L$.
Then in all cases, $a = f'(b) + a_2$, where
$a_2 = (f - f')(b) + a_1 \in I^{r + 1}M$. (Namely, either $c \geq r$
and $(f - f')(b) \in I^{r + 1}M$ by assumption, or $c < r$ and
$b \in I^{r - c}$, whence again $(f - f')(b) \in I^{c + 1} I^{r - c} M =
I^{r + 1}M$.) So we can adjust $a$ by the element $f'(b) \in f'(L)$ to
increase $r$ by $1$.

\medskip\noindent
In fact, the argument above shows that
$(g')^{-1}(I^nN) \subset f'(L) + I^{n - c}M$ for all $n \geq c$.
Hence $S'$ is exact because
$$
(g')^{-1}(0) = (g')^{-1}(\bigcap I^nN) \subset
\bigcap f'(L) + I^{n - c}M = f'(L)
$$
as $I$ is contained in the Jacobson radical of $A$, see Algebra, Lemma
\ref{algebra-lemma-intersection-powers-ideal-module}.
\end{proof}

\noindent
Given an ideal $I \subset A$ of a ring $A$ and an $A$-module $M$
we set
$$
\text{Gr}_I(M) = \bigoplus I^nM/I^{n + 1}M.
$$
We think of this as a graded $\text{Gr}_I(A)$-module.

\begin{lemma}
\label{lemma-approximate-complex-graded}
Assumptions as in Lemma \ref{lemma-approximate-complex}.
Let $Q = \Coker(g)$ and $Q' = \Coker(g')$. Then
$\text{Gr}_I(Q) \cong \text{Gr}_I(Q')$
as graded $\text{Gr}_I(A)$-modules.
\end{lemma}

\begin{proof}
In degree $n$ we have
$\text{Gr}_I(Q)_n = I^nN/(I^{n + 1}N + g(M) \cap I^nN)$
and similarly for $Q'$. We claim that
$$
g(M) \cap I^nN \subset I^{n + 1}N + g'(M) \cap I^nN.
$$
By symmetry (the proof of the claim will only use that $c$ works
for $g$ which also holds for $g'$ by the lemma) this will imply that
$$
I^{n + 1}N + g(M) \cap I^nN = I^{n + 1}N + g'(M) \cap I^nN
$$
whence $\text{Gr}_I(Q)_n$ and $\text{Gr}_I(Q')_n$ agree as subquotients
of $N$, implying the lemma. Observe that the claim is clear for
$n \leq c$ as $g = g' \bmod I^{c + 1}N$. If $n > c$, then suppose
$b \in g(M) \cap I^nN$. Write $b = g(a)$ for $a \in I^{n - c}M$.
Set $b' = g'(a)$. We have $b - b' = (g - g')(a) \in I^{n + 1}N$
as desired.
\end{proof}

\begin{lemma}
\label{lemma-works-flat-extension}
Let $A \to B$ be a flat map of Noetherian rings. Let $I \subset A$ be
an ideal. Let $f : M \to N$ be a homomorphism of finite $A$-modules.
Assume that $c$ works for $f$ in the Artin-Rees lemma. Then $c$ works for
$f \otimes 1 : M \otimes_A B \to N \otimes_A B$ in the Artin-Rees lemma
for the ideal $IB$.
\end{lemma}

\begin{proof}
Note that
$$
(f \otimes 1)(M) \cap I^n N \otimes_A B
= (f \otimes 1)\left((f \otimes 1)^{-1}(I^n N \otimes_A B)\right)
$$
On the other hand,
\begin{align*}
(f \otimes 1)^{-1}(I^n N \otimes_A B) &
= \Ker(M \otimes_A B \to N \otimes_A B/(I^n N \otimes_A B)) \\
& =
\Ker(M \otimes_A B \to (N/I^nN) \otimes_A B)
\end{align*}
As $A \to B$ is flat taking kernels and cokernels commutes with
tensoring with $B$, whence this is equal to
$f^{-1}(I^nN) \otimes_A B$. By assumption $f^{-1}(I^nN)$ is contained in
$\Ker(f) + I^{n - c}M$. Thus the lemma holds.
\end{proof}








\section{Fibre products of rings, I}
\label{section-fibre-products-rings}

\noindent
Fibre products of rings have to do with pushouts of schemes. Some
cases of pushouts of schemes are discussed in
More on Morphisms, Section \ref{more-morphisms-section-pushouts}.

\begin{lemma}
\label{lemma-fibre-product-finite-type}
Let $R$ be a ring. Let $A \to B$ and $C \to B$ be $R$-algebra maps.
Assume
\begin{enumerate}
\item $R$ is Noetherian,
\item $A$, $B$, $C$ are of finite type over $R$,
\item $A \to B$ is surjective, and
\item $B$ is finite over $C$.
\end{enumerate}
Then $A \times_B C$ is of finite type over $R$.
\end{lemma}

\begin{proof}
Set $D = A \times_B C$. There is a commutative diagram
$$
\xymatrix{
0 &
B \ar[l] &
A \ar[l] &
I \ar[l] &
0 \ar[l] \\
0 &
C \ar[l] \ar[u] &
D \ar[l] \ar[u] &
I \ar[l] \ar[u] &
0 \ar[l]
}
$$
with exact rows. Choose $y_1, \ldots, y_n \in B$ which are generators for
$B$ as a $C$-module. Choose $x_i \in A$ mapping to $y_i$.
Then $1, x_1, \ldots, x_n$ are generators for $A$ as a $D$-module.
The map $D \to A \times C$ is injective, and the ring $A \times C$ is finite
as a $D$-module (because it is the direct sum of the finite $D$-modules
$A$ and $C$). Hence the lemma follows from the Artin-Tate lemma
(Algebra, Lemma \ref{algebra-lemma-Artin-Tate}).
\end{proof}

\begin{lemma}
\label{lemma-formal-consequence}
Let $R$ be a Noetherian ring. Let $I$ be a finite set. Suppose given a
cartesian diagram
$$
\xymatrix{
\prod B_i &
\prod A_i \ar[l]^{\prod \varphi_i} \\
Q \ar[u]^{\prod \psi_i} &
P \ar[u] \ar[l]
}
$$
with $\psi_i$ and $\varphi_i$ surjective, and $Q$, $A_i$, $B_i$ of
finite type over $R$. Then $P$ is of finite type over $R$.
\end{lemma}

\begin{proof}
Follows from Lemma \ref{lemma-fibre-product-finite-type}
and induction on the size of $I$.
Namely, let $I = I' \amalg \{i_0\}$. Let $P'$ be the ring defined
by the diagram of the lemma using $I'$. Then $P'$ is of finite type
by induction hypothesis. Finally, $P$ sits in a fibre product diagram
$$
\xymatrix{
B_{i_0} &
A_{i_0} \ar[l] \\
P' \ar[u] &
P \ar[u] \ar[l] &
}
$$
to which the lemma applies.
\end{proof}

\begin{lemma}
\label{lemma-diagram-localize}
Suppose given a cartesian diagram of rings
$$
\xymatrix{
R &
R' \ar[l]^t \\
B \ar[u]_s &
B'\ar[u] \ar[l]
}
$$
i.e., $B' = B \times_R R'$. If $h \in B'$ corresponds to $g \in B$
and $f \in R'$ such that $s(g) = t(f)$, then the diagram
$$
\xymatrix{
R_{s(g)} = R_{t(f)} &
(R')_f \ar[l]^-t \\
B_g \ar[u]_s &
(B')_h \ar[u] \ar[l]
}
$$
is cartesian too.
\end{lemma}

\begin{proof}
The equality $B' = B \times_R R'$ tells us that
$$
0 \to B' \to B \oplus R' \xrightarrow{s, -t} R
$$
is an exact sequence of $B'$-modules. We have $B_g = B_h$,
$R'_f = R'_h$, and $R_{s(g)} = R_{t(f)} = R_h$ as $B'$-modules.
By exactness of localization
(Algebra, Proposition \ref{algebra-proposition-localization-exact})
we find that
$$
0 \to B'_h \to B_g \oplus R'_f \xrightarrow{s, -t} R_{s(g)} = R_{t(f)}
$$
is an exact sequence. This proves the lemma.
\end{proof}

\noindent
Consider a commutative diagram of rings
$$
\xymatrix{
R &
R' \ar[l] \\
B \ar[u] &
B' \ar[u] \ar[l]
}
$$
Consider the functor (where the fibre product of categories is
as constructed in
Categories, Example \ref{categories-example-2-fibre-product-categories})
\begin{equation}
\label{equation-modules}
\text{Mod}_{B'} \longrightarrow
\text{Mod}_B \times_{\text{Mod}_R} \text{Mod}_{R'},\quad
L' \longmapsto (L' \otimes_{B'} B, L' \otimes_{B'} R', can)
\end{equation}
where $can$ is the canonical identification
$L' \otimes_{B'} B \otimes_B R = L' \otimes_{B'} R' \otimes_{R'} R$.
In the following we will write $(N, M', \varphi)$ for an object
of the right hand side, i.e., $N$ is a $B$-module, $M'$ is an $R'$-module
and $\varphi : N \otimes_B R \to M' \otimes_{R'} R$ is an isomorphism.

\begin{lemma}
\label{lemma-modules}
Given a commutative diagram of rings
$$
\xymatrix{
R &
R' \ar[l] \\
B \ar[u] &
B' \ar[u] \ar[l]
}
$$
the functor (\ref{equation-modules}) has a right adjoint,
namely the functor
$$
F : (N, M', \varphi) \longmapsto N \times_\varphi M'
$$
(see proof for elucidation).
\end{lemma}

\begin{proof}
Given an object $(N, M', \varphi)$ of the category
$\text{Mod}_B \times_{\text{Mod}_R} \text{Mod}_{R'}$
we set
$$
N \times_\varphi M' = \{(n, m') \in N \times M' \mid
\varphi(n \otimes 1) = m' \otimes 1\text{ in }M' \otimes_{R'} R\}
$$
viewed as a $B'$-module.
The adjointness statement is that for a $B'$-module $L'$ and
a triple $(N, M', \varphi)$ we have
$$
\Hom_{B'}(L', N \times_\varphi M') =
\Hom_B(L' \otimes_{B'} B, N)
\times_{\Hom_R(L' \otimes_{B'} R, M' \otimes_{R'} R)}
\Hom_{R'}(L' \otimes_{B'} R', M')
$$
By Algebra, Lemma \ref{algebra-lemma-adjoint-tensor-restrict}
the right hand side is equal to
$$
\Hom_{B'}(L', N) \times_{\Hom_{B'}(L', M' \otimes_{R'} R)} \Hom_{B'}(L', M')
$$
Thus it is clear that for a pair $(g, f')$ of elements
of this fibre product we get an $B'$-linear map
$L' \to N \times_\varphi M'$, $l' \mapsto (g(l'), f'(l'))$.
Conversely, given a $B'$ linear map $g' : L' \to N \times_\varphi M'$
we can set $g$ equal to the composition $L' \to N \times_\varphi M' \to N$
and $f'$ equal to the composition $L' \to N \times_\varphi M' \to M'$.
These constructions are mutually inverse to each other and
define the desired isomorphism.
\end{proof}






\section{Fibre products of rings, II}
\label{section-fibre-products-rings-II}

\noindent
In this section we discuss fibre products in the following situation.

\begin{situation}
\label{situation-module-over-fibre-product}
In the following we will consider ring maps
$$
\xymatrix{
B \ar[r] & A & A' \ar[l]
}
$$
where we assume $A' \to A$ is surjective with kernel $I$.
In this situation we set $B' = B \times_A A'$ to
obtain a cartesian square
$$
\xymatrix{
A & A' \ar[l] \\
B \ar[u] & B' \ar[l] \ar[u]
}
$$
\end{situation}

\begin{lemma}
\label{lemma-points-of-fibre-product}
In Situation \ref{situation-module-over-fibre-product}
we have
$$
\Spec(B') = \Spec(B) \amalg_{\Spec(A)} \Spec(A')
$$
as topological spaces.
\end{lemma}

\begin{proof}
Since $B' = B \times_A A'$ we obtain a commutative square of
spectra, which induces a continuous map
$$
can : \Spec(B) \amalg_{\Spec(A)} \Spec(A') \longrightarrow \Spec(B')
$$
as the source is a pushout in the category
of topological spaces (which exists by
Topology, Section \ref{topology-section-colimits}).

\medskip\noindent
To show the map $can$ is surjective, let $\mathfrak q' \subset B'$ be a prime
ideal. If $I \subset \mathfrak q'$ (here and below we take the liberty
of considering $I$ as an ideal of $B'$ as well as an ideal of $A'$),
then $\mathfrak q'$ corresponds to a prime ideal of $B$ and is in the image.
If not, then pick $h \in I$, $h \not \in \mathfrak q'$.
In this case $B_h = A_h = 0$
and the ring map $B'_h \to A'_h$ is an isomorphism, see
Lemma \ref{lemma-diagram-localize}. Thus we see that
$\mathfrak q'$ corresponds to a unique prime ideal $\mathfrak p' \subset A'$
which does not contain $I$.

\medskip\noindent
Since $B' \to B$ is surjective, we see that $can$ is injective on
the summand $\Spec(B)$. We have seen above that $\Spec(A') \to \Spec(B')$
is injective on the complement of $V(I) \subset \Spec(A')$. Since
$V(I) \subset \Spec(A')$ is exactly the image of $\Spec(A) \to \Spec(A')$
a trivial set theoretic argument shows that $can$ is injective.

\medskip\noindent
To finish the proof we have to show that $can$ is open. To do this, observe
that an open of the pushout is of the form $V \amalg U'$ where
$V \subset \Spec(B)$ and $U' \subset \Spec(A')$ are opens whose inverse
images in $\Spec(A)$ agree. Let $v \in V$. We can find a $g \in B$
such that $v \in D(g) \subset V$. Let $f \in A$ be the image.
Pick $f' \in A'$ mapping to $f$. Then
$D(f') \cap U' \cap V(I) = D(f') \cap V(I)$.
Hence $V(I) \cap D(f')$ and $D(f') \cap (U')^c$ are disjoint closed
subsets of $D(f') = \Spec(A'_{f'})$. Write $(U')^c = V(J)$
for some ideal $J \subset A'$. Since
$A'_{f'} \to A'_{f'}/IA'_{f'} \times A'_{f'}/JA'_{f'}$
is surjective by the disjointness just shown, we can find an $a'' \in A'_{f'}$
mapping to $1$ in $A'_{f'}/IA'_{f'}$ and mapping to zero in
$A'_{f'}/JA'_{f'}$. Clearing denominators, we find an element
$a' \in J$ mapping to $f^n$ in $A$. Then $D(a'f') \subset U'$.
Let $h' = (g^{n + 1}, a'f') \in B'$.
Since $B'_{h'} = B_{g^{n + 1}} \times_{A_{f^{n + 1}}} A'_{a'f'}$ by a previously
cited lemma, we see that $D(h')$ pulls back to an open neighbourhood
of $v$ in the pushout, i.e., the image of $V \amalg U'$ contains
an open neighbourhood of the image of $v$. We omit the (easier) proof that
the same thing is true for $u' \in U'$ with $u' \not \in V(I)$.
\end{proof}

\begin{lemma}
\label{lemma-fibre-product-integral}
In Situation \ref{situation-module-over-fibre-product}
if $B \to A$ is integral, then $B' \to A'$ is integral.
\end{lemma}

\begin{proof}
Let $a' \in A'$ with image $a \in A$. Let $x^d + b_1 x^{d - 1} + \ldots + b_d$
be a monical polynomial with coefficients in $B$ satisfied by $a$.
Choose $b'_i \in B'$ mapping to $b_i \in B$ (possible).
Then $(a')^d + b'_1 (a')^{d - 1} + \ldots + b'_d$ is in the kernel
of $A' \to A$. Since $\Ker(B' \to B) = \Ker(A' \to A)$ we can
modify our choice of $b'_d$ to get
$(a')^d + b'_1 (a')^{d - 1} + \ldots + b'_d = 0$ as desired.
\end{proof}

\noindent
In Situation \ref{situation-module-over-fibre-product}
we'd like to understand $B'$-modules in terms of modules over $A'$, $A$,
and $B$. In order to do this we consider the functor (where the
fibre product of categories as constructed in
Categories, Example \ref{categories-example-2-fibre-product-categories})
\begin{equation}
\label{equation-functor}
\text{Mod}_{B'} \longrightarrow
\text{Mod}_B \times_{\text{Mod}_A} \text{Mod}_{A'},\quad
L' \longmapsto (L' \otimes_{B'} B, L' \otimes_{B'} A', can)
\end{equation}
where $can$ is the canonical identification
$L' \otimes_{B'} B \otimes_B A = L' \otimes_{B'} A' \otimes_{A'} A$.
In the following we will write $(N, M', \varphi)$ for an object
of the right hand side, i.e., $N$ is a $B$-module, $M'$ is an $A'$-module
and $\varphi : N \otimes_B A \to M' \otimes_{A'} A$ is an isomorphism.
However, it is often more convenient think of $\varphi$ as a $B$-linear
map $\varphi : N \to M'/IM'$ which induces an isomorphism
$N \otimes_B A \to M' \otimes_{A'} A = M'/IM'$.

\begin{lemma}
\label{lemma-module-over-fibre-product}
In Situation \ref{situation-module-over-fibre-product}
the functor (\ref{equation-functor}) has a right adjoint, namely
the functor
$$
F : (N, M', \varphi) \longmapsto N \times_{\varphi, M} M'
$$
where $M = M'/IM'$. Moreover, the composition of $F$ with
(\ref{equation-functor}) is the identity functor on
$\text{Mod}_B \times_{\text{Mod}_A} \text{Mod}_{A'}$. In other words,
setting $N' = N \times_{\varphi, M} M'$ we have
$N' \otimes_{B'} B = N$ and $N' \otimes_{B'} A' = M'$.
\end{lemma}

\begin{proof}
The adjointness statement follows from the more general
Lemma \ref{lemma-modules}.
To prove the final assertion, recall that
$B' = B \times_A A'$ and $N' = N \times_{\varphi, M} M'$ and extend
these equalities to 
$$
\vcenter{
\xymatrix{
A & A' \ar[l] & I \ar[l] \\
B \ar[u] & B' \ar[l] \ar[u] & J \ar[l] \ar[u]
}
}
\quad\text{and}\quad
\vcenter{
\xymatrix{
M & M' \ar[l] & K \ar[l] \\
N \ar[u]_\varphi & N' \ar[l] \ar[u] & L \ar[l] \ar[u]
}
}
$$
where $I, J, K, L$ are the kernels of the horizontal maps of the original
diagrams. We present the proof as a sequence of observations:
\begin{enumerate}
\item $K = IM'$ (see statement lemma),
\item $B' \to B$ is surjective with kernel $J$ and $J \to I$ is bijective,
\item $N' \to N$ is surjective with kernel $L$ and $L \to K$ is bijective,
\item $JN' \subset L$,
\item $\Im(N \to M)$ generates $M$ as an $A$-module
(because $N \otimes_B A = M$),
\item $\Im(N' \to M')$ generates $M'$ as an $A'$-module
(because it holds modulo $K$ and $L$ maps isomorphically to $K$),
\item $JN' = L$ (because $L \cong K = I M'$
is generated by images of elements $x n'$ with $x \in I$ and
$n' \in N'$ by the previous statement),
\item $N' \otimes_{B'} B = N$ (because $N = N'/L$, $B = B'/J$, and
the previous statement),
\item there is a map $\gamma : N' \otimes_{B'} A' \to M'$,
\item $\gamma$ is surjective (see above),
\item the kernel of the composition $N' \otimes_{B'} A' \to M' \to M$ 
is generated by elements $l \otimes 1$ and $n' \otimes x$ with
$l \in K$, $n' \in N'$, $x \in I$ (because $M = N \otimes_B A$ by assumption
and because $N' \to N$ and $A' \to A$ are surjective with kernels
$L$ and $I$),
\item any element of $N' \otimes_{B'} A'$ in the submodule generated
by the elements $l \otimes 1$ and $n' \otimes x$ with
$l \in L$, $n' \in N'$, $x \in I$ can be written as $l \otimes 1$
for some $l \in L$ (because $J$ maps isomorphically to $I$ we see
that $n' \otimes x = n'x \otimes 1$ in $N' \otimes_{B'} A'$;
similarly $x n' \otimes a' = n' \otimes xa' = n'(xa') \otimes 1$
in $N' \otimes_{B'} A'$ when $n' \in N'$, $x \in J$ and $a' \in A'$;
since we have seen that $JN' = L$ this proves the assertion),
\item the kernel of $\gamma$ is zero (because by (10) and (11) any element of
the kernel is of the form $l \otimes 1$ with $l \in L$ which
is mapped to $l \in K \subset M'$ by $\gamma$).
\end{enumerate}
This finishes the proof.
\end{proof}

\begin{lemma}
\label{lemma-module-over-fibre-product-bis}
In the situation of Lemma \ref{lemma-module-over-fibre-product}
for a $B'$-module $L'$ the adjunction map
$$
L' \longrightarrow 
(L' \otimes_{B'} B) \times_{(L' \otimes_{B'} A)} (L' \otimes_{B'} A')
$$
is surjective but in general not injective.
\end{lemma}

\begin{proof}
As in the proof of Lemma \ref{lemma-module-over-fibre-product}
let $J \subset B'$ be the kernel of the map $B' \to B$.
Then $L' \otimes_{B'} B = L'/JL'$. Hence to prove surjectivity it suffices
to show that elements of the form $(0, z)$ of the fibre product are in the
image of the map of the lemma. The kernel of the map
$L' \otimes_{B'} A' \to L' \otimes_{B'} A$ is the image of
$L' \otimes_{B'} I \to L' \otimes_{B'} A'$. Since the map $J \to I$
induced by $B' \to A'$ is an isomorphism
the composition
$$
L' \otimes_{B'} J \to L' \to
(L' \otimes_{B'} B) \times_{(L' \otimes_{B'} A)} (L' \otimes_{B'} A')
$$
induces a surjection of $L' \otimes_{B'} J$ onto the set of elements
of the form $(0, z)$. To see the map is not injective in general we
present a simple example. Namely, take a field $k$,
set $B' = k[x, y]/(xy)$, $A' = B'/(x)$, $B = B'/(y)$, $A = B'/(x, y)$
and $L' = B'/(x - y)$. In that case the class of $x$ in $L'$ is nonzero
but is mapped to zero under the displayed arrow.
\end{proof}

\begin{lemma}
\label{lemma-surjection-module-over-fibre-product}
In Situation \ref{situation-module-over-fibre-product}
let $(N_1, M'_1, \varphi_1) \to (N_2, M'_2, \varphi_2)$ be a morphism
of $\text{Mod}_B \times_{\text{Mod}_A} \text{Mod}_{A'}$
with $N_1 \to N_2$ and $M'_1 \to M'_2$ surjective. Then
$$
N_1 \times_{\varphi_1, M_1} M'_1 \to N_2 \times_{\varphi_2, M_2} M'_2
$$
where $M_1 = M'_1/IM'_1$ and $M_2 = M'_2/IM'_2$ is surjective.
\end{lemma}

\begin{proof}
Pick $(x_2, y_2) \in N_2 \times_{\varphi_2, M_2} M'_2$. Choose $x_1 \in N_1$
mapping to $x_2$. Since $M'_1 \to M_1$ is surjective we can find
$y_1 \in M'_1$ mapping to $\varphi_1(x_1)$. Then $(x_1, y_1)$
maps to $(x_2, y'_2)$ in $N_2 \times_{\varphi_2, M_2} M'_2$. Thus it suffices
to show that elements of the form $(0, y_2)$ are in the image of the map.
Here we see that $y_2 \in IM'_2$. Write $y_2 = \sum t_i y_{2, i}$
with $t_i \in I$. Choose $y_{1, i} \in M'_1$ mapping to $y_{2, i}$.
Then $y_1 = \sum t_iy_{1, i} \in IM'_1$ and the element $(0, y_1)$
does the job.
\end{proof}

\begin{lemma}
\label{lemma-finite-module-over-fibre-product}
Let $A, A', B, B', I, M, M', N, \varphi$ be as in
Lemma \ref{lemma-module-over-fibre-product}.
If $N$ finite over $B$ and $M'$ finite over $A'$, then
$N' = N \times_{\varphi, M} M'$ is finite over $B'$.
\end{lemma}

\begin{proof}
We will use the results of
Lemma \ref{lemma-module-over-fibre-product}
without further mention. Choose generators $y_1, \ldots, y_r$ of $N$ over $B$
and generators $x_1, \ldots, x_s$ of $M'$ over $A'$. Using that
$N = N' \otimes_{B'} B$ and $B' \to B$ is surjective we can find
$u_1, \ldots, u_r \in N'$ mapping to $y_1, \ldots, y_r$ in $N$.
Using that $M' = N' \otimes_{B'} A'$ we can find $v_1, \ldots, v_t \in N'$
such that $x_i = \sum v_j \otimes a'_{ij}$ for some $a'_{ij} \in A'$.
In particular we see that the images $\overline{v}_j \in M'$
of the $v_j$ generate $M'$ over $A'$.
We claim that $u_1, \ldots, u_r, v_1, \ldots, v_t$
generate $N'$ as a $B'$-module. Namely, pick $\xi \in N'$. We first choose
$b'_1, \ldots, b'_r \in B'$ such that $\xi$ and $\sum b'_i u_i$ map
to the same element of $N$. This is possible because $B' \to B$
is surjective and $y_1, \ldots, y_r$ generate $N$ over $B$.
The difference $\xi - \sum b'_i u_i$ is of the form $(0, \theta)$
for some $\theta$ in $IM'$. Say $\theta$ is $\sum t_j\overline{v}_j$
with $t_j \in I$. As $J = \Ker(B' \to B)$ maps isomorphically to $I$
we can choose $s_j \in J \subset B'$ mapping to $t_j$.
Because $N' = N \times_{\varphi, M} M'$ it follows
that $\xi = \sum b'_i u_i + \sum s_j v_j$ as desired.
\end{proof}

\begin{lemma}
\label{lemma-flat-module-over-fibre-product}
With $A, A', B, B', I$ as in
Situation \ref{situation-module-over-fibre-product}.
\begin{enumerate}
\item Let $(N, M', \varphi)$ be an object of
$\text{Mod}_B \times_{\text{Mod}_A} \text{Mod}_{A'}$.
If $M'$ is flat over $A'$ and $N$ is flat over $B$, then
$N' = N \times_{\varphi, M} M'$ is flat over $B'$.
\item If $L'$ is a flat $B'$-module, then
$L' = (L \otimes_{B'} B) \times_{(L \otimes_{B'} A)} (L \otimes_{B'} A')$.
\item The category of flat $B'$-modules is equivalent to the
full subcategory of $\text{Mod}_B \times_{\text{Mod}_A} \text{Mod}_{A'}$
consisting of triples
$(N, M', \varphi)$ with $N$ flat over $B$ and $M'$ flat over $A'$.
\end{enumerate}
\end{lemma}

\begin{proof}
In the proof we will use Lemma \ref{lemma-module-over-fibre-product}
without further mention.

\medskip\noindent
Proof of (1). Set $J = \Ker(B' \to B)$. This is an ideal of $B'$
mapping isomorphically to $I = \Ker(A' \to A)$.
Let $\mathfrak b' \subset B'$ be an ideal.
We have to show that $\mathfrak b' \otimes_{B'} N' \to N'$
is injective, see Algebra, Lemma \ref{algebra-lemma-flat}.
We know that
$$
\mathfrak b'/(\mathfrak b' \cap J) \otimes_{B'} N' =
\mathfrak b'/(\mathfrak b' \cap J) \otimes_B N \to N
$$
is injective as $N$ is flat over $B$. As
$\mathfrak b' \cap J \to \mathfrak b' \to
\mathfrak b'/(\mathfrak b' \cap J) \to 0$ is exact, we
conclude that it suffices to show that
$(\mathfrak b' \cap J) \otimes_{B'} N' \to N'$
is injective. Thus we may assume that $\mathfrak b' \subset J$.
Next, since $J \to I$ is an isomorphism we have
$$
J \otimes_{B'} N' =
I \otimes_{A'} A' \otimes_{B'} N' =
I \otimes_{A'} M'
$$
which maps injectively into $M'$ as $M'$ is a flat $A'$-module.
Hence $J \otimes_{B'} N' \to N'$ is injective and we conclude
that $\text{Tor}_1^{B'}(B'/J, N') = 0$, see
Algebra, Remark \ref{algebra-remark-Tor-ring-mod-ideal}.
Thus we may apply Algebra, Lemma \ref{algebra-lemma-what-does-it-mean}
to $N'$ over $B'$ and the ideal $J$.
Going back to our ideal $\mathfrak b' \subset J$, let
$\mathfrak b' \subset \mathfrak b'' \subset J$ be the smallest ideal
whose image in $I$ is an $A'$-submodule of $I$. In other words,
we have $\mathfrak b'' = A' \mathfrak b'$ if we view $J = I$
as $A'$-module. Then $\mathfrak b''/\mathfrak b'$ is killed
by $J$ and we get a short exact sequence
$$
0 \to \mathfrak b' \otimes_{B'} N' \to
\mathfrak b'' \otimes_{B'} N' \to
\mathfrak b''/\mathfrak b' \otimes_{B'} N' \to 0
$$
by the vanishing of $\text{Tor}_1^{B'}(\mathfrak b''/\mathfrak b', N')$
we get from the application of the lemma.
Thus we may replace $\mathfrak b'$ by $\mathfrak b''$.
In particular we may assume $\mathfrak b'$ is an $A'$-module
and maps to an ideal of $A'$. Then
$$
\mathfrak b' \otimes_{B'} N' =
\mathfrak b' \otimes_{A'} A' \otimes_{B'} N' =
\mathfrak b' \otimes_{A'} M'
$$
This tensor product maps injectively into $M'$ by our assumption that
$M'$ is flat over $A'$. We conclude that
$\mathfrak b' \otimes_{B'} N' \to N' \to M'$ is injective
and hence the first map is injective as desired.

\medskip\noindent
Proof of (2). This follows by tensoring the short exact sequence
$0 \to B' \to B \oplus A' \to A \to 0$ with $L'$ over $B'$.

\medskip\noindent
Proof of (3). Immediate consequence of (1) and (2).
\end{proof}

\begin{lemma}
\label{lemma-finitely-presented-module-over-fibre-product}
Let $A, A', B, B', I$ be as in
Situation \ref{situation-module-over-fibre-product}.
The category of finite projective $B'$-modules
is equivalent to the full subcategory
of $\text{Mod}_B \times_{\text{Mod}_A} \text{Mod}_{A'}$
consisting of triples $(N, M', \varphi)$
with $N$ finite projective over $B$ and $M'$ finite projective over $A'$.
\end{lemma}

\begin{proof}
Recall that a module is finite projective if and only if
it is finitely presented and flat, see
Algebra, Lemma \ref{algebra-lemma-finite-projective}.
Using Lemmas \ref{lemma-flat-module-over-fibre-product} and
\ref{lemma-finite-module-over-fibre-product}
we reduce to showing that $N' = N \times_{\varphi, M} M'$ is
a $B'$-module of finite presentation if
$N$ finite projective over $B$
and $M'$ finite projective over $A'$.

\medskip\noindent
By Lemma \ref{lemma-finite-module-over-fibre-product}
the module $N'$ is finite over $B'$. Choose a surjection
$(B')^{\oplus n} \to N'$ with kernel $K'$. By base change we obtain maps
$B^{\oplus n} \to N$, $(A')^{\oplus n} \to M'$, and $A^{\oplus n} \to M$
with kernels $K_B$, $K_{A'}$, and $K_A$. There is a canonical map
$$
K' \longrightarrow K_B \times_{K_A} K_{A'}
$$
On the other hand, since $N' = N \times_{\varphi, M} M'$ and
$B' = B \times_A A'$ there is also a
canonical map $K_B \times_{K_A} K_{A'} \to K'$ inverse to the displayed
arrow. Hence the displayed map is an isomorphism. By
Algebra, Lemma \ref{algebra-lemma-extension}
the modules $K_B$ and $K_{A'}$ are finite. We conclude from
Lemma \ref{lemma-finite-module-over-fibre-product}
that $K'$ is a finite $B'$-module provided that $K_B \to K_A$ and
$K_{A'} \to K_A$ induce isomorphisms
$K_B \otimes_B A = K_A = K_{A'} \otimes_{A'} A$.
This is true because the flatness assumptions implies the sequences
$$
0 \to K_B \to B^{\oplus n} \to N \to 0
\quad\text{and}\quad
0 \to K_{A'} \to (A')^{\oplus n} \to M' \to 0
$$
stay exact upon tensoring, see
Algebra, Lemma \ref{algebra-lemma-flat-tor-zero}.
\end{proof}





\section{Fibre products of rings, III}
\label{section-fibre-products-rings-III}

\noindent
In this section we discuss fibre products in the following situation.

\begin{situation}
\label{situation-relative-module-over-fibre-product}
Let $A, A', B, B', I$ be as in
Situation \ref{situation-module-over-fibre-product}.
Let $B' \to D'$ be a ring map. Set $D = D' \otimes_{B'} B$,
$C' = D' \otimes_{B'} A'$, and $C = D' \otimes_{B'} A$. This leads
to a big commutative diagram
$$
\xymatrix{
C & & & C' \ar[lll] \\
& A \ar[ul] & A' \ar[l] \ar[ru] \\
& B \ar[u] \ar[ld] & B' \ar[l] \ar[u] \ar[rd] \\
D \ar[uuu] & & & D' \ar[lll] \ar[uuu]
}
$$
of rings.
Observe that we do {\bf not} assume that the map $D' \to D \times_C C'$
is an isomorphism\footnote{But $D' \to D \times_C C'$ is surjective
by Lemma \ref{lemma-module-over-fibre-product-bis}.}.
In this situation we have the functor
\begin{equation}
\label{equation-relative-functor}
\text{Mod}_{D'} \longrightarrow
\text{Mod}_D \times_{\text{Mod}_C} \text{Mod}_{C'},\quad
L' \longmapsto (L' \otimes_{D'} D, L' \otimes_{D'} C', can)
\end{equation}
analogous to (\ref{equation-functor}). Note that
$L' \otimes_{D'} D = L \otimes_{D'} (D' \otimes_{B'} B) = L \otimes_{B'} B$
and similarly
$L' \otimes_{D'} C' = L \otimes_{D'} (D' \otimes_{B'} A') = L \otimes_{B'} A'$
hence the diagram
$$
\xymatrix{
\text{Mod}_{D'} \ar[r] \ar[d] &
\text{Mod}_D \times_{\text{Mod}_C} \text{Mod}_{C'} \ar[d] \\
\text{Mod}_{B'} \ar[r] &
\text{Mod}_B \times_{\text{Mod}_A} \text{Mod}_{A'}
}
$$
is commutative. In the following we will write $(N, M', \varphi)$ for an
object of $\text{Mod}_D \times_{\text{Mod}_C} \text{Mod}_{C'}$,
i.e., $N$ is a $D$-module, $M'$ is an $C'$-module
and $\varphi : N \otimes_B A \to M' \otimes_{A'} A$ is an isomorphism
of $C$-modules.
However, it is often more convenient think of $\varphi$ as a $D$-linear
map $\varphi : N \to M'/IM'$ which induces an isomorphism
$N \otimes_B A \to M' \otimes_{A'} A = M'/IM'$.
\end{situation}

\begin{lemma}
\label{lemma-relative-module-over-fibre-product}
In Situation \ref{situation-relative-module-over-fibre-product}
the functor (\ref{equation-relative-functor}) has a right adjoint, namely
the functor
$$
F : (N, M', \varphi) \longmapsto N \times_{\varphi, M} M'
$$
where $M = M'/IM'$. Moreover, the composition of $F$ with
(\ref{equation-relative-functor}) is the identity functor on
$\text{Mod}_D \times_{\text{Mod}_C} \text{Mod}_{C'}$. In other words,
setting $N' = N \times_{\varphi, M} M'$ we have
$N' \otimes_{D'} D = N$ and $N' \otimes_{D'} C' = M'$.
\end{lemma}

\begin{proof}
The adjointness statement follows from the more general
Lemma \ref{lemma-modules}.
The final assertion follows from the corresponding assertion of
Lemma \ref{lemma-module-over-fibre-product}
because
$N' \otimes_{D'} D = N' \otimes_{D'} D' \otimes_{B'} B = N' \otimes_{B'} B$
and
$N' \otimes_{D'} C' = N' \otimes_{D'} D' \otimes_{B'} A' = N' \otimes_{B'} A'$.
\end{proof}

\begin{lemma}
\label{lemma-relative-surjection-ideals}
In Situation \ref{situation-relative-module-over-fibre-product}
the map $JD' \to IC'$ is surjective where $J = \Ker(B' \to B)$.
\end{lemma}

\begin{proof}
Since $C' = D' \otimes_{B'} A'$ we have that $IC'$ is the image
of $D' \otimes_{B'} I = C' \otimes_{A'} I \to C'$. As the ring
map $B' \to A'$ induces an isomorphism $J \to I$ the lemma follows.
\end{proof}

\begin{lemma}
\label{lemma-relative-finite-module-over-fibre-product}
Let $A, A', B, B', C, C', D, D', I, M', M, N, \varphi$ be as in
Lemma \ref{lemma-relative-module-over-fibre-product}.
If $N$ finite over $D$ and $M'$ finite over $C'$, then
$N' = N \times_{\varphi, M} M'$ is finite over $D'$.
\end{lemma}

\begin{proof}
Recall that $D' \to D \times_C C'$ is surjective by
Lemma \ref{lemma-module-over-fibre-product-bis}.
Observe that $N' = N \times_{\varphi, M} M'$ is a module
over $D \times_C C'$. We can apply
Lemma \ref{lemma-finite-module-over-fibre-product}
to the data $C, C', D, D', IC', M', M, N, \varphi$
to see that $N' = N \times_{\varphi, M} M'$ is finite
over $D \times_C C'$. Thus it is finite over $D'$.
\end{proof}

\begin{lemma}
\label{lemma-relative-flat-module-over-fibre-product}
With $A, A', B, B', C, C', D, D', I$ as in
Situation \ref{situation-relative-module-over-fibre-product}.
\begin{enumerate}
\item Let $(N, M', \varphi)$ be an object of
$\text{Mod}_D \times_{\text{Mod}_C} \text{Mod}_{C'}$.
If $M'$ is flat over $A'$ and $N$ is flat over $B$, then
$N' = N \times_{\varphi, M} M'$ is flat over $B'$.
\item If $L'$ is a $D'$-module flat over $B'$, then
$L' = (L \otimes_{D'} D) \times_{(L \otimes_{D'} C)} (L \otimes_{D'} C')$.
\item The category of $D'$-modules flat over $B'$
is equivalent to the categories of objects $(N, M', \varphi)$
of $\text{Mod}_D \times_{\text{Mod}_C} \text{Mod}_{C'}$
with $N$ flat over $B$ and $M'$ flat over $A'$.
\end{enumerate}
\end{lemma}

\begin{proof}
Part (1) follows from part (1) of
Lemma \ref{lemma-flat-module-over-fibre-product}.

\medskip\noindent
Part (2) follows from part (2) of
Lemma \ref{lemma-flat-module-over-fibre-product}
using that $L' \otimes_{D'} D = L' \otimes_{B'} B$,
$L' \otimes_{D'} C' = L' \otimes_{B'} A'$, and
$L' \otimes_{D'} C = L' \otimes_{B'} A$, see discussion in
Situation \ref{situation-relative-module-over-fibre-product}.

\medskip\noindent
Part (3) is an immediate consequence of (1) and (2).
\end{proof}

\noindent
The following lemma is a good deal more interesting than its
counter part in the absolute case
(Lemma \ref{lemma-finitely-presented-module-over-fibre-product}),
although the proof is essentially the same.

\begin{lemma}
\label{lemma-relative-finitely-presented-module-over-fibre-product}
Let $A, A', B, B', C, C', D, D', I, M', M, N, \varphi$ be as in
Lemma \ref{lemma-relative-module-over-fibre-product}. If
\begin{enumerate}
\item $N$ is finitely presented over $D$ and flat over $B$,
\item $M'$ finitely presented over $C'$ and flat over $A'$, and
\item the ring map $B' \to D'$ factors as $B' \to D'' \to D'$
with $B' \to D''$ flat and $D'' \to D'$ of finite presentation,
\end{enumerate}
then $N' = N \times_M M'$ is finitely presented over $D'$.
\end{lemma}

\begin{proof}
Choose a surjection $D''' = D''[x_1, \ldots, x_n] \to D'$ with
finitely generated kernel $J$.
By Algebra, Lemma \ref{algebra-lemma-finite-finitely-presented-extension}
it suffices to show that $N'$ is finitely presented as a
$D'''$-module. Moreover, $D''' \otimes_{B'} B \to D' \otimes_{B'} B = D$
and $D''' \otimes_{B'} A' \to D' \otimes_{B'} A' = C'$ are surjections
whose kernels are generated by the image of $J$, hence $N$ is a
finitely presented $D''' \otimes_{B'} B$-module and
$M'$ is a finitely presented $D''' \otimes_{B'} A'$-module by
Algebra, Lemma \ref{algebra-lemma-finite-finitely-presented-extension}
again. Thus we may replace $D'$ by $D'''$ and $D$ by
$D''' \otimes_{B'} B$, etc. Since $D'''$ is
flat over $B'$, it follows that we may assume that $B' \to D'$ is flat.

\medskip\noindent
Assume $B' \to D'$ is flat.
By Lemma \ref{lemma-relative-finite-module-over-fibre-product}
the module $N'$ is finite over $D'$. Choose a surjection
$(D')^{\oplus n} \to N'$ with kernel $K'$. By base change we obtain maps
$D^{\oplus n} \to N$, $(C')^{\oplus n} \to M'$, and $C^{\oplus n} \to M$
with kernels $K_D$, $K_{C'}$, and $K_C$. There is a canonical map
$$
K' \longrightarrow K_D \times_{K_C} K_{C'}
$$
On the other hand, since $N' = N \times_M M'$ and
$D' = D \times_C C'$ (by Lemma \ref{lemma-flat-module-over-fibre-product};
applied to the flat $B'$-module $D'$)
there is also a
canonical map $K_D \times_{K_C} K_{C'} \to K'$ inverse to the displayed
arrow. Hence the displayed map is an isomorphism. By
Algebra, Lemma \ref{algebra-lemma-extension}
the modules $K_D$ and $K_{C'}$ are finite. We conclude from
Lemma \ref{lemma-relative-finite-module-over-fibre-product}
that $K'$ is a finite $D'$-module provided that $K_D \to K_C$ and
$K_{C'} \to K_C$ induce isomorphisms
$K_D \otimes_B A = K_C = K_{C'} \otimes_{A'} A$.
This is true because the flatness assumptions implies the sequences
$$
0 \to K_D \to D^{\oplus n} \to N \to 0
\quad\text{and}\quad
0 \to K_{C'} \to (C')^{\oplus n} \to M' \to 0
$$
stay exact upon tensoring, see
Algebra, Lemma \ref{algebra-lemma-flat-tor-zero}.
\end{proof}

\begin{lemma}
\label{lemma-properties-algebras-over-fibre-product}
Let $A, A', B, B', I$ be as in
Situation \ref{situation-module-over-fibre-product}.
Let $(D, C', \varphi)$ be a system consisting of an $B$-algebra $D$,
a $A'$-algebra $C'$ and an isomorphism $D \otimes_B A \to C'/IC' = C$.
Set $D' = D \times_C C'$ (as in
Lemma \ref{lemma-module-over-fibre-product}). Then
\begin{enumerate}
\item $B' \to D'$ is finite type if and only if $B \to D$ and
$A' \to C'$ are finite type,
\item $B' \to D'$ is flat if and only if $B \to D$ and $A' \to C'$ are flat,
\item $B' \to D'$ is flat and of finite presentation if and only if
$B \to D$ and $A' \to C'$ are flat and of finite presentation,
\item $B' \to D'$ is smooth if and only if $B \to D$ and $A' \to C'$ are smooth,
\item $B' \to D'$ is \'etale if and only if $B \to D$ and $A' \to C'$
are \'etale.
\end{enumerate}
Moreover, if $D'$ is a flat $B'$-algebra, then
$D' \to (D' \otimes_{B'} B) \times_{(D' \otimes_{B'} A)} (D' \otimes_{B'} A')$
is an isomorphism. In this way the category of flat $B'$-algebras
is equivalent to the categories of systems $(D, C', \varphi)$ as above
with $D$ flat over $B$ and $C'$ flat over $A'$.
\end{lemma}

\begin{proof}
The implication ``$\Rightarrow$'' follows from
Algebra, Lemmas \ref{algebra-lemma-base-change-finiteness},
\ref{algebra-lemma-flat-base-change},
\ref{algebra-lemma-base-change-smooth}, and
\ref{algebra-lemma-etale} because we have
$D' \otimes_{B'} B = D$ and $D' \otimes_{B'} A' = C'$
by Lemma \ref{lemma-module-over-fibre-product}.
Thus it suffices to prove the implications in the other direction.

\medskip\noindent
Ad (1). Assume $D$ of finite type over $B$ and $C'$ of finite type over $A'$.
We will use the results of
Lemma \ref{lemma-module-over-fibre-product}
without further mention. Choose generators $x_1, \ldots, x_r$ of $D$ over $B$
and generators $y_1, \ldots, y_s$ of $C'$ over $A'$. Using that
$D = D' \otimes_{B'} B$ and $B' \to B$ is surjective we can find
$u_1, \ldots, u_r \in D'$ mapping to $x_1, \ldots, x_r$ in $D$.
Using that $C' = D' \otimes_{B'} A'$ we can find $v_1, \ldots, v_t \in D'$
such that $y_i = \sum v_j \otimes a'_{ij}$ for some $a'_{ij} \in A'$.
In particular, the images of $v_j$ in $C'$ generate $C'$ as an
$A'$-algebra. Set $N = r + t$ and consider the cube of rings
$$
\xymatrix{
A[x_1, \ldots, x_N] & & A'[x_1, \ldots, x_N] \ar[ll] \\
& A \ar[lu] & & A' \ar[ll] \ar[lu] \\
B[x_1, \ldots, x_N] \ar[uu] & & B'[x_1, \ldots, x_N] \ar[uu] \ar[ll] \\
& B \ar[uu] \ar[lu] & & B' \ar[ll] \ar[uu] \ar[lu]
}
$$
Observe that the back square is cartesian as well.
Consider the ring map
$$
B'[x_1, \ldots, x_N] \to D',\quad
x_i \mapsto u_i \quad\text{and}\quad x_{r + j} \mapsto v_j.
$$
Then we see that the induced maps $B[x_1, \ldots, x_N] \to D$ and
$A'[x_1, \ldots, x_N] \to C'$
are surjective, in particular finite. We conclude from
Lemma \ref{lemma-relative-finite-module-over-fibre-product}
that $B'[x_1, \ldots, x_N] \to D'$ is finite, which implies that $D'$
is of finite type over $B'$ for example by
Algebra, Lemma \ref{algebra-lemma-compose-finite-type}.

\medskip\noindent
Ad (2). The implication ``$\Leftarrow$'' follows from
Lemma \ref{lemma-relative-flat-module-over-fibre-product}.
Moreover, the final statement follows from the final
statement of Lemma \ref{lemma-relative-flat-module-over-fibre-product}.

\medskip\noindent
Ad (3). Assume $B \to D$ and $A' \to C'$ are flat and of finite presentation.
The flatness of $B' \to D'$ we've seen in (2). We know $B' \to D'$
is of finite type by (1). Choose a surjection $B'[x_1, \ldots, x_N] \to D'$.
By Algebra, Lemma \ref{algebra-lemma-finite-presentation-independent}
the ring $D$ is of finite presentation as a $B[x_1, \ldots, x_N]$-module
and the ring $C'$ is of finite presentation as a
$A'[x_1, \ldots, x_N]$-module. By
Lemma \ref{lemma-relative-finitely-presented-module-over-fibre-product}
we see that $D'$ is of finite presentation as a $B'[x_1, \ldots, x_N]$-module,
i.e., $B' \to D'$ is of finite presentation.

\medskip\noindent
Ad (4). Assume $B \to D$ and $A' \to C'$ smooth.
By (3) we see that $B' \to D'$ is flat and of finite presentation.
By Algebra, Lemma \ref{algebra-lemma-flat-fibre-smooth}
it suffices to check that $D' \otimes_{B'} k$ is smooth for any
field $k$ over $B'$. If the composition $J \to B' \to k$ is zero,
then $B' \to k$ factors as $B' \to B \to k$ and we see that
$$
D' \otimes_{B'} k = D' \otimes_{B'} B \otimes_B k
= D \otimes_B k
$$
is smooth as $B \to D$ is smooth. If the composition $J \to B' \to k$
is nonzero, then there exists an $h \in J$ which does not map to zero
in $k$. Then $B' \to k$ factors as $B' \to B'_h \to k$.
Observe that $h$ maps to zero in $B$, hence $B_h = 0$.
Thus by Lemma \ref{lemma-diagram-localize} we have $B'_h = A'_h$ and we get
$$
D' \otimes_{B'} k = D' \otimes_{B'} B'_h \otimes_{B'_h} k
= C'_h \otimes_{A'_h} k
$$
is smooth as $A' \to C'$ is smooth.

\medskip\noindent
Ad (5). Assume $B \to D$ and $A' \to C'$ are \'etale. By (4) we see that
$B' \to D'$ is smooth. As we can read off whether or not a smooth
map is \'etale from the dimension of fibres we see that (5) holds
(argue as in the proof of (4) to identify fibres -- some details omitted).
\end{proof}

\begin{remark}
\label{remark-relative-modules-over-fibre-product}
In Situation \ref{situation-relative-module-over-fibre-product}.
Assume $B' \to D'$ is of finite presentation and
suppose we are given a $D'$-module $L'$.
We claim there is a bijective correspondence between
\begin{enumerate}
\item surjections of $D'$-modules $L' \to Q'$ with $Q'$ of finite presentation
over $D'$ and flat over $B'$, and
\item pairs of surjections of modules
$(L' \otimes_{D'} D \to Q_1, L' \otimes_{D'} C' \to Q_2)$
with
\begin{enumerate}
\item $Q_1$ of finite presentation over $D$ and flat over $B$,
\item $Q_2$ of finite presentation over $C'$ and flat over $A'$,
\item $Q_1 \otimes_D C = Q_2 \otimes_{C'} C$ as quotients of
$L' \otimes_{D'} C$.
\end{enumerate}
\end{enumerate}
The correspondence between these is given by $Q \mapsto (Q_1, Q_2)$ with
$Q_1 = Q \otimes_{D'} D$ and $Q_2 = Q \otimes_{D'} C'$. And for the converse
we use $Q = Q_1 \times_{Q_{12}} Q_2$ where $Q_{12}$ the common quotient
$Q_1 \otimes_D C = Q_2 \otimes_{C'} C$ of $L' \otimes_{D'} C$. As quotient
map we use
$$
L' \longrightarrow
(L' \otimes_{D'} D) \times_{(L' \otimes_{D'} C)} (L' \otimes_{D'} C')
\longrightarrow Q_1 \times_{Q_{12}} Q_2 = Q
$$
where the first arrow is surjective by
Lemma \ref{lemma-module-over-fibre-product-bis}
and the second by Lemma \ref{lemma-surjection-module-over-fibre-product}.
The claim follows by
Lemmas \ref{lemma-relative-flat-module-over-fibre-product} and
\ref{lemma-relative-finitely-presented-module-over-fibre-product}.
\end{remark}








\section{Fitting ideals}
\label{section-fitting-ideals}

\noindent
The Fitting ideals of a finite module are the ideals determined
by the construction of Lemma \ref{lemma-fitting-ideal}.

\begin{lemma}
\label{lemma-ideals-generated-by-minors}
Let $R$ be a ring. Let $A$ be an $n \times m$ matrix with coefficients
in $R$. Let $I_r(A)$ be the ideal generated by the $r \times r$-minors
of $A$ with the convention that $I_0(A) = R$ and $I_r(A) = 0$ if
$r > \min(n, m)$. Then
\begin{enumerate}
\item $I_0(A) \supset I_1(A) \supset I_2(A) \supset \ldots$,
\item if $B$ is an $(n + n') \times m$ matrix, and $A$ is the first
$n$ rows of $B$, then $I_{r + n'}(B) \subset I_r(A)$,
\item if $C$ is an $n \times n$ matrix then $I_r(CA) \subset I_r(A)$.
\item If $A$ is a block matrix
$$
\left(
\begin{matrix}
A_1 & 0 \\
0 & A_2 
\end{matrix}
\right)
$$
then $I_r(A) = \sum_{r_1 + r_2 = r} I_{r_1}(A_1) I_{r_2}(A_2)$.
\item Add more here.
\end{enumerate}
\end{lemma}

\begin{proof}
Omitted. (Hint: Use that a determinant can be computed by expanding
along a column or a row.)
\end{proof}

\begin{lemma}
\label{lemma-fitting-ideal}
Let $R$ be a ring. Let $M$ be a finite $R$-module. Choose a presentation
$$
\bigoplus\nolimits_{j \in J} R \longrightarrow R^{\oplus n}
\longrightarrow M \longrightarrow 0.
$$
of $M$. Let $A = (a_{ij})_{i = 1, \ldots, n, j \in J}$ be the matrix
of the map $\bigoplus_{j \in J} R \to R^{\oplus n}$.
The ideal $\text{Fit}_k(M)$ generated by the
$(n - k) \times (n - k)$ minors of
$A$ is independent of the choice of the presentation.
\end{lemma}

\begin{proof}
Let $K \subset R^{\oplus n}$ be the kernel of the surjection
$R^{\oplus n} \to M$. Pick $z_1, \ldots, z_{n - k} \in K$
and write $z_j = (z_{1j}, \ldots, z_{nj})$.
Another description of the ideal $\text{Fit}_k(M)$
is that it is the ideal generated by the $(n - k) \times (n - k)$ minors of
all the matrices $(z_{ij})$ we obtain in this way.

\medskip\noindent
Suppose we change the surjection into the surjection
$R^{\oplus n + n'} \to M$ with kernel $K'$ where we use the original
map on the first $n$ standard basis elements of $R^{\oplus n + n'}$
and $0$ on the last $n'$ basis vectors. Then the corresponding ideals
are the same. Namely, if $z_1, \ldots, z_{n - k} \in K$ as above,
let $z'_j = (z_{1j}, \ldots, z_{nj}, 0, \ldots, 0) \in K'$ for
$j = 1, \ldots, n - k$ and
$z'_{n + j'} = (0, \ldots, 0, 1, 0, \ldots, 0) \in K'$. Then we see that
the ideal of $(n - k) \times (n - k)$ minors of $(z_{ij})$ agrees
with the ideal of $(n + n' - k) \times (n + n' - k)$ minors of
$(z'_{ij})$. This gives one of the inclusions.
Conversely, given  $z'_1, \ldots, z'_{n + n' - k}$
in $K'$ we can project these to $R^{\oplus n}$ to get
$z_1, \ldots, z_{n + n' - k}$ in $K$. By
Lemma \ref{lemma-ideals-generated-by-minors}
we see that the ideal generated by the
$(n + n' - k) \times (n + n' - k)$ minors of
$(z'_{ij})$ is contained in the ideal generated by the
$(n - k) \times (n - k)$ minors of $(z_{ij})$. This gives the
other inclusion.

\medskip\noindent
Let $R^{\oplus m} \to M$ be another surjection with kernel $L$.
By Schanuel's lemma (Algebra, Lemma \ref{algebra-lemma-Schanuel})
and the results of the previous paragraph, we may assume $m = n$
and that there is an isomorphism $R^{\oplus n} \to R^{\oplus m}$
commuting with the surjections to $M$. Let $C = (c_{li})$ be the
(invertible) matrix of this map (it is a square matrix as $n = m$).
Then given $z'_1, \ldots, z'_{n - k} \in L$ as above we can find
$z_1, \ldots, z_{n - k} \in K$ with
$z_1' = Cz_1, \ldots, z'_{n - k} = Cz_{n - k}$. By
Lemma \ref{lemma-ideals-generated-by-minors} we get one of the
inclusions. By symmetry we get the other.
\end{proof}

\begin{definition}
\label{definition-fitting-ideal}
Let $R$ be a ring. Let $M$ be a finite $R$-module. Let $k \geq 0$.
The {\it $k$th Fitting ideal} of $M$ is the ideal $\text{Fit}_k(M)$
constructed in Lemma \ref{lemma-fitting-ideal}. Set $\text{Fit}_{-1}(M) = 0$.
\end{definition}

\noindent
Since the Fitting ideals are the ideals of minors of a big matrix
(numbered in reverse ordering from the ordering in
Lemma \ref{lemma-ideals-generated-by-minors})
we see that
$$
0 = \text{Fit}_{-1}(M) \subset \text{Fit}_0(M) \subset \text{Fit}_1(M)
\subset \ldots \subset \text{Fit}_t(M) = R
$$
for some $t \gg 0$. Here are some basic properties of Fitting ideals.

\begin{lemma}
\label{lemma-fitting-ideal-basics}
Let $R$ be a ring. Let $M$ be a finite $R$-module.
\begin{enumerate}
\item If $M$ can be generated by $n$ elements, then
$\text{Fit}_n(M) = R$.
\item Given a second finite $R$-module $M'$ we have
$\text{Fit}_0(M \oplus M') = \text{Fit}_0(M)\text{Fit}_0(M')$ and
more generally
$$
\text{Fit}_l(M \oplus M') =
\sum\nolimits_{k + k' = l} \text{Fit}_k(M)\text{Fit}_{k'}(M')
$$
\item If $R \to R'$ is a ring map, then $\text{Fit}_k(M \otimes_R R')$
is the ideal of $R'$ generated by the image of $\text{Fit}_k(M)$.
\item If $M$ is of finite presentation, then $\text{Fit}_k(M)$
is a finitely generated ideal.
\item If $M \to M'$ is a surjection, then
$\text{Fit}_k(M) \subset \text{Fit}_k(M')$.
\item We have $\text{Fit}_0(M) \subset \text{Ann}_R(M)$.
\item We have $V(\text{Fit}_0(M)) = \text{Supp}(M)$.
\item If $I$ is an ideal of $R$, then $\text{Fit}_0(R/I) = I$.
\item Add more here.
\end{enumerate}
\end{lemma}

\begin{proof}
Part (1) follows from the fact that $I_0(A) = R$ in
Lemma \ref{lemma-ideals-generated-by-minors}.

\medskip\noindent
Part (2) follows form the corresponding statement in
Lemma \ref{lemma-ideals-generated-by-minors}.

\medskip\noindent
Part (3) follows from the fact that $\otimes_R R'$ is right exact,
so the base change of a presentation of $M$ is a presentation of
$M \otimes_R R'$.

\medskip\noindent
Proof of (4). Let $R^{\oplus m} \xrightarrow{A} R^{\oplus n} \to M \to 0$
be a presentation. Then $\text{Fit}_k(M)$ is the ideal generated by the
$(n - k) \times (n - k)$ minors of the matrix $A$.

\medskip\noindent
Part (5) is immediate from the definition.

\medskip\noindent
Proof of (6). Choose a presentation of $M$ with matrix $A$
as in Lemma \ref{lemma-fitting-ideal}.
Let $J' \subset J$ be a subset of cardinality $n$.
It suffices to show that
$f = \det(a_{ij})_{i = 1, \ldots, n, j \in J'}$
annihilates $M$.
This is clear because the cokernel of
$$
R^{\oplus n} \xrightarrow{A' = (a_{ij})_{i = 1, \ldots, n, j \in J'}}
R^{\oplus n} \to M \to 0
$$
is killed by $f$ as there is a matrix $B$ with $A' B = f1_{n \times n}$.

\medskip\noindent
Proof of (7). Choose a presentation of $M$ with matrix $A$
as in Lemma \ref{lemma-fitting-ideal}.
By Nakayama's lemma (Algebra, Lemma \ref{algebra-lemma-NAK})
we have
$$
M_\mathfrak p \not = 0
\Leftrightarrow
M \otimes_R \kappa(\mathfrak p) \not = 0
\Leftrightarrow
\text{rank}(\text{image }A\text{ in }\kappa(\mathfrak p)) < n
$$
Clearly $\text{Fit}_0(M)$ exactly cuts out the set of primes
with this property.
\end{proof}

\begin{example}
\label{example-fitting-free}
Let $R$ be a ring.
The Fitting ideals of the finite free module $M = R^{\oplus n}$
are $\text{Fit}_k(M) = 0$ for $k < n$ and $\text{Fit}_k(M) = R$
for $k \geq n$.
\end{example}

\begin{example}
\label{example-laurent}
Let $R$ be a ring and let $I, J\subset R$ be ideals.
Then $\text{Fit}_0(R/I) = I$, $\text{Fit}_0(R/I \oplus R/J) = IJ$,
$\text{Fit}_1(R/I\oplus R/J) = I+J$.
\end{example}

\begin{lemma}
\label{lemma-fitting-ideal-generate-locally}
Let $R$ be a ring. Let $M$ be a finite $R$-module. Let $k \geq 0$.
Let $\mathfrak p \subset R$ be a prime ideal. The following
are equivalent
\begin{enumerate}
\item $\text{Fit}_k(M) \not \subset \mathfrak p$,
\item $\dim_{\kappa(\mathfrak p)} M \otimes_R \kappa(\mathfrak p) \leq k$,
\item $M_\mathfrak p$ can be generated by $k$ elements over $R_\mathfrak p$, and
\item $M_f$ can be generated by $k$ elements over $R_f$
for some $f \in R$, $f \not \in \mathfrak p$.
\end{enumerate}
\end{lemma}

\begin{proof}
By Nakayama's lemma (Algebra, Lemma \ref{algebra-lemma-NAK}) we see that
$M_f$ can be generated by $k$ elements over $R_f$ for some
$f \in R$, $f \not \in \mathfrak p$ if $M \otimes_R \kappa(\mathfrak p)$
can be generated by $k$ elements. Hence (2), (3), and (4)
are equivalent. Using
Lemma \ref{lemma-fitting-ideal-basics} part (3)
this reduces the problem to the
case where $R$ is a field and $\mathfrak p = (0)$. In this case
the result follows from Example \ref{example-fitting-free}.
\end{proof}

\begin{lemma}
\label{lemma-fitting-ideal-finite-locally-free}
Let $R$ be a ring. Let $M$ be a finite $R$-module. Let $r \geq 0$.
The following are equivalent
\begin{enumerate}
\item $M$ is finite locally free of rank $r$
(Algebra, Definition \ref{algebra-definition-locally-free}),
\item $\text{Fit}_{r - 1}(M) = 0$ and $\text{Fit}_r(M) = R$, and
\item $\text{Fit}_k(M) = 0$ for $k < r$ and $\text{Fit}_k(M) = R$
for $k \geq r$.
\end{enumerate}
\end{lemma}

\begin{proof}
It is immediate that (2) is equivalent to (3) because the Fitting ideals
form an increasing sequence of ideals.
Since the formation of $\text{Fit}_k(M)$ commutes with base change
(Lemma \ref{lemma-fitting-ideal-basics}) we see that (1) implies (2) by
Example \ref{example-fitting-free}
and glueing results (Algebra, Section \ref{algebra-section-more-glueing}).
Conversely, assume (2). By
Lemma \ref{lemma-fitting-ideal-generate-locally} we may assume that $M$
is generated by $r$ elements. Thus a presentation
$\bigoplus_{j \in J} R \to R^{\oplus r} \to M \to 0$.
But now the assumption that $\text{Fit}_{r - 1}(M) = 0$ implies
that all entries of the matrix of the map
$\bigoplus_{j \in J} R \to R^{\oplus r}$ are zero.
Thus $M$ is free.
\end{proof}

\begin{lemma}
\label{lemma-principal-fitting-ideal}
Let $R$ be a local ring. Let $M$ be a finite $R$-module. Let $k \geq 0$.
Assume that $\text{Fit}_k(M) = (f)$ for some $f \in R$.
Let $M'$ be the quotient of $M$ by $\{x \in M \mid fx = 0\}$. Then
$M'$ can be generated by $k$ elements.
\end{lemma}

\begin{proof}
Choose generators $x_1, \ldots, x_n \in M$ corresponding to the
surjection $R^{\oplus n} \to M$. Since $R$ is local if a set
of elements $E \subset (f)$ generates $(f)$, then some $e \in E$ generates
$(f)$, see Algebra, Lemma \ref{algebra-lemma-NAK}. Hence we may pick
$z_1, \ldots, z_{n - k}$ in the kernel of $R^{\oplus n} \to M$ such
that some $(n - k) \times (n - k)$ minor of the $n \times (n - k)$
matrix $A = (z_{ij})$ generates $(f)$. After renumbering the $x_i$ we may
assume the first minor $\det(z_{ij})_{1 \leq i, j \leq n - k}$
generates $(f)$, i.e., $\det(z_{ij})_{1 \leq i, j \leq n - k} = uf$
for some unit $u \in R$. Every other minor is a multiple of $f$.
By Algebra, Lemma \ref{algebra-lemma-matrix-right-inverse} there exists a
$n - k \times n - k$ matrix $B$ such that
$$
AB = f
\left(
\begin{matrix}
u 1_{n - k \times n - k} \\
C
\end{matrix}
\right)
$$
for some matrix $C$ with coefficients in $R$. This implies that for every
$i \leq n - k$ the element $y_i = ux_i + \sum_j c_{ji}x_j$ is annihilated
by $f$. Since $M/\sum Ry_i$ is generated by the images of
$x_{n - k + 1}, \ldots, x_n$ we win.
\end{proof}

\begin{lemma}
\label{lemma-fitting-ideals-and-pd1}
Let $R$ be a ring. Let $M$ be a finitely presented $R$-module. Let $k \geq 0$.
Assume that $\text{Fit}_k(M) = (f)$ for some nonzerodivisor $f \in R$
and $\text{Fit}_{k - 1}(M) = 0$. Then
\begin{enumerate}
\item $M$ has projective dimension $\leq 1$,
\item $M' = \Ker(f : M \to M)$ is the $f$-power torsion submodule of $M$,
\item $M'$ has projective dimension $\leq 1$,
\item $M/M'$ is finite locally free of rank $k$, and
\item $M \cong M/M' \oplus M'$.
\end{enumerate}
\end{lemma}

\begin{proof}
Choose a presentation
$$
R^{\oplus m} \xrightarrow{A} R^{\oplus n} \to M \to 0
$$
for some matrix $A$ with coefficients in $R$.

\medskip\noindent
We first prove the lemma when $R$ is local.
Set $M' = \{x \in M \mid fx = 0\}$ as in the statement. By
Lemma \ref{lemma-principal-fitting-ideal} we can choose
$x_1, \ldots, x_k \in M$ which generate $M/M'$. Then
$x_1, \ldots, x_k$ generate $M_f = (M/M')_f$. Hence,
if there is a relation $\sum a_ix_i = 0$ in $M$,
then we see that $a_1, \ldots, a_k$ map to zero in $R_f$
since otherwise $\text{Fit}_{k - 1}(M) R_f = \text{Fit}_{k - 1}(M_f)$
would be nonzero.
Since $f$ is a nonzerodivisor, we conclude $a_1 = \ldots = a_k = 0$.
Thus $M \cong R^{\oplus k} \oplus M'$. After a change of basis
in our presentation above, we may assume the first $n - k$
basis vectors of $R^{\oplus n}$ map into the summand $M'$ of $M$
and the last $k$-basis vectors of $R^{\oplus n}$ map to
basis elements of the summand $R^{\oplus k}$ of $M$.
Having done so, the last $k$ rows of the matrix $A$ vanish.
In this way we see that, replacing $M$ by $M'$, $k$ by $0$,
$n$ by $n - k$, and $A$ by the submatrix where we
delete the last $k$ rows, we reduce to the case discussed
in the next paragraph.

\medskip\noindent
Assume $R$ is local, $k = 0$, and $M$ annihilated by $f$.
Now the $0$th Fitting ideal of $M$ is $(f)$ and is generated by the
$n \times n$ minors of the matrix $A$ of size $n \times m$.
(This in particular implies $m \geq n$.)
Since $R$ is local, some $n \times n$ minor of $A$ is $uf$
for a unit $u \in R$.
After renumbering we may assume this minor is the first one.
Moreover, we know all other $n \times n$ minors of $A$ are
divisible by $f$. Write $A = (A_1 A_2)$ in block form where
$A_1$ is an $n \times n$ matrix and $A_2$ is an
$n \times (m - n)$ matrix. By
Algebra, Lemma \ref{algebra-lemma-matrix-right-inverse}
applied to the transpose of $A$ (!) we find there exists an
$n \times n$ matrix $B$ such that
$$
BA = B(A_1 A_2) = f
\left(
\begin{matrix}
u 1_{n \times n} & C
\end{matrix}
\right)
$$
for some $n \times (m - n)$ matrix $C$ with coefficients in $R$.
Then we first conclude $BA_1 = fu 1_{n \times n}$.
Thus
$$
BA_2 = fC = u^{-1}fuC = u^{-1}BA_1C
$$
Since the determinant of $B$ is a nonzerodivisor we conclude
that $A_2 = u^{-1}A_1C$. Therefore the image of $A$ is equal
to the image of $A_1$ which is isomorphic to $R^{\oplus n}$
because the determinant of $A_1$ is a nonzerodivisor.
Hence $M$ has projective dimension $\leq 1$.

\medskip\noindent
We return to the case of a general ring $R$. By the local case
we see that $M/M'$ is a finite locally free module of rank $k$, see
Algebra, Lemma \ref{algebra-lemma-finite-projective}. Hence the extension
$0 \to M' \to M \to M/M' \to 0$ splits. It follows that $M'$
is a finitely presented module. Choose a short exact sequence
$0 \to K \to R^{\oplus a} \to M' \to 0$. Then $K$ is a finite
$R$-module, see Algebra, Lemma \ref{algebra-lemma-extension}.
By the local case we see that $K_\mathfrak p \cong R_\mathfrak p^{\oplus a}$
for all primes. Hence by
Algebra, Lemma \ref{algebra-lemma-finite-projective}
again we see that $K$ is finite locally free of rank $a$.
It follows that $M'$ has projective dimension $\leq 1$
and the lemma is proved.
\end{proof}










\section{Lifting}
\label{section-lifting}

\noindent
In this section we collection some lemmas concerning lifting
statements of the following kind: If $A$ is a ring and $I \subset A$
is an ideal, and $\overline{\xi}$ is some kind of structure over
$A/I$, then we can lift $\overline{\xi}$ to a similar kind of structure
$\xi$ over $A$ or over some \'etale extension of $A$. Here are some types
of structure for which we have already proved some results:
\begin{enumerate}
\item idempotents, see
Algebra, Lemmas \ref{algebra-lemma-lift-idempotents} and
\ref{algebra-lemma-lift-idempotents-noncommutative},
\item projective modules, see
Algebra, Lemmas \ref{algebra-lemma-lift-projective-module} and
\ref{algebra-lemma-lift-finite-projective-module},
\item finite stably free modules, see Lemma \ref{lemma-lift-stably-free},
\item basis elements, see
Algebra, Lemmas \ref{algebra-lemma-local-artinian-basis-when-flat} and
\ref{algebra-lemma-lift-basis},
\item ring maps, i.e., proving certain algebras are formally smooth, see
Algebra, Lemma \ref{algebra-lemma-polynomial-ring-formally-smooth},
Proposition \ref{algebra-proposition-smooth-formally-smooth}, and
Lemma \ref{algebra-lemma-smooth-strong-lift},
\item syntomic ring maps, see
Algebra, Lemma \ref{algebra-lemma-lift-syntomic},
\item smooth ring maps, see
Algebra, Lemma \ref{algebra-lemma-lift-smooth},
\item \'etale ring maps, see
Algebra, Lemma \ref{algebra-lemma-lift-etale},
\item factoring polynomials, see
Algebra, Lemma \ref{algebra-lemma-factor-mod-lift-etale}, and
\item Algebra, Section \ref{algebra-section-henselian} discusses henselian
local rings.
\end{enumerate}
The interested reader will find more results of this nature in
Smoothing Ring Maps, Section \ref{smoothing-section-presentations}
in particular
Smoothing Ring Maps, Proposition \ref{smoothing-proposition-lift-smooth}.

\medskip\noindent
Let $A$ be a ring and let $I \subset A$ be an ideal. Let $\overline{\xi}$
be some kind of structure over $A/I$. In the following lemmas we look for
\'etale ring maps $A \to A'$ which induce isomorphisms $A/I \to A'/IA'$
and objects $\xi'$ over $A'$ lifting $\overline{\xi}$. A general remark is
that given \'etale ring maps $A \to A' \to A''$ such that
$A/I \cong A'/IA'$ and $A'/IA' \cong A''/IA''$ the composition
$A \to A''$ is also \'etale (Algebra, Lemma \ref{algebra-lemma-etale})
and also satisfies $A/I \cong A''/IA''$.
We will frequently use this in the following lemmas without further mention.
Here is a trivial example of the type of result we are looking for.

\begin{lemma}
\label{lemma-lift-invertible-element}
Let $A$ be a ring, let $I \subset A$ be an ideal, let $\overline{u} \in A/I$
be an invertible element. There exists an \'etale ring map $A \to A'$ which
induces an isomorphism $A/I \to A'/IA'$ and an invertible element $u' \in A'$
lifting $\overline{u}$.
\end{lemma}

\begin{proof}
Choose any lift $f \in A$ of $\overline{u}$ and set $A' = A_f$ and $u$
the image of $f$ in $A'$.
\end{proof}

\begin{lemma}
\label{lemma-lift-idempotent}
Let $A$ be a ring, let $I \subset A$ be an ideal, let $\overline{e} \in A/I$
be an idempotent. There exists an \'etale ring map $A \to A'$ which
induces an isomorphism $A/I \to A'/IA'$ and an idempotent $e' \in A'$
lifting $\overline{e}$.
\end{lemma}

\begin{proof}
Choose any lift $x \in A$ of $\overline{e}$. Set
$$
A' = A[t]/(t^2 - t)\left[\frac{1}{t - 1 + x}\right].
$$
The ring map $A \to A'$ is \'etale because $(2t - 1)\text{d}t = 0$
and $(2t - 1)(2t - 1) = 1$ which is invertible. We have
$A'/IA' = A/I[t]/(t^2 - t)[\frac{1}{t - 1 + \overline{e}}] \cong A/I$
the last map sending $t$ to $\overline{e}$ which works as
$\overline{e}$ is a root of $t^2 - t$. This also shows that setting
$e'$ equal to the class of $t$ in $A'$ works.
\end{proof}

\begin{lemma}
\label{lemma-lift-open-covering}
Let $A$ be a ring, let $I \subset A$ be an ideal. Let
$\Spec(A/I) = \coprod_{j \in J} \overline{U}_j$ be a finite disjoint open
covering. Then there exists an \'etale ring map $A \to A'$ which
induces an isomorphism $A/I \to A'/IA'$ and a finite disjoint open covering
$\Spec(A') = \coprod_{j \in J} U'_j$ lifting the given covering.
\end{lemma}

\begin{proof}
This follows from Lemma \ref{lemma-lift-idempotent} and
the fact that open and closed subsets of Spectra correspond
to idempotents, see Algebra, Lemma \ref{algebra-lemma-disjoint-decomposition}.
\end{proof}

\begin{lemma}
\label{lemma-localize-upstairs}
Let $A \to B$ be a ring map and $J \subset B$ an ideal. If
$A \to B$ is \'etale at every prime of $V(J)$, then there exists
a $g \in B$ mapping to an invertible element
of $B/J$ such that $A' = B_g$ is \'etale over $A$.
\end{lemma}

\begin{proof}
The set of points of $\Spec(B)$ where $A \to B$ is not \'etale is a
closed subset of $\Spec(B)$, see
Algebra, Definition \ref{algebra-definition-etale}.
Write this as $V(J')$ for some ideal $J' \subset B$. Then
$V(J') \cap V(J) = \emptyset$ hence $J + J' = B$ by
Algebra, Lemma \ref{algebra-lemma-Zariski-topology}.
Write $1 = f + g$ with $f \in J$ and $g \in J'$.
Then $g$ works.
\end{proof}

\noindent
Next we have three lemmas saying we can lift factorizations of
polynomials.

\begin{lemma}
\label{lemma-lift-factorization-monic}
Let $A$ be a ring, let $I \subset A$ be an ideal. Let $f \in A[x]$ be a
monic polynomial. Let $\overline{f} = \overline{g} \overline{h}$ be a
factorization of $f$ in $A/I[x]$ such that $\overline{g}$ and $\overline{h}$
are monic and generate the unit ideal in $A/I[x]$. Then there exists an
\'etale ring map $A \to A'$ which induces an isomorphism $A/I \to A'/IA'$
and a factorization $f = g' h'$ in $A'[x]$ with $g'$, $h'$ monic
lifting the given factorization over $A/I$.
\end{lemma}

\begin{proof}
We will deduce this from results on the universal factorization proved
earlier; however, we encourage the reader to find their own proof not
using this trick. Say $\deg(\overline{g}) = n$ and $\deg(\overline{h}) = m$ so
that $\deg(f) = n + m$. Write $f = x^{n + m} + \sum \alpha_i x^{n + m - i}$
for some $\alpha_1, \ldots, \alpha_{n + m} \in A$. Consider the ring map
$$
R = \mathbf{Z}[a_1, \ldots, a_{n + m}]
\longrightarrow
S = \mathbf{Z}[b_1, \ldots, b_n, c_1, \ldots, c_m]
$$
of Algebra, Example \ref{algebra-example-factor-polynomials-etale}.
Let $R \to A$ be the ring map which sends $a_i$ to $\alpha_i$.
Set
$$
B = A \otimes_R S
$$
By construction the image $f_B$ of $f$ in $B[x]$ factors, say
$f_B = g_B h_B$ with $g_B = x^n + \sum (1 \otimes b_i) x^{n - i}$
and similarly for $h_B$.
Write $\overline{g} = x^n + \sum \overline{\beta}_i x^{n - i}$ and
$\overline{h} = x^m + \sum \overline{\gamma}_i x^{m - i}$.
The $A$-algebra map
$$
B \longrightarrow A/I, \quad
1 \otimes b_i \mapsto \overline{\beta}_i, \quad
1 \otimes c_i \mapsto \overline{\gamma}_i
$$
maps $g_B$ and $h_B$ to $\overline{g}$ and $\overline{h}$ in $A/I[x]$.
The displayed map is surjective; denote $J \subset B$ its kernel.
From the discussion in 
Algebra, Example \ref{algebra-example-factor-polynomials-etale}
it is clear that $A \to B$ is etale at all points of $V(J) \subset \Spec(B)$.
Choose $g \in B$ as in Lemma \ref{lemma-localize-upstairs} and
consider the $A$-algebra $B_g$. Since $g$ maps to a unit
in $B/J = A/I$ we obtain also a map $B_g/I B_g \to A/I$ of $A/I$-algebras.
Since $A/I \to B_g/I B_g$ is \'etale, also $B_g/IB_g \to A/I$ is \'etale
(Algebra, Lemma \ref{algebra-lemma-map-between-etale}). Hence there exists an
idempotent $e \in B_g/I B_g$ such that $A/I = (B_g/I B_g)_e$
(Algebra, Lemma \ref{algebra-lemma-surjective-flat-finitely-presented}).
Choose a lift $h \in B_g$ of $e$. Then $A \to A' = (B_g)_h$ with
factorization given by the image of the factorization $f_B = g_B h_B$
in $A'$ is a solution to the problem posed by the lemma.
\end{proof}

\begin{lemma}
\label{lemma-lift-factorization-easy}
Let $A$ be a ring, let $I \subset A$ be an ideal. Let $f \in A[x]$ be a
monic polynomial. Let $\overline{f} = \overline{g} \overline{h}$ be a
factorization of $f$ in $A/I[x]$ and assume
\begin{enumerate}
\item the leading coefficient of $\overline{g}$ is an invertible element
of $A/I$, and
\item $\overline{g}$, $\overline{h}$ generate the unit ideal in $A/I[x]$.
\end{enumerate}
Then there exists an \'etale ring map $A \to A'$ which induces an
isomorphism $A/I \to A'/IA'$ and a factorization $f = g' h'$ in $A'[x]$
lifting the given factorization over $A/I$.
\end{lemma}

\begin{proof}
Applying Lemma \ref{lemma-lift-invertible-element} we may assume that
the leading coefficient of $\overline{g}$ is the reduction of an
invertible element $u \in A$. Then we may replace $\overline{g}$ by
$\overline{u}^{-1}\overline{g}$ and $\overline{h}$ by
$\overline{u}\overline{h}$. Thus we may assume that $\overline{g}$
is monic. Since $f$ is monic we conclude that $\overline{h}$ is monic
too. In this case the result follows from
Lemma \ref{lemma-lift-factorization-monic}.
\end{proof}

\begin{example}
\label{example-cannot-lift-factorization}
This example shows that one cannot drop the condition on the leading
coefficient in Lemma \ref{lemma-lift-factorization-easy}. Namely,
let $A = \mathbf{Z}$, $I = 4\mathbf{Z}$, $f = 1$ and
$\bar{g} = \bar{h} = 2x^2 + 2x + 1$ in $A/I[x]$. If $A \to A'$ is \'etale
and $A'/I A' = A/I$, then the $2$-adic completion of $A'$ is isomorphic
to $\mathbf{Z}_2$. Now any polynomial $g'$ in $\mathbf{Z}_2[x]$ congruent
to $2x^2 + 2x + 1$ mod $4$ will not have a polynomial multiplicative
inverse in $\mathbf{Z}_2[x]$, whence the conclusion of the lemma does
not hold.
\end{example}

\begin{lemma}
\label{lemma-separate-image-closed-from-closed}
Let $R \to S$ be a ring map. Let $I \subset R$ be an ideal of $R$
and let $J \subset S$ be an ideal of $S$. If the closure of the image
of $V(J)$ in $\Spec(R)$ is disjoint from $V(I)$, then there exists
an element $f \in R$ which maps to $1$ in $R/I$ and to an element
of $J$ in $S$.
\end{lemma}

\begin{proof}
Let $I' \subset R$ be an ideal such that $V(I')$ is the closure of
the image of $V(J)$. Then $V(I) \cap V(I') = \emptyset$ by assumption
and hence $I + I' = R$ by
Algebra, Lemma \ref{algebra-lemma-Zariski-topology}.
Write $1 = g + f$ with $g \in I$ and $f \in I'$.
We have $V(f') \supset V(J)$ where $f'$ is the image of $f$ in $S$.
Hence $(f')^n \in J$ for some $n$, see
Algebra, Lemma \ref{algebra-lemma-Zariski-topology}.
Replacing $f$ by $f^n$ we win.
\end{proof}

\begin{lemma}
\label{lemma-helper-integral}
Let $I$ be an ideal of a ring $A$. Let $A \to B$ be an integral ring map.
Let $b \in B$ map to an idempotent in $B/IB$. Then there exists a
monic $f \in A[x]$ with $f(b) = 0$ and $f \bmod I = x^d(x - 1)^d$
for some $d \geq 1$.
\end{lemma}

\begin{proof}
Observe that $z = b^2 - b$ is an element of $IB$. By
Algebra, Lemma \ref{algebra-lemma-integral-integral-over-ideal}
there exist a monic polynomial
$g(x) = x^d + \sum a_j x^j$ of degree $d$ with $a_j \in I$ such that
$g(z) = 0$ in $B$. Hence $f(x) = g(x^2 - x) \in A[x]$ is a monic
polynomial such that $f(x) \equiv x^d(x - 1)^d \bmod I$
and such that $f(b) = 0$ in $B$.
\end{proof}

\begin{lemma}
\label{lemma-lift-idempotent-upstairs}
Let $A$ be a ring, let $I \subset A$ be an ideal.
Let $A \to B$ be an integral ring map.
Let $\overline{e} \in B/IB$ be an idempotent.
Then there exists an \'etale ring map $A \to A'$
which induces an isomorphism $A/I \to A'/IA'$ and an idempotent
$e' \in B \otimes_A A'$ lifting $\overline{e}$.
\end{lemma}

\begin{proof}
Choose an element $y \in B$ lifting $\overline{e}$.
Choose $f \in A[x]$ as in Lemma \ref{lemma-helper-integral} for $y$.
By Lemma \ref{lemma-lift-factorization-easy}
we can find an \'etale ring map $A \to A'$ which induces
an isomorphism $A/I \to A'/IA'$ and such that $f = gh$
in $A[x]$ with $g(x) = x^d \bmod IA'$ and $h(x) = (x - 1)^d \bmod IA'$.
After replacing $A$ by $A'$ we may assume that the factorization
is defined over $A$. In that case we see that
$b_1 = g(y) \in B$ is a lift of $\overline{e}^d = \overline{e}$ and
$b_2 = h(y) \in B$ is a lift of
$(\overline{e} - 1)^d = (-1)^d (1 - \overline{e})^d = (-1)^d(1 - \overline{e})$
and moreover $b_1b_2 = 0$. Thus $(b_1, b_2)B/IB = B/IB$ and
$V(b_1, b_2) \subset \Spec(B)$ is disjoint from $V(IB)$. Since
$\Spec(B) \to \Spec(A)$ is closed (see
Algebra, Lemmas \ref{algebra-lemma-integral-going-up} and
\ref{algebra-lemma-going-up-closed})
we can find an $a \in A$ which maps to an invertible element
of $A/I$ whose image in $B$ lies in $(b_1, b_2)$, see
Lemma \ref{lemma-separate-image-closed-from-closed}.
After replacing $A$ by the localization $A_a$ we get that
$(b_1, b_2) = B$. Then $\Spec(B) = D(b_1) \amalg D(b_2)$;
disjoint union because $b_1b_2 = 0$ and covers
$\Spec(B)$ because $(b_1, b_2) = B$. Let $e \in B$ be the idempotent
corresponding to the open and closed subset $D(b_1)$, see
Algebra, Lemma \ref{algebra-lemma-disjoint-decomposition}.
Since $b_1$ is a lift of $\overline{e}$ and $b_2$ is a
lift of $\pm (1 - \overline{e})$ we conclude that $e$ is
a lift of $\overline{e}$ by the uniqueness statement in
Algebra, Lemma \ref{algebra-lemma-disjoint-decomposition}.
\end{proof}

\begin{lemma}
\label{lemma-lift-projective-module}
Let $A$ be a ring, let $I \subset A$ be an ideal.
Let $\overline{P}$ be a finite projective $A/I$-module.
Then there exists an \'etale ring map $A \to A'$ which induces
an isomorphism $A/I \to A'/IA'$ and a finite projective
$A'$-module $P'$ lifting $\overline{P}$.
\end{lemma}

\begin{proof}
We can choose an integer $n$ and a direct sum decomposition
$(A/I)^{\oplus n} = \overline{P} \oplus \overline{K}$
for some $R/I$-module $\overline{K}$. Choose a lift
$\varphi : A^{\oplus n} \to A^{\oplus n}$ of the projector $\overline{p}$
associated to the direct summand $\overline{P}$.
Let $f \in A[x]$ be the characteristic polynomial of $\varphi$.
Set $B = A[x]/(f)$. By Cayley-Hamilton
(Algebra, Lemma \ref{algebra-lemma-charpoly}) there is a map
$B \to \text{End}_A(A^{\oplus n})$ mapping $x$ to $\varphi$.
For every prime $\mathfrak p \supset I$ the image of $f$ in
$\kappa(\mathfrak p)$ is $(x - 1)^rx^{n - r}$ where $r$ is the
dimension of $\overline{P} \otimes_{A/I} \kappa(\mathfrak p)$.
Hence $(x - 1)^nx^n$ maps to zero in $B \otimes_A \kappa(\mathfrak p)$
for all $\mathfrak p \supset I$. Thus $x(1 - x)$ is contained
in every prime ideal of $B/IB$. Hence $x^N(1 - x)^N$ is
contained in $IB$ for some $N \geq 1$.
It follows that $x^N + (1 - x)^N$ is a unit in $B/IB$ and that
$$
\overline{e} = \text{image of }\frac{x^N}{x^N + (1 - x)^N}\text{ in }B/IB
$$
is an idempotent as both assertions hold in $\mathbf{Z}[x]/(x^N(x - 1)^N)$.
The image of $\overline{e}$ in $\text{End}_{A/I}((A/I)^{\oplus n})$ is
$$
\frac{\overline{p}^N}{\overline{p}^N + (1 - \overline{p})^N} = \overline{p}
$$
as $\overline{p}$ is an idempotent. After replacing $A$ by an \'etale
extension $A'$ as in the lemma, we may assume there exists an idempotent
$e \in B$ which maps to $\overline{e}$ in $B/IB$, see
Lemma \ref{lemma-lift-idempotent-upstairs}.
Then the image of $e$ under the map
$$
B = A[x]/(f) \longrightarrow \text{End}_A(A^{\oplus n}).
$$
is an idempotent element $p$ which lifts $\overline{p}$.
Setting $P = \Im(p)$ we win.
\end{proof}

\begin{lemma}
\label{lemma-cotangent-complex-symmetric-algebra}
Let $A$ be a ring. Let $0 \to K \to A^{\oplus m} \to M \to 0$
be a sequence of $A$-modules. Consider the $A$-algebra
$C = \text{Sym}^*_A(M)$ with its presentation
$\alpha : A[y_1, \ldots, y_m] \to C$
coming from the surjection $A^{\oplus m} \to M$. Then
$$
\NL(\alpha) =
(K \otimes_A C \to \bigoplus\nolimits_{j = 1, \ldots, m} C \text{d}y_j)
$$
(see Algebra, Section \ref{algebra-section-netherlander})
in particular $\Omega_{C/A} = M \otimes_A C$.
\end{lemma}

\begin{proof}
Let $J = \Ker(\alpha)$. The lemma asserts that
$J/J^2 \cong K \otimes_A C$. Note that $\alpha$ is a homomorphism
of graded algebras. We will prove that in degree $d$ we have
$(J/J^2)_d = K \otimes_A C_{d - 1}$. Note that
$$
J_d = \Ker(\text{Sym}^d_A(A^{\oplus m}) \to \text{Sym}^d_A(M))
= \Im(K \otimes_A \text{Sym}^{d - 1}_A(A^{\oplus m})
\to \text{Sym}^d_A(A^{\oplus m})),
$$
see Algebra, Lemma \ref{algebra-lemma-presentation-sym-exterior}.
It follows that $(J^2)_d = \sum_{a + b = d} J_a \cdot J_b$ is the image of
$$
K \otimes_A K \otimes_A \text{Sym}^{d - 2}_A(A^{\otimes m})
\to \text{Sym}^d_A(A^{\oplus m}).
$$
The cokernel of the map $K \otimes_A \text{Sym}^{d - 2}_A(A^{\otimes m}) \to
\text{Sym}^{d - 1}_A(A^{\oplus m})$ is $\text{Sym}^{d - 1}_A(M)$ by
the lemma referenced above.
Hence it is clear that $(J/J^2)_d = J_d/(J^2)_d$ is equal to
\begin{align*}
\Coker(
K \otimes_A K \otimes_A \text{Sym}^{d - 2}_A(A^{\otimes m})
\to K \otimes_A \text{Sym}^{d - 1}_A(A^{\otimes m}))
& = K \otimes_A \text{Sym}^{d - 1}_A(M) \\
& = K \otimes_A C_{d -1}
\end{align*}
as desired.
\end{proof}

\begin{lemma}
\label{lemma-symmetric-algebra-smooth}
Let $A$ be a ring. Let $M$ be an $A$-module. Then $C = \text{Sym}_A^*(M)$
is smooth over $A$ if and only if $M$ is a finite projective $A$-module.
\end{lemma}

\begin{proof}
Let $\sigma : C \to A$ be the projection onto the degree $0$ part of $C$.
Then $J = \Ker(\sigma)$ is the part of degree $> 0$ and we see that
$J/J^2 = M$ as an $A$-module. Hence if $A \to C$ is smooth then $M$ is
a finite projective $A$-module by
Algebra, Lemma \ref{algebra-lemma-section-smooth}.

\medskip\noindent
Conversely, assume that $M$ is finite projective and choose a surjection
$A^{\oplus n} \to M$ with kernel $K$. Of course the sequence
$0 \to K \to A^{\oplus n} \to M \to 0$ is split as $M$ is projective.
In particular we see that $K$ is a finite $A$-module and hence
$C$ is of finite presentation over $A$ as $C$ is a quotient of
$A[x_1, \ldots, x_n]$ by the ideal generated by $K \subset \bigoplus Ax_i$.
The computation of Lemma \ref{lemma-cotangent-complex-symmetric-algebra}
shows that $\NL_{C/A}$ is homotopy equivalent to $(K \to M) \otimes_A C$.
Hence $\NL_{C/A}$ is quasi-isomorphic to $C \otimes_A M$ placed in degree
$0$ which means that $C$ is smooth over $A$ by
Algebra, Definition \ref{algebra-definition-smooth}.
\end{proof}

\begin{lemma}
\label{lemma-lift-section-smooth-morphism}
Let $A$ be a ring, let $I \subset A$ be an ideal. Consider a commutative
diagram
$$
\xymatrix{
B \ar[rd] \\
A \ar[u] \ar[r] & A/I
}
$$
where $B$ is a smooth $A$-algebra. Then there exists an \'etale ring
map $A \to A'$ which induces an isomorphism $A/I \to A'/IA'$ and an
$A$-algebra map $B \to A'$ lifting the ring map $B \to A/I$.
\end{lemma}

\begin{proof}
Let $J \subset B$ be the kernel of $B \to A/I$ so that $B/J = A/I$. By
Algebra, Lemma \ref{algebra-lemma-application-NL-smooth} the sequence
$$
0 \to I/I^2 \to J/J^2 \to \Omega_{B/A} \otimes_B B/J \to 0
$$
is split exact. Thus $\overline{P} = J/(J^2 + IB) = \Omega_{B/A} \otimes_B B/J$
is a finite projective $A/I$-module. Choose an integer $n$ and a direct sum
decomposition $A/I^{\oplus n} = \overline{P} \oplus \overline{K}$.
By Lemma \ref{lemma-lift-projective-module} we can find an
\'etale ring map $A \to A'$ which induces an isomorphism
$A/I \to A'/IA'$ and a finite projective $A$-module $K$ which
lifts $\overline{K}$. We may and do replace $A$ by $A'$.
Set $B' = B \otimes_A \text{Sym}_A^*(K)$. Since $A \to \text{Sym}_A^*(K)$
is smooth by Lemma \ref{lemma-symmetric-algebra-smooth} we see that
$B \to B'$ is smooth which in turn implies that $A \to B'$ is smooth (see
Algebra, Lemmas \ref{algebra-lemma-base-change-smooth} and
\ref{algebra-lemma-locally-smooth}).
Moreover the section $\text{Sym}^*_A(K) \to A$ determines a section
$B' \to B$ and we let $B' \to A/I$ be the composition $B' \to B \to A/I$.
Let $J' \subset B'$ be the kernel of $B' \to A/I$. We have
$JB' \subset J'$ and $B \otimes_A K \subset J'$. These maps combine
to give an isomorphism
$$
(A/I)^{\oplus n} \cong J/(J^2 + IB) \oplus \overline{K}
\longrightarrow
J'/((J')^2 + IB')
$$
Thus, after replacing $B$ by $B'$ we may assume that
$J/(J^2 + IB) = \Omega_{B/A} \otimes_B B/J$ is a free
$A/I$-module of rank $n$.

\medskip\noindent
In this case, choose $f_1, \ldots, f_n \in J$ which map to a
basis of $J/(J^2 + IB)$. Consider the finitely presented $A$-algebra
$C = B/(f_1, \ldots, f_n)$. Note that we have an exact sequence
$$
0 \to H_1(L_{C/A}) \to (f_1, \ldots, f_n)/(f_1, \ldots, f_n)^2
\to \Omega_{B/A} \otimes_B C \to \Omega_{C/A} \to 0
$$
see Algebra, Lemma \ref{algebra-lemma-exact-sequence-NL} (note that
$H_1(L_{B/A}) = 0$ and that $\Omega_{B/A}$ is finite projective,
in particular flat so the Tor group vanishes). For any prime
$\mathfrak q \supset J$ of $B$ the module $\Omega_{B/A, \mathfrak q}$
is free of rank $n$ because $\Omega_{B/A}$ is finite projective and
because $\Omega_{B/A} \otimes_B B/J$ is free of rank $n$
(see Algebra, Lemma \ref{algebra-lemma-finite-projective}). By our choice
of $f_1, \ldots, f_n$ the map
$$
\left((f_1, \ldots, f_n)/(f_1, \ldots, f_n)^2\right)_{\mathfrak q}
\to
\Omega_{B/A, \mathfrak q}
$$
is surjective modulo $J$. Hence we see that this map of modules over
the local ring $C_{\mathfrak q}$ has to be an isomorphism
(by Algebra, Lemma \ref{algebra-lemma-NAK} the map is surjective
and for example by Algebra, Lemma \ref{algebra-lemma-fun}
because $((f_1, \ldots, f_n)/(f_1, \ldots, f_n)^2)_{\mathfrak q}$
is generated by $n$ elements the map is injective). Thus
$H_1(L_{C/A})_{\mathfrak q} = 0$ and $\Omega_{C/A, \mathfrak q} = 0$. By
Algebra, Lemma \ref{algebra-lemma-smooth-at-point}
we see that $A \to C$ is smooth at the prime $\overline{\mathfrak q}$
of $C$ corresponding to $\mathfrak q$. Since
$\Omega_{C/A, \mathfrak q} = 0$ it is actually \'etale at
$\overline{\mathfrak q}$. Thus $A \to C$ is \'etale at all primes
of $C$ containing $JC$. By Lemma \ref{lemma-localize-upstairs}
we can find an $f \in C$ mapping to an invertible element of $C/JC$ such that
$A \to C_f$ is \'etale. By our choice of $f$ it is still true that
$C_f/JC_f = A/I$. The map $C_f/IC_f \to A/I$ is surjective and
\'etale by Algebra, Lemma \ref{algebra-lemma-map-between-etale}.
Hence $A/I$ is isomorphic to the localization of $C_f/IC_f$ at
some element $g \in C$, see
Algebra, Lemma \ref{algebra-lemma-surjective-flat-finitely-presented}.
Set $A' = C_{fg}$ to conclude the proof.
\end{proof}








\section{Zariski pairs}
\label{section-zariski-pairs}

\noindent
In this section and the next a {\it pair} is a pair $(A, I)$ where $A$
is a ring and $I \subset A$ is an ideal. A {\it morphism of pairs}
$(A, I) \to (B, J)$ is a ring map $\varphi : A \to B$ with
$\varphi(I) \subset J$.

\begin{definition}
\label{definition-zariski-pair}
A {\it Zariski pair} is a pair $(A, I)$ such that
$I$ is contained in the Jacobson radical of $A$.
\end{definition}

\begin{lemma}
\label{lemma-idempotents-determined-modulo-radical}
Let $(A, I)$ be a Zariski pair. Then the map from
idempotents of $A$ to idempotents of $A/I$ is injective.
\end{lemma}

\begin{proof}
An idempotent of a local ring is either $0$ or $1$.
Thus an idempotent is determined by the set of maximal ideals
where it vanishes, by
Algebra, Lemma \ref{algebra-lemma-characterize-zero-local}.
\end{proof}

\begin{lemma}
\label{lemma-check-isomorphism-zariski}
Let $(A, I)$ be a Zariski pair. Let $A \to B$ be a flat,
integral, finitely presented ring map such that $A/I \to B/IB$
is an isomorphism. Then $A \to B$ is an isomorphism.
\end{lemma}

\begin{proof}
The ring map $A \to B$ is finite by
Algebra, Lemma \ref{algebra-lemma-characterize-finite-in-terms-of-integral}.
Hence $B$ is finitely presented as an $A$-module by
Algebra, Lemma \ref{algebra-lemma-finite-finitely-presented-extension}.
Hence $B$ is a finite locally free $A$-module by
Algebra, Lemma \ref{algebra-lemma-finite-projective}.
Since the module $B$ has rank $1$
along $V(I)$ (see rank function described in
Algebra, Lemma \ref{algebra-lemma-finite-projective}),
and as $(A, I)$ is a Zariski pair, we conclude that the
rank is $1$ everywhere.
It follows that $A \to B$ is an isomorphism:
it is a pleasant exercise to show that
a ring map $R \to S$ such that $S$ is a locally free $R$-module
of rank $1$ is an isomorphism (hint: look at local rings).
\end{proof}

\begin{lemma}
\label{lemma-helper-finite}
Let $(A, I)$ be a Zariski pair. Let $A \to B$ be a finite ring map.
Assume
\begin{enumerate}
\item $B/IB = B_1 \times B_2$ is a product of $A/I$-algebras
\item $A/I \to B_1/IB_1$ is surjective,
\item $b \in B$ maps to $(1, 0)$ in the product.
\end{enumerate}
Then there exists a monic $f \in A[x]$ with $f(b) = 0$ and
$f \bmod I = (x - 1)x^d$ for some $d \geq 1$.
\end{lemma}

\begin{proof}
By Lemma \ref{lemma-lift-idempotent-upstairs} we can find an
\'etale ring map $A \to A'$ inducing an isomorphism $A/I \to A'/IA'$
such that $B' = B \otimes_A A'$ contains
an idempotent $e'$ lifting the image of $b$ in $B'/IB'$.
Consider the corresponding $A'$-algebra decomposition
$$
B' = B'_1 \times B'_2
$$
which is compatible with the one given in the lemma upon
reduction modulo $I$.
The map $A' \to B'_1$ is surjective modulo $IA'$. By Nakayama's lemma
(Algebra, Lemma \ref{algebra-lemma-NAK})
we can find $i \in IA'$ such that after replacing $A'$
by $A'_{1 + i}$ the map $A' \to B'_1$ is surjective.
Observe that the image $b'_1 \in B'_1$ of $b$ 
satisfies $b'_1 - 1 \in IB'_1$.
Thus we may pick $a' \in IA'$ mapping to $b'_1 - 1$.
On the other hand, the image $b'_2 \in B'_2$ of $b$ is in $IB'_2$. By
Algebra, Lemma \ref{algebra-lemma-integral-integral-over-ideal}
there exist a monic polynomial
$g(x) = x^d + \sum a'_j x^j$ of degree $d$ with $a'_j \in IA'$ such that
$g(b'_2) = 0$ in $B'_2$. Thus the image $b' = (b'_1, b'_2) \in B'$
of $b$ is a root of the polynomial $(x - 1 - a')g(x)$. We conclude that
$$
(b' - 1)(b')^d \in \sum\nolimits_{j = 0, \ldots, d} IA' \cdot (b')^j
$$
We claim that this implies
$$
(b - 1)b^d \in \sum\nolimits_{j = 0, \ldots, d} I \cdot b^j
$$
in $B$. For this it is enough to see that the ring map
$A \to A'$ is faithfully flat, because the condition is that
the image of $(b - 1)b^d$ is zero in $B/\sum_{j = 0, \ldots, d} Ib^j$
(use Algebra, Lemma \ref{algebra-lemma-faithfully-flat-universally-injective}).
The map $A \to A'$ flat because it is \'etale
(Algebra, Lemma \ref{algebra-lemma-etale}).
On the other hand, the induced map on spectra is open
(see Algebra, Proposition \ref{algebra-proposition-fppf-open} and use
previous lemma referenced) and the image
contains $V(I)$. Since $I$ is contained in the Jacobson radical of $A$
we conclude.
\end{proof}

\begin{lemma}
\label{lemma-noetherian-zariski-jacobson-complement}
Let $(A, I)$ be a Zariski pair with $A$ Noetherian. Let $f \in I$.
Then $A_f$ is a Jacobson ring.
\end{lemma}

\begin{proof}
We will use the criterion of
Algebra, Lemma \ref{algebra-lemma-noetherian-dim-1-Jacobson}.
Let $\mathfrak p \subset A$ be a prime ideal such that
$\mathfrak p_f = \mathfrak p A_f$ is prime and not maximal. We have to
show that $A_f/\mathfrak p_f = (A/\mathfrak p)_f$
has infinitely many prime ideals. After replacing $A$ by $A/\mathfrak p$
we may assume $A$ is a domain, $\dim A_f > 0$, and our goal is to
show that $\Spec(A_f)$ is infinite. Since $\dim A_f > 0$ we can
find a nonzero prime ideal $\mathfrak q \subset A$ not containing $f$.
Choose a maximal ideal $\mathfrak m \subset A$ containing $\mathfrak q$.
Since $(A, I)$ is a Zariski pair, we see $I \subset \mathfrak m$.
Hence $\mathfrak m \not = \mathfrak q$ and $\dim(A_\mathfrak m) > 1$.
Hence $\Spec((A_\mathfrak m)_f) \subset \Spec(A_f)$ is infinite by
Algebra, Lemma \ref{algebra-lemma-Noetherian-local-domain-dim-2-infinite-opens}
and we win.
\end{proof}











\section{Henselian pairs}
\label{section-henselian-pairs}

\noindent
Some of the results of Section \ref{section-lifting} may be viewed as results
about henselian pairs. In this section a {\it pair} is a pair $(A, I)$
where $A$ is a ring and $I \subset A$ is an ideal. A {\it morphism of pairs}
$(A, I) \to (B, J)$ is a ring map $\varphi : A \to B$ with
$\varphi(I) \subset J$. As in
Section \ref{section-lifting} given an object $\xi$ over $A$ we denote
$\overline{\xi}$ the ``base change'' of $\xi$ to an object over $A/I$
(provided this makes sense).

\begin{definition}
\label{definition-henselian-pair}
A {\it henselian pair} is a pair $(A, I)$ satisfying
\begin{enumerate}
\item $I$ is contained in the Jacobson radical of $A$, and
\item for any monic polynomial $f \in A[T]$ and factorization
$\overline{f} = g_0h_0$ with $g_0, h_0 \in A/I[T]$ monic
generating the unit ideal in $A/I[T]$, there
exists a factorization $f = gh$ in $A[T]$ with $g, h$ monic
and $g_0 = \overline{g}$ and $h_0 = \overline{h}$.
\end{enumerate}
\end{definition}

\noindent
Observe that if $A$ is a local ring and $I = \mathfrak m$ is the maximal
ideal, then $(A, I)$ is a henselian pair if and only if $A$ is a henselian
local ring, see
Algebra, Lemma \ref{algebra-lemma-characterize-henselian}.
In Lemma \ref{lemma-characterize-henselian-pair} we give a number of
equivalent characterizations of
henselian pairs (and we will add more as time goes on).

\begin{lemma}
\label{lemma-locally-nilpotent-henselian}
Let $(A, I)$ be a pair with $I$ locally nilpotent. Then the functor
$B \mapsto B/IB$ induces an equivalence between the category of
\'etale algebras over $A$ and the category of \'etale algebras over $A/I$.
Moreover, the pair is henselian.
\end{lemma}

\begin{proof}
Essential surjectivity holds by Algebra, Lemma \ref{algebra-lemma-lift-etale}.
If $B$, $B'$ are \'etale over $A$ and $B/IB \to B'/IB'$ is a morphism
of $A/I$-algebras, then we can lift this by
Algebra, Lemma \ref{algebra-lemma-smooth-strong-lift}.
Finally, suppose that $f, g : B \to B'$ are two $A$-algebra
maps with $f \bmod I = g \bmod I$. Choose an idempotent $e \in B \otimes_A B$
generating the kernel of the multiplication map $B \otimes_A B \to B$,
see Algebra, Lemmas \ref{algebra-lemma-diagonal-unramified}
and \ref{algebra-lemma-unramified} (to see that \'etale is unramified).
Then $(f \otimes g)(e) \in IB'$. Since $IB'$ is locally nilpotent
(Algebra, Lemma \ref{algebra-lemma-locally-nilpotent}) this implies
$(f \otimes g)(e) = 0$ by Algebra, Lemma \ref{algebra-lemma-lift-idempotents}.
Thus $f = g$.

\medskip\noindent
It is clear that $I$ is contained in the Jacobson radical of $A$.
Let $f \in A[T]$ be a monic polynomial and let
$\overline{f} = g_0h_0$ be a factorization
of $\overline{f} = f \bmod I$ with $g_0, h_0 \in A/I[T]$ monic
generating the unit ideal in $A/I[T]$. By
Lemma \ref{lemma-lift-factorization-monic}
there exists an \'etale ring map $A \to A'$ which
induces an isomorphism $A/I \to A'/IA'$ such that
the factorization lifts to a factorization into monic polynomials
over $A'$. By the above we have $A = A'$ and the factorization
is over $A$.
\end{proof}

\begin{lemma}
\label{lemma-limit-henselian}
Let $A = \lim A_n$ where $(A_n)$ is an inverse system of rings
whose transition maps are surjective and have locally nilpotent kernels.
Then $(A, I_n)$ is a henselian pair, where $I_n = \Ker(A \to A_n)$.
\end{lemma}

\begin{proof}
Fix $n$. Let $a \in A$ be an element which maps to $1$ in $A_n$.
By Algebra, Lemma \ref{algebra-lemma-locally-nilpotent-unit}
we see that $a$ maps to a unit in $A_m$ for all $m \geq n$.
Hence $a$ is a unit in $A$. Thus by
Algebra, Lemma \ref{algebra-lemma-contained-in-radical}
the ideal $I_n$ is contained in the Jacobson radical of $A$.
Let $f \in A[T]$ be a monic polynomial and let
$\overline{f} = g_nh_n$ be a factorization
of $\overline{f} = f \bmod I_n$ with $g_n, h_n \in A_n[T]$ monic
generating the unit ideal in $A_n[T]$. By
Lemma \ref{lemma-locally-nilpotent-henselian}
we can successively lift this factorization to
$f \bmod I_m = g_m h_m$ with $g_m, h_m$ monic
in $A_m[T]$ for all $m \geq n$. At each step we have
to verify that our lifts $g_m, h_m$ generate the unit ideal in
$A_n[T]$; this follows from the corresponding fact for
$g_n, h_n$ and the fact that $\Spec(A_n[T]) = \Spec(A_m[T])$
because the kernel of $A_m \to A_n$ is locally nilpotent.
As $A = \lim A_m$ this finishes the proof.
\end{proof}

\begin{lemma}
\label{lemma-complete-henselian}
Let $(A, I)$ be a pair. If $A$ is $I$-adically complete, then
the pair is henselian.
\end{lemma}

\begin{proof}
By Algebra, Lemma \ref{algebra-lemma-radical-completion}
the ideal $I$ is contained in the Jacobson radical of $A$.
Let $f \in A[T]$ be a monic polynomial and let
$\overline{f} = g_0h_0$ be a factorization
of $\overline{f} = f \bmod I$ with $g_0, h_0 \in A/I[T]$ monic
generating the unit ideal in $A/I[T]$. By
Lemma \ref{lemma-locally-nilpotent-henselian}
we can successively lift this factorization to
$f \bmod I^n = g_n h_n$ with $g_n, h_n$ monic
in $A/I^n[T]$ for all $n \geq 1$.
As $A = \lim A/I^n$ this finishes the proof.
\end{proof}

\begin{lemma}
\label{lemma-helper-finite-type}
Let $(A, I)$ be a pair. Let $A \to B$ be a finite type ring map
such that $B/IB = C_1 \times C_2$ with $A/I \to C_1$ finite.
Let $B'$ be the integral closure of $A$ in $B$.
Then we can write $B'/IB' = C_1 \times C'_2$ such that
the map $B'/IB' \to B/IB$ preserves product decompositions
and there exists a $g \in B'$ mapping to $(1, 0)$ in
$C_1 \times C'_2$ with $B'_g \to B_g$ an isomorphism.
\end{lemma}

\begin{proof}
Observe that $A \to B$ is quasi-finite at every prime of the
closed subset $T = \Spec(C_1) \subset \Spec(B)$ (this follows
by looking at fibre rings, see
Algebra, Definition \ref{algebra-definition-quasi-finite}).
Consider the diagram of topological spaces
$$
\xymatrix{
\Spec(B) \ar[rr]_\phi \ar[rd]_\psi & & \Spec(B') \ar[ld]^{\psi'} \\
& \Spec(A)
}
$$
By Algebra, Theorem \ref{algebra-theorem-main-theorem}
for every $\mathfrak p \in T$ there is a $h_\mathfrak p \in B'$,
$h_\mathfrak p \not \in \mathfrak p$ such that $B'_h \to B_h$ is
an isomorphism. The union $U = \bigcup D(h_\mathfrak p)$ gives an open
$U \subset \Spec(B')$ such that $\phi^{-1}(U) \to U$ is a homeomorphism
and $T \subset \phi^{-1}(U)$. Since $T$ is open in $\psi^{-1}(V(I))$
we conclude that $\phi(T)$ is open in $U \cap (\psi')^{-1}(V(I))$.
Thus $\phi(T)$ is open in $(\psi')^{-1}(V(I))$.
On the other hand, since $C_1$ is finite over $A/I$ it is
finite over $B'$. Hence $\phi(T)$ is a closed subset of $\Spec(B')$
by Algebra, Lemmas \ref{algebra-lemma-going-up-closed} and
\ref{algebra-lemma-integral-going-up}. We conclude that
$\Spec(B'/IB') \supset \phi(T)$ is open and closed. By
Algebra, Lemma \ref{algebra-lemma-disjoint-implies-product}
we get a corresponding product decomposition $B'/IB' = C'_1 \times C'_2$.
The map $B'/IB' \to B/IB$ maps $C'_1$ into $C_1$ and $C'_2$ into $C_2$
as one sees by looking at what happens on spectra (hint: the inverse
image of $\phi(T)$ is exactly $T$; some details omitted).
Pick a $g \in B'$ mapping to $(1, 0)$ in $C'_1 \times C'_2$
such that $D(g) \subset U$; this is possible because $\Spec(C'_1)$
and $\Spec(C'_2)$ are disjoint and closed in $\Spec(B')$ and
$\Spec(C'_1)$ is contained in $U$. Then $B'_g \to B_g$ defines a homeomorphism
on spectra and an isomorphism on local rings (by our choice of $U$ above).
Hence it is an isomorphism, as follows for example from
Algebra, Lemma \ref{algebra-lemma-characterize-zero-local}.
Finally, it follows that $C'_1 = C_1$ and the proof is complete.
\end{proof}

\begin{lemma}
\label{lemma-characterize-henselian-pair}
\begin{reference}
\cite[Chapter XI]{Henselian} and \cite[Proposition 1]{Gabber-henselian}
\end{reference}
Let $(A, I)$ be a pair. The following are equivalent
\begin{enumerate}
\item $(A, I)$ is a henselian pair,
\item given an \'etale ring map $A \to A'$ and an $A$-algebra map
$\sigma : A' \to A/I$, there exists an $A$-algebra map $A' \to A$
lifting $\sigma$,
\item for any finite $A$-algebra $B$ the map $B \to B/IB$ induces
a bijection on idempotents,
\item for any integral $A$-algebra $B$ the map $B \to B/IB$ induces
a bijection on idempotents, and
\item (Gabber) $I$ is contained in the Jacobson radical of $A$ and
every monic polynomial $f(T) \in A[T]$ of the form
$$
f(T) = T^n(T - 1) + a_n T^n + \ldots + a_1 T + a_0
$$
with $a_n, \ldots, a_0 \in I$ and $n \ge 1$ has a root $\alpha \in 1 + I$.
\end{enumerate}
Moreover, in part (5) the root is unique.
\end{lemma}

\begin{proof}
Assume (2) holds. Then $I$ is contained in the Jacobson radical of $A$, since
otherwise there would be a nonunit $f \in A$ congruent to $1$ modulo $I$
and the map $A \to A_f$ would contradict (2). Hence $IB \subset B$
is contained in the Jacobson radical of $B$ for $B$ integral over $A$
because $\Spec(B) \to \Spec(A)$ is closed by
Algebra, Lemmas \ref{algebra-lemma-going-up-closed} and
\ref{algebra-lemma-integral-going-up}.
Thus the map from idempotents of $B$ to idempotents of $B/IB$
is injective by Lemma \ref{lemma-idempotents-determined-modulo-radical}.
On the other hand, since (2) holds, every idempotent
of $B/IB$ lifts to an idempotent of $B$
by Lemma \ref{lemma-lift-idempotent-upstairs}.
In this way we see that (2) implies (4).

\medskip\noindent
The implication (4) $\Rightarrow$ (3) is trivial.

\medskip\noindent
Assume (3). Let $\mathfrak m$ be a maximal ideal and consider the
finite map $A \to B = A/(I \cap \mathfrak m)$. The condition that
$B \to B/IB$ induces a bijection on idempotents implies that
$I \subset \mathfrak m$ (if not, then $B = A/I \times A/\mathfrak m$
and $B/IB = A/I$). Thus we see that $I$ is contained in the Jacobson
radical of $A$. Let $f \in A[T]$ be monic and suppose given a
factorization $\overline{f} = g_0h_0$ with $g_0, h_0 \in A/I[T]$ monic
generating the unit ideal in $A/I[T]$.
Set $B = A[T]/(f)$. Let $\overline{e}$ be the idempotent
of $B/IB$ corresponding to the decomposition
$$
B/IB = A/I[T]/(g_0) \times A/I[T]/(h_0)
$$
of $A$-algebras. Let $e \in B$ be an idempotent lifting $\overline{e}$
which exists as we assumed (3). This gives a product decomposition
$$
B = eB \times (1 - e)B
$$
Note that $B$ is free of rank $\deg(f)$ as an $A$-module.
Hence $eB$ and $(1 - e)B$ are finite locally free $A$-modules.
However, since $eB$ and $(1 - e)B$ have constant rank
$\deg(g_0)$ and $\deg(h_0)$ over $A/I$ we find that the same
is true over $\Spec(A)$. We conclude that
\begin{align*}
f & = \text{CharPol}_A(T : B \to B) \\
& =
\text{CharPol}_A(T : eB \to eB)
\text{CharPol}_A(T : (1 - e)B \to (1 - e)B)
\end{align*}
is a factorization into monic polynomials reducing to the given
factorization modulo $I$. Here $\text{CharPol}_A$ denotes the characteristic
polynomial of an endomorphism of a finite locally free module over $A$.
If the module is free the $\text{CharPol}_A$ is defined as the
characteristic polynomial of the corresponding matrix and in
general one uses Algebra, Lemma \ref{algebra-lemma-standard-covering} to
glue. Details omitted. Thus (3) implies (1).

\medskip\noindent
Assume (1). Let $f$ be as in (5). The factorization of $f \bmod I$
as $T^n$ times $T - 1$ lifts to a factorization $f = gh$ with $g$
and $h$ monic by Definition \ref{definition-henselian-pair}.
Then $h$ has to have degree $1$
and we see that $f$ has a root reducing to $1$ modulo $I$.
Finally, $I$ is contained in the Jacobson radical by
the definition of a henselian pair.
Thus (1) implies (5).

\medskip\noindent
Before we give the proof of the last step, let us show that
the root $\alpha$ in (5), if it exists, is unique. Namely,
due to the explicit shape of $f(T)$, we have
$f'(\alpha) \in 1 + I$ where $f'$ is the derivative of $f$ with
respect to $T$. An elementary argument shows that
$$
f(T) = f(\alpha + T - \alpha) =
f(\alpha) + f'(\alpha) \cdot (T - \alpha) \bmod
(T - \alpha)^2 A[T]
$$
This shows that any other root $\alpha' \in 1 + I$ of $f(T)$
satisfies $0 = f(\alpha') - f(\alpha) = (\alpha' - \alpha)(1 + i)$
for some $i \in I$, so that, since $1 + i$ is a unit in $A$,
we have $\alpha = \alpha'$.

\medskip\noindent
Assume (5). We will show that (2) holds, in other words, that
for every \'etale map $A \to A'$, every section $\sigma : A' \to A/I$
modulo $I$ lifts to a section $A' \to A$.
Since $A \to A'$ is \'etale, the section $\sigma$
determines a decomposition
\begin{equation}
\label{equation-GCHP}
A'/IA' \cong A/I \times C
\end{equation}
of $A/I$-algebras. Namely, the surjective ring map
$A'/IA' \to A/I$ is \'etale by
Algebra, Lemma \ref{algebra-lemma-map-between-etale}
and then we get the desired idempotent by
Algebra, Lemma \ref{algebra-lemma-surjective-flat-finitely-presented}.
We will show that this decomposition lifts to a decomposition 
\begin{equation}
\label{equation-GCHP-want}
A' \cong A'_1 \times A'_2
\end{equation}
of $A$-algebras with $A'_1$ integral over $A$. Then $A \to A'_1$
is integral and \'etale and $A/I \to A'_1/IA'_1$ is an isomorphism,
thus $A \to A'_1$ is an isomorphism by
Lemma \ref{lemma-check-isomorphism-zariski}
(here we also use that an \'etale ring map is flat
and of finite presentation, see Algebra, Lemma \ref{algebra-lemma-etale}).

\medskip\noindent
Let $B'$ be the integral closure of $A$ in $A'$. By
Lemma \ref{lemma-helper-finite-type}
we may decompose 
\begin{equation}
\label{equation-dec-mod-I}
B'/IB' \cong A/I \times C'
\end{equation}
as $A/I$-algebras compatibly with (\ref{equation-GCHP})
and we may find $b \in B'$ that lifts $(1, 0)$ such that
$B'_b \to A'_b$ is an isomorphism. If the decomposition
(\ref{equation-dec-mod-I}) lifts to a decomposition
\begin{equation}
\label{equation-want-2}
B' \cong B'_1 \times B'_2
\end{equation}
of $A$-algebras, then the induced decomposition
$A' = A'_1 \times A'_2$ will give the
desired (\ref{equation-GCHP-want}): indeed, since $b$ is a unit in $B'_1$
(details omitted),
we will have $B'_1 \cong A'_1$, so that $A'_1$ will be integral over $A$. 

\medskip\noindent
Choose a finite $A$-subalgebra $B'' \subset B'$ containing $b$
(observe that any finitely generated $A$-subalgebra of $B'$
is finite over $A$). After enlarging $B''$ we may assume
$b$ maps to an idempotent in $B''/IB''$ producing
\begin{equation}
\label{equation-again-dec-mod-I}
B''/IB'' \cong C''_1 \times C''_2
\end{equation}
Since $B'_b \cong A'_b$ we see that $B'_b$ is of finite type over $A$.
Say $B'_b$ is generated by $b_1/b^n, \ldots, b_t/b^n$ over $A$
and enlarge $B''$ so that $b_1, \ldots, b_t \in B''$. Then
$B''_b \to B'_b$ is surjective as well as injective, hence an isomorphism.
In particular, we see that $C''_1 = A/I$! Therefore $A/I \to C''_1$
is an isomorphism, in particular surjective.
By Lemma \ref{lemma-helper-finite} we can find
an $f(T) \in A[T]$ of the form
$$
f(T) = T^n(T - 1) + a_n T^n + \ldots + a_1 T + a_0
$$
with $a_n, \ldots, a_0 \in I$ and $n \ge 1$ such that $f(b) = 0$.
In particular, we find that $B'$ is a $A[T]/(f)$-algebra.
By (5) we deduce there is a root $a \in 1 + I$ of $f$.
This produces a product decomposition $A[T]/(f) = A[T]/(T - a) \times D$
compatible with the splitting (\ref{equation-dec-mod-I}) of $B'/IB'$.
The induced splitting of $B'$ is then a desired (\ref{equation-want-2}).
\end{proof}

\begin{lemma}
\label{lemma-change-ideal-henselian-pair}
Let $A$ be a ring. Let $I, J \subset A$ be ideals with $V(I) = V(J)$.
Then $(A, I)$ is henselian if and only if $(A, J)$ is henselian.
\end{lemma}

\begin{proof}
For any integral ring map $A \to B$ we see that $V(IB) = V(JB)$.
Hence idempotents of $B/IB$ and $B/JB$ are in bijective correspondence
(Algebra, Lemma \ref{algebra-lemma-disjoint-decomposition}).
It follows that $B \to B/IB$ induces a bijection on sets of
idempotents if and only if $B \to B/JB$ induces a bijection on sets
of idempotents. Thus we conclude by
Lemma \ref{lemma-characterize-henselian-pair}.
\end{proof}

\begin{lemma}
\label{lemma-integral-over-henselian-pair}
Let $(A, I)$ be a henselian pair and let $A \to B$ be an integral ring
map. Then $(B, IB)$ is a henselian pair.
\end{lemma}

\begin{proof}
Immediate from the fourth characterization of henselian pairs in
Lemma \ref{lemma-characterize-henselian-pair} and the fact that the
composition of integral ring maps is integral.
\end{proof}

\begin{lemma}
\label{lemma-henselian-henselian-pair}
Let $I \subset J \subset A$ be ideals of a ring $A$.
The following are equivalent
\begin{enumerate}
\item $(A, I)$ and $(A/I, J/I)$ are henselian pairs, and
\item $(A, J)$ is an henselian pair.
\end{enumerate}
\end{lemma}

\begin{proof}
Assume (1). Let $B$ be an integral $A$-algebra. Consider the ring maps
$$
B \to B/IB \to B/JB
$$
By Lemma \ref{lemma-characterize-henselian-pair} we find that
both arrows induce bijections on idempotents. Hence so does the
composition. Whence $(A, J)$ is a henselian pair
by Lemma \ref{lemma-characterize-henselian-pair}.

\medskip\noindent
Conversely, assume (2) holds. Then $(A/I, J/I)$ is a henselian pair
by Lemma \ref{lemma-integral-over-henselian-pair}. Let $B$ be an
integral $A$-algebra. Consider the ring maps
$$
B \to B/IB \to B/JB
$$
By Lemma \ref{lemma-characterize-henselian-pair} we find that the composition
and the second arrow induce bijections on idempotents.
Hence so does the first arrow. It follows that $(A, I)$ is a henselian
pair (by the lemma again).
\end{proof}

\begin{lemma}
\label{lemma-sum-henselian}
Let $A$ be a ring and let $(A, I)$ and $(A, I')$ be henselian pairs.
Then $(A, I + I')$ is an henselian pair.
\end{lemma}

\begin{proof}
By Lemma \ref{lemma-integral-over-henselian-pair} the pair
$(A/I, (I' + I)/I)$ is henselian. Thus we get the conclusion from
Lemma \ref{lemma-henselian-henselian-pair}.
\end{proof}

\begin{lemma}
\label{lemma-product-henselian-pairs}
Let $J$ be a set and let $\{ (A_j, I_j)\}_{j \in J}$ be a collection
of pairs. Then $(\prod_{j \in J} A_j, \prod_{j\in J} I_j)$ is Henselian
if and only if so is each $(A_j, I_j)$.
\end{lemma}
 
\begin{proof}
For every $j \in J$, the projection $\prod_{j \in J} A_j \rightarrow A_j$
is an integral ring map, so Lemma \ref{lemma-integral-over-henselian-pair}
proves that each $(A_j, I_j)$ is Henselian if
$(\prod_{j \in J} A_j, \prod_{j\in J} I_j)$ is Henselian.

\medskip\noindent
Conversely, suppose that each $(A_j, I_j)$ is a Henselian pair.
Then every $1 + x$ with $x \in \prod_{j \in J} I_j$  is a unit
in $\prod_{j \in J} A_j$ because it is so componentwise by
Algebra, Lemma \ref{algebra-lemma-contained-in-radical} and
Definition \ref{definition-henselian-pair}.
Thus, by Algebra, Lemma \ref{algebra-lemma-contained-in-radical}
again, $\prod_{j \in J} I_j$ is contained in the Jacobson radical
of $\prod_{j \in J} A_j$. Continuing to work componentwise, it
likewise follows that for every monic $f \in (\prod_{j \in J} A_j)[T]$
and every factorization $\overline{f} = g_0h_0$ with monic
$g_0, h_0 \in (\prod_{j \in J} A_j / \prod_{j \in J} I_j)[T] =
(\prod_{j \in J} A_j/I_j)[T]$ that generate the unit ideal in
$(\prod_{j \in J} A_j / \prod_{j \in J} I_j)[T]$, there exists a
factorization $f = gh$ in $(\prod_{j \in J} A_j)[T]$ with $g$, $h$ monic
and reducing to $g_0$, $h_0$. In conclusion, according to
Definition \ref{definition-henselian-pair}
$(\prod_{j \in J} A_j, \prod_{j\in J} I_j)$ is a Henselian pair.
\end{proof}

\begin{lemma}
\label{lemma-limits-henselian}
The property of being Henselian is preserved under limits of pairs.
More precisely, let $J$ be a preordered set and let $(A_j, I_j)$
be an inverse system of henselian pairs over $J$.
Then $A = \lim A_j$ equipped with the ideal $I = \lim I_j$
is a henselian pair $(A, I)$.
\end{lemma}

\begin{proof}
By Categories, Lemma \ref{categories-lemma-limits-products-equalizers},
we only need to consider products and equalizers.
For products, the claim follows from
Lemma \ref{lemma-product-henselian-pairs}.
Thus, consider an equalizer diagram
$$
\xymatrix{
(A, I) \ar[r] & (A', I')
\ar@<1ex>[r]^{\varphi} \ar@<-1ex>[r]_{\psi}
&
(A'', I'')  
}
$$
in which the pairs $(A', I')$ and $(A'', I'')$ are henselian.
To check that the pair $(A, I)$ is also henselian, we will use the
Gabber's criterion in Lemma \ref{lemma-characterize-henselian-pair}.
Every element of $1 + I$ is a unit in $A$ because,
due to the uniqueness of the inverses of units,
this may be checked in $(A', I')$.
Thus $I$ is contained in the Jacobson radical of $A$, see
Algebra, Lemma \ref{algebra-lemma-contained-in-radical}.
Thus, let 
$$
f(T) = T^{N - 1}(T - 1) + a_{N - 1} T^{N - 1} + \dotsb + a_1 T + a_0
$$
be a polynomial in $A[T]$ with $a_{N - 1}, \dotsc, a_0 \in I$ and $N \ge 1$.
The image of $f(T)$ in $A'[T]$ has a unique root $\alpha' \in 1 + I'$
and likewise for the further image in $A''[T]$.
Thus, due to the uniqueness, $\varphi(\alpha') = \psi(\alpha')$,
to the effect that $\alpha'$ defines a root of $f(T)$ in $1 + I$, as desired.
\end{proof}

\begin{lemma}
\label{lemma-filtered-colimits-henselian}
The property of being Henselian is preserved under filtered colimits of pairs.
More precisely, let $J$ be a directed set and let $(A_j, I_j)$
be a system of henselian pairs over $J$.
Then $A = \colim A_j$ equipped with the ideal $I = \colim I_j$
is a henselian pair $(A, I)$.
\end{lemma}

\begin{proof}
If $u \in 1 + I$ then for some $j \in J$ we see that $u$ is the
image of some $u_j \in 1 + I_j$. Then $u_j$ is invertible in $A_j$ by
Algebra, Lemma \ref{algebra-lemma-contained-in-radical}
and the assumption that $I_j$ is contained in the Jacobson radical of $A_j$.
Hence $u$ is invertible in $A$. Thus $I$
is contained in the Jacobson radical of $A$ (by the lemma).

\medskip\noindent
Let $f \in A[T]$ be a monic polynomial and let $\overline{f} = g_0 h_0$
be a factorization with $g_0, h_0 \in A/I[T]$ monic generating the unit
ideal in $A/I[T]$. Write $1 = g_0 g'_0 + h_0 h'_0$ for some
$g'_0, h'_0 \in A/I[T]$. Since $A = \colim A_j$ and $A/I = \colim A_j/I_j$
are filtered colimits we can find a $j \in J$ and $f_j \in A_j$ and
a factorization $\overline{f}_j = g_{j, 0} h_{j, 0}$
with $g_{j, 0}, h_{j, 0} \in A_j/I_j[T]$ monic
and $1 = g_{j, 0} g'_{j, 0} + h_{j, 0} h'_{j, 0}$
for some $g'_{j, 0}, h'_{j, 0} \in A_j/I_j[T]$ with
$f_j, g_{j, 0}, h_{j, 0}, g'_{j, 0}, h'_{j, 0}$
mapping to $f, g_0, h_0, g'_0, h'_0$.
Since $(A_j, I_j)$ is a henselian pair, we can lift
$\overline{f}_j = g_{j, 0} h_{j, 0}$ to a factorization
over $A_j$ and taking the image in $A$ we obtain a
corresponding factorization in $A$. Hence $(A, I)$ is henselian.
\end{proof}

\begin{example}[Moret-Bailly]
\label{example-moret-bailly}
Lemma \ref{lemma-filtered-colimits-henselian} is wrong if the colimit
isn't filtered. For example, if we take the coproduct of the
henselian pairs $(\mathbf{Z}_p, (p))$ and $(\mathbf{Z}_p, (p))$, then we 
obtain $(A, pA)$ with $A = \mathbf{Z}_p \otimes_\mathbf{Z} \mathbf{Z}_p$.
This isn't a henselian pair: $A/pA = \mathbf{F}_p$ hence if $(A, pA)$
where henselian, then $A$ would have to be local. However, $\Spec(A)$
is disconnected; for example for odd primes $p$ we have the nontrivial
idempotent
$$
(1/2 \otimes 1)
\left(1 \otimes 1 - (1 + p)^{-1}u \otimes u\right)
$$
where $u \in \mathbf{Z}_p$ is a square root of $1 + p$. Some details omitted.
\end{example}

\begin{lemma}
\label{lemma-largest-ideal-henselian}
Let $A$ be a ring. There exists a largest ideal $I \subset A$ such that
$(A, I)$ is a henselian pair.
\end{lemma}

\begin{proof}
Combine Lemmas \ref{lemma-henselian-henselian-pair}, \ref{lemma-sum-henselian},
and \ref{lemma-filtered-colimits-henselian}.
\end{proof}

\begin{lemma}
\label{lemma-irreducible-henselian-pair-connected}
Let $(A, I)$ be a henselian pair. Let $\mathfrak p \subset A$
be a prime ideal. Then $V(\mathfrak p + I)$ is connected.
\end{lemma}

\begin{proof}
By Lemma \ref{lemma-integral-over-henselian-pair} we see that
$(A/\mathfrak p, I + \mathfrak p/\mathfrak p)$ is a henselian pair.
Thus it suffices to prove: If $(A, I)$ is a henselian pair and
$A$ is a domain, then $\Spec(A/I) = V(I)$ is connected. If not,
then $A/I$ has a nontrivial idempotent by
Algebra, Lemma \ref{algebra-lemma-characterize-spec-connected}.
By Lemma \ref{lemma-characterize-henselian-pair}
this would imply $A$ has a nontrivial idempotent. This is a contradiction.
\end{proof}









\section{Henselization of pairs}
\label{section-henselization}

\noindent
We continue the discussion started in
Section \ref{section-henselian-pairs}.

\begin{lemma}
\label{lemma-henselization}
The inclusion functor
$$
\text{category of henselian pairs}
\longrightarrow
\text{category of pairs}
$$
has a left adjoint $(A, I) \mapsto (A^h, I^h)$.
\end{lemma}

\begin{proof}
Let $(A, I)$ be a pair. Consider the category $\mathcal{C}$ consisting
of \'etale ring maps $A \to B$ such that $A/I \to B/IB$ is an isomorphism.
We will show that the category $\mathcal{C}$ is directed and that
$A^h = \colim_{B \in \mathcal{C}} B$ with ideal $I^h = IA^h$ gives
the desired adjoint.

\medskip\noindent
We first prove that $\mathcal{C}$ is directed
(Categories, Definition \ref{categories-definition-directed}).
It is nonempty because $\text{id} : A \to A$ is an object.
If $B$ and $B'$ are two objects of $\mathcal{C}$, then
$B'' = B \otimes_A B'$ is an object of $\mathcal{C}$
(use Algebra, Lemma \ref{algebra-lemma-etale})
and there are morphisms $B \to B''$ and $B' \to B''$.
Suppose that $f, g : B \to B'$ are two maps between
objects of $\mathcal{C}$. Then a coequalizer is
$$
(B' \otimes_{f, B, g} B') \otimes_{(B' \otimes_A B')} B'
$$
which is \'etale over $A$ by
Algebra, Lemmas \ref{algebra-lemma-etale} and
\ref{algebra-lemma-map-between-etale}.
Thus the category $\mathcal{C}$ is directed.

\medskip\noindent
Since $B/IB = A/I$ for all objects $B$ of $\mathcal{C}$ we
see that $A^h/I^h = A^h/IA^h = \colim B/IB = \colim A/I = A/I$.

\medskip\noindent
Next, we show that $A^h = \colim_{B \in \mathcal{C}} B$ with
$I^h = IA^h$ is a henselian pair. To do this we will verify
condition (2) of Lemma \ref{lemma-characterize-henselian-pair}.
Namely, suppose given an \'etale ring map $A^h \to A'$
and $A^h$-algebra map $\sigma : A' \to A^h/I^h$. Then there exists a
$B \in \mathcal{C}$ and an \'etale ring map $B \to B'$ such that
$A' = B' \otimes_B A^h$. See Algebra, Lemma \ref{algebra-lemma-etale}.
Since $A^h/I^h = A/I$, the map $\sigma$ induces an $A$-algebra
map $s : B' \to A/I$. Then $B'/IB' = A/I \times C$ as $A/I$-algebra,
where $C$ is the kernel of the map $B'/IB' \to A/I$ induced by $s$.
Let $g \in B'$ map to $(1, 0) \in A/I \times C$. Then $B \to B'_g$
is \'etale and $A/I \to B'_g/IB'_g$ is an isomorphism, i.e.,
$B'_g$ is an object of $\mathcal{C}$. Thus we obtain a canonical
map $B'_g \to A^h$ such that
$$
\vcenter{
\xymatrix{
B'_g \ar[r] & A^h \\
B \ar[u] \ar[ur]
}
}
\quad\text{and}\quad
\vcenter{
\xymatrix{
B' \ar[r] \ar[rrd]_s & B'_g \ar[r] & A^h \ar[d] \\
& & A/I
}
}
$$
commute. This induces a map $A' = B' \otimes_B A^h \to A^h$
compatible with $\sigma$ as desired.

\medskip\noindent
Let $(A, I) \to (A', I')$ be a morphism of pairs with $(A', I')$ henselian.
We will show there is a unique factorization $A \to A^h \to A'$ which will
finish the proof. Namely, for each $A \to B$ in $\mathcal{C}$
the ring map $A' \to B' = A' \otimes_A B$ is \'etale and induces
an isomorphism $A'/I' \to B'/I'B'$. Hence there is a section
$\sigma_B : B' \to A'$ by Lemma \ref{lemma-characterize-henselian-pair}.
Given a morphism $B_1 \to B_2$ in $\mathcal{C}$ we claim the diagram
$$
\xymatrix{
B'_1 \ar[rr] \ar[rd]_{\sigma_{B_1}} & &
B'_2 \ar[ld]^{\sigma_{B_2}} \\
& A'
}
$$
commutes. This follows once we prove that for every $B$ in $\mathcal{C}$
the section $\sigma_B$ is the unique $A'$-algebra map $B' \to A'$.
We have $B' \otimes_{A'} B' = B' \times R$ for some ring $R$, see
Algebra, Lemma \ref{algebra-lemma-diagonal-unramified}. In our case
$R/I'R = 0$ as $B'/I'B' = A'/I'$. Thus given two $A'$-algebra maps
$\sigma_B, \sigma_B' : B' \to A'$ then
$e = (\sigma_B \otimes \sigma_B')(0, 1) \in A'$
is an idempotent contained in $I'$. We conclude that $e = 0$
by Lemma \ref{lemma-idempotents-determined-modulo-radical}.
Hence $\sigma_B = \sigma_B'$ as desired.
Using the commutativity we obtain
$$
A^h = \colim_{B \in \mathcal{C}} B \to
\colim_{B \in \mathcal{C}} A' \otimes_A B \xrightarrow{\colim \sigma_B} A'
$$
as desired. The uniqueness of the maps $\sigma_B$ also guarantees that
this map is unique. Hence $(A, I) \mapsto (A^h, I^h)$ is the desired adjoint.
\end{proof}

\begin{lemma}
\label{lemma-henselization-flat}
Let $(A, I)$ be a pair. Let $(A^h, I^h)$ be as in 
Lemma \ref{lemma-henselization}. Then $A \to A^h$ is flat,
$I^h = IA^h$ and $A/I^n \to A^h/I^nA^h$ is an isomorphism
for all $n$.
\end{lemma}

\begin{proof}
In the proof of Lemma \ref{lemma-henselization} we have seen that
$A^h$ is a filtered colimit of \'etale $A$-algebras $B$ such that
$A/I \to B/IB$ is an isomorphism and we have seen that
$I^h = IA^h$. As an \'etale ring map is flat
(Algebra, Lemma \ref{algebra-lemma-etale}) we conclude that
$A \to A^h$ is flat by Algebra, Lemma \ref{algebra-lemma-colimit-flat}.
Since each $A \to B$ is flat we find that the maps
$A/I^n \to B/I^nB$ are isomorphisms as well (for example by
Algebra, Lemma \ref{algebra-lemma-lift-basis}).
Taking the colimit we find that $A/I^n = A^h/I^nA^h$
as desired.
\end{proof}

\begin{lemma}
\label{lemma-henselization-local-ring}
\begin{slogan}
Compatibility henselization of pairs and of local rings.
\end{slogan}
The functor of Lemma \ref{lemma-henselization} associates to a local ring
$(A, \mathfrak m)$ its henselization.
\end{lemma}

\begin{proof}
Let $(A^h, \mathfrak m^h)$ be the henselization of the pair $(A, \mathfrak m)$
constructed in Lemma \ref{lemma-henselization}. Then
$\mathfrak m^h = \mathfrak m A^h$ is a maximal ideal
by Lemma \ref{lemma-henselization-flat} and since it
is contained in the Jacobson radical, we conclude $A^h$ is
local with maximal ideal $\mathfrak m^h$. Having said this
there are two ways to finish the proof.

\medskip\noindent
First proof: observe that the construction in the proof of
Algebra, Lemma \ref{algebra-lemma-henselization} as a colimit
is the same as the colimit used to construct $A^h$ in
Lemma \ref{lemma-henselization}.
Second proof: Both the henselization $A \to S$ and $A \to A^h$ of
Lemma \ref{lemma-henselization} are local ring homomorphisms,
both $S$ and $A^h$ are filtered colimits of \'etale $A$-algebras,
both $S$ and $A^h$ are henselian local rings, and both $S$ and $A^h$
have residue fields equal to $\kappa(\mathfrak m)$ (by
Lemma \ref{lemma-henselization-flat} for the second case).
Hence they are canonically isomorphic by
Algebra, Lemma \ref{algebra-lemma-uniqueness-henselian}.
\end{proof}

\begin{lemma}
\label{lemma-henselization-Noetherian-pair}
\begin{slogan}
The henselization of a Noetherian pair is a Noetherian pair
with the same completion
\end{slogan}
Let $(A, I)$ be a pair with $A$ Noetherian.  Let $(A^h, I^h)$ be as in 
Lemma \ref{lemma-henselization}. Then the map of $I$-adic completions
$$
A^\wedge \to (A^h)^\wedge
$$
is an isomorphism. Moreover, $A^h$ is Noetherian, the maps
$A \to A^h \to A^\wedge$ are flat, and $A^h \to A^\wedge$ is
faithfully flat.
\end{lemma}

\begin{proof}
The first statement is an immediate consequence of
Lemma \ref{lemma-henselization-flat}
and in fact holds without assuming $A$ is Noetherian.
In the proof of Lemma \ref{lemma-henselization} we have seen that
$A^h$ is a filtered colimit of \'etale $A$-algebras $B$ such that
$A/I \to B/IB$ is an isomorphism. For each such $A \to B$
the induced map $A^\wedge \to B^\wedge$ is an isomorphism
(see proof of Lemma \ref{lemma-henselization-flat}).
By Algebra, Lemma \ref{algebra-lemma-completion-flat} the ring map
$B \to A^\wedge = B^\wedge = (A^h)^\wedge$ is flat for each $B$.
Thus $A^h \to A^\wedge = (A^h)^\wedge$ is flat by
Algebra, Lemma \ref{algebra-lemma-colimit-rings-flat}.
Since $I^h = IA^h$ is contained in the Jacobson radical of $A^h$
and since $A^h \to A^\wedge$ induces an isomorphism $A^h/I^h \to A/I$
we see that $A^h \to A^\wedge$ is faithfully flat by
Algebra, Lemma \ref{algebra-lemma-ff}.
By Algebra, Lemma \ref{algebra-lemma-completion-Noetherian-Noetherian}
the ring $A^\wedge$ is Noetherian.
Hence we conclude that $A^h$ is Noetherian by
Algebra, Lemma \ref{algebra-lemma-descent-Noetherian}.
\end{proof}

\begin{lemma}
\label{lemma-henselization-colimit}
Let $(A, I) = \colim (A_i, I_i)$ be a filtered colimit of pairs. The functor of
Lemma \ref{lemma-henselization} gives
$A^h = \colim A_i^h$ and $I^h = \colim I_i^h$.
\end{lemma}

\noindent
This lemma is false for non-filtered colimits, see
Example \ref{example-moret-bailly}.

\begin{proof}
By Categories, Lemma \ref{categories-lemma-adjoint-exact}
we see that $(A^h, I^h)$ is the colimit of the system $(A_i^h, I_i^h)$
in the category of henselian pairs. Thus for a henselian pair $(B, J)$
we have
$$
\Mor((A^h, I^h), (B, J)) =
\lim \Mor((A_i^h, I_i^h), (B, J)) =
\Mor(\colim (A_i^h, I_i^h), (B, J))
$$
Here the colimit is in the category of pairs. Since the colimit is
filtered we obtain $\colim (A_i^h, I_i^h) = (\colim A_i^h, \colim I_i^h)$
in the category of pairs; details omitted. Again using the colimit is filtered,
this is a henselian pair (Lemma \ref{lemma-filtered-colimits-henselian}).
Hence by the Yoneda lemma we
find $(A^h, I^h) = (\colim A_i^h, \colim I_i^h)$.
\end{proof}

\begin{lemma}
\label{lemma-henselization-change-ideal}
\begin{slogan}
The henselization of a pair only depends on the radical of the ideal
\end{slogan}
Let $A$ be a ring with ideals $I$ and $J$. If $V(I) = V(J)$ then the functor
of Lemma \ref{lemma-henselization} produces the same ring for the pair
$(A, I)$ as for the pair $(A, J)$.
\end{lemma}

\begin{proof}
Let $(A', IA')$ be the pair produced by Lemma \ref{lemma-henselization}
starting with the pair $(A, I)$, see Lemma \ref{lemma-henselization-flat}.
Let $(A'', JA'')$ be the pair produced by Lemma \ref{lemma-henselization}
starting with the pair $(A, J)$. By
Lemma \ref{lemma-change-ideal-henselian-pair} we see that
$(A', JA')$ is a henselian pair and $(A'', IA'')$ is a henselian pair.
By the universal property of the construction we obtain
unique $A$-algebra maps $A'' \to A'$ and $A' \to A''$. The
uniqueness shows that these are mutually inverse.
\end{proof}

\begin{lemma}
\label{lemma-henselization-integral}
\begin{slogan}
Henselization commutes with integral base change
\end{slogan}
Let $(A, I) \to (B, J)$ be a map of pairs such that $V(J) = V(IB)$.
Let $(A^h , I^h) \to (B^h, J^h)$ be the induced map
on henselizations (Lemma \ref{lemma-henselization}).
If $A \to B$ is integral, then the induced map
$A^h \otimes_A B \to B^h$ is an isomorphism.
\end{lemma}

\begin{proof}
By Lemma \ref{lemma-henselization-change-ideal} we may assume $J = IB$.
By Lemma \ref{lemma-integral-over-henselian-pair}
the pair $(A^h \otimes_A B, I^h(A^h \otimes_A B))$ is henselian.
By the universal property of $(B^h, IB^h)$ we obtain a map
$B^h \to A^h \otimes_A B$. We omit the proof
that this map is the inverse of the map in the lemma.
\end{proof}


\begin{lemma}
\label{lemma-CRT-henselization}
Let $I_1, I_2, \dots, I_n \subset A$ be ideals in a ring $A$.
Suppose $I_i + I_j = A$ for $i \not = j$, $1 \leq i, j \leq n$.
Denote $A^h$ the henselization of the pair $(A, I_1 \cap \ldots \cap I_n)$
and $A_i^h$ the henselization of the pair $(A, I_i)$.
Then the natural morphism
$$
A^h \to \prod\nolimits_{i = 1, \ldots, n} A^h_i
$$
is an isomorphism.
\end{lemma}

\begin{proof}
The map comes from the functoriality of the construction.
If $n = 1$ the result is immediate. If $n > 1$, observe that
$(I_1 \cap \ldots \cap I_{n - 1}) + I_n = A$ by our condition
on the ideals. Thus if we prove the lemma for $n = 2$, then
the lemma follows for all $n$ by induction.

\medskip\noindent
The case $n = 2$. Choose $f_i \in I_i$ such that $f_1 + f_2 = 1$.
Recall that $A^h$ is constructed as the colimit of $A$-algebras $B$
such that $A \to B$ is \'etale and $A/I_1 \cap I_2 \to B/(I_1 \cap I_2)B$
is an isomorphism. See proof of Lemma \ref{lemma-henselization}.
Denote $Hens$ the category of such $A$-algebras $B$.
Similarly, denote $Hens_i$, $i = 1, 2$ the category of \'etale $A$-algebras
$B$ such that $A/I_i \to B/I_iB$ is an isomorphism.

\medskip\noindent
Consider the subcategory $Hens_1(f_2) \subset Hens_1$ consisting of
those $B$ such that $f_2$ is invertible. This is a cofinal subcategory
since if $B$ is in $Hens_1$, then $B_{f_2}$ is in $Hens_1$ as well
(hint: $B_{f_2}/I_1B_{f_2} = (B/I_1B)_{f_2} = B/I_1B$ because
$f_2$ is a unit modulo $I_1$). Similarly, the subcategory
$Hens_2(f_1) \subset Hens_2$ is cofinal. Hence to compute $A_1^h$ and
$A_2^h$ we may work with this smaller categories, see
Categories, Lemma \ref{categories-lemma-cofinal}.

\medskip\noindent
Consider the functor
$$
Hens_1(f_2) \times Hens_2(f_1) \to Hens,\quad (B_1, B_2) \mapsto B_1 \times B_2
$$
This functor makes sense as
\begin{align*}
(B_1 \times B_2)/(I_1 \cap I_2)(B_1 \times B_2)
& =
B_1/(I_1 \cap I_2)B_1 \times B_2/(I_1 \cap I_2)B_2 \\
& =
B_1/I_1B_1 \times B_2/I_2B_2 \\
& = A/I_1 \times A/I_2 \\
& = A/(I_1 \cap I_2)
\end{align*}
The second equality holds because
$(I_1 \cap I_2)B_1 = I_1B_1 \cap I_2B_1 = I_1B_1$ as $f_2 \in I_2$
is invertible in $B_1$
(we have also used Algebra, Lemma \ref{algebra-lemma-flat-intersect-ideals}).
The last equality holds by Algebra, Lemma \ref{algebra-lemma-chinese-remainder}.
To finish the proof, it suffices to prove that the displayed functor
is cofinal (see lemma quoted above).
The first condition of
Categories, Definition \ref{categories-definition-cofinal}
holds because an object $B$ of $Hens$ maps to $B_{f_2} \times B_{f_1}$
which is in the image of the functor. The second one always holds because
if we have $B \to B_1 \times B_2$ and $B \to B'_1 \times B'_2$ with
$B_1, B'_1 \in Hens_1(f_2)$ and $B_2, B'_2 \in Hens_2(f_1)$, then
we can set $B''_1 = B_1 \otimes_B B'_1$ which is in $Hens_1(f_2)$
and similarly $B''_2 = B_2 \otimes_B B''_2$ in $Hens_2(f_1)$
to find a map $B \to B''_1 \times B''_2$ fitting into the commutative diagram
$$
\xymatrix{
& B \ar[ld] \ar[d] \ar[rd] \\
B_1 \times B_2 \ar[r] & B_1'' \times B_2'' & B_1' \times B_2' \ar[l]
}
$$
This finishes the proof.
\end{proof}









\section{Lifting and henselian pairs}
\label{section-lifting-henselian-pairs}

\noindent
In this section we mostly combine results from
Sections \ref{section-lifting} and \ref{section-henselian-pairs}.

\begin{lemma}
\label{lemma-lift-finite-projective-module}
Let $(R, I)$ be a henselian pair. The map
$$
P \longrightarrow P/IP
$$
induces a bijection between the sets of isomorphism classes of finite
projective $R$-modules and finite projective $R/I$-modules. In particular,
any finite projective $R/I$-module is isomorphic to $P/IP$
for some finite projective $R$-module $P$.
\end{lemma}

\begin{proof}
We first prove the final statement. Let $\overline{P}$ be a finite
projective $R/I$-module. We can find a finite projective module $P'$
over some $R'$ \'etale over $R$ with $R/I = R'/IR'$ such that
$P'/IP'$ is isomorphic to $\overline{P}$, see
Lemma \ref{lemma-lift-projective-module}.
Then, since $(R, I)$ is a henselian pair, the \'etale ring map $R \to R'$ has
a section $\tau : R' \to R$ (Lemma \ref{lemma-characterize-henselian-pair}).
Setting $P = P' \otimes_{R', \tau} R$ we conclude that $P/IP$
is isomorphic to $\overline{P}$. Of course, this tells us that
the map in the statement of the lemma is surjective.

\medskip\noindent
Injectivity. Suppose that $P_1$ and $P_2$ are finite projective $R$-modules such
that $P_1/IP_1 \cong P_2/IP_2$ as $R/I$-modules. Since $P_1$ is projective,
we can find an $R$-module map $u : P_1 \to P_2$ lifting the given isomorphism.
Then $u$ is surjective by Nakayama's lemma
(Algebra, Lemma \ref{algebra-lemma-NAK}).
We similarly find a surjection $v : P_2 \to P_1$.
By Algebra, Lemma \ref{algebra-lemma-fun} the map $v \circ u$
is an isomorphism and we conclude $u$ is an isomorphism.
\end{proof}

\begin{lemma}
\label{lemma-finite-etale-equivalence}
Let $(A, I)$ be a henselian pair. The functor $B \to B/IB$ determines
an equivalence between finite \'etale $A$-algebras and finite \'etale
$A/I$-algebras.
\end{lemma}

\begin{proof}
Let $B, B'$ be two $A$-algebras finite \'etale over $A$.
Then $B' \to B'' = B \otimes_A B'$ is finite \'etale as well
(Algebra, Lemmas \ref{algebra-lemma-etale} and
\ref{algebra-lemma-base-change-integral}).
Now we have $1$-to-$1$ correspondences between
\begin{enumerate}
\item $A$-algebra maps $B \to B'$,
\item sections of $B' \to B''$, and
\item idempotents $e$ of $B''$ such that $B' \to B'' \to eB''$ is
an isomorphism.
\end{enumerate}
The bijection between (2) and (3) sends $\sigma : B'' \to B'$
to $e$ such that $(1 - e)$ is the idempotent
that generates the kernel of $\sigma$ which exists by
Algebra, Lemmas \ref{algebra-lemma-map-between-etale} and
\ref{algebra-lemma-surjective-flat-finitely-presented}.
There is a similar correspondence between
$A/I$-algebra maps $B/IB \to B'/IB'$ and idempotents
$\overline{e}$ of $B''/IB''$ such that
$B'/IB' \to B''/IB'' \to \overline{e}(B''/IB'')$ is
an isomorphism. However every idempotent $\overline{e}$ of $B''/IB''$
lifts uniquely to an idempotent $e$ of $B''$
(Lemma \ref{lemma-characterize-henselian-pair}).
Moreover, if $B'/IB' \to \overline{e}(B''/IB'')$ is an isomorphism,
then $B' \to eB''$ is an isomorphism too by Nakayama's lemma
(Algebra, Lemma \ref{algebra-lemma-NAK}).
In this way we see that the functor is fully faithful.

\medskip\noindent
Essential surjectivity. Let $A/I \to C$ be a finite \'etale map.
By Algebra, Lemma \ref{algebra-lemma-lift-etale}
there exists an \'etale map $A \to B$ such that $B/IB \cong C$.
Let $B'$ be the integral closure of $A$ in $B$. 
By Lemma \ref{lemma-helper-finite-type} we have
$B'/IB' = C \times C'$ for some ring $C'$
and $B'_g \cong B_g$ for some $g \in B'$ mapping to $(1, 0) \in C \times C'$.
Since idempotents lift
(Lemma \ref{lemma-characterize-henselian-pair})
we get $B' = B'_1 \times B'_2$ with $C = B'_1/IB'_1$ and $C' = B'_2/IB'_2$.
The image of $g$ in $B'_1$ is invertible.
Then $B_g = B'_g = B'_1 \times (B_2)_g$ and this implies
that $A \to B'_1$ is \'etale.
We conclude that $B'_1$ is finite \'etale over $A$
(integral \'etale implies finite \'etale by
Algebra, Lemma \ref{algebra-lemma-characterize-finite-in-terms-of-integral}
for example)
and the proof is done.
\end{proof}

\begin{lemma}
\label{lemma-lift-smooth-henselian}
Let $R$ be a ring and $S$ a smooth $R$-algebra. Assume that $A$
is an $R$-algebra and $(A,I)$ is a henselian pair. Then any $R$-algebra
map $S \to A/I$ can be lifted to an $R$-algebra map $S \to A$.
\end{lemma}

\begin{proof}
Let $\tau : S \to A/I$ be an $R$-algebra map.
Observe that $S \otimes_R A$ is a smooth $A$-algebra by
Algebra, Lemma \ref{algebra-lemma-base-change-smooth}.
Thus by Lemma \ref{lemma-lift-section-smooth-morphism} we can lift
the induced map $S \otimes_R A \to A/I$ to an $A$-algebra
homorphism $S \otimes_R A \to A'$ where $A \to A'$ is \'etale
and induces an isomorphism $A/I \to A'/IA'$. Since $(A, I)$
is henselian there is an $A$-algebra map $A' \to A$, see
Lemma \ref{lemma-characterize-henselian-pair}.
The composition $S \to S \otimes_R A \to A' \to A$ is the
desired lift.
\end{proof}

\begin{lemma}
\label{lemma-lim-finite-projective-gives-finite-projective}
Let $A = \lim A_n$ be a limit of an inverse system $(A_n)$ of rings.
Suppose given $A_n$-modules $M_n$ and $A_{n + 1}$-module maps
$M_{n + 1} \to M_n$. Assume
\begin{enumerate}
\item the transition maps $A_{n + 1} \to A_n$ are surjective
with locally nilpotent kernels,
\item $M_1$ is a finite projective $A_1$-module,
\item $M_n$ is a finite flat $A_n$-module, and
\item the maps induce isomorphisms
$M_{n + 1} \otimes_{A_{n + 1}} A_n \to M_n$.
\end{enumerate}
Then $M = \lim M_n$ is a finite projective $A$-module
and $M \otimes_A A_n \to M_n$ is an isomorphism for all $n$.
\end{lemma}

\begin{proof}
By Lemma \ref{lemma-limit-henselian} the pair $(A, \Ker(A \to A_1))$ is
henselian. By Lemma \ref{lemma-lift-finite-projective-module}
we can choose a finite projective $A$-module $P$ and an isomorphism
$P \otimes_A A_1 \to M_1$. Since $P$ is projective, we can successively lift
the $A$-module map $P \to M_1$ to $A$-module maps
$P \to M_2$, $P \to M_3$, and so on. Thus we obtain a map
$$
P \longrightarrow M
$$
Since $P$ is finite projective, we can write $A^{\oplus m} = P \oplus Q$
for some $m \geq 0$ and $A$-module $Q$. Since $A = \lim A_n$ we conclude
that $P = \lim P \otimes_A A_n$. Hence, in order to show that the displayed
$A$-module map is an isomorphism, it suffices to show that the maps
$P \otimes_A A_n \to M_n$ are isomorphisms.
From Lemma \ref{lemma-lift-projective} we see that
$M_n$ is a finite projective module.
By Lemma \ref{lemma-isomorphic-finite-projective-lifts}
the maps $P \otimes_A A_n \to M_n$ are isomorphisms.
\end{proof}






\section{Absolute integral closure}
\label{section-absolute-integral-closure}

\noindent
Here is our definition.

\begin{definition}
\label{definition-absolutely-integrally-closed}
A ring $A$ is {\it absolutely integrally closed} if every
monic $f \in A[T]$ is a product of linear factors.
\end{definition}

\noindent
Be careful: it may be possible to write $f$ as a product of
linear factors in many different ways.

\begin{lemma}
\label{lemma-absolutely-integrally-closed}
Let $A$ be a ring. The following are equivalent
\begin{enumerate}
\item $A$ is absolutely integrally closed, and
\item any monic $f \in A[T]$ has a root in $A$.
\end{enumerate}
\end{lemma}

\begin{proof}
Omitted.
\end{proof}

\begin{lemma}
\label{lemma-absolutely-integrally-closed-quotient-localization}
Let $A$ be absolutely integrally closed.
\begin{enumerate}
\item Any quotient ring $A/I$ of $A$ is absolutely integrally closed.
\item Any localization $S^{-1}A$ is absolutely integrally closed.
\end{enumerate}
\end{lemma}

\begin{proof}
Omitted.
\end{proof}

\begin{lemma}
\label{lemma-integrally-closed-in-absolutely-integrally-closed}
Let $A$ be a ring. Let $S \subset A$ be a multiplicative subset
consisting of nonzerodivisors.
If $S^{-1}A$ is absolutely integrally closed and $A \subset S^{-1}A$
is integrally closed in $S^{-1}A$, then $A$ is absolutely integrally closed.
\end{lemma}

\begin{proof}
Omitted.
\end{proof}

\begin{lemma}
\label{lemma-normal-domain-absolutely-integrally-closed}
Let $A$ be a normal domain. Then $A$ is absolutely integrally closed
if and only if its fraction field is algebraically closed.
\end{lemma}

\begin{proof}
Observe that a field is algebraically closed if and only if
it is absolutely integrally closed as a ring. Hence the lemma follows
from Lemmas \ref{lemma-absolutely-integrally-closed-quotient-localization} and
\ref{lemma-integrally-closed-in-absolutely-integrally-closed}.
\end{proof}

\begin{lemma}
\label{lemma-construct-absolute-integral-closure}
For any ring $A$ there exists an extension $A \subset B$ such that
\begin{enumerate}
\item $B$ is a filtered colimit of finite free $A$-algebras,
\item $B$ is free as an $A$-module, and
\item $B$ is absolutely integrally closed.
\end{enumerate}
\end{lemma}

\begin{proof}
Let $I$ be the set of monic polynomials over $A$. For $i \in I$
denote $x_i$ a variable and $P_i$ the corresponding monic polynomial
in the variable $x_i$. Then we set
$$
F(A) = A[x_i; i \in I]/(P_i; i \in I)
$$
As the notation suggests $F$ is a functor from the category
of rings to itself. Note that $A \subset F(A)$, that $F(A)$ is free
as an $A$-module, and that $F(A)$ is a filtered colimit of
finite free $A$-algebras. Then we take
$$
B = \colim F^n(A)
$$
where the transition maps are the inclusions
$F^n(A) \subset F(F^n(A)) = F^{n + 1}(A)$.
Any monic polynomial with coefficients in $B$
actually has coefficients in $F^n(A)$ for some $n$
and hence has a solution in $F^{n + 1}(A)$ by construction.
This implies that $B$ is absolutely integrally closed by
Lemma \ref{lemma-absolutely-integrally-closed}.
We omit the proof of the other properties.
\end{proof}

\begin{lemma}
\label{lemma-absolutely-integrally-closed-strictly-henselian}
Let $A$ be absolutely integrally closed. Let $\mathfrak p \subset A$
be a prime. Then the local ring $A_\mathfrak p$ is strictly henselian.
\end{lemma}

\begin{proof}
By Lemma \ref{lemma-absolutely-integrally-closed-quotient-localization}
we may assume $A$ is a local ring and $\mathfrak p$ is its maximal ideal.
The residue field is algebraically closed by
Lemma \ref{lemma-absolutely-integrally-closed-quotient-localization}.
Every monic polynomial decomposes completely into linear factors
hence Algebra, Definition \ref{algebra-definition-henselian} applies directly.
\end{proof}

\begin{lemma}
\label{lemma-absolutely-integrally-closed-henselian-pair}
Let $A$ be absolutely integrally closed. Let $I \subset A$ be an ideal.
Then $(A, I)$ is a henselian pair if (and only if) the following
conditions hold
\begin{enumerate}
\item $I$ is contained in the Jacobson radical of $A$,
\item $A \to A/I$ induces a bijection on idempotents.
\end{enumerate}
\end{lemma}

\begin{proof}
Let $f \in A[T]$ be a monic polynomial and let
$f \bmod I = g_0 h_0$ be a factorization over $A/I$
with $g_0$, $h_0$ monic such that
$g_0$ and $h_0$ generate the unit ideal of $A/I[T]$.
This means that
$$
A/I[T]/(f) = A/I[T]/(g_0) \times A/I[T]/(h_0)
$$
Denote $e \in A/I[T]/(f)$ the element correspoing to the
idempotent $(1, 0)$ in the ring on the right.
Write $f = (T - a_1) \ldots (T - a_d)$ with $a_i \in A$.
For each $i \in \{1, \ldots, d\}$ we obtain an $A$-algebra map
$\varphi_i : A[T]/(f) \to A$, $T \mapsto a_i$ which induces
a similar $A/I$-algebra map $\overline{\varphi}_i : A/I[T]/(f) \to A/I$.
Denote $e_i = \overline{\varphi}_i(e) \in A/I$. These are idempotents.
By our assumption (2) we can lift $e_i$ to an idempotent in $A$.
This means we can write $A = \prod A_j$ as a finite product of
rings such that in $A_j/IA_j$ each $e_i$ is either $0$ or $1$.
Some details omitted. Observe that $A_j$ is absolutely integrally
closed as a factor ring of $A$. It suffices to lift the factorization
of $f$ over $A_j/IA_j$ to $A_j$.
This reduces us to the situation discussed in the next paragraph.

\medskip\noindent
Assume $e_i = 1$ for $i = 1, \ldots, r$ and $e_i = 0$ for
$i = r + 1, \ldots, d$. From $(g_0, h_0) = A/I[T]$ we have that there are
$k_0, l_0 \in A/I[T]$ such that $g_0 k_0 + h_0 l_0 = 1$.
We see that $e = h_0 l_0$ and $e_i = h_0(a_i) l_0(a_i)$.
We conclude that $h_0(a_i)$ is a unit for $i = 1, \dots ,r$.
Since $f(a_i) = 0$ we find $0 = h_0(a_i)g_0(a_i)$
and we conclude that $g_0(a_i) = 0$ for $i = 1, \ldots, r$.
Thus $(T - a_1)$ divides $g_0$ in $A/I[T]$, say
$g_0 = (T - a_1) g_0'$.
Set $f' = (T - a_2) \ldots (T - a_d)$ and $h'_0 = h_0$.
By induction on $d$ we can lift the factorization $f' \bmod I = g'_0 h'_0$
to a factorization of $f' = g' h'$ over over $A$ which
gives the factorization $f = (T - a_1) g' h'$
lifting the factorization $f \bmod I = g_0 h_0$
as desired.
\end{proof}









\section{Auto-associated rings}
\label{section-auto-ass}

\noindent
Some of this material is in \cite{Autour}.

\begin{definition}
\label{definition-auto-ass}
A ring $R$ is said to be {\it auto-associated} if $R$ is local and its
maximal ideal $\mathfrak m$ is weakly associated to $R$.
\end{definition}

\begin{lemma}
\label{lemma-auto-ass-implies-P}
An auto-associated ring $R$ has the following property: (P)
Every proper finitely generated ideal $I \subset R$ has a nonzero
annihilator.
\end{lemma}

\begin{proof}
By assumption there exists a nonzero element $x \in R$ such that for every
$f \in \mathfrak m$ we have $f^n x = 0$. Say $I = (f_1, \ldots, f_r)$.
Then $x$ is in the kernel of $R \to \bigoplus R_{f_i}$. Hence we see
that there exists a nonzero $y \in R$ such that $f_i y = 0$ for all $i$, see
Algebra, Lemma \ref{algebra-lemma-when-injective-covering}.
As $y \in \text{Ann}_R(I)$ we win.
\end{proof}

\begin{lemma}
\label{lemma-P-universally-injective}
Let $R$ be a ring having property (P) of
Lemma \ref{lemma-auto-ass-implies-P}.
Let $u : N \to M$ be a homomorphism of projective $R$-modules.
Then $u$ is universally injective if and only if $u$ is injective.
\end{lemma}

\begin{proof}
Assume $u$ is injective. Our goal is to show $u$ is universally injective.
First we choose a module $Q$ such that $N \oplus Q$ is free. On considering
the map $N \oplus Q \to M \oplus Q$ we see that it suffices to prove
the lemma in case $N$ is free. In this case $N$ is a directed colimit of
finite free $R$-modules. Thus we reduce to the case that $N$ is a finite
free $R$-module, say $N = R^{\oplus n}$. We prove the lemma by induction
on $n$. The case $n = 0$ is trivial.

\medskip\noindent
Let $u : R^{\oplus n} \to M$ be an injective module map with $M$ projective.
Choose an $R$-module $Q$ such that $M \oplus Q$ is free. After replacing
$u$ by the composition $R^{\oplus n} \to M \to M \oplus Q$ we see that
we may assume that $M$ is free. Then we can find a direct summand
$R^{\oplus m} \subset M$ such that $u(R^{\oplus n}) \subset R^{\oplus m}$.
Hence we may assume that $M = R^{\oplus m}$.
In this case $u$ is given by a matrix $A = (a_{ij})$ so that
$u(x_1, \ldots, x_n) = (\sum x_i a_{i1}, \ldots, \sum x_i a_{im})$.
As $u$ is injective, in particular
$u(x, 0, \ldots, 0) = (xa_{11}, xa_{12}, \ldots, xa_{1m}) \not = 0$ if
$x \not = 0$, and as $R$ has property (P) we see that
$a_{11}R + a_{12}R + \ldots + a_{1m}R = R$. Hence see that
$R(a_{11}, \ldots, a_{1m}) \subset R^{\oplus m}$ is a direct summand
of $R^{\oplus m}$, in particular $R^{\oplus m}/R(a_{11}, \ldots, a_{1m})$
is a projective $R$-module. We get a commutative diagram
$$
\xymatrix{
0 \ar[r] &
R \ar[rr] \ar[d]^1 & & R^{\oplus n} \ar[r] \ar[d]^u &
R^{\oplus n - 1} \ar[r] \ar[d] & 0 \\
0 \ar[r] & R \ar[rr]^{(a_{11}, \ldots, a_{1m})} & &
R^{\oplus m} \ar[r] & R^{\oplus m}/R(a_{11}, \ldots, a_{1m}) \ar[r] & 0
}
$$
with split exact rows. Thus the right vertical arrow is injective
and we may apply the induction hypothesis to conclude that
the right vertical arrow is universally injective. It follows that the
middle vertical arrow is universally injective.
\end{proof}

\begin{lemma}
\label{lemma-P-fPD-zero}
Let $R$ be a ring. The following are equivalent
\begin{enumerate}
\item $R$ has property (P) of
Lemma \ref{lemma-auto-ass-implies-P},
\item any injective map of projective $R$-modules is
universally injective,
\item if $u : N \to M$ is injective and $N$, $M$ are finite projective
$R$-modules then $\Coker(u)$ is a finite projective $R$-module,
\item if $N \subset M$ and $N$, $M$ are finite projective as $R$-modules, then
$N$ is a direct summand of $M$, and
\item any injective map $R \to R^{\oplus n}$ is a split injection.
\end{enumerate}
\end{lemma}

\begin{proof}
The implication (1) $\Rightarrow$ (2) is
Lemma \ref{lemma-P-universally-injective}.
It is clear that (3) and (4) are equivalent.
We have (2) $\Rightarrow$ (3), (4) by
Algebra, Lemma \ref{algebra-lemma-universally-exact-split}.
Part (5) is a special case of (4).
Assume (5). Let $I = (a_1, \ldots, a_n)$ be a proper finitely generated
ideal of $R$. As $I \not = R$ we see that
$R \to R^{\oplus n}$, $x \mapsto (xa_1, \ldots, xa_n)$
is not a split injection. Hence it has a nonzero kernel and we conclude
that $\text{Ann}_R(I) \not = 0$. Thus (1) holds.
\end{proof}

\begin{example}
\label{example-auto-ass-weird-flat}
If the equivalent conditions of
Lemma \ref{lemma-P-fPD-zero}
hold, then it is not always the case that every injective map of
free $R$-modules is a split injection. For example suppose that
$R = k[x_1, x_2, x_3, \ldots]/(x_i^2)$. This is an auto-associated ring.
Consider the map of free $R$-modules
$$
u :
\bigoplus\nolimits_{i \geq 1} Re_i
\longrightarrow
\bigoplus\nolimits_{i \geq 1} Rf_i, \quad
e_i \longmapsto f_i - x_if_{i + 1}.
$$
For any integer $n$ the restriction of $u$ to
$\bigoplus_{i = 1, \ldots, n} Re_i$ is injective as the images
$u(e_1), \ldots, u(e_n)$ are $R$-linearly independent. Hence
$u$ is injective and hence universally injective by the lemma.
Since $u \otimes \text{id}_k$ is bijective we see that if
$u$ were a split injection then $u$ would be surjective. But $u$ is not
surjective because the inverse image of $f_1$ would be the element
$$
\sum\nolimits_{i \geq 0} x_1 \ldots x_ie_{i + 1} =
e_1 + x_1e_2 + x_1x_2e_3 + \ldots
$$
which is not an element of the direct sum. A side remark is that
$\Coker(u)$ is a flat (because $u$ is universally injective),
countably generated $R$-module which is not projective (as $u$ is not
split), hence not Mittag-Leffler (see
Algebra, Lemma \ref{algebra-lemma-countgen-projective}).
\end{example}

\noindent
The following lemma is a special case of
Algebra, Proposition \ref{algebra-proposition-what-exact}
in case the local ring is Noetherian.

\begin{lemma}
\label{lemma-exact-length-1}
Let $(R, \mathfrak m)$ be a local ring.
Suppose that $\varphi : R^m \to R^n$ is a map
of finite free modules. The following are equivalent
\begin{enumerate}
\item $\varphi$ is injective,
\item the rank of $\varphi$ is $m$ and the annihilator of
$I(\varphi)$ in $R$ is zero.
\end{enumerate}
If $R$ is Noetherian these are also equivalent to
\begin{enumerate}
\item[(3)] the rank of $\varphi$ is $m$ and
either $I(\varphi) = R$ or it contains a nonzerodivisor.
\end{enumerate}
Here the rank of $\varphi$ and $I(\varphi)$ are defined
as in Algebra, Definition \ref{algebra-definition-rank}.
\end{lemma}

\begin{proof}
If any matrix coefficient of $\varphi$ is not in $\mathfrak m$, then we apply
Algebra, Lemma \ref{algebra-lemma-add-trivial-complex}
to write $\varphi$ as the sum of $1 : R \to R$ and a map
$\varphi' : R^{m-1} \to R^{n-1}$. It is easy to see that
the lemma for $\varphi'$ implies the lemma for $\varphi$.
Thus we may assume from the outset that all the matrix
coefficients of $\varphi$ are in $\mathfrak m$.

\medskip\noindent
Suppose $\varphi$ is injective. We may assume $m > 0$.
Let $\mathfrak q \in \text{WeakAss}(R)$ so that $R_\mathfrak q$
is an auto-associated ring. Then $\varphi$
induces an injective map $R_\mathfrak q^m \to R_\mathfrak q^n$
which is universally injective by
Lemmas \ref{lemma-auto-ass-implies-P} and \ref{lemma-P-universally-injective}.
Thus $\varphi : \kappa(\mathfrak q)^m \to \kappa(\mathfrak q)^n$
is injective. Hence the rank of $\varphi \bmod \mathfrak q$ is $m$ and
$I(\varphi \otimes \kappa(\mathfrak q))$ is not the zero ideal.
Since $m$ is the maximum rank $\varphi$ can have, we conclude
that $\varphi$ has rank $m$ as well (ranks of matrices
can only drop after base change). Hence
$I(\varphi) \cdot \kappa(\mathfrak q) =
I(\varphi \otimes \kappa(\mathfrak q))$ is not zero.
Thus $I(\varphi)$ is not contained in $\mathfrak q$.
Thus none of the weakly associated primes of $R$
are weakly associated primes of the $R$-module
$\text{Ann}_R I(\varphi)$. Thus $\text{Ann}_R I(\varphi)$
has no weakly associated primes, see
Algebra, Lemma \ref{algebra-lemma-weakly-ass}.
It follows from Algebra, Lemma
\ref{algebra-lemma-weakly-ass-zero}
that $\text{Ann}_R I(\varphi)$ is zero.

\medskip\noindent
Conversely, assume (2). The rank being $m$ implies $n \geq m$.
Write $I(\varphi) = (f_1, \ldots, f_r)$ which is possible as
$I(\varphi)$ is finitely generated. By
Algebra, Lemma \ref{algebra-lemma-matrix-left-inverse}
we can find maps $\psi_i : R^n \to R^m$ such that
$\psi \circ \varphi = f_i \text{id}_{R^m}$. Thus
$\varphi(x) = 0$ implies $f_i x = 0$ for $i = 1, \ldots, r$.
This implies $x = 0$ and hence $\varphi$ is injective.

\medskip\noindent
For the equivalence of (1) and (3) in the Noetherian local case
we refer to Algebra, Proposition \ref{algebra-proposition-what-exact}.
If the ring $R$ is Noetherian but not local, then the reader can
deduce it from the local case; details omitted.
Another option is to redo the argument above using associated
primes, using that there are finitely many of these, using
prime avoidance, and using the characterization of nonzerodivisors
as elements of a Noetherian ring not contained in any associated prime.
\end{proof}

\begin{lemma}
\label{lemma-coker-injective-free}
Let $R$ be a ring. Suppose that $\varphi : R^n \to R^n$ be an injective
map of finite free modules of the same rank. Then
$\Hom_R(\Coker(\varphi), R) = 0$.
\end{lemma}

\begin{proof}
Let $\varphi^t : R^n \to R^n$ be the transpose of $\varphi$.
The lemma claims that $\varphi^t$ is injective.
With notation as in Lemma \ref{lemma-exact-length-1}
we see that the rank of $\varphi^t$ is $n$ and that
$I(\varphi) = I(\varphi^t)$. Thus we conclude by
the equivalence of (1) and (2) of the lemma.
\end{proof}









\section{Flattening stratification}
\label{section-flattening}

\noindent
Let $R \to S$ be a ring map and let $M$ be an $S$-module.
For any $R$-algebra $R'$ we can consider the base changes
$S' = S \otimes_R R'$ and $M' = M \otimes_R R'$.
We say $R \to R'$ {\it flattens} $M$ if the module $M'$ is flat over $R'$.
We would like to understand the structure of the collection of ring maps
$R \to R'$ which flatten $M$. In particular we would like to know
if there exists a {\it universal flattening $R \to R_{univ}$ of $M$},
i.e., a ring map $R \to R_{univ}$ which flattens $M$ and has the property
that any ring map $R \to R'$ which flattens $M$ factors through
$R \to R_{univ}$. It turns out that such a universal solution usually
does not exist.

\medskip\noindent
We will discuss {\it universal flattenings} and
{\it flattening stratifications} in a scheme theoretic
setting $\mathcal{F}/X/S$ in
More on Flatness, Section \ref{flat-section-flattening}.
If the universal flattening $R \to R_{univ}$ exists then the
morphism of schemes $\Spec(R_{univ}) \to \Spec(R)$ is the
universal flattening of the quasi-coherent module
$\widetilde{M}$ on $\Spec(S)$.

\medskip\noindent
In this and the next few sections we prove some basic algebra facts
related to this. The most basic result is perhaps the following.

\begin{lemma}
\label{lemma-intersection-flat}
Let $R$ be a ring. Let $M$ be an $R$-module. Let $I_1$, $I_2$ be ideals of $R$.
If $M/I_1M$ is flat over $R/I_1$ and $M/I_2M$ is flat over $R/I_2$,
then $M/(I_1 \cap I_2)M$ is flat over $R/(I_1 \cap I_2)$.
\end{lemma}

\begin{proof}
By replacing $R$ with $R/(I_1 \cap I_2)$ and $M$ by $M/(I_1 \cap I_2)M$
we may assume that $I_1 \cap I_2 = 0$. Let $J \subset R$ be an ideal.
To prove that $M$ is flat over $R$ we have to show that
$J \otimes_R M \to M$ is injective, see
Algebra, Lemma \ref{algebra-lemma-flat}.
By flatness of $M/I_1M$ over $R/I_1$ the map
$$
J/(J \cap I_1) \otimes_R M =
(J + I_1)/I_1 \otimes_{R/I_1} M/I_1M
\longrightarrow M/I_1M
$$
is injective. As $0 \to (J \cap I_1) \to J \to J/(J \cap I_1) \to 0$
is exact we obtain a diagram
$$
\xymatrix{
(J \cap I_1) \otimes_R M \ar[r] \ar[d] &
J \otimes_R M \ar[r] \ar[d] &
J/(J \cap I_1) \otimes_R M \ar[r] \ar[d] & 0 \\
M \ar@{=}[r] &
M \ar[r] &
M/I_1M
}
$$
hence it suffices to show that $(J \cap I_1) \otimes_R M \to M$ is
injective. Since $I_1 \cap I_2 = 0$ the ideal $J \cap I_1$
maps isomorphically to an ideal $J' \subset R/I_2$ and we see that
$(J \cap I_1) \otimes_R M = J' \otimes_{R/I_2} M/I_2M$. By flatness
of $M/I_2M$ over $R/I_2$ the map $J' \otimes_{R/I_2} M/I_2M \to M/I_2M$
is injective, which clearly implies that $(J \cap I_1) \otimes_R M \to M$ is
injective.
\end{proof}





\section{Flattening over an Artinian ring}
\label{section-flattening-artinian}

\noindent
A universal flattening exists when the base ring is an Artinian local
ring. It exists for an arbitrary module. Hence, as we will see later,
a flatting stratification exists when the base scheme is the spectrum
of an Artinian local ring.

\begin{lemma}
\label{lemma-flattening-artinian}
Let $R$ be an Artinian ring.
Let $M$ be an $R$-module.
Then there exists a smallest ideal $I \subset R$ such that
$M/IM$ is flat over $R/I$.
\end{lemma}

\begin{proof}
This follows directly from
Lemma \ref{lemma-intersection-flat}
and the Artinian property.
\end{proof}

\noindent
This ideal has the following universal property.

\begin{lemma}
\label{lemma-flattening-artinian-universal-property}
Let $R$ be an Artinian ring. Let $M$ be an $R$-module.
Let $I \subset R$ be the smallest ideal $I \subset R$ such that
$M/IM$ is flat over $R/I$.
Then $I$ has the following universal property:
For every ring map $\varphi : R \to R'$ we have
$$
R' \otimes_R M\text{ is flat over }R'
\Leftrightarrow
\text{we have }\varphi(I) = 0.
$$
\end{lemma}

\begin{proof}
Note that $I$ exists by
Lemma \ref{lemma-flattening-artinian}.
The implication $\Rightarrow$ follows from
Algebra, Lemma \ref{algebra-lemma-flat-base-change}.
Let $\varphi : R \to R'$ be such that $M \otimes_R R'$ is flat over $R'$.
Let $J = \Ker(\varphi)$. By
Algebra,
Lemma \ref{algebra-lemma-descent-flatness-injective-map-artinian-rings}
and as $R' \otimes_R M = R' \otimes_{R/J} M/JM$ is
flat over $R'$ we conclude that $M/JM$ is flat over $R/J$.
Hence $I \subset J$ as desired.
\end{proof}











\section{Flattening over a closed subset of the base}
\label{section-flattening-local-base}

\noindent
Let $R \to S$ be a ring map.
Let $I \subset R$ be an ideal.
Let $M$ be an $S$-module.
In the following we will consider the following condition
\begin{equation}
\label{equation-flat-at-primes-over}
\forall \mathfrak q \in V(IS) \subset \Spec(S) :
M_{\mathfrak q}\text{ is flat over }R.
\end{equation}
Geometrically, this means that $M$ is flat over $R$ along the inverse
image of $V(I)$ in $\Spec(S)$. If $R$ and $S$ are Noetherian rings and
$M$ is a finite $S$-module, then (\ref{equation-flat-at-primes-over})
is equivalent to the condition that $M/I^nM$ is flat over $R/I^n$
for all $n \geq 1$, see
Algebra, Lemma \ref{algebra-lemma-flat-module-powers}.

\begin{lemma}
\label{lemma-base-change-flat-at-primes-over}
Let $R \to S$ be a ring map.
Let $I \subset R$ be an ideal.
Let $M$ be an $S$-module.
Let $R \to R'$ be a ring map and $IR' \subset I' \subset R'$ an ideal.
If (\ref{equation-flat-at-primes-over}) holds for
$(R \to S, I, M)$, then (\ref{equation-flat-at-primes-over})
holds for $(R' \to S \otimes_R R', I', M \otimes_R R')$.
\end{lemma}

\begin{proof}
Assume (\ref{equation-flat-at-primes-over}) holds for
$(R \to S, I \subset R, M)$.
Let $I'(S \otimes_R R') \subset \mathfrak q'$ be a prime of $S \otimes_R R'$.
Let $\mathfrak q \subset S$ be the corresponding prime of $S$.
Then $IS \subset \mathfrak q$. Note that $(M \otimes_R R')_{\mathfrak q'}$
is a localization of the base change $M_{\mathfrak q} \otimes_R R'$.
Hence $(M \otimes_R R')_{\mathfrak q'}$ is flat over $R'$ as a localization
of a flat module, see
Algebra, Lemmas \ref{algebra-lemma-flat-base-change} and
\ref{algebra-lemma-flat-localization}.
\end{proof}

\begin{lemma}
\label{lemma-flat-descent-flat-at-primes-over}
Let $R \to S$ be a ring map.
Let $I \subset R$ be an ideal.
Let $M$ be an $S$-module.
Let $R \to R'$ be a ring map and $IR' \subset I' \subset R'$ an ideal
such that
\begin{enumerate}
\item the map $V(I') \to V(I)$ induced by
$\Spec(R') \to \Spec(R)$ is surjective, and
\item $R'_{\mathfrak p'}$ is flat over $R$ for all primes
$\mathfrak p' \in V(I')$.
\end{enumerate}
If (\ref{equation-flat-at-primes-over}) holds for
$(R' \to S \otimes_R R', I', M \otimes_R R')$, then
(\ref{equation-flat-at-primes-over}) holds for $(R \to S, I, M)$.
\end{lemma}

\begin{proof}
Assume (\ref{equation-flat-at-primes-over}) holds for
$(R' \to S \otimes_R R', IR', M \otimes_R R')$. Pick a prime
$IS \subset \mathfrak q \subset S$. Let $I \subset \mathfrak p \subset R$
be the corresponding prime of $R$. By assumption there exists
a prime $\mathfrak p' \in V(I')$ of $R'$ lying over $\mathfrak p$ and
$R_{\mathfrak p} \to R'_{\mathfrak p'}$ is flat. Choose a prime
$\overline{\mathfrak q}' \subset
\kappa(\mathfrak q) \otimes_{\kappa(\mathfrak p)} \kappa(\mathfrak p')$
which corresponds to a prime $\mathfrak q' \subset S \otimes_R R'$ which
lies over $\mathfrak q$ and over $\mathfrak p'$. Note that
$(S \otimes_R R')_{\mathfrak q'}$ is a localization of
$S_{\mathfrak q} \otimes_{R_{\mathfrak p}} R'_{\mathfrak p'}$.
By assumption the module $(M \otimes_R R')_{\mathfrak q'}$ is
flat over $R'_{\mathfrak p'}$. Hence
Algebra, Lemma \ref{algebra-lemma-base-change-flat-up-down}
implies that $M_{\mathfrak q}$ is flat over $R_{\mathfrak p}$
which is what we wanted to prove.
\end{proof}

\begin{lemma}
\label{lemma-limit-preserving-flat-at-primes-over}
Let $R \to S$ be a ring map of finite presentation.
Let $M$ be an $S$-module of finite presentation.
Let $R' = \colim_{\lambda \in \Lambda} R_\lambda$
be a directed colimit of $R$-algebras. Let $I_\lambda \subset R_\lambda$
be ideals such that $I_\lambda R_\mu \subset I_\mu$ for all $\mu \geq \lambda$
and set $I' = \colim_\lambda I_\lambda$.
If (\ref{equation-flat-at-primes-over}) holds for
$(R' \to S \otimes_R R', I', M \otimes_R R')$, then there exists
a $\lambda \in \Lambda$ such that
(\ref{equation-flat-at-primes-over}) holds for
$(R_\lambda \to S \otimes_R R_\lambda, I_\lambda, M \otimes_R R_\lambda)$.
\end{lemma}

\begin{proof}
We are going to write $S_\lambda = S \otimes_R R_\lambda$,
$S' = S \otimes_R R'$, $M_\lambda = M \otimes_R R_\lambda$, and
$M' = M \otimes_R R'$.
The base change $S'$ is of finite presentation over $R'$ and
$M'$ is of finite presentation over $S'$ and similarly for the
versions with subscript $\lambda$, see
Algebra, Lemma \ref{algebra-lemma-base-change-finiteness}.
By
Algebra, Theorem \ref{algebra-theorem-openness-flatness}
the set
$$
U' = \{\mathfrak q' \in \Spec(S') \mid
M'_{\mathfrak q'}\text{ is flat over }R'\}
$$
is open in $\Spec(S')$. Note that $V(I'S')$ is a quasi-compact space
which is contained in $U'$ by assumption. Hence there exist finitely many
$g'_j \in S'$, $j = 1, \ldots, m$ such that $D(g'_j) \subset U'$ and such
that $V(I'S') \subset \bigcup D(g'_j)$.
Note that in particular $(M')_{g'_j}$ is a flat module over $R'$.

\medskip\noindent
We are going to pick increasingly large elements $\lambda \in \Lambda$.
First we pick it large enough so that we can find
$g_{j, \lambda} \in S_{\lambda}$ mapping to $g'_j$.
The inclusion $V(I'S') \subset \bigcup D(g'_j)$ means that
$I'S' + (g'_1, \ldots, g'_m) = S'$ which can be expressed as
$1 = \sum z_sh_s + \sum f_jg'_j$ for some $z_s \in I'$, $h_s, f_j \in S'$.
After increasing $\lambda$ we may assume such an equation holds in
$S_\lambda$. Hence we may assume that
$V(I_\lambda S_\lambda) \subset \bigcup D(g_{j, \lambda})$. By
Algebra, Lemma \ref{algebra-lemma-flat-finite-presentation-limit-flat}
we see that for some sufficiently large $\lambda$ the modules
$(M_\lambda)_{g_{j, \lambda}}$ are flat over $R_\lambda$.
In particular the module $M_\lambda$ is flat over $R_\lambda$
at all the primes lying over the ideal $I_\lambda$.
\end{proof}









\section{Flattening over a closed subsets of source and base}
\label{section-flattening-local-source-base}

\noindent
In this section we slightly generalize the discussion in
Section \ref{section-flattening-local-base}.
We strongly suggest the reader first read and understand
that section.

\begin{situation}
\label{situation-flattening-general}
Let $R \to S$ be a ring map. Let $J \subset S$ be an ideal.
Let $M$ be an $S$-module.
\end{situation}

\noindent
In this situation, given an $R$-algebra $R'$ and an ideal $I' \subset R'$
we set $S' = S \otimes_R R'$ and $M' = M \otimes_R R'$.
We will consider the condition
\begin{equation}
\label{equation-flat-at-primes}
\forall \mathfrak q' \in V(I'S' + JS') \subset \Spec(S') :
M'_{\mathfrak q'}\text{ is flat over }R'.
\end{equation}
Geometrically, this means that $M'$ is flat over $R'$ along the intersection
of the inverse image of $V(I')$ with the inverse image of $V(J)$.
Since $(R \to S, J, M)$ are fixed, condition (\ref{equation-flat-at-primes})
only depends on the pair $(R', I')$ where $R'$ is viewed as an $R$-algebra.

\begin{lemma}
\label{lemma-base-change-flat-at-primes}
In Situation \ref{situation-flattening-general}
let $R' \to R''$ be an $R$-algebra map.
Let $I' \subset R'$ and $I'R'' \subset I'' \subset R''$ be ideals.
If (\ref{equation-flat-at-primes}) holds for
$(R', I')$, then (\ref{equation-flat-at-primes})
holds for $(R'', I'')$.
\end{lemma}

\begin{proof}
Assume (\ref{equation-flat-at-primes}) holds for $(R', I')$.
Let $I''S'' + JS'' \subset \mathfrak q''$ be a prime of $S''$.
Let $\mathfrak q' \subset S'$ be the corresponding prime of $S'$.
Then both $I'S' \subset \mathfrak q'$ and $JS' \subset \mathfrak q'$
because the corresponding conditions hold for $\mathfrak q''$.
Note that $(M'')_{\mathfrak q''}$
is a localization of the base change $M'_{\mathfrak q'} \otimes_R R''$.
Hence $(M'')_{\mathfrak q''}$ is flat over $R''$ as a localization
of a flat module, see
Algebra, Lemmas \ref{algebra-lemma-flat-base-change} and
\ref{algebra-lemma-flat-localization}.
\end{proof}

\begin{lemma}
\label{lemma-flat-descent-flat-at-primes}
In Situation \ref{situation-flattening-general}
let $R' \to R''$ be an $R$-algebra map.
Let $I' \subset R'$ and $I'R'' \subset I'' \subset R''$ be ideals.
Assume
\begin{enumerate}
\item the map $V(I'') \to V(I')$ induced by
$\Spec(R'') \to \Spec(R')$ is surjective, and
\item $R''_{\mathfrak p''}$ is flat over $R'$ for all primes
$\mathfrak p'' \in V(I'')$.
\end{enumerate}
If (\ref{equation-flat-at-primes}) holds for
$(R'', I'')$, then (\ref{equation-flat-at-primes}) holds for $(R', I')$.
\end{lemma}

\begin{proof}
Assume (\ref{equation-flat-at-primes}) holds for $(R'', I'')$. Pick a prime
$I'S' + JS' \subset \mathfrak q' \subset S'$. Let
$I' \subset \mathfrak p' \subset R'$ be the corresponding prime of $R'$.
By assumption there exists a prime $\mathfrak p'' \in V(I'')$ of $R''$
lying over $\mathfrak p'$ and $R'_{\mathfrak p'} \to R''_{\mathfrak p''}$
is flat. Choose a prime
$\overline{\mathfrak q}'' \subset
\kappa(\mathfrak q') \otimes_{\kappa(\mathfrak p')} \kappa(\mathfrak p'')$.
This corresponds to a prime $\mathfrak q'' \subset S'' = S' \otimes_{R'} R''$
which lies over $\mathfrak q'$ and over $\mathfrak p''$. In particular
we see that $I''S'' \subset \mathfrak q''$ and that
$JS'' \subset \mathfrak q''$. Note that
$(S' \otimes_{R'} R'')_{\mathfrak q''}$ is a localization of
$S'_{\mathfrak q'} \otimes_{R'_{\mathfrak p'}} R''_{\mathfrak p''}$.
By assumption the module $(M' \otimes_{R'} R'')_{\mathfrak q''}$ is
flat over $R''_{\mathfrak p''}$. Hence
Algebra, Lemma \ref{algebra-lemma-base-change-flat-up-down}
implies that $M'_{\mathfrak q'}$ is flat over $R'_{\mathfrak p'}$
which is what we wanted to prove.
\end{proof}

\begin{lemma}
\label{lemma-limit-preserving-flat-at-primes}
In Situation \ref{situation-flattening-general}
assume $R \to S$ is essentially of finite presentation
and $M$ is an $S$-module of finite presentation. Let
$R' = \colim_{\lambda \in \Lambda} R_\lambda$
be a directed colimit of $R$-algebras. Let $I_\lambda \subset R_\lambda$
be ideals such that $I_\lambda R_\mu \subset I_\mu$ for all
$\mu \geq \lambda$ and set $I' = \colim_\lambda I_\lambda$.
If (\ref{equation-flat-at-primes}) holds for
$(R', I')$, then there exists a $\lambda \in \Lambda$ such that
(\ref{equation-flat-at-primes}) holds for $(R_\lambda, I_\lambda)$.
\end{lemma}

\begin{proof}
We first prove the lemma in case $R \to S$ is of finite presentation
and then we explain what needs to be changed in the general case.
We are going to write $S_\lambda = S \otimes_R R_\lambda$,
$S' = S \otimes_R R'$, $M_\lambda = M \otimes_R R_\lambda$, and
$M' = M \otimes_R R'$.
The base change $S'$ is of finite presentation over $R'$ and
$M'$ is of finite presentation over $S'$ and similarly for the
versions with subscript $\lambda$, see
Algebra, Lemma \ref{algebra-lemma-base-change-finiteness}.
By
Algebra, Theorem \ref{algebra-theorem-openness-flatness}
the set
$$
U' = \{\mathfrak q' \in \Spec(S') \mid
M'_{\mathfrak q'}\text{ is flat over }R'\}
$$
is open in $\Spec(S')$. Note that $V(I'S' + JS')$
is a quasi-compact space which is contained in $U'$ by assumption.
Hence there exist finitely many $g'_j \in S'$, $j = 1, \ldots, m$
such that $D(g'_j) \subset U'$ and such
that $V(I'S' + JS') \subset \bigcup D(g'_j)$.
Note that in particular $(M')_{g'_j}$ is a flat module over $R'$.

\medskip\noindent
We are going to pick increasingly large elements $\lambda \in \Lambda$.
First we pick it large enough so that we can find
$g_{j, \lambda} \in S_{\lambda}$ mapping to $g'_j$.
The inclusion $V(I'S' + JS') \subset \bigcup D(g'_j)$ means that
$I'S' + JS' + (g'_1, \ldots, g'_m) = S'$ which can be expressed as
$$
1 = \sum y_tk_t + \sum z_sh_s + \sum f_jg'_j
$$
for some $z_s \in I'$, $y_t \in J$, $k_t, h_s, f_j \in S'$.
After increasing $\lambda$ we may assume such an equation holds in
$S_\lambda$. Hence we may assume that
$V(I_\lambda S_\lambda + J S_\lambda) \subset \bigcup D(g_{j, \lambda})$. By
Algebra, Lemma \ref{algebra-lemma-flat-finite-presentation-limit-flat}
we see that for some sufficiently large $\lambda$ the modules
$(M_\lambda)_{g_{j, \lambda}}$ are flat over $R_\lambda$.
In particular the module $M_\lambda$ is flat over $R_\lambda$
at all the primes corresponding to points of
$V(I_\lambda S_\lambda + J S_\lambda)$.

\medskip\noindent
In the case that $S$ is essentially of finite presentation, we can write
$S = \Sigma^{-1}C$ where $R \to C$ is of finite presentation and
$\Sigma \subset C$ is a multiplicative subset. We can also write
$M = \Sigma^{-1}N$ for some finitely presented $C$-module $N$, see
Algebra, Lemma \ref{algebra-lemma-construct-fp-module}.
At this point we introduce $C_\lambda$, $C'$, $N_\lambda$, $N'$. Then in
the discussion above we obtain an open $U' \subset \Spec(C')$
over which $N'$ is flat over $R'$. The assumption that
(\ref{equation-flat-at-primes}) is true means that $V(I'S' + JS')$ maps
into $U'$, because for a prime $\mathfrak q' \subset S'$, corresponding
to a prime $\mathfrak r' \subset C'$ we have
$M'_{\mathfrak q'} = N'_{\mathfrak r'}$. Thus we can find
$g'_j \in C'$ such that $\bigcup D(g'_j)$ contains the image of
$V(I'S' + JS')$. The rest of the proof is exactly the same as before.
\end{proof}

\begin{lemma}
\label{lemma-flat-module-powers-variant}
In Situation \ref{situation-flattening-general}.
Let $I \subset R$ be an ideal. Assume
\begin{enumerate}
\item $R$ is a Noetherian ring,
\item $S$ is a Noetherian ring,
\item $M$ is a finite $S$-module, and
\item for each $n \geq 1$ and any prime
$\mathfrak q \in V(J + IS)$ the module $(M/I^n M)_{\mathfrak q}$
is flat over $R/I^n$.
\end{enumerate}
Then (\ref{equation-flat-at-primes}) holds for $(R, I)$, i.e.,
for every prime $\mathfrak q \in V(J + IS)$
the localization $M_{\mathfrak q}$ is flat over $R$.
\end{lemma}

\begin{proof}
Let $\mathfrak q \in V(J + IS)$. Then
Algebra, Lemma \ref{algebra-lemma-flat-module-powers}
applied to $R \to S_{\mathfrak q}$ and $M_{\mathfrak q}$
implies that $M_{\mathfrak q}$ is flat over $R$.
\end{proof}










\section{Flattening over a Noetherian complete local ring}
\label{section-flattening-Noetherian-complete-local}

\noindent
The following three lemmas give a completely algebraic proof of the existence
of the ``local'' flattening stratification when the base is a complete
local Noetherian ring $R$ and the given module is finite over a finite
type $R$-algebra $S$.

\begin{lemma}
\label{lemma-flattening-complete-local-noetherian}
Let $R \to S$ be a ring map.
Let $M$ be an $S$-module.
Assume
\begin{enumerate}
\item $(R, \mathfrak m)$ is a complete local Noetherian ring,
\item $S$ is a Noetherian ring, and
\item $M$ is finite over $S$.
\end{enumerate}
Then there exists an ideal $I \subset \mathfrak m$ such that
\begin{enumerate}
\item $(M/IM)_{\mathfrak q}$ is flat over $R/I$ for all
primes $\mathfrak q$ of $S/IS$ lying over $\mathfrak m$, and
\item if $J \subset R$ is an ideal such that $(M/JM)_{\mathfrak q}$
is flat over $R/J$ for all primes $\mathfrak q$ lying over
$\mathfrak m$, then $I \subset J$.
\end{enumerate}
In other words, $I$ is the smallest ideal of $R$ such that
(\ref{equation-flat-at-primes-over}) holds for
$(\overline{R} \to \overline{S}, \overline{\mathfrak m}, \overline{M})$
where $\overline{R} = R/I$, $\overline{S} = S/IS$,
$\overline{\mathfrak m} = \mathfrak m/I$ and $\overline{M} = M/IM$.
\end{lemma}

\begin{proof}
Let $J \subset R$ be an ideal. Apply
Algebra, Lemma \ref{algebra-lemma-flat-module-powers}
to the module $M/JM$ over the ring $R/J$.
Then we see that $(M/JM)_{\mathfrak q}$
is flat over $R/J$ for all primes $\mathfrak q$ of $S/JS$
if and only if $M/(J + \mathfrak m^n)M$ is flat over
$R/(J + \mathfrak m^n)$ for all $n \geq 1$.
We will use this remark below.

\medskip\noindent
For every $n \geq 1$ the local ring $R/\mathfrak m^n$ is Artinian.
Hence, by
Lemma \ref{lemma-flattening-artinian}
there exists a smallest ideal $I_n \supset \mathfrak m^n$ such that
$M/I_nM$ is flat over $R/I_n$. It is clear that $I_{n + 1} + \mathfrak m^n$
is contains $I_n$ and applying
Lemma \ref{lemma-intersection-flat}
we see that $I_n = I_{n + 1} + \mathfrak m^n$. Since
$R = \lim_n\ R/\mathfrak m^n$ we see that $I = \lim_n\ I_n/\mathfrak m^n$
is an ideal in $R$ such that $I_n = I + \mathfrak m^n$ for all $n \geq 1$.
By the initial remarks of the proof we see that $I$ verifies (1)
and (2). Some details omitted.
\end{proof}

\begin{lemma}
\label{lemma-flattening-complete-local-noetherian-property-by-finite-type}
With notation $R \to S$, $M$, and $I$ and assumptions as in
Lemma \ref{lemma-flattening-complete-local-noetherian}.
Consider a local homomorphism of local rings
$\varphi : (R, \mathfrak m) \to (R', \mathfrak m')$
such that $R'$ is Noetherian. Then the following are equivalent
\begin{enumerate}
\item condition (\ref{equation-flat-at-primes-over}) holds
for $(R' \to S \otimes_R R', \mathfrak m', M \otimes_R R')$, and
\item $\varphi(I) = 0$.
\end{enumerate}
\end{lemma}

\begin{proof}
The implication (2) $\Rightarrow$ (1) follows from
Lemma \ref{lemma-base-change-flat-at-primes-over}.
Let $\varphi : R \to R'$ be as in the lemma satisfying (1).
We have to show that $\varphi(I) = 0$.
This is equivalent to the condition that $\varphi(I)R' = 0$.
By Artin-Rees in the Noetherian local ring $R'$ (see
Algebra, Lemma \ref{algebra-lemma-intersect-powers-ideal-module-zero})
this is equivalent to the condition that
$\varphi(I)R' + (\mathfrak m')^n = (\mathfrak m')^n$ for all $n > 0$.
Hence this is equivalent to the condition that the composition
$\varphi_n : R \to R' \to R'/(\mathfrak m')^n$ annihilates $I$ for each $n$.
Now assumption (1) for $\varphi$ implies assumption (1) for
$\varphi_n$ by
Lemma \ref{lemma-base-change-flat-at-primes-over}.
This reduces us to the case where $R'$ is Artinian local.

\medskip\noindent
Assume $R'$ Artinian. Let $J = \Ker(\varphi)$. We have to show that
$I \subset J$. By the construction of $I$ in
Lemma \ref{lemma-flattening-complete-local-noetherian}
it suffices to show that $(M/JM)_{\mathfrak q}$ is flat over $R/J$
for every prime $\mathfrak q$ of $S/JS$ lying over $\mathfrak m$.
As $R'$ is Artinian, condition (1) signifies that $M \otimes_R R'$
is flat over $R'$. As $R'$ is Artinian and $R/J \to R'$ is a local
injective ring map, it follows that $R/J$ is Artinian
too. Hence the flatness of $M \otimes_R R' = M/JM \otimes_{R/J} R'$ over
$R'$ implies that $M/JM$ is flat over $R/J$ by
Algebra,
Lemma \ref{algebra-lemma-descent-flatness-injective-map-artinian-rings}.
This concludes the proof.
\end{proof}

\begin{lemma}
\label{lemma-flattening-complete-local-universal-property}
With notation $R \to S$, $M$, and $I$ and assumptions as in
Lemma \ref{lemma-flattening-complete-local-noetherian}.
In addition assume that $R \to S$ is of finite type.
Then for any local homomorphism of local rings
$\varphi : (R, \mathfrak m) \to (R', \mathfrak m')$
the following are equivalent
\begin{enumerate}
\item condition (\ref{equation-flat-at-primes-over}) holds
for $(R' \to S \otimes_R R', \mathfrak m', M \otimes_R R')$, and
\item $\varphi(I) = 0$.
\end{enumerate}
\end{lemma}

\begin{proof}
The implication (2) $\Rightarrow$ (1) follows from
Lemma \ref{lemma-base-change-flat-at-primes-over}.
Let $\varphi : R \to R'$ be as in the lemma satisfying (1).
As $R$ is Noetherian we see that $R \to S$ is of finite presentation
and $M$ is an $S$-module of finite presentation.
Write $R' = \colim_\lambda R_\lambda$
as a directed colimit of local $R$-subalgebras $R_\lambda \subset R'$,
with maximal ideals $\mathfrak m_\lambda = R_\lambda \cap \mathfrak m'$
such that each $R_\lambda$ is essentially of finite type over $R$. By
Lemma \ref{lemma-limit-preserving-flat-at-primes-over}
we see that condition (\ref{equation-flat-at-primes-over}) holds for
$(R_\lambda \to S \otimes_R R_\lambda, \mathfrak m_\lambda,
M \otimes_R R_\lambda)$ for some $\lambda$. Hence
Lemma \ref{lemma-flattening-complete-local-noetherian-property-by-finite-type}
applies to the ring map $R \to R_\lambda$ and we see that
$I$ maps to zero in $R_\lambda$, a fortiori it maps to zero in $R'$.
\end{proof}




\section{Descent of flatness along integral maps}
\label{section-descent-flatness-integral}

\noindent
First a few simple lemmas.

\begin{lemma}
\label{lemma-have-one-root}
Let $R$ be a ring. Let $P(T)$ be a monic polynomial with coefficients
in $R$. Let $\alpha \in R$ be such that $P(\alpha) = 0$. Then
$P(T) = (T - \alpha)Q(T)$ for some monic polynomial $Q(T) \in R[T]$.
\end{lemma}

\begin{proof}
By induction on the degree of $P$. If $\deg(P) = 1$, then
$P(T) = T - \alpha$ and the result is true. If $\deg(P) > 1$, then
we can write $P(T) = (T - \alpha)Q(T) + r$ for some polynomial
$Q \in R[T]$ of degree $< \deg(P)$ and some $r \in R$ by long
division. By assumption $0 = P(\alpha) = (\alpha - \alpha)Q(\alpha) + r = r$
and we conclude that $r = 0$ as desired.
\end{proof}

\begin{lemma}
\label{lemma-adjoin-one-root}
Let $R$ be a ring. Let $P(T)$ be a monic polynomial with coefficients
in $R$. There exists a finite free ring map $R \to R'$ such that
$P(T) = (T - \alpha)Q(T)$ for some $\alpha \in R'$ and some
monic polynomial $Q(T) \in R'[T]$.
\end{lemma}

\begin{proof}
Write $P(T) = T^d + a_1T^{d - 1} + \ldots + a_0$.
Set $R' = R[x]/(x^d + a_1x^{d - 1} + \ldots + a_0)$.
Set $\alpha$ equal to the congruence class of $x$.
Then it is clear that $P(\alpha) = 0$. Thus we win by
Lemma \ref{lemma-have-one-root}.
\end{proof}

\begin{lemma}
\label{lemma-finite-split}
Let $R \to S$ be a finite ring map.
There exists a finite free ring extension $R \subset R'$ such
that $S \otimes_R R'$ is a quotient of a ring of the form
$$
R'[T_1, \ldots, T_n]/(P_1(T_1), \ldots, P_n(T_n))
$$
with $P_i(T) = \prod_{j = 1, \ldots, d_i} (T - \alpha_{ij})$ for some
$\alpha_{ij} \in R'$.
\end{lemma}

\begin{proof}
Let $x_1, \ldots, x_n \in S$ be generators of $S$ over $R$.
For each $i$ we can choose a monic polynomial $P_i(T) \in R[T]$
such that $P_i(x_i) = 0$ in $S$, see
Algebra, Lemma \ref{algebra-lemma-finite-is-integral}.
Say $\deg(P_i) = d_i$. By
Lemma \ref{lemma-adjoin-one-root}
(applied $\sum d_i$ times) there exists a finite free ring
extension $R \subset R'$ such that each $P_i$ splits completely:
$$
P_i(T) = \prod\nolimits_{j = 1, \ldots, d_i} (T - \alpha_{ij})
$$
for certain $\alpha_{ik} \in R'$. Let
$R'[T_1, \ldots, T_n] \to S \otimes_R R'$ be the $R'$-algebra map
which maps $T_i$ to $x_i \otimes 1$. As this maps $P_i(T_i)$ to zero,
this induces the desired surjection.
\end{proof}

\begin{lemma}
\label{lemma-split-image}
Let $R$ be a ring.
Let $S = R[T_1, \ldots, T_n]/J$.
Assume $J$ contains elements of the form $P_i(T_i)$
with $P_i(T) = \prod_{j = 1, \ldots, d_i} (T - \alpha_{ij})$ for some
$\alpha_{ij} \in R$. For $\underline{k} = (k_1, \ldots, k_n)$
with $1 \leq k_i \leq d_i$ consider the ring map
$$
\Phi_{\underline{k}} : R[T_1, \ldots, T_n] \to R,
\quad
T_i \longmapsto \alpha_{ik_i}
$$
Set $J_{\underline{k}} = \Phi_{\underline{k}}(J)$.
Then the image of $\Spec(S) \to \Spec(R)$ is equal to
$V(\bigcap J_{\underline{k}})$.
\end{lemma}

\begin{proof}
This lemma proves itself. Hint:
$V(\bigcap J_{\underline{k}}) = \bigcup V(J_{\underline{k}})$.
\end{proof}

\noindent
The following result is due to Ferrand, see \cite{Ferrand}.

\begin{lemma}
\label{lemma-descent-flatness-injective-finite-Noetherian-rings}
Let $R \to S$ be a finite injective homomorphism of Noetherian rings.
Let $M$ be an $R$-module. If $M \otimes_R S$ is a flat $S$-module,
then $M$ is a flat $R$-module.
\end{lemma}

\begin{proof}
Let $M$ be an $R$-module such that $M \otimes_R S$ is flat over $S$. By
Algebra, Lemma \ref{algebra-lemma-flatness-descends}
in order to prove that $M$ is flat we may replace $R$ by any faithfully flat
ring extension. By
Lemma \ref{lemma-finite-split}
we can find a finite locally free ring extension $R \subset R'$ such
that $S' = S \otimes_R R' = R'[T_1, \ldots, T_n]/J$ for some ideal
$J \subset R'[T_1, \ldots, T_n]$ which contains the  elements of the form
$P_i(T_i)$ with $P_i(T) = \prod_{j = 1, \ldots, d_i} (T - \alpha_{ij})$
for some $\alpha_{ij} \in R'$. Note that $R'$ is Noetherian
and that $R' \subset S'$ is a finite extension of rings. Hence we may
replace $R$ by $R'$ and assume that $S$ has a presentation as in
Lemma \ref{lemma-split-image}.
Note that $\Spec(S) \to \Spec(R)$ is surjective, see
Algebra, Lemma \ref{algebra-lemma-integral-overring-surjective}.
Thus, using
Lemma \ref{lemma-split-image}
we conclude that $I = \bigcap J_{\underline{k}}$ is an ideal
such that $V(I) = \Spec(R)$. This means that
$I \subset \sqrt{(0)}$, and since $R$ is Noetherian that $I$
is nilpotent. The maps $\Phi_{\underline{k}}$ induce commutative
diagrams
$$
\xymatrix{
S \ar[rr] & & R/J_{\underline{k}} \\
& R \ar[lu] \ar[ru]
}
$$
from which we conclude that $M/J_{\underline{k}}M$ is flat over
$R/J_{\underline{k}}$. By
Lemma \ref{lemma-intersection-flat}
we see that $M/IM$ is flat over $R/I$. Finally, applying
Algebra, Lemma \ref{algebra-lemma-lift-flatness}
we conclude that $M$ is flat over $R$.
\end{proof}

\begin{lemma}
\label{lemma-descent-flatness-injective-integral}
Let $R \to S$ be an injective integral ring map.
Let $M$ be a finitely presented module over $R[x_1, \ldots, x_n]$.
If $M \otimes_R S$ is flat over $S$, then $M$ is flat over $R$.
\end{lemma}

\begin{proof}
Choose a presentation
$$
R[x_1, \ldots, x_n]^{\oplus t} \to R[x_1, \ldots, x_n]^{\oplus r} \to M \to 0.
$$
Let's say that the first map is given by the $r \times t$-matrix
$T = (f_{ij})$ with $f_{ij} \in R[x_1, \ldots, x_n]$. Write
$f_{ij} = \sum f_{ij, I} x^I$ with $f_{ij, I} \in R$ (multi-index notation).
Consider diagrams
$$
\xymatrix{
R \ar[r] & S \\
R_\lambda \ar[u] \ar[r] & S_\lambda \ar[u]
}
$$
where $R_\lambda$ is a finitely generated $\mathbf{Z}$-subalgebra of
$R$ containing all $f_{ij, I}$ and $S_\lambda$ is a finite
$R_\lambda$-subalgebra of $S$. Let $M_\lambda$ be the finite
$R_\lambda[x_1, \ldots, x_n]$-module defined by a presentation
as above, using the same matrix $T$ but now viewed as a matrix
over $R_\lambda[x_1, \ldots, x_n]$. Note that $S$ is the directed colimit
of the $S_\lambda$ (details omitted). By
Algebra, Lemma \ref{algebra-lemma-flat-finite-presentation-limit-flat}
we see that for some $\lambda$ the module
$M_\lambda \otimes_{R_\lambda} S_\lambda$ is flat over $S_\lambda$. By
Lemma \ref{lemma-descent-flatness-injective-finite-Noetherian-rings}
we conclude that $M_\lambda$ is flat over $R_\lambda$. Since
$M = M_\lambda \otimes_{R_\lambda} R$ we win by
Algebra, Lemma \ref{algebra-lemma-flat-base-change}.
\end{proof}

\begin{lemma}
\label{lemma-descent-projective-injective-finite-Noetherian-rings}
Let $R \to S$ be a finite injective homomorphism of Noetherian rings.
Let $P$ be an $R$-module. If $P \otimes_R S$ is a projective $S$-module,
then $P$ is a projective $R$-module.
\end{lemma}

\begin{proof}
Let $P$ be an $R$-module such that $P \otimes_R S$ is projective over $S$. By
Algebra, Theorem \ref{algebra-theorem-ffdescent-projectivity}
in order to prove that $P$ is projective we may replace $R$
by any faithfully flat ring extension. By
Lemma \ref{lemma-finite-split}
we can find a finite locally free ring extension $R \subset R'$ such
that $S' = S \otimes_R R' = R'[T_1, \ldots, T_n]/J$ for some ideal
$J \subset R'[T_1, \ldots, T_n]$ which contains the  elements of the form
$P_i(T_i)$ with $P_i(T) = \prod_{j = 1, \ldots, d_i} (T - \alpha_{ij})$
for some $\alpha_{ij} \in R'$. Note that $R'$ is Noetherian
and that $R' \subset S'$ is a finite extension of rings. Hence we may
replace $R$ by $R'$ and assume that $S$ has a presentation as in
Lemma \ref{lemma-split-image}.
Note that $\Spec(S) \to \Spec(R)$ is surjective, see
Algebra, Lemma \ref{algebra-lemma-integral-overring-surjective}.
Thus, using Lemma \ref{lemma-split-image}
we conclude that $I = \bigcap J_{\underline{k}}$ is an ideal
such that $V(I) = \Spec(R)$. This means that
$I \subset \sqrt{(0)}$, and since $R$ is Noetherian that $I$
is nilpotent. The maps $\Phi_{\underline{k}}$ induce commutative
diagrams
$$
\xymatrix{
S \ar[rr] & & R/J_{\underline{k}} \\
& R \ar[lu] \ar[ru]
}
$$
from which we conclude that $P/J_{\underline{k}}P$ is projective over
$R/J_{\underline{k}}$. By
Algebra, Lemma \ref{algebra-lemma-projective-modulo-two-ideals}
we see that $P/IP$ is projective over $R/I$. Since $P$ is flat
over $R$ by
Lemma \ref{lemma-descent-flatness-injective-finite-Noetherian-rings},
we may apply Algebra, Lemma \ref{algebra-lemma-lift-projective} to
conclude that $P$ is projective over $R$.
\end{proof}









\section{Torsion free modules}
\label{section-torsion-flat}

\noindent
In this section we discuss torsion free modules and the relationship
with flatness (especially over dimension 1 rings).

\begin{definition}
\label{definition-torsion}
Let $R$ be a domain. Let $M$ be an $R$-module.
\begin{enumerate}
\item We say an element $x \in M$ is {\it torsion} if there exists
a nonzero $f \in R$ such that $fx = 0$.
\item We say $M$ is {\it torsion free} if the only torsion element of $M$
is $0$.
\item We say $M$ is a {\it torsion module} if every element of $M$
is torsion.
\end{enumerate}
\end{definition}

\noindent
Let $R$ be a domain and let $S = R \setminus \{0\}$ be the multiplicative
set of nonzero elements of $R$. Then an $R$-module $M$ is torsion free if
and only if $M \to S^{-1}M$ is injective. In other words, if and only if
the map $M \to M \otimes_R K$ is injective where $K = S^{-1}R$ is the fraction
field of $R$.

\begin{lemma}
\label{lemma-torsion}
Let $R$ be a domain. Let $M$ be an $R$-module. The set of torsion elements
$M_{tors}$ of $M$ is the kernel of the map $M \to M \otimes_R K$.
Thus $M_{tors}$ is an $R$-submodule of $M$.
The quotient module $M/M_{tors}$ is torsion free.
\end{lemma}

\begin{proof}
See discussion above.
\end{proof}

\begin{lemma}
\label{lemma-localize-torsion}
Let $R$ be a domain. Let $M$ be a torsion free $R$-module.
For any multiplicative set $S \subset R$ the module
$S^{-1}M$ is a torsion free $S^{-1}R$-module.
\end{lemma}

\begin{proof}
Omitted.
\end{proof}

\begin{lemma}
\label{lemma-flat-pullback-torsion}
Let $R \to R'$ be a flat homomorphism of domains. If $M$ is a torsion
free $R$-module, then $M \otimes_R R'$ is a torsion free $R'$-module.
\end{lemma}

\begin{proof}
If $M$ is torsion free, then $M \subset M \otimes_R K$ is injective
where $K$ is the fraction field of $R$. Since $R'$ is flat over $R$
we see that $M \otimes_R R' \to (M \otimes_R K) \otimes_R R'$ is injective.
Since $M \otimes_R K$ is isomorphic to a direct sum of copies of $K$,
it suffices to see that $K \otimes_R R'$ is torsion free. This is true
because it is a localization of $R'$.
\end{proof}

\begin{lemma}
\label{lemma-extension-torsion-free}
Let $R$ be a domain. Let $0 \to M \to M' \to M'' \to 0$
be a short exact sequence of $R$-modules. If $M$ and $M''$
are torsion free, then $M'$ is torsion free.
\end{lemma}

\begin{proof}
Omitted.
\end{proof}

\begin{lemma}
\label{lemma-check-torsion}
Let $R$ be a domain. Let $M$ be an $R$-module.
Then $M$ is torsion free if and only if $M_\mathfrak m$ is a
torsion free $R_\mathfrak m$-module for all maximal ideals
$\mathfrak m$ of $R$.
\end{lemma}

\begin{proof}
Omitted. Hint: Use Lemma \ref{lemma-localize-torsion} and
Algebra, Lemma \ref{algebra-lemma-characterize-zero-local}.
\end{proof}

\begin{lemma}
\label{lemma-finite-torsion-free-submodule-free}
Let $R$ be a domain. Let $M$ be a finite $R$-module.
Then $M$ is torsion free if and only if $M$ is a
submodule of a finite free module.
\end{lemma}

\begin{proof}
If $M$ is a submodule of $R^{\oplus n}$, then $M$ is torsion free.
For the converse, assume $M$ is torsion free. Let $K$ be the
fraction field of $R$. Then $M \otimes_R K$ is a finite dimensional
$K$-vector space. Choose a basis $e_1, \ldots, e_r$ for this vector
space. Let $x_1, \ldots, x_n$ be generators of $M$. Write
$x_i = \sum (a_{ij}/b_{ij}) e_j$ for some $a_{ij}, b_{ij} \in R$
with $b_{ij} \not = 0$. Set $b = \prod_{i, j} b_{ij}$.
Since $M$ is torsion free the map
$M \to M \otimes_R K$ is injective and the image is contained
in $R^{\oplus r} = R e_1/b \oplus \ldots \oplus Re_r/b$.
\end{proof}

\begin{lemma}
\label{lemma-torsion-free-finite-noetherian-domain}
Let $R$ be a Noetherian domain. Let $M$ be a nonzero finite $R$-module.
The following are equivalent
\begin{enumerate}
\item $M$ is torsion free,
\item $M$ is a submodule of a finite free module,
\item $(0)$ is the only associated prime of $M$,
\item $(0)$ is in the support of $M$ and $M$ has property $(S_1)$, and
\item $(0)$ is in the support of $M$ and $M$ has no embedded associated prime.
\end{enumerate}
\end{lemma}

\begin{proof}
We have seen the equivalence of (1) and (2) in
Lemma \ref{lemma-finite-torsion-free-submodule-free}.
We have seen the equivalence of (4) and (5) in
Algebra, Lemma \ref{algebra-lemma-criterion-no-embedded-primes}.
The equivalence between (3) and (5) is immediate from the definition.
A localization of a torsion free module is torsion free
(Lemma \ref{lemma-localize-torsion}), hence it is clear that a
$M$ has no associated primes different from $(0)$. Thus (1)
implies (5). Conversely, assume (5). If $M$ has torsion,
then there exists an embedding $R/I \subset M$ for some nonzero
ideal $I$ of $R$. Hence $M$ has an associated prime different
from $(0)$
(see Algebra, Lemmas \ref{algebra-lemma-ass} and \ref{algebra-lemma-ass-zero}).
This is an embedded associated prime which contradicts the assumption.
\end{proof}

\begin{lemma}
\label{lemma-flat-torsion-free}
Let $R$ be a domain. Any flat $R$-module is torsion free.
\end{lemma}

\begin{proof}
If $x \in R$ is nonzero, then $x : R \to R$ is injective, and hence if $M$
is flat over $R$, then $x : M \to M$ is injective. Thus if $M$ is flat over
$R$, then $M$ is torsion free.
\end{proof}

\begin{lemma}
\label{lemma-valuation-ring-torsion-free-flat}
Let $A$ be a valuation ring.
An $A$-module $M$ is flat over $A$ if and only if $M$ is torsion free.
\end{lemma}

\begin{proof}
The implication ``flat $\Rightarrow$ torsion free'' is
Lemma \ref{lemma-flat-torsion-free}.
For the converse, assume $M$ is torsion free.
By the equational criterion of flatness (see
Algebra, Lemma \ref{algebra-lemma-flat-eq})
we have to show that every relation in $M$ is trivial. To do this assume that
$\sum_{i = 1, \ldots, n} a_i x_i = 0$ with $x_i \in M$ and $a_i \in A$.
After renumbering we may assume that $v(a_1) \leq v(a_i)$ for all $i$.
Hence we can write $a_i = a'_i a_1$ for some $a'_i \in A$. Note that
$a'_1 = 1$. As $M$ is torsion free we see that
$x_1 = - \sum_{i \geq 2} a'_i x_i$. Thus, if we choose
$y_i = x_i$, $i = 2, \ldots, n$ then
$$
x_1 = \sum\nolimits_{j \geq 2} -a'_j y_j, \quad
x_i = y_i, (i \geq 2)\quad
0 = a_1 \cdot (-a'_j) + a_j \cdot 1 (j \geq 2)
$$
shows that the relation was trivial (to be explicit the elements
$a_{ij}$ are defined by setting $a_{11} = 0$, $a_{1j} = -a'_j$
for $j > 1$, and $a_{ij} = \delta_{ij}$ for $i, j \geq 2$).
\end{proof}

\begin{lemma}
\label{lemma-dedekind-torsion-free-flat}
Let $A$ be a Dedekind domain (for example a discrete valuation ring
or more generally a PID).
\begin{enumerate}
\item An $A$-module is flat if and only if it is torsion free.
\item A finite torsion free $A$-module is finite locally free.
\item A finite torsion free $A$-module is finite free if
$A$ is a PID.
\end{enumerate}
\end{lemma}

\begin{proof}
(For the parenthetical remark in the statement of the lemma, see
Algebra, Lemma \ref{algebra-lemma-PID-dedekind}.) Proof of (1).
By Lemma \ref{lemma-check-torsion}
and Algebra, Lemma \ref{algebra-lemma-flat-localization}
it suffices to check the statement over $A_\mathfrak m$
for $\mathfrak m \subset A$ maximal. Since $A_\mathfrak m$
is a discrete valuation ring
(Algebra, Lemma \ref{algebra-lemma-characterize-Dedekind})
we win by Lemma \ref{lemma-valuation-ring-torsion-free-flat}.

\medskip\noindent
Proof of (2). Follows from
Algebra, Lemma \ref{algebra-lemma-finite-projective}
and (1).

\medskip\noindent
Proof of (3). Let $A$ be a PID and let $M$ be a finite
torsion free module. By Lemma \ref{lemma-finite-torsion-free-submodule-free}
we see that $M \subset A^{\oplus n}$ for some $n$. We argue that
$M$ is free by induction on $n$. The case $n = 1$ expresses exactly the
fact that $A$ is a PID. If $n > 1$ let $M' \subset A^{\oplus n - 1}$
be the image of the projection onto the last $n - 1$ summands of
$A^{\oplus n}$. Then we obtain a short exact sequence
$0 \to I \to M \to M' \to 0$ where $I$ is the intersection of $M$
with the first summand $A$ of $A^{\oplus n}$. By induction we
see that $M$ is an extension of finite free $A$-modules, whence
finite free.
\end{proof}

\begin{lemma}
\label{lemma-hom-into-torsion-free}
Let $R$ be a domain. Let $M$, $N$ be $R$-modules.
If $N$ is torsion free, so is $\Hom_R(M, N)$.
\end{lemma}

\begin{proof}
Choose a surjection $\bigoplus_{i \in I} R \to M$.
Then $\Hom_R(M, N) \subset \prod_{i \in I} N$.
\end{proof}




\section{Ranks of modules}
\label{section-rank}

\noindent
Here is our definition.

\begin{definition}
\label{definition-rank}
Let $R$ be a domain with fraction field $K$.
The {\it rank} of an $R$-module $M$ is $\dim_K(M \otimes_R K)$.
\end{definition}

\noindent
This definition does not conflict with the notion of a locally free module
of rank $r$ in Algebra, Definition \ref{algebra-definition-locally-free}
when both definitions apply.

\begin{lemma}
\label{lemma-torsion-does-not-matter-rank}
Let $R$ be a domain. If $M \to M'$ is a map of $R$-modules
whose kernel and cokernel are torsion, then the rank of $M$
equals the rank of $M'$.
\end{lemma}

\begin{proof}
Omitted. Hint: the induced map $M \otimes_R K \to M' \otimes_R K$
is an isomorphism if $K$ is the fraction field of $R$.
\end{proof}

\begin{lemma}
\label{lemma-rank-additive}
Let $R$ be a domain. Let $0 \to M \to M' \to M'' \to 0$ be a short 
exact sequence of $R$-modules. Then the rank of $M'$ is the
sum of the ranks of $M$ and $M''$.
\end{lemma}

\begin{proof}
Omitted.
\end{proof}

\begin{lemma}
\label{lemma-rank-tensor}
Let $R$ be a domain. Let $M$ and $N$ be $R$-modules.
\begin{enumerate}
\item The rank of $M \otimes_R N$ is the product of the
ranks of $M$ and $N$.
\item If $M$ is a finitely presented $R$-module, then
the rank of $\Hom_R(M, N)$ is the product of the
ranks of $M$ and $N$.
\end{enumerate}
\end{lemma}

\begin{proof}
Part (1) follows from the same fact for vector spaces.
Assume $M$ is finitely presented as an $R$-module.
Then $\Hom_R(M, N) \otimes_R K$ is isomorphic to
$\Hom_K(M \otimes_R K, N \otimes_R K)$ as a $K$-vector space by
Algebra, Lemma \ref{algebra-lemma-hom-from-finitely-presented}.
Since also $M \otimes_R K$ is finite dimensional, we conclude
from the corresponding fact for vector spaces.
\end{proof}

\begin{lemma}
\label{lemma-dominant-pullback-rank}
Let $R \subset R'$ be an extension of domains. If $M$ is an $R$-module,
then the rank of $M$ over $R$ is equal to the rank of $M \otimes_R R'$
over $R'$.
\end{lemma}

\begin{proof}
This is true because the dimension of a vector space is invariant
under extension of ground field.
\end{proof}

\begin{lemma}
\label{lemma-finite-module-rank}
Let $R$ be a domain. Let $M$ be a finite $R$-module. Then
\begin{enumerate}
\item $M$ has finite rank $r \geq 0$,
\item there is a map $M \to R^{\oplus r}$ whose kernel and cokernel
are torsion modules,
\item there is an $f \in R$, $f \not = 0$ such that $M_f$
is free of rank $r$, and
\item there is an injective map $R^{\oplus r} \to M$
whose cokernel is a torsion module.
\end{enumerate}
\end{lemma}

\begin{proof}
Since $M$ is finite, $M \otimes_R K$ is a finite dimensional
vector space over the fraction field $K$ of $R$. Thus the rank $r$ is finite.

\medskip\noindent
Choose generators $x_1, \ldots, x_n$ of $M$. Choose a basis
$e_1, \ldots, e_r$ for $M \otimes_R K$. We may write
$x_i \otimes 1 = \sum (a_{ij}/b_{ij}) e_j$ for  some
$a_{ij} , b_{ij} \in R$ with $b_{ij} \not = 0$. Setting
$b = \prod b_{ij}$ and replacing $a_{ij}$ by $a_{ij}b/b_{ij}$
we may assume that $x_i \otimes 1 = \sum (a_{ij}/b) e_j$.
Note that if we have a relation $\sum r_i x_i = 0$ in $M$
then $\sum r_ia_{ij} = 0$ in $R$ for all $j = 1, \ldots, r$.

\medskip\noindent
Proof of (2). The last comment above tells us the map
$M \to \bigoplus Re_j$ sending $x_i$ to $\sum a_{ij} e_j$.
is well defined. The fact that $e_1, \ldots, e_r$ forms a basis
of $M \otimes_R K$ tells us that this map becomes an
isomorphism upon tensoring by $K$. Hence this is a map as in (2).

\medskip\noindent
Proof of (3). The cokernel of the map $M \to R^{\oplus r}$
constructed in (2) is finite and torsion. Hence we can find
a $g \in R$, $g \not = 0$ which annihilates this cokernel.
In other words, such that the map $M_g \to R_g^{\oplus r}$
is surjective. Then we can choose a splitting
$M_g = N \oplus R_g^{\oplus r}$. Then $N$ is a finite
$R_g$-module with $N \otimes_{R_g} K = 0$. In other words,
$N$ is torsion and finite over $R_g$. Thus we can find
a $g' \in R$ which annihilates $N$. Setting $f= gg'$ we
see that (3) holds.

\medskip\noindent
To prove (4) choose a basis $y_1, \ldots, y_r$
of $M_f$ where $f$ is as in (3). Then write
$y_j = m_j/f^{t_j}$ for some $m_j \in M$ and
$t_j \geq 0$. The map $R^{\oplus r} \to M$ sending
the $j$th basis vector to $m_j$ will be as in (4).
Details omitted.
\end{proof}




\section{Reflexive modules}
\label{section-reflexive}

\noindent
Here is our definition.

\begin{definition}
\label{definition-reflexive}
Let $R$ be a domain. We say an $R$-module $M$ is {\it reflexive} if
the natural map
$$
j : M \longrightarrow \Hom_R(\Hom_R(M, R), R)
$$
which sends $m \in M$ to the map sending $\varphi \in \Hom_R(M, R)$
to $\varphi(m) \in R$ is an isomorphism.
\end{definition}

\noindent
We can make this definition for more general rings, but already the
definition above has drawbacks. It would be wise to restrict
to Noetherian domains and finite torsion free modules and (perhaps)
impose some regularity conditions on $R$ (e.g., $R$ is normal).

\begin{lemma}
\label{lemma-reflexive-torsion-free}
Let $R$ be a domain and let $M$ be an $R$-module.
\begin{enumerate}
\item If $M$ is reflexive, then $M$ is torsion free.
\item If $M$ is finite, then the kernel and cokernel of
$j : M \to \Hom_R(\Hom_R(M, R), R)$ are torsion modules.
\item If $M$ is finite, then $j$ is injective
if and only if $M$ is torsion free.
\end{enumerate}
\end{lemma}

\begin{proof}
To see (1) use that the double dual of $M$ is torsion free by
Lemma \ref{lemma-hom-into-torsion-free}. Assume that $M$ is finite.
Choose generators $x_1, \ldots, x_n$ of $M$ over $R$.
Let $K$ be the fraction field of $R$. After renumbering we
may assume that for some $0 \leq r \leq n$ the elements
$x_1, \ldots, x_r$ map to a basis of $M \otimes_R K$ over $K$.
For $r < i \leq n$ write
$x_i \otimes 1 = \sum_{1 \leq j \leq r} x_j \otimes k_{ij}$ with
$k_{ij} \in K$. Choose $c \in R$ such that $a_{ij} = ck_{ij} \in R$
and such that $cx_i =  \sum_{1 \leq j \leq r} a_{ij} x_j$ in $M$
(small detail omitted). Then we get maps
$$
R^{\oplus r} \xrightarrow{\alpha} M \xrightarrow{\beta} R^{\oplus r}
\xrightarrow{\alpha} M
$$
with both $\alpha \circ \beta$ and $\beta \circ \alpha$
given by multiplication with $c$. Here $\alpha$ uses
$x_1, \ldots, x_r$ and $\beta$ sends $x_i$ to $c$ times the $i$th
basis vector for $1 \leq i \leq r$ and
to the vector $(a_{i1}, \ldots, a_{ir})$ for $r < i \leq n$.
Consider the commutative diagram
$$
\xymatrix{
R^{\oplus r} \ar[r] \ar[d]^{\text{id}} &
M \ar[r] \ar[d]^j &
R^{\oplus r} \ar[d]^{\text{id}} \ar[r] &
M \ar[d]^j \\
R^{\oplus r} \ar[r] &
\Hom_R(\Hom_R(M, R), R) \ar[r] &
R^{\oplus r} \ar[r] &
\Hom_R(\Hom_R(M, R), R)
}
$$
Of course the compositions of consecutive arrows on the bottom row
are equal to multiplication by $c$. The diagram shows that the kernel
and cokernel of $j$ are annihilated by $c$ which proves (2).
For (3) note that if $M$ is torsion free, then
$\beta$ is injective which implies that $j$ is injective.
The converse is clear because we've seen already that the double
dual of $M$ is torsion free.
\end{proof}

\begin{lemma}
\label{lemma-cokernel-map-double-dual-dvr}
Let $R$ be a discrete valuation ring and let $M$ be a finite $R$-module.
Then the map $j : M \to \Hom_R(\Hom_R(M, R), R)$ is surjective.
\end{lemma}

\begin{proof}
Let $M_{tors} \subset M$ be the torsion submodule. Then we have
$\Hom_R(M, R) = \Hom_R(M/M_{tors}, R)$ (holds over any domain).
Hence we may assume that $M$ is torsion free. Then $M$ is free
by Lemma \ref{lemma-dedekind-torsion-free-flat} and the lemma is clear.
\end{proof}

\begin{lemma}
\label{lemma-check-reflexive}
Let $R$ be a Noetherian domain. Let $M$ be a finite $R$-module.
The following are equivalent:
\begin{enumerate}
\item $M$ is reflexive,
\item $M_\mathfrak p$ is a reflexive $R_\mathfrak p$-module
for all primes $\mathfrak p \subset R$, and
\item $M_\mathfrak m$ is a reflexive $R_\mathfrak m$-module
for all maximal ideals $\mathfrak m$ of $R$.
\end{enumerate}
\end{lemma}

\begin{proof}
The localization of $j : M \to \Hom_R(\Hom_R(M, R), R)$
at a prime $\mathfrak p$ is the corresponding map for the module
$M_\mathfrak p$ over the Noetherian local domain $R_\mathfrak p$.
See Algebra, Lemma \ref{algebra-lemma-hom-from-finitely-presented}.
Thus the lemma holds by
Algebra, Lemma \ref{algebra-lemma-characterize-zero-local}.
\end{proof}

\begin{lemma}
\label{lemma-sequence-reflexive}
Let $R$ be a Noetherian domain. Let $0 \to M \to M' \to M''$
an exact sequence of finite $R$-modules. If $M'$ is reflexive
and $M''$ is torsion free, then $M$ is reflexive.
\end{lemma}

\begin{proof}
We will use without further mention that $\Hom_R(N, N')$ is a finite
$R$-module for any finite $R$-modules $N$ and $N'$, see
Algebra, Lemma \ref{algebra-lemma-ext-noetherian}.
We take duals to get a sequence
$$
\Hom_R(M, R) \leftarrow \Hom_R(M', R) \leftarrow \Hom_R(M'', R)
$$
Dualizing again we obtain a commutative diagram
$$
\xymatrix{
\Hom_R(\Hom_R(M, R), R) \ar[r]_j &
\Hom_R(\Hom_R(M', R), R) \ar[r] &
\Hom_R(\Hom_R(M'', R), R) \\
M \ar[u] \ar[r] & M' \ar[u] \ar[r] & M'' \ar[u]
}
$$
We do not know the top row is exact. But, by assumption
the middle vertical arrow is an isomorphism and the right
vertical arrow is injective (Lemma \ref{lemma-reflexive-torsion-free}).
We claim $j$ is injective. Assuming the claim
a diagram chase shows that the left vertical
arrow is an isomorphism, i.e., $M$ is reflexive.

\medskip\noindent
Proof of the claim. Consider the exact sequence
$\Hom_R(M',R)\to \Hom_R(M,R)\to Q \to 0$
defining $Q$. One applies
Algebra, Lemma \ref{algebra-lemma-hom-from-finitely-presented}
to obtain
$$
\Hom_K(M'\otimes_R K,K) \to
\Hom_K(M\otimes_R K,K) \to
Q\otimes_R K\to 0
$$
But $M \otimes_R K \to M' \otimes_R K$ is an injective map of vector spaces,
hence split injective, so $Q \otimes_R K = 0$, that is, $Q$ is torsion.
Then one gets the exact sequence 
$$
0 \to \Hom_R(Q,R) \to
\Hom_R(\Hom_R(M,R),R) \to
\Hom_R(\Hom_R(M',R),R)
$$
and $\Hom_R(Q,R)=0$ because $Q$ is torsion.
\end{proof}

\begin{lemma}
\label{lemma-characterize-reflexive}
Let $R$ be a Noetherian domain. Let $M$ be a finite $R$-module.
The following are equivalent
\begin{enumerate}
\item $M$ is reflexive,
\item there exists a short exact sequence $0 \to M \to F \to N \to 0$
with $F$ finite free and $N$ torsion free.
\end{enumerate}
\end{lemma}

\begin{proof}
Observe that a finite free module is reflexive.
By Lemma \ref{lemma-sequence-reflexive} we see
that (2) implies (1). Assume $M$ is reflexive. Choose a presentation
$R^{\oplus m} \to R^{\oplus n} \to \Hom_R(M, R) \to 0$.
Dualizing we get an exact sequence
$$
0 \to \Hom_R(\Hom_R(M, R), R) \to R^{\oplus n} \to N \to 0
$$
with $N = \Im(R^{\oplus n} \to R^{\oplus m})$ a torsion free module.
As $M = \Hom_R(\Hom_R(M, R), R)$ we get an exact sequence as in (2).
\end{proof}

\begin{lemma}
\label{lemma-flat-pullback-reflexive}
Let $R \to R'$ be a flat homomorphism of Noetherian domains.
If $M$ is a finite reflexive $R$-module, then $M \otimes_R R'$
is a finite reflexive $R'$-module.
\end{lemma}

\begin{proof}
Choose a short exact sequence $0 \to M \to F \to N \to 0$
with $F$ finite free and $N$ torsion free, see
Lemma \ref{lemma-characterize-reflexive}. Since $R \to R'$ is flat
we obtain a short exact sequence
$0 \to M \otimes_R R' \to F \otimes_R R' \to N \otimes_R R' \to 0$
with $F \otimes_R R'$ finite free and $N \otimes_R R'$ torsion
free (Lemma \ref{lemma-flat-pullback-torsion}). Thus $M \otimes_R R'$
is reflexive by Lemma \ref{lemma-characterize-reflexive}.
\end{proof}

\begin{lemma}
\label{lemma-dual-reflexive}
Let $R$ be a Noetherian domain. Let $M$ be a finite $R$-module.
Let $N$ be a finite reflexive $R$-module. Then $\Hom_R(M, N)$ is reflexive.
\end{lemma}

\begin{proof}
Choose a presentation $R^{\oplus m} \to R^{\oplus n} \to M \to 0$.
Then we obtain
$$
0 \to \Hom_R(M, N) \to N^{\oplus n} \to N' \to 0
$$
with $N' = \Im(N^{\oplus n} \to N^{\oplus m})$ torsion free.
We conclude by Lemma \ref{lemma-sequence-reflexive}.
\end{proof}

\begin{definition}
\label{definition-reflexive-hull}
Let $R$ be a Noetherian domain. Let $M$ be a finite $R$-module.
The module $M^{**} = \Hom_R(\Hom_R(M, R), R)$ is called the
{\it reflexive hull} of $M$.
\end{definition}

\noindent
This makes sense because the reflexive hull is reflexive by
Lemma \ref{lemma-dual-reflexive}. The assignment $M \mapsto M^{**}$
is a functor. If $\varphi : M \to N$ is an $R$-module map into a reflexive
$R$-module $N$, then $\varphi$ factors $M \to M^{**} \to N$
through the reflexive hull of $M$. Another way to say this is that
taking the reflexive hull is the left adjoint to the inclusion functor
$$
\text{finite reflexive modules} \subset
\text{finite modules}
$$
over a Noetherian domain $R$.

\begin{lemma}
\label{lemma-hom-into-depth}
Let $R$ be a Noetherian local ring. Let $M$, $N$ be finite $R$-modules.
\begin{enumerate}
\item If $N$ has depth $\geq 1$, then $\Hom_R(M, N)$ has depth $\geq 1$.
\item If $N$ has depth $\geq 2$, then $\Hom_R(M, N)$ has depth $\geq 2$.
\end{enumerate}
\end{lemma}

\begin{proof}
Choose a presentation $R^{\oplus m} \to R^{\oplus n} \to M \to 0$.
Dualizing we get an exact sequence
$$
0 \to \Hom_R(M, N) \to N^{\oplus n} \to N' \to 0
$$
with $N' = \Im(N^{\oplus n} \to N^{\oplus m})$. A submodule of a module
with depth $\geq 1$ has depth $\geq 1$; this follows immediately from
the definition. Thus part (1) is clear. For (2) note that here the
assumption and the previous remark implies $N'$ has depth $\geq 1$.
The module $N^{\oplus n}$ has depth $\geq 2$.
From Algebra, Lemma \ref{algebra-lemma-depth-in-ses}
we conclude $\Hom_R(M, N)$ has depth $\geq 2$.
\end{proof}

\begin{lemma}
\label{lemma-hom-into-S2}
Let $R$ be a Noetherian ring. Let $M$, $N$ be finite $R$-modules.
\begin{enumerate}
\item If $N$ has property $(S_1)$, then $\Hom_R(M, N)$ has property $(S_1)$.
\item If $N$ has property $(S_2)$, then $\Hom_R(M, N)$ has property $(S_2)$.
\item If $R$ is a domain, $N$ is torsion free and $(S_2)$, then
$\Hom_R(M, N)$ is torsion free and has property $(S_2)$.
\end{enumerate}
\end{lemma}

\begin{proof}
Since localizing at primes commutes with taking $\Hom_R$ for finite
$R$-modules (Algebra, Lemma \ref{algebra-lemma-ext-noetherian})
parts (1) and (2) follow immediately from Lemma \ref{lemma-hom-into-depth}.
Part (3) follows from (2) and
Lemma \ref{lemma-hom-into-torsion-free}.
\end{proof}

\begin{lemma}
\label{lemma-check-injective-on-ass}
Let $R$ be a Noetherian ring. Let $\varphi : M \to N$ be a map of
$R$-modules. Assume that for every prime $\mathfrak p$
of $R$ at least one of the following happens
\begin{enumerate}
\item $M_\mathfrak p \to N_\mathfrak p$ is injective, or
\item $\mathfrak p \not \in \text{Ass}(M)$.
\end{enumerate}
Then $\varphi$ is injective.
\end{lemma}

\begin{proof}
Let $\mathfrak p$ be an associated prime of $\Ker(\varphi)$.
Then there exists an element $x \in M_\mathfrak p$ which is
in the kernel of $M_\mathfrak p \to N_\mathfrak p$ and whose
annihilator is $\mathfrak pR_\mathfrak p$
(Algebra, Lemma \ref{algebra-lemma-associated-primes-localize}).
This is impossible in both cases. Hence
$\text{Ass}(\Ker(\varphi)) = \emptyset$ and we conclude $\Ker(\varphi) = 0$ by
Algebra, Lemma \ref{algebra-lemma-ass-zero}.
\end{proof}

\begin{lemma}
\label{lemma-check-isomorphism-via-depth-and-ass}
Let $R$ be a Noetherian ring. Let $\varphi : M \to N$ be a map of
$R$-modules. Assume $M$ is finite and that for every prime $\mathfrak p$
of $R$ one of the following happens
\begin{enumerate}
\item $M_\mathfrak p \to N_\mathfrak p$ is an isomorphism, or
\item $\text{depth}(M_\mathfrak p) \geq 2$ and
$\mathfrak p \not \in \text{Ass}(N)$.
\end{enumerate}
Then $\varphi$ is an isomorphism.
\end{lemma}

\begin{proof}
By Lemma \ref{lemma-check-injective-on-ass} we see that $\varphi$ is injective.
Let $N' \subset N$ be an finitely generated $R$-module containing
the image of $M$. Then $\text{Ass}(N_\mathfrak p) = \emptyset$ implies
$\text{Ass}(N'_\mathfrak p) = \emptyset$.
Hence the assumptions of the lemma hold for $M \to N'$.
In order to prove that $\varphi$ is an isomorphism, it suffices
to prove the same thing for every such $N' \subset N$. Thus we may
assume $N$ is a finite $R$-module. In this case,
$\mathfrak p \not \in \text{Ass}(N) \Rightarrow
\text{depth}(N_\mathfrak p) \geq 1$, see
Algebra, Lemma \ref{algebra-lemma-ideal-nonzerodivisor}.
Consider the short exact sequence
$$
0 \to M \to N \to Q \to 0
$$
defining $Q$.
Looking at the conditions we see that either $Q_\mathfrak p = 0$
in case (1) or $\text{depth}(Q_\mathfrak p) \geq 1$ in case (2)
by Algebra, Lemma \ref{algebra-lemma-depth-in-ses}.
This implies that $Q$ does not have any associated primes, hence $Q = 0$ by
Algebra, Lemma \ref{algebra-lemma-ass-zero}.
\end{proof}

\begin{lemma}
\label{lemma-isom-depth-2-torsion-free}
Let $R$ be a Noetherian domain. Let $\varphi : M \to N$ be a map of
$R$-modules. Assume $M$ is finite, $N$ is torsion free, and
that for every prime $\mathfrak p$ of $R$ one of the following happens
\begin{enumerate}
\item $M_\mathfrak p \to N_\mathfrak p$ is an isomorphism, or
\item $\text{depth}(M_\mathfrak p) \geq 2$.
\end{enumerate}
Then $\varphi$ is an isomorphism.
\end{lemma}

\begin{proof}
This is a special case of
Lemma \ref{lemma-check-isomorphism-via-depth-and-ass}.
\end{proof}

\begin{lemma}
\label{lemma-reflexive-depth-2}
Let $R$ be a Noetherian domain. Let $M$ be a finite $R$-module.
The following are equivalent
\begin{enumerate}
\item $M$ is reflexive,
\item for every prime $\mathfrak p$ of $R$ one of the following happens
\begin{enumerate}
\item $M_\mathfrak p$ is a reflexive $R_\mathfrak p$-module, or
\item $\text{depth}(M_\mathfrak p) \geq 2$.
\end{enumerate}
\end{enumerate}
\end{lemma}

\begin{proof}
If (1) is true, then $M_\mathfrak p$ is a reflexive module
for all primes of $\mathfrak p$ by Lemma \ref{lemma-check-reflexive}.
Thus (1) $\Rightarrow$ (2). Assume (2). Set $N = \Hom_R(\Hom_R(M, R), R)$
so that
$$
N_\mathfrak p =
\Hom_{R_\mathfrak p}(
\Hom_{R_\mathfrak p}(M_\mathfrak p, R_\mathfrak p), R_\mathfrak p)
$$
for every prime $\mathfrak p$ of $R$.
See Algebra, Lemma \ref{algebra-lemma-hom-from-finitely-presented}.
We apply Lemma \ref{lemma-isom-depth-2-torsion-free} to the map
$j : M \to N$. This is allowed because $M$ is finite and
$N$ is torsion free by Lemma \ref{lemma-hom-into-torsion-free}.
In case (2)(a) the map $M_\mathfrak p \to N_\mathfrak p$ is an
isomorphism and in case (2)(b) we have $\text{depth}(M_\mathfrak p) \geq 2$.
\end{proof}

\begin{lemma}
\label{lemma-reflexive-S2}
Let $R$ be a Noetherian domain. Let $M$ be a finite reflexive $R$-module.
Let $\mathfrak p \subset R$ be a prime ideal.
\begin{enumerate}
\item If $\text{depth}(R_\mathfrak p) \geq 2$, then
$\text{depth}(M_\mathfrak p) \geq 2$.
\item If $R$ is $(S_2)$, then $M$ is $(S_2)$.
\end{enumerate}
\end{lemma}

\begin{proof}
Since formation of reflexive hull $\Hom_R(\Hom_R(M, R), R)$
commutes with localization
(Algebra, Lemma \ref{algebra-lemma-hom-from-finitely-presented})
part (1) follows from Lemma \ref{lemma-hom-into-depth}.
Part (2) is immediate from Lemma \ref{lemma-hom-into-S2}.
\end{proof}

\begin{example}
\label{example-ring-not-S2}
The results above and below suggest reflexivity is related to the $(S_2)$
condition; here is an example to prevent too optimistic conjectures.
Let $k$ be a field. Let $R$ be the $k$-subalgebra of $k[x, y]$
generated by $1, y, x^2, xy, x^3$. Then $R$ is not $(S_2)$. So
$R$ as an $R$-module is an example of a reflexive $R$-module
which is not $(S_2)$. Let $M = k[x, y]$ viewed as an $R$-module.
Then $M$ is a reflexive $R$-module because
$$
\Hom_R(M, R) = \mathfrak m = (y, x^2, xy, x^3)
\quad\text{and}\quad
\Hom_R(\mathfrak m, R) = M
$$
and $M$ is $(S_2)$ as an $R$-module (computations omitted).
Thus $R$ is a Noetherian domain possessing
a reflexive $(S_2)$ module but $R$ is not $(S_2)$ itself.
\end{example}

\begin{lemma}
\label{lemma-reflexive-over-normal}
Let $R$ be a Noetherian normal domain with fraction field $K$.
Let $M$ be a finite $R$-module. The following are equivalent
\begin{enumerate}
\item $M$ is reflexive,
\item $M$ is torsion free and has property $(S_2)$,
\item $M$ is torsion free and
$M = \bigcap_{\text{height}(\mathfrak p) = 1} M_{\mathfrak p}$
where the intersection happens in $M_K = M \otimes_R K$.
\end{enumerate}
\end{lemma}

\begin{proof}
By Algebra, Lemma \ref{algebra-lemma-criterion-normal}
we see that $R$ satisfies $(R_1)$ and $(S_2)$.

\medskip\noindent
Assume (1). Then $M$ is torsion free by Lemma \ref{lemma-reflexive-torsion-free}
and satisfies $(S_2)$ by Lemma \ref{lemma-reflexive-S2}.
Thus (2) holds.

\medskip\noindent
Assume (2). By definition
$M' = \bigcap_{\text{height}(\mathfrak p) = 1} M_{\mathfrak p}$
is the kernel of the map
$$
M_K
\longrightarrow
\bigoplus\nolimits_{\text{height}(\mathfrak p) = 1} M_K/M_\mathfrak p
\subset
\prod\nolimits_{\text{height}(\mathfrak p) = 1} M_K/M_\mathfrak p
$$
Observe that our map indeed factors through the direct sum as indicated
since given
$a/b \in K$ there are at most finitely many height $1$ primes $\mathfrak p$
with $b \in \mathfrak p$. Let $\mathfrak p_0$ be a prime of height $1$.
Then $(M_K/M_\mathfrak p)_{\mathfrak p_0} = 0$ unless
$\mathfrak p = \mathfrak p_0$ in which case we get
$(M_K/M_\mathfrak p)_{\mathfrak p_0} = M_K/M_{\mathfrak p_0}$.
Thus by exactness of localization and the fact that localization
commutes with direct sums, we see that
$M'_{\mathfrak p_0} = M_{\mathfrak p_0}$.
Since $M$ has depth $\geq 2$ at primes of height $> 1$,
we see that $M \to M'$ is an isomorphism by
Lemma \ref{lemma-isom-depth-2-torsion-free}. Hence (3) holds.

\medskip\noindent
Assume (3). Let $\mathfrak p$ be a prime of height $1$.
Then $R_\mathfrak p$ is a discrete valuation ring by $(R_1)$.
By Lemma \ref{lemma-dedekind-torsion-free-flat}
we see that $M_\mathfrak p$ is finite free, in particular
reflexive. Hence the map $M \to M^{**}$ induces an isomorphism at all
the primes $\mathfrak p$ of height $1$. Thus the condition
$M = \bigcap_{\text{height}(\mathfrak p) = 1} M_{\mathfrak p}$
implies that $M = M^{**}$ and (1) holds.
\end{proof}

\begin{lemma}
\label{lemma-describe-reflexive-hull}
Let $R$ be a Noetherian normal domain. Let $M$ be a finite $R$-module.
Then the reflexive hull of $M$ is the intersection
$$
M^{**} =
\bigcap\nolimits_{\text{height}(\mathfrak p) = 1}
M_{\mathfrak p}/(M_\mathfrak p)_{tors} =
\bigcap\nolimits_{\text{height}(\mathfrak p) = 1}
(M/M_{tors})_\mathfrak p
$$
taken in $M \otimes_R K$.
\end{lemma}

\begin{proof}
Let $\mathfrak p$ be a prime of height $1$.
The kernel of $M_\mathfrak p \to M \otimes_R K$ is the
torsion submodule $(M_\mathfrak p)_{tors}$ of $M_\mathfrak p$.
Moreover, we have
$(M/M_{tors})_\mathfrak p = M_\mathfrak p/(M_\mathfrak p)_{tors}$
and this is a finite free module over the discrete valuation ring
$R_\mathfrak p$
(Lemma \ref{lemma-dedekind-torsion-free-flat}).
Then $M_\mathfrak p/(M_\mathfrak p)_{tors} \to
(M_\mathfrak p)^{**} = (M^{**})_\mathfrak p$
is an isomorphism, hence the lemma is a consequence of
Lemma \ref{lemma-reflexive-over-normal}.
\end{proof}

\begin{lemma}
\label{lemma-integral-closure-reflexive}
Let $A$ be a Noetherian normal domain with fraction field $K$.
Let $L$ be a finite extension of $K$. If the integral closure
$B$ of $A$ in $L$ is finite over $A$, then $B$ is reflexive as an $A$-module.
\end{lemma}

\begin{proof}
It suffices to show that $B = \bigcap B_\mathfrak p$ where the intersection
is over height $1$ primes $\mathfrak p \subset A$, see
Lemma \ref{lemma-reflexive-over-normal}.
Let $b \in \bigcap B_\mathfrak p$. Let $x^d + a_1x^{d - 1} + \ldots + a_d$
be the minimal polynomial of $b$ over $K$.
We want to show $a_i \in A$.
By Algebra, Lemma \ref{algebra-lemma-minimal-polynomial-normal-domain}
we see that $a_i \in A_\mathfrak p$ for all $i$ and all height one primes
$\mathfrak p$. Hence we get what we want from
Algebra, Lemma
\ref{algebra-lemma-normal-domain-intersection-localizations-height-1}
(or the lemma already cited as $A$ is a reflexive module over itself).
\end{proof}








\section{Content ideals}
\label{section-content}

\noindent
The definition may not be what you expect.

\begin{definition}
\label{definition-content-ideal}
Let $A$ be a ring. Let $M$ be a flat $A$-module. Let $x \in M$.
If the set of ideals $I$ in $A$ such that $x \in IM$ has a
smallest element, we call it the {\it content ideal of $x$}.
\end{definition}

\noindent
Note that since $M$ is flat over $A$, for a pair of ideals $I, I'$
of $A$ we have $IM \cap I'M = (I \cap I')M$
as can be seen by tensoring the exact sequence
$0 \to I \cap I' \to I \oplus I' \to I + I' \to 0$ by $M$.

\begin{lemma}
\label{lemma-content-finitely-generated}
Let $A$ be a ring. Let $M$ be a flat $A$-module. Let $x \in M$.
The content ideal of $x$, if it exists, is finitely generated.
\end{lemma}

\begin{proof}
Say $x \in IM$. Then we can write $x = \sum_{i = 1, \ldots, n} f_i x_i$ with
$f_i \in I$ and $x_i \in M$. Hence $x \in I'M$ with
$I' = (f_1, \ldots, f_n)$.
\end{proof}

\begin{lemma}
\label{lemma-equal-content}
Let $(A, \mathfrak m)$ be a local ring. Let $u : M \to N$ be a map of flat
$A$-modules such that $\overline{u} : M/\mathfrak m M \to N/\mathfrak m N$
is injective. If $x \in M$ has content ideal $I$, then $u(x)$ has content
ideal $I$ as well.
\end{lemma}

\begin{proof}
It is clear that $u(x) \in IN$. If $u(x) \in I'N$, then
$u(x) \in (I' \cap I)N$, see discussion following
Definition \ref{definition-content-ideal}. Hence it suffices to
show: if $x \in I'N$ and $I' \subset I$, $I' \not =  I$, then
$u(x) \not \in I'N$. Since $I/I'$ is a nonzero finite $A$-module
(Lemma \ref{lemma-content-finitely-generated}) there is a nonzero map
$\chi : I/I' \to A/\mathfrak m$ of $A$-modules
by Nakayama's lemma (Algebra, Lemma \ref{algebra-lemma-NAK}).
Since $I$ is the content ideal of $x$ we see that
$x \not \in I''M$ where $I'' = \Ker(\chi)$.
Hence $x$ is not in the kernel of the map
$$
IM = I \otimes_A M \xrightarrow{\chi \otimes 1}
A/\mathfrak m \otimes M \cong M/\mathfrak m M
$$
Applying our hypothesis on $\overline{u}$ we conclude that
$u(x)$ does not map to zero under the map
$$
IN = I \otimes_A N \xrightarrow{\chi \otimes 1}
A/\mathfrak m \otimes N \cong N/\mathfrak m N
$$
and we conclude.
\end{proof}

\begin{lemma}
\label{lemma-content-exists-flat-Mittag-Leffler}
Let $A$ be a ring. Let $M$ be a flat Mittag-Leffler module.
Then every element of $M$ has a content ideal.
\end{lemma}

\begin{proof}
This is a special case of Algebra, Lemma \ref{algebra-lemma-flat-ML-criterion}.
\end{proof}








\section{Flatness and finiteness conditions}
\label{section-flat-finite}

\noindent
In this section we discuss some implications of the type
``flat $+$ finite type $\Rightarrow$ finite presentation''.
We will revisit this result in the chapter on flatness, see
More on Flatness, Section \ref{flat-section-introduction}.
A first result of this type was proved in
Algebra, Lemma \ref{algebra-lemma-finite-flat-module-finitely-presented}.

\begin{lemma}
\label{lemma-flat-finite-type-finite-presentation-local-module}
Let $R$ be a ring. Let $S = R[x_1, \ldots, x_n]$ be a polynomial
ring over $R$. Let $M$ be an $S$-module.
Assume
\begin{enumerate}
\item there exist finitely many primes $\mathfrak p_1, \ldots, \mathfrak p_m$
of $R$ such that the map $R \to \prod R_{\mathfrak p_j}$ is injective,
\item $M$ is a finite $S$-module,
\item $M$ flat over $R$, and
\item for every prime $\mathfrak p$ of $R$ the module $M_{\mathfrak p}$
is of finite presentation over $S_{\mathfrak p}$.
\end{enumerate}
Then $M$ is of finite presentation over $S$.
\end{lemma}

\begin{proof}
Choose a presentation
$$
0 \to K \to S^{\oplus r} \to M \to 0
$$
of $M$ as an $S$-module. Let $\mathfrak q$ be a prime ideal of $S$
lying over a prime $\mathfrak p$ of $R$. By assumption there exist
finitely many elements $k_1, \ldots, k_t \in K$ such that if we set
$K' = \sum Sk_j \subset K$ then
$K'_{\mathfrak p} = K_{\mathfrak p}$ and
$K'_{\mathfrak p_j} = K_{\mathfrak p_j}$ for $j = 1, \ldots, m$.
Setting $M' = S^{\oplus r}/K'$ we deduce that in particular
$M'_{\mathfrak q} = M_{\mathfrak q}$. By openness of flatness, see
Algebra, Theorem \ref{algebra-theorem-openness-flatness}
we conclude that there exists a $g \in S$, $g \not \in \mathfrak q$
such that $M'_g$ is flat over $R$. Thus $M'_g \to M_g$ is a surjective
map of flat $R$-modules. Consider the commutative diagram
$$
\xymatrix{
M'_g \ar[r] \ar[d] & M_g \ar[d] \\
\prod (M'_g)_{\mathfrak p_j} \ar[r] & \prod (M_g)_{\mathfrak p_j}
}
$$
The bottom arrow is an isomorphism by choice of $k_1, \ldots, k_t$.
The left vertical arrow is an injective map as
$R \to \prod R_{\mathfrak p_j}$ is injective and $M'_g$ is flat over $R$.
Hence the top horizontal arrow is injective, hence an isomorphism.
This proves that $M_g$ is of finite presentation over $S_g$.
We conclude by applying
Algebra, Lemma \ref{algebra-lemma-cover}.
\end{proof}

\begin{lemma}
\label{lemma-flat-finite-type-finite-presentation-local}
Let $R \to S$ be a ring homomorphism.
Assume
\begin{enumerate}
\item there exist finitely many primes
$\mathfrak p_1, \ldots, \mathfrak p_m$ of $R$ such that
the map $R \to \prod R_{\mathfrak p_j}$ is injective,
\item $R \to S$ is of finite type,
\item $S$ flat over $R$, and
\item for every prime $\mathfrak p$ of $R$ the ring $S_{\mathfrak p}$
is of finite presentation over $R_{\mathfrak p}$.
\end{enumerate}
Then $S$ is of finite presentation over $R$.
\end{lemma}

\begin{proof}
By assumption $S$ is a quotient of a polynomial ring over $R$.
Thus the result follows directly from
Lemma \ref{lemma-flat-finite-type-finite-presentation-local-module}.
\end{proof}

\begin{lemma}
\label{lemma-flat-graded-finite-type-finite-presentation-module}
Let $R$ be a ring.
Let $S = R[x_1, \ldots, x_n]$ be a graded polynomial algebra over $R$,
i.e., $\deg(x_i) > 0$ but not necessarily equal to $1$.
Let $M$ be a graded $S$-module.
Assume
\begin{enumerate}
\item $R$ is a local ring,
\item $M$ is a finite $S$-module, and
\item $M$ is flat over $R$.
\end{enumerate}
Then $M$ is finitely presented as an $S$-module.
\end{lemma}

\begin{proof}
Let $M = \bigoplus M_d$ be the grading on $M$.
Pick homogeneous generators $m_1, \ldots, m_r \in M$ of $M$.
Say $\deg(m_i) = d_i \in \mathbf{Z}$. This gives us a presentation
$$
0 \to K \to \bigoplus\nolimits_{i = 1, \ldots, r} S(-d_i) \to M \to 0
$$
which in each degree $d$ leads to the short exact sequence
$$
0 \to K_d \to \bigoplus\nolimits_{i = 1, \ldots, r} S_{d - d_i} \to
M_d \to 0.
$$
By assumption each $M_d$ is a finite flat $R$-module. By
Algebra, Lemma \ref{algebra-lemma-finite-flat-local}
this implies each $M_d$ is a finite free $R$-module. Hence
we see each $K_d$ is a finite $R$-module. Also each $K_d$ is flat
over $R$ by
Algebra, Lemma \ref{algebra-lemma-flat-ses}.
Hence we conclude that each $K_d$ is finite free by
Algebra, Lemma \ref{algebra-lemma-finite-flat-local}
again.

\medskip\noindent
Let $\mathfrak m$ be the maximal ideal of $R$. By the flatness of $M$
over $R$ the short exact sequences above remain short exact after tensoring
with $\kappa = \kappa(\mathfrak m)$. As the ring $S \otimes_R \kappa$ is
Noetherian we see that there exist homogeneous elements
$k_1, \ldots, k_t \in K$ such that the images $\overline{k}_j$
generate $K \otimes_R \kappa$ over $S \otimes_R \kappa$. Say $\deg(k_j) = e_j$.
Thus for any $d$ the map
$$
\bigoplus\nolimits_{j = 1, \ldots, t} S_{d - e_j}
\longrightarrow
K_d
$$
becomes surjective after tensoring with $\kappa$. By
Nakayama's lemma
(Algebra, Lemma \ref{algebra-lemma-NAK})
this implies the map is surjective over $R$. Hence $K$ is generated
by $k_1, \ldots, k_t$ over $S$ and we win.
\end{proof}

\begin{lemma}
\label{lemma-flat-graded-finite-type-finite-presentation}
Let $R$ be a ring. Let $S = \bigoplus_{n \geq 0} S_n$ be a graded $R$-algebra.
Let $M = \bigoplus_{d \in \mathbf{Z}} M_d$ be a graded $S$-module.
Assume $S$ is finitely generated as an $R$-algebra, assume $S_0$ is a finite
$R$-algebra, and assume there exist finitely many primes
$\mathfrak p_j$, $i = 1, \ldots, m$ such that
$R \to \prod R_{\mathfrak p_j}$ is injective.
\begin{enumerate}
\item If $S$ is flat over $R$, then $S$ is a finitely presented $R$-algebra.
\item If $M$ is flat as an $R$-module and finite as an $S$-module,
then $M$ is finitely presented as an $S$-module.
\end{enumerate}
\end{lemma}

\begin{proof}
As $S$ is finitely generated as an $R$-algebra, it is finitely generated
as an $S_0$ algebra, say by homogeneous elements $t_1, \ldots, t_n \in S$
of degrees $d_1, \ldots, d_n > 0$. Set $P = R[x_1, \ldots, x_n]$ with
$\deg(x_i) = d_i$. The ring map $P \to S$, $x_i \to t_i$ is finite
as $S_0$ is a finite $R$-module. To prove (1) it suffices to prove
that $S$ is a finitely presented $P$-module.  To prove (2) it suffices
to prove that $M$ is a finitely presented $P$-module. Thus it suffices
to prove that if $S = P$ is a graded polynomial ring and $M$ is a finite
$S$-module flat over $R$, then $M$ is finitely presented as an $S$-module. By
Lemma \ref{lemma-flat-graded-finite-type-finite-presentation-module}
we see $M_{\mathfrak p}$ is a finitely presented $S_{\mathfrak p}$-module
for every prime $\mathfrak p$ of $R$. Thus the result follows from
Lemma \ref{lemma-flat-finite-type-finite-presentation-local-module}.
\end{proof}

\begin{remark}
\label{remark-when-does-condition-hold}
Let $R$ be a ring. When does $R$ satisfy the condition mentioned in
Lemmas \ref{lemma-flat-finite-type-finite-presentation-local-module},
\ref{lemma-flat-finite-type-finite-presentation-local}, and
\ref{lemma-flat-graded-finite-type-finite-presentation}?
This holds if
\begin{enumerate}
\item $R$ is local,
\item $R$ is Noetherian,
\item $R$ is a domain,
\item $R$ is a reduced ring with finitely many minimal primes, or
\item $R$ has finitely many weakly associated primes, see
Algebra, Lemma \ref{algebra-lemma-zero-at-weakly-ass-zero}.
\end{enumerate}
Thus these lemmas hold in all cases listed above.
\end{remark}

\noindent
The following lemma will be improved on in
More on Flatness, Proposition
\ref{flat-proposition-flat-finite-type-finite-presentation-domain}.

\begin{lemma}
\label{lemma-flat-finite-type-valuation-ring-finite-presentation}
\begin{reference}
\cite[Theorem 3]{Nagata-Finitely}
\end{reference}
Let $A$ be a valuation ring. Let $A \to B$ be a ring map of finite type.
Let $M$ be a finite $B$-module.
\begin{enumerate}
\item If $B$ is flat over $A$, then $B$ is a finitely presented $A$-algebra.
\item If $M$ is flat as an $A$-module, then $M$ is finitely presented
as a $B$-module.
\end{enumerate}
\end{lemma}

\begin{proof}
We are going to use that an $A$-module is flat if and only if it is
torsion free, see
Lemma \ref{lemma-valuation-ring-torsion-free-flat}.
By
Algebra, Lemma \ref{algebra-lemma-homogenize}
we can find a graded $A$-algebra $S$ with $S_0 = A$ and generated
by finitely many elements in degree $1$, an element $f \in S_1$ and a
finite graded $S$-module $N$ such that $B \cong S_{(f)}$ and
$M \cong N_{(f)}$. If $M$ is torsion free, then we can take $N$ torsion
free by replacing it by $N/N_{tors}$, see
Lemma \ref{lemma-torsion}.
Similarly, if $B$ is torsion free, then we can take
$S$ torsion free by replacing it by $S/S_{tors}$.
Hence in case (1), we may apply
Lemma \ref{lemma-flat-graded-finite-type-finite-presentation}
to see that $S$ is a finitely presented
$A$-algebra, which implies that $B = S_{(f)}$ is a finitely
presented $A$-algebra. To see (2) we may first replace $S$ by
a graded polynomial ring, and then we may apply
Lemma \ref{lemma-flat-graded-finite-type-finite-presentation-module}
to conclude.
\end{proof}

\begin{lemma}
\label{lemma-valuation-ring-flat-essentially-finite-type}
Let $A$ be a valuation ring. Let $A \to B$ be a local homomorphism
which is essentially of finite type. Let $M$ be a finite $B$-module.
\begin{enumerate}
\item If $B$ is flat over $A$, then $B$ is essentially of finite
presentation over $A$.
\item If $M$ is flat as an $A$-module, then $M$ is finitely presented
as a $B$-module.
\end{enumerate}
\end{lemma}

\begin{proof}
By assumption we can write $B$ as a quotient of the localization of
a polynomial algebra $P = A[x_1, \ldots, x_n]$ at a prime ideal $\mathfrak q$.
In case (1) we consider $M = B$ as a finite module over
$P_\mathfrak q$ and in case (2) we consider $M$ as a finite module over
$P_\mathfrak q$. In both cases, we have to show that this is a
finitely presented $P_\mathfrak q$-module, see
Algebra, Lemma \ref{algebra-lemma-finitely-presented-over-subring}
for case (2).

\medskip\noindent
Choose a presentation $0 \to K \to P_\mathfrak q^{\oplus r} \to M \to 0$
which is possible because $M$ is finite over $P_\mathfrak q$.
Let $L = P^{\oplus r} \cap K$. Then $K = L_\mathfrak q$, see
Algebra, Lemma \ref{algebra-lemma-submodule-localization}.
Then $N = P^{\oplus r}/L$ is a submodule of $M$ and hence flat by
Lemma \ref{lemma-valuation-ring-torsion-free-flat}.
Since also $N$ is a finite $P$-module, we see that $N$ is finitely
presented as a $P$-module by
Lemma \ref{lemma-flat-finite-type-valuation-ring-finite-presentation}.
Since localization is exact
(Algebra, Proposition \ref{algebra-proposition-localization-exact})
we see that $N_\mathfrak q = M$ and we conclude.
\end{proof}











\section{Blowing up and flatness}
\label{section-blowup-flat}

\noindent
In this section we begin our discussion of results of the form: ``After a
blowup the strict transform becomes flat''. More results of this type may
be found in Divisors, Section \ref{divisors-section-blowup-flat} and
More on Flatness, Section \ref{flat-section-blowup-flat}.

\begin{definition}
\label{definition-strict-transform}
Let $R$ be a ring. Let $I \subset R$ be an ideal and $a \in I$.
Let $R[\frac{I}{a}]$ be the affine blowup algebra, see
Algebra, Definition \ref{algebra-definition-blow-up}.
Let $M$ be an $R$-module.
The {\it strict transform of $M$ along $R \to R[\frac{I}{a}]$} is
the $R[\frac{I}{a}]$-module
$$
M' = \left(M \otimes_R R[\textstyle{\frac{I}{a}}]\right)/a\text{-power torsion}
$$
\end{definition}

\noindent
The following is a very weak version of flattening by blowing up, but
it is already sometimes a useful result.

\begin{lemma}
\label{lemma-flatten-on-affine-blowup}
Let $(R, \mathfrak m)$ be a local domain with fraction field $K$.
Let $S$ be a finite type $R$-algebra.
Let $M$ be a finite $S$-module.
For every valuation ring $A \subset K$ dominating $R$
there exists an ideal $I \subset \mathfrak m$ and a nonzero
element $a \in I$ such that
\begin{enumerate}
\item $I$ is finitely generated,
\item $A$ has center on $R[\frac{I}{a}]$,
\item the fibre ring of $R \to R[\frac{I}{a}]$ at $\mathfrak m$
is not zero, and
\item the strict transform $S_{I, a}$ of $S$ along $R \to R[\frac{I}{a}]$
is flat and of finite presentation over $R$, and the strict transform
$M_{I, a}$ of $M$ along $R \to R[\frac{I}{a}]$ is flat over $R$ and
finitely presented over $S_{I, a}$.
\end{enumerate}
\end{lemma}

\begin{proof}
Write $S = R[x_1, \ldots, x_n]/J$ and denote $N = S \oplus M$
viewed as a module over $P = R[x_1, \ldots, x_n]$. If we can prove the
lemma in case $S$ is a polynomial algebra over $R$, then we can
find $I, a$ satisfying (1), (2), (3) such that the strict transform
$N_{I, a}$ of $N$ along $R \to R[\frac{I}{a}]$ is flat over
$R$ and finitely presented as a module over the strict transform $P_{I, a}$
of $P$. Since $P_{I, a} = R[\frac{I}{a}][x_1, \ldots, x_n]$
(small detail omitted)
we find that the summand $S_{I, a} \subset N_{I, a}$ is flat over $R$
and finitely presented as a module over $R[\frac{I}{a}][x_1, \ldots, x_n]$.
Hence $S_{I, a}$ is finitely presented as an $R[\frac{I}{a}]$-algebra.
Moreover, the summand $M_{I, a} \subset N_{I, a}$ is flat over $R$
and finitely presented as a module over $P_{I, a}$ hence also
finitely presented as a module over $S_{I, a}$, see
Algebra, Lemma \ref{algebra-lemma-finitely-presented-over-subring}.
This reduces us to the case discussed in the next paragraph.

\medskip\noindent
Assume $S = R[x_1, \ldots, x_n]$. Choose a presentation
$$
0 \to K \to S^{\oplus r} \to M \to 0.
$$
Let $M_A$ be the quotient of $M \otimes_R A$ by its torsion submodule, see
Lemma \ref{lemma-torsion}. Then $M_A$ is a finite module over
$S_A = A[x_1, \ldots, x_n]$. By
Lemma \ref{lemma-valuation-ring-torsion-free-flat}
we see that $M_A$ is flat over $A$. By
Lemma \ref{lemma-flat-finite-type-valuation-ring-finite-presentation}
we see that $M_A$ is finitely presented. Hence there exist finitely many
elements $k_1, \ldots, k_t \in S_A^{\oplus r}$ which generate the
kernel of the presentation $S_A^{\oplus r} \to M_A$ as
an $S_A$-module. For any choice of $a \in I \subset \mathfrak m$
satisfying (1), (2), and (3) we denote $M_{I, a}$ the strict transform of
$M$ along $R \to R[\frac{I}{a}]$. It is a finite module over
$S_{I, a} = R[\frac{I}{a}][x_1, \ldots, x_n]$. By
Algebra, Lemma \ref{algebra-lemma-valuation-ring-colimit-affine-blowups}
we have $A = \colim_{I, a} R[\frac{I}{a}]$.
This implies that $S_A = \colim S_{I, a}$ and
$$
\colim M \otimes_R R[\textstyle{\frac{I}{a}}] = M \otimes_R A
$$
Choose $I, a$ and lifts $k_1, \ldots, k_t \in S_{I, a}^{\oplus r}$.
Since $M_A$ is the quotient of $M \otimes_R A$
by torsion, we see that the images of $k_1, \ldots, k_t$ in $M \otimes_R A$
are annihilated by a nonzero element $\alpha \in A$.
After replacing $I, a$ by a different pair
(recall that the colimit is filtered),
we may assume $\alpha = x/a^n$ for some $x \in I^n$ nonzero.
Then we find that $x k_1, \ldots, x k_t$ map to zero in $M \otimes_R A$.
Hence after replacing $I, a$ by a different pair
we may assume $x k_1, \ldots, x k_t$ map to zero in
$M \otimes_R R[\frac{I}{a}]$ for some nonzero $x \in R$.
Then finally replacing $I, a$ by $xI, xa$ we find that we may assume
$k_1, \ldots, k_t$ map to $a$-power torsion elements of
$M \otimes_R R[\frac{I}{a}]$. For any such pair $(I, a)$ we set
$$
M'_{I, a} = S_{I, a}^{\oplus r}/ \sum S_{I, a}k_j.
$$
Since $M_A = S_A^{\oplus r}/ \sum S_Ak_j$ we see that
$M_A = \colim_{I, a} M'_{I, a}$. At this point we finally apply
Algebra, Lemma \ref{algebra-lemma-flat-finite-presentation-limit-flat} (3)
to conclude that $M'_{I, a}$ is flat for some pair $(I, a)$ as above.
This lemma does not apply a priori to the system of strict transforms
$$
M_{I, a} = (M \otimes_R R[\textstyle{\frac{I}{a}}])/a\text{-power torsion}
$$
as the transition maps may not satisfy the assumptions of the lemma.
But now, since flatness implies torsion free
(Lemma \ref{lemma-flat-torsion-free}) and since $M_{I, a}$
is the quotient of $M'_{I, a}$ (because we arranged it so the
elements $k_1, \ldots, k_t$ map to zero in $M_{I, a}$)
by the $a$-power torsion submodule
we also conclude that $M'_{I, a} = M_{I, a}$ for such a pair and we win.
\end{proof}

\begin{lemma}
\label{lemma-blowup-fitting-ideal}
Let $R$ be a ring. Let $M$ be a finite $R$-module.
Let $k \geq 0$ and $I = \text{Fit}_k(M)$. For every $a \in I$
with $R' = R[\frac{I}{a}]$ the strict transform
$$
M' = (M \otimes_R R')/a\text{-power torsion}
$$
has $\text{Fit}_k(M') = R'$.
\end{lemma}

\begin{proof}
First observe that $\text{Fit}_k(M \otimes_R R') = IR' = aR'$.
The first equality by Lemma \ref{lemma-fitting-ideal-basics} part (3)
and the second equality by
Algebra, Lemma \ref{algebra-lemma-affine-blowup}.
From Lemma \ref{lemma-principal-fitting-ideal}
and exactness of localization
we see that $M'_{\mathfrak p'}$
can be generated by $\leq k$ elements for every prime $\mathfrak p'$
of $R'$. Then $\text{Fit}_k(M') = R'$ for example by
Lemma \ref{lemma-fitting-ideal-generate-locally}.
\end{proof}

\begin{lemma}
\label{lemma-blowup-fitting-ideal-locally-free}
Let $R$ be a ring. Let $M$ be a finite $R$-module.
Let $k \geq 0$ and $I = \text{Fit}_k(M)$. Assume that
$M_\mathfrak p$ is free of rank $k$ for every
$\mathfrak p \not \in V(I)$. Then for every $a \in I$
with $R' = R[\frac{I}{a}]$ the strict transform
$$
M' = (M \otimes_R R')/a\text{-power torsion}
$$
is locally free of rank $k$.
\end{lemma}

\begin{proof}
By Lemma \ref{lemma-blowup-fitting-ideal} we have
$\text{Fit}_k(M') = R'$. By Lemma \ref{lemma-fitting-ideal-finite-locally-free}
it suffices to show that $\text{Fit}_{k - 1}(M') = 0$.
Recall that $R' \subset R'_a = R_a$, see
Algebra, Lemma \ref{algebra-lemma-affine-blowup}.
Hence it suffices to prove that $\text{Fit}_{k - 1}(M')$
maps to zero in $R'_a = R_a$.
Since clearly $(M')_a = M_a$ this reduces us to showing
that $\text{Fit}_{k - 1}(M_a) = 0$
because formation of Fitting ideals commutes with base
change according to Lemma \ref{lemma-fitting-ideal-basics} part (3).
This is true by our assumption that
$M_a$ is finite locally free of rank $k$
(see Algebra, Lemma \ref{algebra-lemma-finite-projective})
and the already cited Lemma \ref{lemma-fitting-ideal-finite-locally-free}.
\end{proof}

\begin{lemma}
\label{lemma-blowup-module}
Let $R$ be a ring. Let $M$ be a finite $R$-module. Let $f \in R$
be an element such that $M_f$ is finite locally free of rank $r$.
Then there exists a finitely generated ideal $I \subset R$ with
$V(f) = V(I)$ such that for all $a \in I$ with $R' = R[\frac{I}{a}]$
the strict transform
$$
M' = (M \otimes_R R')/a\text{-power torsion}
$$
is locally free of rank $r$.
\end{lemma}

\begin{proof}
Choose a surjection $R^{\oplus n} \to M$. Choose a finite submodule
$K \subset \Ker(R^{\oplus n} \to M)$ such that $R^{\oplus n}/K \to M$
becomes an isomorphism after inverting $f$. This is possible because
$M_f$ is of finite presentation for example by
Algebra, Lemma \ref{algebra-lemma-finite-projective}.
Set $M_1 = R^{\oplus n}/K$
and suppose we can prove the lemma for $M_1$. Say $I \subset R$ is the
corresponding ideal. Then for $a \in I$ the map
$$
M_1' = (M_1 \otimes_R R')/a\text{-power torsion}
\longrightarrow
M' = (M \otimes_R R')/a\text{-power torsion}
$$
is surjective. It is also an isomorphism after inverting $a$ in $R'$
as $R'_a = R_f$, see Algebra, Lemma \ref{algebra-lemma-blowup-in-principal}.
But $a$ is a nonzerodivisor on $M'_1$, whence the displayed map is an
isomorphism. Thus it suffices to prove the lemma in case $M$ is a finitely
presented $R$-module.

\medskip\noindent
Assume $M$ is a finitely presented $R$-module.
Then $J = \text{Fit}_r(M) \subset S$ is a finitely generated ideal.
We claim that $I = fJ$ works.

\medskip\noindent
We first check that $V(f) = V(I)$. The inclusion $V(f) \subset V(I)$ is
clear. Conversely, if $f \not \in \mathfrak p$, then
$\mathfrak p$ is not an element of $V(J)$ by 
Lemma \ref{lemma-fitting-ideal-generate-locally}.
Thus $\mathfrak p \not \in V(fJ) = V(I)$.

\medskip\noindent
Let $a \in I$ and set $R' = R[\frac{I}{a}]$. We may write $a = fb$
for some $b \in J$. By Algebra, Lemmas \ref{algebra-lemma-affine-blowup} and
\ref{algebra-lemma-blowup-add-principal} we see that $J R' = b R'$
and $b$ is a nonzerodivisor in $R'$. 
Let $\mathfrak p' \subset R' = R[\frac{I}{a}]$ be
a prime ideal. Then $JR'_{\mathfrak p'}$ is generated by $b$.
It follows from
Lemma \ref{lemma-principal-fitting-ideal}
that $M'_{\mathfrak p'}$ can be generated by $r$ elements.
Since $M'$ is finite, there exist $m_1, \ldots, m_r \in M'$ and
$g \in R'$, $g \not \in \mathfrak p'$ such that the corresponding map
$(R')^{\oplus r} \to M'$ becomes surjective after inverting $g$.

\medskip\noindent
Finally, consider the ideal $J' = \text{Fit}_{k - 1}(M')$.
Note that $J' R'_g$ is generated by the coefficients of relations between
$m_1, \ldots, m_r$ (compatibility of Fitting ideal with base change).
Thus it suffices to show that $J' = 0$, see
Lemma \ref{lemma-fitting-ideal-finite-locally-free}.
Since $R'_a = R_f$ (Algebra, Lemma \ref{algebra-lemma-blowup-in-principal})
and $M'_a = M_f$ is free of rank $r$ we see that $J'_a = 0$.
Since $a$ is a nonzerodivisor in $R'$ we
conclude that $J' = 0$ and we win.
\end{proof}












\section{Completion and flatness}
\label{section-completion-flat}

\noindent
In this section we discuss when the completion of a ``big'' flat module
is flat.

\begin{lemma}
\label{lemma-ui-completion-direct-sum-into-product}
Let $R$ be a ring.
Let $I \subset R$ be an ideal.
Let $A$ be a set.
Assume $R$ is Noetherian and complete with respect to $I$. There is a
canonical map
$$
\left(\bigoplus\nolimits_{\alpha \in A} R\right)^\wedge
\longrightarrow
\prod\nolimits_{\alpha \in A} R
$$
from the $I$-adic completion of the direct sum into the product
which is universally injective.
\end{lemma}

\begin{proof}
By definition an element $x$ of the left hand side is $x = (x_n)$ where
$x_n = (x_{n, \alpha}) \in \bigoplus\nolimits_{\alpha \in A} R/I^n$
such that $x_{n, \alpha} = x_{n + 1, \alpha} \bmod I^n$.
As $R = R^\wedge$ we see that for any $\alpha$ there exists a $y_\alpha \in R$
such that $x_{n, \alpha} = y_\alpha \bmod I^n$. Note that for each $n$ there
are only finitely many $\alpha$ such that the elements $x_{n, \alpha}$ are
nonzero. Conversely, given $(y_\alpha) \in \prod_\alpha R$ such that for each
$n$ there are only finitely many $\alpha$ such that $y_{\alpha} \bmod I^n$
is nonzero, then this defines an element of the left hand side.
Hence we can think of an element of the left hand side as infinite
``convergent sums'' $\sum_\alpha y_\alpha$ with $y_\alpha \in R$
such that for each $n$ there are only finitely many $y_\alpha$
which are nonzero modulo $I^n$. The displayed map maps this element
to the element to $(y_\alpha)$ in the product.
In particular the map is injective.

\medskip\noindent
Let $Q$ be a finite $R$-module. We have to show that the map
$$
Q \otimes_R \left(\bigoplus\nolimits_{\alpha \in A} R\right)^\wedge
\longrightarrow
Q \otimes_R \left(\prod\nolimits_{\alpha \in A} R\right)
$$
is injective, see
Algebra, Theorem \ref{algebra-theorem-universally-exact-criteria}.
Choose a presentation $R^{\oplus k} \to R^{\oplus m} \to Q \to 0$
and denote $q_1, \ldots, q_m \in Q$ the corresponding generators for $Q$.
By Artin-Rees
(Algebra, Lemma \ref{algebra-lemma-Artin-Rees})
there exists a constant $c$ such that
$\Im(R^{\oplus k} \to R^{\oplus m}) \cap (I^N)^{\oplus m}
\subset \Im((I^{N - c})^{\oplus k} \to R^{\oplus m})$.
Let us contemplate the diagram
$$
\xymatrix{
\bigoplus_{l = 1}^k \left(\bigoplus\nolimits_{\alpha \in A} R\right)^\wedge
\ar[r] \ar[d] &
\bigoplus_{j = 1}^m \left(\bigoplus\nolimits_{\alpha \in A} R\right)^\wedge
\ar[r] \ar[d] &
Q \otimes_R \left(\bigoplus\nolimits_{\alpha \in A} R\right)^\wedge
\ar[r] \ar[d] &
0 \\
\bigoplus_{l = 1}^k \left(\prod\nolimits_{\alpha \in A} R\right)
\ar[r] &
\bigoplus_{j = 1}^m \left(\prod\nolimits_{\alpha \in A} R\right)
\ar[r] &
Q \otimes_R \left(\prod\nolimits_{\alpha \in A} R\right)
\ar[r] &
0
}
$$
with exact rows. Pick an element $\sum_j \sum_\alpha y_{j, \alpha}$ of
$\bigoplus_{j = 1, \ldots, m}
\left(\bigoplus\nolimits_{\alpha \in A} R\right)^\wedge$.
If this element maps to zero in the module
$Q \otimes_R \left(\prod\nolimits_{\alpha \in A} R\right)$,
then we see in particular that
$\sum_j q_j \otimes y_{j, \alpha} = 0$ in $Q$ for each $\alpha$.
Thus we can find an element
$(z_{1, \alpha}, \ldots, z_{k, \alpha}) \in \bigoplus_{l = 1, \ldots, k} R$
which maps to
$(y_{1, \alpha}, \ldots, y_{m, \alpha}) \in \bigoplus_{j = 1, \ldots, m} R$.
Moreover, if $y_{j, \alpha} \in I^{N_\alpha}$ for $j = 1, \ldots, m$, then
we may assume that $z_{l, \alpha} \in I^{N_\alpha - c}$ for
$l = 1, \ldots, k$.
Hence the sum $\sum_l \sum_\alpha z_{l, \alpha}$ is ``convergent'' and
defines an element of
$\bigoplus_{l = 1, \ldots, k}
\left(\bigoplus\nolimits_{\alpha \in A} R\right)^\wedge$
which maps to the element $\sum_j \sum_\alpha y_{j, \alpha}$ we started
out with. Thus the right vertical arrow is injective and we win.
\end{proof}

\noindent
The following lemma can also be deduced from
Lemma \ref{lemma-limit-flat} below.

\begin{lemma}
\label{lemma-completed-direct-sum-flat}
Let $R$ be a ring. Let $I \subset R$ be an ideal. Let $A$ be a set.
Assume $R$ is Noetherian. The completion
$(\bigoplus\nolimits_{\alpha \in A} R)^\wedge$
is a flat $R$-module.
\end{lemma}

\begin{proof}
Denote $R^\wedge$ the completion of $R$ with respect to $I$. As
$R \to R^\wedge$ is flat by
Algebra, Lemma \ref{algebra-lemma-completion-flat}
it suffices to prove that
$(\bigoplus\nolimits_{\alpha \in A} R)^\wedge$ is a flat
$R^\wedge$-module (use
Algebra, Lemma \ref{algebra-lemma-composition-flat}).
Since
$$
(\bigoplus\nolimits_{\alpha \in A} R)^\wedge
=
(\bigoplus\nolimits_{\alpha \in A} R^\wedge)^\wedge
$$
we may replace $R$ by $R^\wedge$ and assume that $R$ is complete with
respect to $I$ (see
Algebra, Lemma \ref{algebra-lemma-completion-complete}).
In this case
Lemma \ref{lemma-ui-completion-direct-sum-into-product}
tells us the map
$(\bigoplus\nolimits_{\alpha \in A} R)^\wedge \to \prod_{\alpha \in A} R$
is universally injective. Thus, by
Algebra, Lemma \ref{algebra-lemma-ui-flat-domain}
it suffices to show that $\prod_{\alpha \in A} R$ is flat. By
Algebra, Proposition \ref{algebra-proposition-characterize-coherent}
(and
Algebra, Lemma \ref{algebra-lemma-Noetherian-coherent})
we see that $\prod_{\alpha \in A} R$ is flat.
\end{proof}

\begin{lemma}
\label{lemma-tor-strictly-pro-zero}
\begin{reference}
This is \cite[Lemma 9.9]{quillenhomology}; note that
the author forgot the word ``strict'' in the statement
although it was clearly intended.
\end{reference}
Let $A$ be a Noetherian ring. Let $I$ be an ideal of $A$.
Let $M$ be a finite $A$-module. For every $p > 0$ there exists a $c > 0$
such that $\text{Tor}_p^A(M, A/I^n) \to \text{Tor}_p^A(M, A/I^{n - c})$
is zero for all $n \geq c$.
\end{lemma}

\begin{proof}
Proof for $p = 1$. Choose a short exact sequence
$0 \to K \to A^{\oplus t} \to M \to 0$. Then
$\text{Tor}_1^A(M, A/I^n) = K \cap (I^n)^{\oplus t}/I^nK$.
By Artin-Rees (Algebra, Lemma \ref{algebra-lemma-Artin-Rees})
there is a constant $c \geq 0$ such that
$K \cap (I^n)^{\oplus t} \subset I^{n - c}K$ for $n \geq c$. Thus the result
for $p = 1$. For $p > 1$ we have
$\text{Tor}_p^A(M, A/I^n) = \text{Tor}^A_{p - 1}(K, A/I^n)$.
Thus the lemma follows by induction.
\end{proof}

\begin{lemma}
\label{lemma-limit-flat}
Let $A$ be a Noetherian ring. Let $I$ be an ideal of $A$. Let
$(M_n)$ be an inverse system of $A$-modules such that
\begin{enumerate}
\item $M_n$ is a flat $A/I^n$-module,
\item $M_{n + 1} \to M_n$ is surjective.
\end{enumerate}
Then $M = \lim M_n$ is a flat $A$-module and
$Q \otimes_A M = \lim Q \otimes_A M_n$ for every finite $A$-module $Q$.
\end{lemma}

\begin{proof}
We first show that $Q \otimes_A M = \lim Q \otimes_A M_n$ for every finite
$A$-module $Q$. Choose a resolution $F_2 \to F_1 \to F_0 \to Q \to 0$
by finite free $A$-modules $F_i$. Then
$$
F_2 \otimes_A M_n \to F_1 \otimes_A M_n \to F_0 \otimes_A M_n
$$
is a chain complex whose homology in degree $0$ is $Q \otimes_A M_n$
and whose homology in degree $1$ is
$$
\text{Tor}_1^A(Q, M_n) = \text{Tor}_1^A(Q, A/I^n) \otimes_{A/I^n} M_n
$$
as $M_n$ is flat over $A/I^n$. By Lemma \ref{lemma-tor-strictly-pro-zero}
we see that this system is essentially constant (with value $0$).
It follows from Homology, Lemma \ref{homology-lemma-apply-Mittag-Leffler-again}
that $\lim Q \otimes_A M_n =
\Coker(\lim F_1 \otimes_A M_n \to \lim F_0 \otimes_A M_n)$.
Since $F_i$ is finite free this equals
$\Coker(F_1 \otimes_A M \to F_0 \otimes_A M) = Q \otimes_A M$.

\medskip\noindent
Next, let $Q \to Q'$ be an injective map of finite $A$-modules.
We have to show that $Q \otimes_A M \to Q' \otimes_A M$ is injective
(Algebra, Lemma \ref{algebra-lemma-flat}). By the above we see
$$
\Ker(Q \otimes_A M \to Q' \otimes_A M) =
\Ker(\lim Q \otimes_A M_n \to \lim Q' \otimes_A M_n).
$$
For each $n$ we have an exact sequence
$$
\text{Tor}_1^A(Q', M_n) \to \text{Tor}_1^A(Q'', M_n) \to
Q \otimes_A M_n \to Q' \otimes_A M_n
$$
where $Q'' = \Coker(Q \to Q')$. Above we have seen that the
inverse systems of Tor's are essentially constant with value $0$.
It follows from
Homology, Lemma \ref{homology-lemma-apply-Mittag-Leffler-again}
that the inverse limit of the right most maps is injective.
\end{proof}

\begin{lemma}
\label{lemma-flat-after-completion}
Let $R$ be a ring. Let $I \subset R$ be an ideal. Let $M$ be
an $R$-module. Assume
\begin{enumerate}
\item $I$ is finitely generated,
\item $R/I$ is Noetherian,
\item $M/IM$ is flat over $R/I$,
\item $\text{Tor}_1^R(M, R/I) = 0$.
\end{enumerate}
Then the $I$-adic completion $R^\wedge$
is a Noetherian ring and $M^\wedge$ is flat over $R^\wedge$.
\end{lemma}

\begin{proof}
By Algebra, Lemma \ref{algebra-lemma-what-does-it-mean}
the modules $M/I^nM$ are flat over $R/I^n$ for all $n$.
By Algebra, Lemma \ref{algebra-lemma-hathat-finitely-generated} we have
(a) $R^\wedge$ and $M^\wedge$ are $I$-adically complete and
(b) $R/I^n = R^\wedge/I^nR^\wedge$ for all $n$.
By Algebra, Lemma \ref{algebra-lemma-completion-Noetherian}
the ring $R^\wedge$ is Noetherian.
Applying Lemma \ref{lemma-limit-flat} we conclude that
$M^\wedge = \lim M/I^nM$ is flat as an $R^\wedge$-module.
\end{proof}








\section{The Koszul complex}
\label{section-koszul}

\noindent
We define the Koszul complex as follows.

\begin{definition}
\label{definition-koszul}
Let $R$ be a ring. Let $\varphi : E \to R$ be an $R$-module map. The
{\it Koszul complex} $K_\bullet(\varphi)$ associated to $\varphi$
is the commutative differential graded algebra defined as follows:
\begin{enumerate}
\item the underlying graded algebra is the exterior algebra
$K_\bullet(\varphi) = \wedge(E)$,
\item the differential $d : K_\bullet(\varphi) \to K_\bullet(\varphi)$
is the unique derivation such that $d(e) = \varphi(e)$ for all
$e \in E = K_1(\varphi)$.
\end{enumerate}
\end{definition}

\noindent
Explicitly, if $e_1 \wedge \ldots \wedge e_n$ is one of the generators of
degree $n$ in $K_\bullet(\varphi)$, then
$$
d(e_1 \wedge \ldots \wedge e_n) =
\sum\nolimits_{i = 1, \ldots, n} (-1)^{i + 1}
\varphi(e_i)e_1 \wedge \ldots \wedge \widehat{e_i} \wedge \ldots \wedge e_n.
$$
It is straightforward to see that this gives a well defined derivation
on the tensor algebra, which annihilates $e \otimes e$ and hence factors
through the exterior algebra.

\medskip\noindent
We often assume that $E$ is a finite free module, say $E = R^{\oplus n}$.
In this case the map $\varphi$ is given by a sequence of elements
$f_1, \ldots, f_n \in R$.

\begin{definition}
\label{definition-koszul-complex}
Let $R$ be a ring and let $f_1, \ldots, f_r \in R$. The
{\it Koszul complex on $f_1, \ldots, f_r$} is the Koszul complex
associated to the map $(f_1, \ldots, f_r) : R^{\oplus r} \to R$.
Notation $K_\bullet(f_\bullet)$, $K_\bullet(f_1, \ldots, f_r)$,
$K_\bullet(R, f_1, \ldots, f_r)$, or $K_\bullet(R, f_\bullet)$.
\end{definition}

\noindent
Of course, if $E$ is finite locally free, then $K_\bullet(\varphi)$ is
locally on $\Spec(R)$ isomorphic to a Koszul complex
$K_\bullet(f_1, \ldots, f_r)$.
This complex has many interesting formal properties.

\begin{lemma}
\label{lemma-functorial}
Let $\varphi : E \to R$ and $\varphi' : E' \to R$ be $R$-module maps.
Let $\psi : E \to E'$ be an $R$-module map such that
$\varphi' \circ \psi = \varphi$. Then $\psi$ induces a
homomorphism of differential graded algebras
$K_\bullet(\varphi) \to K_\bullet(\varphi')$.
\end{lemma}

\begin{proof}
This is immediate from the definitions.
\end{proof}

\begin{lemma}
\label{lemma-change-basis}
Let $f_1, \ldots, f_r \in R$ be a sequence.
Let $(x_{ij})$ be an invertible $r \times r$-matrix with
coefficients in $R$. Then the complexes
$K_\bullet(f_\bullet)$ and
$$
K_\bullet(\sum x_{1j}f_j, \sum x_{2j}f_j, \ldots, \sum x_{rj}f_j)
$$
are isomorphic.
\end{lemma}

\begin{proof}
Set $g_i = \sum x_{ij}f_j$.
The matrix $(x_{ji})$ gives an isomorphism $x : R^{\oplus r} \to R^{\oplus r}$
such that $(g_1, \ldots, g_r) = (f_1, \ldots, f_r) \circ x$.
Hence this follows from the functoriality of the Koszul complex
described in
Lemma \ref{lemma-functorial}.
\end{proof}

\begin{lemma}
\label{lemma-homotopy-koszul-abstract}
Let $R$ be a ring. Let $\varphi : E \to R$ be an $R$-module map.
Let $e \in E$ with image $f = \varphi(e)$ in $R$. Then
$$
f = de + ed
$$
as endomorphisms of $K_\bullet(\varphi)$.
\end{lemma}

\begin{proof}
This is true because $d(ea) = d(e)a - ed(a) = fa - ed(a)$.
\end{proof}

\begin{lemma}
\label{lemma-homotopy-koszul}
Let $R$ be a ring. Let $f_1, \ldots, f_r \in R$ be a sequence.
Multiplication by $f_i$ on $K_\bullet(f_\bullet)$ is homotopic to
zero, and in particular the cohomology modules $H_i(K_\bullet(f_\bullet))$
are annihilated by the ideal $(f_1, \ldots, f_r)$.
\end{lemma}

\begin{proof}
Special case of
Lemma \ref{lemma-homotopy-koszul-abstract}.
\end{proof}

\noindent
In
Derived Categories, Section \ref{derived-section-cones}
we defined the cone of a morphism of cochain complexes.
The cone $C(f)_\bullet$ of a morphism of chain complexes
$f : A_\bullet \to B_\bullet$ is the complex $C(f)_\bullet$ given by
$C(f)_n = B_n \oplus A_{n - 1}$ and differential
\begin{equation}
\label{equation-differential-cone}
d_{C(f), n} =
\left(
\begin{matrix}
d_{B, n} & f_{n - 1} \\
0 & -d_{A, n - 1}
\end{matrix}
\right)
\end{equation}
It comes equipped with canonical morphisms of complexes
$i : B_\bullet \to C(f)_\bullet$ and
$p : C(f)_\bullet \to A_\bullet[-1]$
induced by the obvious maps $B_n \to C(f)_n \to A_{n - 1}$.

\begin{lemma}
\label{lemma-cone-koszul-abstract}
Let $R$ be a ring. Let $\varphi : E \to R$ be an $R$-module map.
Let $f \in R$. Set $E' = E \oplus R$ and define $\varphi' : E' \to R$
by $\varphi$ on $E$ and multiplication by $f$ on $R$.
The complex $K_\bullet(\varphi')$ is isomorphic to the
cone of the map of complexes
$$
f :
K_\bullet(\varphi)
\longrightarrow
K_\bullet(\varphi).
$$
\end{lemma}

\begin{proof}
Denote $e_0 \in E'$ the element $1 \in R \subset R \oplus E$.
By our definition of the cone above we see that
$$
C(f)_n = K_n(\varphi) \oplus K_{n - 1}(\varphi) =
\wedge^n(E) \oplus \wedge^{n - 1}(E) = \wedge^n(E')
$$
where in the last $=$ we map $(0, e_1 \wedge \ldots \wedge e_{n - 1})$
to $e_0 \wedge e_1 \wedge \ldots \wedge e_{n - 1}$ in $\wedge^n(E')$.
A computation shows that this isomorphism is compatible with
differentials. Namely, this is clear for elements of the first
summand as $\varphi'|_E = \varphi$ and $d_{C(f)}$ restricted to
the first summand is just $d_{K_\bullet(\varphi)}$.
On the other hand, if $e_1 \wedge \ldots \wedge e_{n - 1}$
is in the second summand, then
$$
d_{C(f)}(0, e_1 \wedge \ldots \wedge e_{n - 1}) =
fe_1 \wedge \ldots \wedge e_{n - 1}
- d_{K_\bullet(\varphi)}(e_1 \wedge \ldots \wedge e_{n - 1})
$$
and on the other hand
\begin{align*}
& d_{K_\bullet(\varphi')}(0, e_0 \wedge e_1 \wedge \ldots \wedge e_{n - 1}) \\
& =
\sum\nolimits_{i = 0, \ldots, n - 1}
(-1)^i \varphi'(e_i)e_0 \wedge \ldots \wedge \widehat{e_i}
\wedge \ldots \wedge e_{n - 1} \\
& =
fe_1 \wedge \ldots \wedge e_{n - 1} +
\sum\nolimits_{i = 1, \ldots, n - 1}
(-1)^i \varphi(e_i)e_0 \wedge \ldots \wedge \widehat{e_i}
\wedge \ldots \wedge e_{n - 1} \\
& =
fe_1 \wedge \ldots \wedge e_{n - 1} -
e_0 \left(\sum\nolimits_{i = 1, \ldots, n - 1}
(-1)^{i + 1} \varphi(e_i)e_1 \wedge \ldots \wedge \widehat{e_i}
\wedge \ldots \wedge e_{n - 1}\right)
\end{align*}
which is the image of the result of the previous computation.
\end{proof}

\begin{lemma}
\label{lemma-cone-koszul}
Let $R$ be a ring. Let $f_1, \ldots, f_r$ be a sequence of elements
of $R$. The complex $K_\bullet(f_1, \ldots, f_r)$ is isomorphic to the
cone of the map of complexes
$$
f_r :
K_\bullet(f_1, \ldots, f_{r - 1})
\longrightarrow
K_\bullet(f_1, \ldots, f_{r - 1}).
$$
\end{lemma}

\begin{proof}
Special case of
Lemma \ref{lemma-cone-koszul-abstract}.
\end{proof}

\begin{lemma}
\label{lemma-cone-squared}
Let $R$ be a ring. Let $A_\bullet$ be a complex of $R$-modules.
Let $f, g \in R$. Let $C(f)_\bullet$ be the cone of
$f : A_\bullet \to A_\bullet$. Define similarly $C(g)_\bullet$ and
$C(fg)_\bullet$. Then $C(fg)_\bullet$ is homotopy equivalent to the
cone of a map
$$
C(f)_\bullet[1] \longrightarrow C(g)_\bullet
$$
\end{lemma}

\begin{proof}
We first prove this if $A_\bullet$ is the complex consisting of $R$ placed
in degree $0$. In this case the complex $C(f)_\bullet$ is the
complex
$$
\ldots \to 0 \to R \xrightarrow{f} R \to 0 \to \ldots
$$
with $R$ placed in (homological) degrees $1$ and $0$. The map
of complexes we use is
$$
\xymatrix{
0 \ar[r] \ar[d] &
0 \ar[r] \ar[d] &
R \ar[r]^f \ar[d]^1 &
R \ar[r] \ar[d] & 0 \ar[d] \\
0 \ar[r] & R \ar[r]^g & R \ar[r] & 0 \ar[r] & 0
}
$$
The cone of this is the chain complex consisting of $R^{\oplus 2}$ placed in
degrees $1$ and $0$ and differential (\ref{equation-differential-cone})
$$
\left(
\begin{matrix}
g & 1 \\
0 & -f
\end{matrix}
\right) :
R^{\oplus 2} \longrightarrow R^{\oplus 2}
$$
To see this chain complex is homotopic to
$C(fg)_\bullet$, i.e., to $R \xrightarrow{fg} R$,
consider the maps of complexes
$$
\xymatrix{
R \ar[d]_{(1, -g)} \ar[r]_{fg} & R \ar[d]^{(0, 1)} \\
R^{\oplus 2} \ar[r] & R^{\oplus 2}
}
\quad\quad
\xymatrix{
R^{\oplus 2} \ar[d]_{(1, 0)} \ar[r] & R^{\oplus 2} \ar[d]^{(f, 1)} \\
R \ar[r]^{fg} & R
}
$$
with obvious notation.
The composition of these two maps in one direction is the
identity on $C(fg)_\bullet$, but in the other direction
it isn't the identity. We omit writing out the required homotopy.

\medskip\noindent
To see the result holds in general, we use that we have a functor
$K_\bullet \mapsto \text{Tot}(A_\bullet \otimes_R K_\bullet)$
on the category of complexes which is compatible with homotopies
and cones. Then we write $C(f)_\bullet$ and $C(g)_\bullet$
as the total complex of the double complexes
$$
(R \xrightarrow{f} R) \otimes_R A_\bullet
\quad\text{and}\quad
(R \xrightarrow{g} R) \otimes_R A_\bullet
$$
and in this way we deduce the result from the special case discussed above.
Some details omitted.
\end{proof}

\begin{lemma}
\label{lemma-koszul-mult-abstract}
Let $R$ be a ring. Let $\varphi : E \to R$ be an $R$-module map.
Let $f, g \in R$. Set $E' = E \oplus R$ and define
$\varphi'_f, \varphi'_g, \varphi'_{fg} : E' \to R$
by $\varphi$ on $E$ and multiplication by $f, g, fg$ on $R$.
The complex $K_\bullet(\varphi'_{fg})$ is homotopy equivalent to the
cone of a map of complexes
$$
K_\bullet(\varphi'_f)[1]
\longrightarrow
K_\bullet(\varphi'_g).
$$
\end{lemma}

\begin{proof}
By
Lemma \ref{lemma-cone-koszul-abstract}
the complex $K_\bullet(\varphi'_f)$ is isomorphic to the cone of
multiplication by $f$ on $K_\bullet(\varphi)$ and similarly
for the other two cases. Hence the lemma follows from
Lemma \ref{lemma-cone-squared}.
\end{proof}

\begin{lemma}
\label{lemma-koszul-mult}
Let $R$ be a ring. Let $f_1, \ldots, f_{r - 1}$ be a sequence of elements
of $R$. Let $f, g \in R$. The complex
$K_\bullet(f_1, \ldots, f_{r - 1}, fg)$
is homotopy equivalent to the cone of a map of complexes
$$
K_\bullet(f_1, \ldots, f_{r - 1}, f)[1]
\longrightarrow
K_\bullet(f_1, \ldots, f_{r - 1}, g)
$$
\end{lemma}

\begin{proof}
Special case of
Lemma \ref{lemma-koszul-mult-abstract}.
\end{proof}

\begin{lemma}
\label{lemma-join-sequences-koszul-complex}
Let $R$ be a ring.
Let $f_1, \ldots, f_r$, $g_1, \ldots, g_s$ be elements of $R$.
Then there is an isomorphism of Koszul complexes
$$
K_\bullet(R, f_1, \ldots, f_r, g_1, \ldots, g_s) =
\text{Tot}(K_\bullet(R, f_1, \ldots, f_r) \otimes_R
K_\bullet(R, g_1, \ldots, g_s)).
$$
\end{lemma}

\begin{proof}
Omitted. Hint: If $K_\bullet(R, f_1, \ldots, f_r)$ is generated as a
differential graded algebra by $x_1, \ldots, x_r$ with $\text{d}(x_i) = f_i$
and $K_\bullet(R, g_1, \ldots, g_s)$ is generated as a
differential graded algebra by $y_1, \ldots, y_s$ with $\text{d}(y_j) = g_j$,
then we can think of $K_\bullet(R, f_1, \ldots, f_r, g_1, \ldots, g_s)$
as the differential graded algebra generated by the sequence of elements
$x_1, \ldots, x_r, y_1, \ldots, y_s$ with $\text{d}(x_i) = f_i$
and $\text{d}(y_j) = g_j$.
\end{proof}





\section{The extended alternating {\v C}ech complex}
\label{section-alternating-cech}

\noindent
Let $R$ be a ring. Let $f_1, \ldots, f_r \in R$. The extended alternating
{\v C}ech complex of $R$ is the cochain complex
$$
R \to \bigoplus\nolimits_{i_0} R_{f_{i_0}} \to
\bigoplus\nolimits_{i_0 < i_1} R_{f_{i_0}f_{i_1}} \to
\ldots \to R_{f_1\ldots f_r}
$$
where $R$ is in degree $0$, the term $\bigoplus_{i_0} R_{f_{i_0}}$
is in degre $1$, and so on. The maps are defined as follows
\begin{enumerate}
\item The map $R \to \bigoplus\nolimits_{i_0} R_{f_{i_0}}$
is given by the canonical maps $R \to R_{f_{i_0}}$.
\item Given $1 \leq i_0 < \ldots < i_{p + 1} \leq r$ and
$0 \leq j \leq p + 1$ we have the canonical localization map
$$
R_{f_{i_0} \ldots \hat f_{i_j} \ldots f_{i_{p + 1}}} \to
R_{f_{i_0} \ldots f_{i_{p + 1}}}
$$
\item The differentials use the canonical maps of (2) with sign $(-1)^j$.
\end{enumerate}
If $M$ is any $R$-module, the extended alternating {\v C}ech complex
of $M$ is the similarly constructed cochain complex
$$
M \to \bigoplus\nolimits_{i_0} M_{f_{i_0}} \to
\bigoplus\nolimits_{i_0 < i_1} M_{f_{i_0}f_{i_1}} \to
\ldots \to M_{f_1\ldots f_r}
$$
where $M$ is in degree $0$ as before.

\begin{lemma}
\label{lemma-extended-alternating-is-complex}
The extended alternating {\v C}ech complexes defined above
are complexes of $R$-modules.
\end{lemma}

\begin{proof}
Omitted.
\end{proof}

\begin{lemma}
\label{lemma-extended-alternating-form-module}
Let $R$ be a ring. Let $f_1, \ldots, f_r \in R$. Let $M$ be an $R$-module.
The extended alternating {\v C}ech complex of $M$ is the tensor product
over $R$ of $M$ with the extended alternating {\v C}ech complex of $R$.
\end{lemma}

\begin{proof}
Omitted.
\end{proof}

\begin{lemma}
\label{lemma-extended-alternating-base-change}
Let $R$ be a ring. Let $f_1, \ldots, f_r \in R$. Let $M$ be an $R$-module.
Let $R \to S$ be a ring map, denote $g_1, \ldots, g_r \in S$ the images
of $f_1, \ldots, f_r$, and set $N = M \otimes_R S$.
The extended alternating {\v C}ech complex constructed using
$S$, $g_1, \ldots, g_r$, and $N$ is the tensor product of the
extended alternating {\v C}ech complex of $M$ with $S$ over $R$.
\end{lemma}

\begin{proof}
Omitted.
\end{proof}

\begin{lemma}
\label{lemma-extended-alternating-homotopy-zero}
Let $R$ be a ring. Let $f_1, \ldots, f_r \in R$. Let $M$ be an $R$-module.
If there exists an $i \in \{1, \ldots, r\}$ such that $f_i$ is a unit, then
the extended alternating {\v C}ech
complex of $M$ is homotopy equivalent to $0$.
\end{lemma}

\begin{proof}
We will use the following notation: a cochain $x$ of degree $p + 1$
in the extended alternating {\v C}ech complex of $M$ is
$x = (x_{i_0 \ldots i_p})$ where $x_{i_0 \ldots i_p}$ is in
$M_{f_{i_0} \ldots f_{i_p}}$. With this notation we have
$$
d(x)_{i_0 \ldots i_{p + 1}} =
\sum\nolimits_j (-1)^j x_{i_0 \ldots \hat i_j \ldots i_{p + 1}}
$$
As homotopy we use the maps
$$
h : \text{cochains of degree }p + 2 \to \text{cochains of degree }p + 1
$$
given by the rule
$$
h(x)_{i_0 \ldots i_p} = 0 \text{ if } i \in \{i_0, \ldots, i_p\}
\text{ and }
h(x)_{i_0 \ldots i_p} = 
(-1)^j x_{i_0 \ldots i_j i i_{j + 1} \ldots i_p} \text{ if not}
$$
Here $j$ is the unique index such that $i_j < i < i_{j + 1}$ in the
second case; also, since $f_i$ is a unit we have the equality
$$
M_{f_{i_0} \ldots f_{i_p}} =
M_{f_{i_0} \ldots f_{i_j} f_i f_{i_{j + 1}} \ldots f_{i_p}}
$$
which we can use to make sense of thinking of
$(-1)^j x_{i_0 \ldots i_j i i_{j + 1} \ldots i_p}$
as an element of $M_{f_{i_0} \ldots f_{i_p}}$.
We will show by a computation that $d h + h d$ equals
the negative of the identity map which finishes the proof.
To do this fix $x$ a cochain of degree $p + 1$ and let
$1 \leq i_0 < \ldots < i_p \leq r$.

\medskip\noindent
Case I: $i \in \{i_0, \ldots, i_p\}$. Say $i = i_t$. Then we have
$h(d(x))_{i_0 \ldots i_p} = 0$. On the other hand we have
$$
d(h(x))_{i_0 \ldots i_p} =
\sum (-1)^j h(x)_{i_0 \ldots \hat i_j \ldots i_p} =
(-1)^t h(x)_{i_0 \ldots \hat i \ldots i_p} =
(-1)^t (-1)^{t - 1} x_{i_0 \ldots i_p}
$$
Thus $(dh + hd)(x)_{i_0 \ldots i_p} = -x_{i_0 \ldots i_p}$ as desired.

\medskip\noindent
Case II: $i \not \in \{i_0, \ldots, i_p\}$. Let $j$ be such that
$i_j < i < i_{j + 1}$. Then we see that
\begin{align*}
h(d(x))_{i_0 \ldots i_p}
& =
(-1)^j d(x)_{i_0 \ldots i_j i i_{j + 1} \ldots i_p} \\
& =
\sum\nolimits_{j' \leq j} (-1)^{j + j'}
x_{i_0 \ldots \hat i_{j'} \ldots i_j i i_{j + 1} \ldots i_p} -
x_{i_0 \ldots i_p} \\
&
+ \sum\nolimits_{j' > j} (-1)^{j + j' + 1}
x_{i_0 \ldots i_j i i_{j + 1} \ldots \hat i_{j'} \ldots i_p}
\end{align*}
On the other hand we have
\begin{align*}
d(h(x))_{i_0 \ldots i_p}
& =
\sum\nolimits_{j'} (-1)^{j'} h(x)_{i_0 \ldots \hat i_{j'} \ldots i_p} \\
& =
\sum\nolimits_{j' \leq j} (-1)^{j' + j - 1}
x_{i_0 \ldots \hat i_{j'} \ldots i_j i i_{j + 1} \ldots i_p} \\
& +
\sum\nolimits_{j' > j} (-1)^{j' + j}
x_{i_0 \ldots i_j i i_{j + 1} \ldots \hat i_{j'} \ldots i_p}
\end{align*}
Adding these up we obtain
$(dh + hd)(x)_{i_0 \ldots i_p} = - x_{i_0 \ldots i_p}$
as desired.
\end{proof}

\begin{lemma}
\label{lemma-extended-alternating-torsion}
Let $R$ be a ring. Let $f_1, \ldots, f_r \in R$. Let $M$ be an $R$-module.
Let $H^q$ be the $q$th cohomology module of the extended alternation
{\v C}ech complex of $M$. Then
\begin{enumerate}
\item $H^q = 0$ if $q \not \in [0, r]$,
\item for $x \in H^i$ there exists an $n \geq 1$ such that $f_i^n x = 0$
for $i = 1, \ldots, r$,
\item the support of $H^q$ is contained in $V(f_1, \ldots, f_r)$,
\item if there is an $f \in (f_1, \ldots, f_r)$ which acts invertibly
on $M$, then $H^q = 0$.
\end{enumerate}
\end{lemma}

\begin{proof}
Part (1) follows from the fact that the extended alternating {\v C}ech complex
is zero in degrees $< 0$ and $> r$. To prove (2) it suffices to show that
for each $i$ there exists an $n \geq 1$ such that $f_i^n x = 0$. To see
this it suffices to show that $(H^q)_{f_i} = 0$. Since localization is
exact, $(H^q)_{f_i}$ is the $q$th cohomology module of the localization
of the extended alternating complex of $M$ at $f_i$. By
Lemma \ref{lemma-extended-alternating-base-change}
this localization is the extended alternating {\v C}ech complex
of $M_{f_i}$ over $R_{f_i}$ with respect to the images of
$f_1, \ldots, f_r$ in $R_{f_i}$. Thus
we reduce to showing that $H^q$ is zero if $f_i$ is invertible, which
follows from Lemma \ref{lemma-extended-alternating-homotopy-zero}.
Part (3) follows from the observation that $(H^q)_{f_i} = 0$
for all $i$ that we just proved. To see part (4) note that in this
case $f$ acts invertibly on $H^q$ and $H^q$ is supported on $V(f)$
by (3). This forces $H^q$ to be zero (small detail omitted).
\end{proof}

\begin{lemma}
\label{lemma-extended-alternating-Cech-is-colimit-koszul}
Let $R$ be a ring. Let $f_1, \ldots, f_r \in R$. The extended alternating
{\v C}ech complex
$$
R \to \bigoplus\nolimits_{i_0} R_{f_{i_0}} \to
\bigoplus\nolimits_{i_0 < i_1} R_{f_{i_0}f_{i_1}} \to
\ldots \to R_{f_1\ldots f_r}
$$
is a colimit of the Koszul complexes $K(R, f_1^n, \ldots, f_r^n)$; see
proof for a precise statement.
\end{lemma}

\begin{proof}
We urge the reader to prove this for themselves.
Denote $K(R, f_1^n, \ldots, f_r^n)$ the Koszul complex of
Definition \ref{definition-koszul-complex} viewed as a cochain complex
sitting in degrees $0, \ldots, r$. Thus we have
$$
K(R, f_1^n, \ldots, f_r^n) :
0 \to \wedge^r(R^{\oplus r}) \to
\wedge^{r - 1}(R^{\oplus r}) \to \ldots \to
R^{\oplus r} \to R \to 0
$$
with the term $\wedge^r(R^{\oplus r})$ sitting in degree $0$.
Let $e^n_1, \ldots, e^n_r$ be the standard basis of $R^{\oplus r}$.
Then the elements $e^n_{j_1} \wedge \ldots \wedge e^n_{j_{r - p}}$ for
$1 \leq j_1 < \ldots < j_{r - p} \leq r$ form a basis for the term in
degree $p$ of the Koszul complex. Further, observe that
$$
d(e^n_{j_1} \wedge \ldots \wedge e^n_{j_{r - p}}) =
\sum (-1)^{a + 1} f_{j_a}^n e^n_{j_1} \wedge \ldots \wedge \hat e^n_{j_a} \wedge
\ldots \wedge e^n_{j_{r - p}}
$$
by our construction of the Koszul complex in Section \ref{section-koszul}.
The transition maps of our system
$$
K(R, f_1^n, \ldots, f_r^n) \to K(R, f_1^{n + 1}, \ldots, f_r^{n + 1})
$$
are given by the rule
$$
e^n_{j_1} \wedge \ldots \wedge e^n_{j_{r - p}}
\longmapsto
f_{i_0} \ldots f_{i_{p - 1}}
e^{n + 1}_{j_1} \wedge \ldots \wedge e^{n + 1}_{j_{r - p}}
$$
where the indices $1 \leq i_0 < \ldots < i_{p - 1} \leq r$ are such that
$\{1, \ldots r\} =
\{i_0, \ldots, i_{p - 1}\} \amalg \{j_1, \ldots, j_{r - p}\}$.
We omit the short computation that shows this is compatible with differentials.
Observe that the transition maps are always $1$ in degree $0$
and equal to $f_1 \ldots f_r$ in degree $r$.

\medskip\noindent
Denote $K^p(R, f_1^n, \ldots, f_r^n)$ the term of degree $p$ in the Koszul
complex. Observe that for any $f \in R$ we have
$$
R_f = \colim (R \xrightarrow{f} R \xrightarrow{f} R \to \ldots )
$$
Hence we see that in degree $p$ we obtain
$$
\colim K^p(R, f_1^n, \ldots f_r^n) =
\bigoplus\nolimits_{1 \leq i_0 < \ldots < i_{p - 1} \leq r}
R_{f_{i_0} \ldots f_{i_{p - 1}}}
$$
Here the element $e^n_{j_1} \wedge \ldots \wedge e^n_{j_{r - p}}$
of the Koszul complex above maps in the colimit
to the element $(f_{i_0} \ldots f_{i_{p - 1}})^{-n}$ in the
summand $R_{f_{i_0} \ldots f_{i_{p - 1}}}$
where the indices are chosen such that
$\{1, \ldots r\} = \{i_0, \ldots, i_{p - 1}\}
\amalg \{j_1, \ldots, j_{r - p}\}$.
Thus the differential on this complex is given by
$$
d(1\text{ in }R_{f_{i_0} \ldots f_{i_{p - 1}}}) =
\sum\nolimits_{i \not \in \{i_0, \ldots, i_{p - 1}\}}
(-1)^{i - t}\text{ in }
R_{f_{i_0} \ldots f_{i_t} f_i f_{i_{t + 1}} \ldots f_{i_{p - 1}}}
$$
Thus if we consider the map of complexes given in degree $p$
by the map
$$
\bigoplus\nolimits_{1 \leq i_0 < \ldots < i_{p - 1} \leq r}
R_{f_{i_0} \ldots f_{i_{p - 1}}}
\longrightarrow
\bigoplus\nolimits_{1 \leq i_0 < \ldots < i_{p - 1} \leq r}
R_{f_{i_0} \ldots f_{i_{p - 1}}}
$$
determined by the rule
$$
1\text{ in }R_{f_{i_0} \ldots f_{i_{p - 1}}}
\longmapsto
(-1)^{i_0 + \ldots + i_{p - 1} + p}\text{ in }R_{f_{i_0} \ldots f_{i_{p - 1}}}
$$
then we get an isomorphism of complexes from
$\colim K(R, f_1^n, \ldots, f_r^n)$ to the
extended alternating {\v C}ech complex defined in this section.
We omit the verification that the signs work out.
\end{proof}








\section{Koszul regular sequences}
\label{section-koszul-regular}

\noindent
Please take a look at
Algebra, Sections \ref{algebra-section-regular-sequences},
\ref{algebra-section-quasi-regular}, and
\ref{algebra-section-depth}
before looking at this one.

\begin{definition}
\label{definition-koszul-regular-sequence}
Let $R$ be a ring. Let $r \geq 0$ and let $f_1, \ldots, f_r \in R$
be a sequence of elements. Let $M$ be an $R$-module.
The sequence $f_1, \ldots, f_r$ is called
\begin{enumerate}
\item {\it $M$-Koszul-regular} if
$H_i(K_\bullet(f_1, \ldots, f_r) \otimes_R M) = 0$ for
all $i \not = 0$,
\item {\it $M$-$H_1$-regular} if
$H_1(K_\bullet(f_1, \ldots, f_r) \otimes_R M) = 0$,
\item {\it Koszul-regular} if $H_i(K_\bullet(f_1, \ldots, f_r)) = 0$ for
all $i \not = 0$, and
\item {\it $H_1$-regular} if $H_1(K_\bullet(f_1, \ldots, f_r)) = 0$.
\end{enumerate}
\end{definition}

\noindent
We will see in Lemmas \ref{lemma-regular-koszul-regular},
\ref{lemma-koszul-regular-H1-regular}, and
\ref{lemma-H1-regular-quasi-regular} that for elements
$f_1, \ldots, f_r$ of a ring $R$ we have the following implications
\begin{align*}
f_1, \ldots, f_r\text{ is a regular sequence}
& \Rightarrow f_1, \ldots, f_r\text{ is a Koszul-regular sequence} \\
& \Rightarrow f_1, \ldots, f_r\text{ is an }H_1\text{-regular sequence} \\
& \Rightarrow f_1, \ldots, f_r\text{ is a quasi-regular sequence.}
\end{align*}
In general none of these implications can be reversed, but if $R$ is
a Noetherian local ring and $f_1, \ldots, f_r \in \mathfrak m_R$,
then the four conditions are all equivalent
(Lemma \ref{lemma-noetherian-finite-all-equivalent}).
If $f = f_1 \in R$ is a length $1$ sequence and $f$ is not a unit of $R$
then it is clear that the following are all equivalent
\begin{enumerate}
\item $f$ is a regular sequence of length one,
\item $f$ is a Koszul-regular sequence of length one, and
\item $f$ is a $H_1$-regular sequence of length one.
\end{enumerate}
It is also clear that these imply that $f$ is a quasi-regular sequence
of length one. But there do exist quasi-regular sequences of length $1$
which are not regular sequences. Namely, let
$$
R = k[x, y_0, y_1, \ldots]/(xy_0, xy_1 - y_0, xy_2 - y_1, \ldots)
$$
and let $f$ be the image of $x$ in $R$. Then $f$ is a zerodivisor, but
$\bigoplus_{n \geq 0} (f^n)/(f^{n + 1}) \cong k[x]$ is a polynomial ring.

\begin{lemma}
\label{lemma-regular-koszul-regular}
Let $R$ be a ring, $M$ an $R$-module, and $f_1, \ldots, f_r \in R$
such that for $i = 1, \ldots, r$ multiplication by $f_i$
is injective on $M/(f_1, \ldots, f_{i - 1})M$. Then $f_1, \ldots, f_r$
is $M$-Koszul regular. In particular, an $M$-regular sequence
is $M$-Koszul-regular and any regular sequence is Koszul-regular.
\end{lemma}

\begin{proof}
Let $R$, $M$, $f_1, \ldots, f_r$ be as in the first sentence of the lemma.
If $r = 1$, it is immediate that $f_1$ is $M$-Koszul-regular.
Assume $r > 1$. Since $f_1$ is a nonzerodivisor on $M$, we obtain
a short exact sequence of complexes:
$$
0 \to K_\bullet(f_2, \ldots, f_r) \otimes M
\xrightarrow{f_1}
K_\bullet(f_2, \ldots, f_r) \otimes M \to
K_\bullet(\overline{f}_2, \ldots, \overline{f}_r) \otimes M/f_1M \to 0
$$
Here $\overline{f}_i$ is the image of $f_i$ in $R/(f_1)$. By
Lemma \ref{lemma-cone-koszul}
the complex $K_\bullet(f_1, \ldots, f_r)$
is isomorphic to the cone of multiplication by $f_1$
on $K_\bullet(f_2, \ldots, f_r)$. Thus
$K_\bullet(R, f_1, \ldots, f_r) \otimes M$ is isomorphic
to the cone on the first map. Hence
$K_\bullet(\overline{f}_2, \ldots, \overline{f}_r) \otimes M/f_1M$
is quasi-isomorphic to $K_\bullet(f_1, \ldots, f_r) \otimes M$.
As $R/(f_1)$, $M/f_1M$, $\overline{f}_2, \ldots, \overline{f}_r$
satisfy the conditions of the lemma, by induction we conclude
this complex is acyclic in postive degrees. This finishes the proof
of the first statement. The second statement immediately follows
from the first.
\end{proof}

\begin{lemma}
\label{lemma-koszul-regular-H1-regular}
A $M$-Koszul-regular sequence is $M$-$H_1$-regular.
A Koszul-regular sequence is $H_1$-regular.
\end{lemma}

\begin{proof}
This is immediate from the definition.
\end{proof}

\begin{lemma}
\label{lemma-mult-koszul-regular}
Let $f_1, \ldots, f_{r - 1} \in R$ be a sequence and $f, g \in R$.
Let $M$ be an $R$-module.
\begin{enumerate}
\item If $f_1, \ldots, f_{r - 1}, f$ and $f_1, \ldots, f_{r - 1}, g$
are $M$-$H_1$-regular then $f_1, \ldots, f_{r - 1}, fg$ is
$M$-$H_1$-regular too.
\item If $f_1, \ldots, f_{r - 1}, f$ and $f_1, \ldots, f_{r - 1}, g$
are $M$-Koszul-regular then $f_1, \ldots, f_{r - 1}, fg$ is
$M$-Koszul-regular too.
\end{enumerate}
\end{lemma}

\begin{proof}
By
Lemma \ref{lemma-koszul-mult}
we have exact sequences
$$
H_i(K_\bullet(f_1, \ldots, f_{r - 1}, f) \otimes M) \to
H_i(K_\bullet(f_1, \ldots, f_{r - 1}, fg) \otimes M) \to
H_i(K_\bullet(f_1, \ldots, f_{r - 1}, g) \otimes M)
$$
for all $i$.
\end{proof}

\begin{lemma}
\label{lemma-koszul-regular-flat-base-change}
Let $\varphi : R \to S$ be a flat ring map. Let $f_1, \ldots, f_r \in R$.
Let $M$ be an $R$-module and set $N = M \otimes_R S$.
\begin{enumerate}
\item If $f_1, \ldots, f_r$ in $R$ is an $M$-$H_1$-regular sequence, then
$\varphi(f_1), \ldots, \varphi(f_r)$ is an $N$-$H_1$-regular
sequence in $S$.
\item If $f_1, \ldots, f_r$ is an $M$-Koszul-regular sequence in $R$, then
$\varphi(f_1), \ldots, \varphi(f_r)$ is an $N$-Koszul-regular
sequence in $S$.
\end{enumerate}
\end{lemma}

\begin{proof}
This is true because
$K_\bullet(f_1, \ldots, f_r) \otimes_R S =
K_\bullet(\varphi(f_1), \ldots, \varphi(f_r))$
and therefore
$(K_\bullet(f_1, \ldots, f_r) \otimes_R M) \otimes_R S =
K_\bullet(\varphi(f_1), \ldots, \varphi(f_r)) \otimes_S N$.
\end{proof}

\begin{lemma}
\label{lemma-H1-regular-quasi-regular}
An $M$-$H_1$-regular sequence is $M$-quasi-regular.
\end{lemma}

\begin{proof}
Let $R$ be a ring and let $M$ be an $R$-module.
Let $f_1, \ldots, f_r$ be an $M$-$H_1$-regular sequence.
Denote $J = (f_1, \ldots, f_r)$. The assumption means that we have
an exact sequence
$$
\wedge^2(R^r) \otimes M \to R^{\oplus r} \otimes M \to JM \to 0
$$
where the first arrow is given by
$e_i \wedge e_j \otimes m \mapsto (f_ie_j - f_je_i) \otimes m$.
Tensoring the sequence with $R/J$ we see that
$$
JM/J^2M = (R/J)^{\oplus r} \otimes_R M = (M/JM)^{\oplus r}
$$
is a finite free module. To finish the proof we have to prove
for every $n \geq 2$ the following: if
$$
\xi = \sum\nolimits_{|I| = n, I = (i_1, \ldots, i_r)}
m_I f_1^{i_1} \ldots f_r^{i_r} \in J^{n + 1}M
$$
then $m_I \in JM$ for all $I$. In the next paragraph, we prove
$m_I \in JM$ for $I = (0, \ldots, 0, n)$ and in the last paragraph
we deduce the general case from this special case.

\medskip\noindent
Let $I = (0, \ldots, 0, n)$. Let $\xi$ be as above. We can write
$\xi = m_1 f_1 + \ldots + m_{r - 1}f_{r - 1} + m_I f_r^n$.
As we have assumed $\xi \in J^{n + 1}M$, we can also write
$\xi = \sum_{1 \leq i \leq j \leq r - 1} m_{ij}f_if_j +
\sum_{1 \leq i \leq r - 1}m'_i f_if_r^n + m'' f_r^{n + 1}$.
Then we see that
$$
\begin{matrix}
(m_1 - m_{11}f_1 - m'_1f_r^n)f_1 + \\
(m_2 - m_{12}f_1 - m_{22}f_2 - m'_2f_r^n)f_2 + \\
\ldots + \\
(m_{r - 1} - m_{1 r - 1}f_1 - \ldots - m_{r - 1 r - 1}f_{r - 1}
- m'_{r - 1}f_r^n)f_{r - 1} + \\
(m_I - m'' f_r)f_r^n = 0
\end{matrix}
$$
Since $f_1, \ldots, f_{r - 1}, f_r^n$ is $M$-$H_1$-regular by
Lemma \ref{lemma-mult-koszul-regular} we see that $m_I - m'' f_r$
is in the submodule $f_1M + \ldots + f_{r - 1}M + f_r^nM$.
Thus $m_I \in f_1M + \ldots + f_rM$.

\medskip\noindent
Let $S = R[x_1, x_2, \ldots, x_r, 1/x_r]$. The ring map $R \to S$ is faithfully
flat, hence $f_1, \ldots, f_r$ is an $M$-$H_1$-regular sequence in $S$, see
Lemma \ref{lemma-koszul-regular-flat-base-change}.
By
Lemma \ref{lemma-change-basis}
we see that
$$
g_1 = f_1 - \frac{x_1}{x_r} f_r,
\ \ldots,
\ g_{r - 1} = f_{r - 1} - \frac{x_{r - 1}}{x_r} f_r,
\ g_r = \frac{1}{x_r}f_r
$$
is an $M$-$H_1$-regular sequence in $S$. Finally, note that our element
$\xi$ can be rewritten
$$
\xi = \sum\nolimits_{|I| = n, I = (i_1, \ldots, i_r)}
m_I (g_1 + x_1 g_r)^{i_1} \ldots (g_{r - 1} + x_{r - 1} g_r)^{i_{r - 1}}
(x_rg_r)^{i_r}
$$
and the coefficient of $g_r^n$ in this expression is
$$
\sum m_I x_1^{i_1} \ldots x_r^{i_r}
$$
By the case discussed in the previous paragraph this sum is in
$J(M \otimes_R S)$. Since the monomials $x_1^{i_1} \ldots x_r^{i_r}$
form part of an $R$-basis of $S$ over $R$ we conclude that $m_I \in J$
for all $I$ as desired.
\end{proof}

\noindent
For nonzero finite modules over Noetherian local rings all of the types of
regular sequences introduced so far are equivalent.

\begin{lemma}
\label{lemma-noetherian-finite-all-equivalent}
Let $(R, \mathfrak m)$ be a Noetherian local ring. Let $M$ be a nonzero
finite $R$-module. Let $f_1, \ldots, f_r \in \mathfrak m$. The following
are equivalent
\begin{enumerate}
\item $f_1, \ldots, f_r$ is an $M$-regular sequence,
\item $f_1, \ldots, f_r$ is a $M$-Koszul-regular sequence,
\item $f_1, \ldots, f_r$ is an $M$-$H_1$-regular sequence,
\item $f_1, \ldots, f_r$ is an $M$-quasi-regular sequence.
\end{enumerate}
In particular the sequence $f_1, \ldots, f_r$ is a regular sequence
in $R$ if and only if it is a Koszul regular sequence, if and only if
it is a $H_1$-regular sequence, if and only if it is a quasi-regular sequence.
\end{lemma}

\begin{proof}
The implication (1) $\Rightarrow$ (2) is 
Lemma \ref{lemma-regular-koszul-regular}.
The implication (2) $\Rightarrow$ (3) is
Lemma \ref{lemma-koszul-regular-H1-regular}.
The implication (3) $\Rightarrow$ (4) is 
Lemma \ref{lemma-H1-regular-quasi-regular}.
The implication (4) $\Rightarrow$ (1) is
Algebra, Lemma \ref{algebra-lemma-quasi-regular-regular}.
\end{proof}

\begin{lemma}
\label{lemma-H1-regular-in-quotient}
Let $A$ be a ring. Let $I \subset A$ be an ideal.
Let $g_1, \ldots, g_m$ be a sequence in $A$ whose image in
$A/I$ is $H_1$-regular. Then $I \cap (g_1, \ldots, g_m) =
I(g_1, \ldots, g_m)$.
\end{lemma}

\begin{proof}
Consider the exact sequence of complexes
$$
0 \to I \otimes_A K_\bullet(A, g_1, \ldots, g_m)
\to K_\bullet(A, g_1, \ldots, g_m) \to
K_\bullet(A/I, g_1, \ldots, g_m) \to 0
$$
Since the complex on the right has $H_1 = 0$ by assumption we
see that
$$
\Coker(I^{\oplus m} \to I)
\longrightarrow
\Coker(A^{\oplus m} \to A)
$$
is injective. This is equivalent to the assertion of the lemma.
\end{proof}

\begin{lemma}
\label{lemma-conormal-sequence-H1-regular}
Let $A$ be a ring. Let $I \subset J \subset A$ be ideals.
Assume that $J/I \subset A/I$ is generated by an $H_1$-regular sequence.
Then $I \cap J^2 = IJ$.
\end{lemma}

\begin{proof}
To prove this choose $g_1, \ldots, g_m \in J$
whose images in $A/I$ form a $H_1$-regular sequence which generates $J/I$.
In particular $J = I + (g_1, \ldots, g_m)$.
Suppose that $x \in I \cap J^2$. Because $x \in J^2$ can write
$$
x =
\sum a_{ij} g_ig_j +
\sum a_j g_j +
a
$$
with $a_{ij} \in A$, $a_j \in I$ and $a \in I^2$.
Then $\sum a_{ij}g_ig_j \in I \cap (g_1, \ldots, g_m)$
hence by
Lemma \ref{lemma-H1-regular-in-quotient}
we see that $\sum a_{ij}g_ig_j \in I(g_1, \ldots, g_m)$.
Thus $x \in IJ$ as desired.
\end{proof}

\begin{lemma}
\label{lemma-join-quasi-regular-H1-regular}
Let $A$ be a ring. Let $I$ be an ideal generated by a quasi-regular
sequence $f_1, \ldots, f_n$ in $A$. Let $g_1, \ldots, g_m \in A$ be
elements whose images $\overline{g}_1, \ldots, \overline{g}_m$ form an
$H_1$-regular sequence in $A/I$. Then $f_1, \ldots, f_n, g_1, \ldots, g_m$
is a quasi-regular sequence in $A$.
\end{lemma}

\begin{proof}
We claim that $g_1, \ldots, g_m$ forms an $H_1$-regular sequence in
$A/I^d$ for every $d$. By induction assume that this holds in
$A/I^{d - 1}$. We have a short exact sequence of complexes
$$
0 \to K_\bullet(A, g_\bullet) \otimes_A I^{d - 1}/I^d
\to K_\bullet(A/I^d, g_\bullet) \to
K_\bullet(A/I^{d - 1}, g_\bullet) \to 0
$$
Since $f_1, \ldots, f_n$ is quasi-regular we see that the first complex
is a direct sum of copies of $K_\bullet(A/I, g_1, \ldots, g_m)$
hence acyclic in degree $1$. By induction hypothesis the last complex is
acyclic in degree $1$. Hence also the middle complex is.
In particular, the sequence $g_1, \ldots, g_m$ forms a quasi-regular
sequence in $A/I^d$ for every $d \geq 1$, see
Lemma \ref{lemma-H1-regular-quasi-regular}.
Now we are ready to prove that $f_1, \ldots, f_n, g_1, \ldots, g_m$
is a quasi-regular sequence in $A$.
Namely, set $J = (f_1, \ldots, f_n, g_1, \ldots, g_m)$ and suppose
that (with multinomial notation)
$$
\sum\nolimits_{|N| + |M| = d} a_{N, M} f^N g^M \in J^{d + 1}
$$
for some $a_{N, M} \in A$. We have to show that $a_{N, M} \in J$
for all $N, M$. Let $e \in \{0, 1, \ldots, d\}$. Then
$$
\sum\nolimits_{|N| = d - e, \ |M| = e} a_{N, M} f^N g^M \in
(g_1, \ldots, g_m)^{e + 1} + I^{d - e + 1}
$$
Because $g_1, \ldots, g_m$ is a quasi-regular sequence in $A/I^{d - e + 1}$
we deduce
$$
\sum\nolimits_{|N| = d - e} a_{N, M} f^N \in
(g_1, \ldots, g_m) + I^{d - e + 1}
$$
for each $M$ with $|M| = e$. By
Lemma \ref{lemma-H1-regular-in-quotient}
applied to $I^{d - e}/I^{d - e + 1}$ in the ring $A/I^{d - e + 1}$
this implies $\sum_{|N| = d - e} a_{N, M} f^N \in I^{d - e}(g_1, \ldots, g_m)$.
Since $f_1, \ldots, f_n$ is quasi-regular in $A$ this implies
that $a_{N, M} \in J$ for each $N, M$ with $|N| = d - e$ and $|M| = e$.
This proves the lemma.
\end{proof}

\begin{lemma}
\label{lemma-join-H1-regular-sequences}
Let $A$ be a ring. Let $I$ be an ideal generated by an
$H_1$-regular sequence $f_1, \ldots, f_n$ in $A$.
Let $g_1, \ldots, g_m \in A$ be elements whose images
$\overline{g}_1, \ldots, \overline{g}_m$ form an $H_1$-regular sequence
in $A/I$. Then $f_1, \ldots, f_n, g_1, \ldots, g_m$ is an $H_1$-regular
sequence in $A$.
\end{lemma}

\begin{proof}
We have to show that $H_1(A, f_1, \ldots, f_n, g_1, \ldots, g_m) = 0$.
To do this consider the commutative diagram
$$
\xymatrix{
\wedge^2(A^{\oplus n + m}) \ar[r] \ar[d] &
A^{\oplus n + m} \ar[r] \ar[d] &
A \ar[r] \ar[d] & 0 \\
\wedge^2(A/I^{\oplus m}) \ar[r] &
A/I^{\oplus m} \ar[r] &
A/I \ar[r] & 0
}
$$
Consider an element $(a_1, \ldots, a_{n + m}) \in A^{\oplus n + m}$
which maps to zero in $A$. Because $\overline{g}_1, \ldots, \overline{g}_m$
form an $H_1$-regular sequence in $A/I$ we see that
$(\overline{a}_{n + 1}, \ldots, \overline{a}_{n + m})$ is the image
of some element $\overline{\alpha}$ of $\wedge^2(A/I^{\oplus m})$.
We can lift $\overline{\alpha}$ to an element
$\alpha \in \wedge^2(A^{\oplus n + m})$ and subtract the image of it
in $A^{\oplus n + m}$ from our element $(a_1, \ldots, a_{n + m})$.
Thus we may assume that $a_{n + 1}, \ldots, a_{n + m} \in I$.
Since $I = (f_1, \ldots, f_n)$ we can modify our element
$(a_1, \ldots, a_{n + m})$ by linear combinations of the elements
$$
(0, \ldots, g_j, 0, \ldots, 0, f_i, 0, \ldots, 0)
$$
in the image of the top left horizontal arrow to reduce to the case
that $a_{n + 1}, \ldots, a_{n + m}$ are zero. In this case
$(a_1, \ldots, a_n, 0, \ldots, 0)$ defines an element of
$H_1(A, f_1, \ldots, f_n)$ which we assumed to be zero.
\end{proof}

\begin{lemma}
\label{lemma-truncate-H1-regular}
Let $A$ be a ring. Let $f_1, \ldots, f_n, g_1, \ldots, g_m \in A$
be an $H_1$-regular sequence. Then the images
$\overline{g}_1, \ldots, \overline{g}_m$ in $A/(f_1, \ldots, f_n)$
form an $H_1$-regular sequence.
\end{lemma}

\begin{proof}
Set $I = (f_1, \ldots, f_n)$. We have to show that any relation
$\sum_{j = 1, \ldots, m} \overline{a}_j \overline{g}_j$ in $A/I$
is a linear combination of trivial relations. Because
$I = (f_1, \ldots, f_n)$ we can lift this relation to a relation
$$
\sum\nolimits_{j = 1, \ldots, m} a_j g_j +
\sum\nolimits_{i = 1, \ldots, n} b_if_i = 0
$$
in $A$. By assumption this relation in $A$ is a linear combination of
trivial relations. Taking the image in $A/I$ we obtain what we want.
\end{proof}

\begin{lemma}
\label{lemma-join-koszul-regular-sequences}
Let $A$ be a ring. Let $I$ be an ideal generated by a Koszul-regular
sequence $f_1, \ldots, f_n$ in $A$. Let $g_1, \ldots, g_m \in A$ be
elements whose images $\overline{g}_1, \ldots, \overline{g}_m$ form a
Koszul-regular sequence in $A/I$. Then $f_1, \ldots, f_n, g_1, \ldots, g_m$
is a Koszul-regular sequence in $A$.
\end{lemma}

\begin{proof}
Our assumptions say that $K_\bullet(A, f_1, \ldots, f_n)$ is a finite free
resolution of $A/I$ and
$K_\bullet(A/I, \overline{g}_1, \ldots, \overline{g}_m)$ is a
finite free resolution of $A/(f_i, g_j)$ over $A/I$. Then
\begin{align*}
K_\bullet(A, f_1, \ldots, f_n, g_1, \ldots, g_m)
& = \text{Tot}(K_\bullet(A, f_1, \ldots, f_n) \otimes_A
K_\bullet(A, g_1, \ldots, g_m)) \\
& \cong A/I \otimes_A K_\bullet(A, g_1, \ldots, g_m) \\
& = K_\bullet(A/I, \overline{g}_1, \ldots, \overline{g}_m) \\
& \cong A/(f_i, g_j)
\end{align*}
The first equality by
Lemma \ref{lemma-join-sequences-koszul-complex}.
The first quasi-isomorphism $\cong$ by (the dual of)
Homology, Lemma \ref{homology-lemma-double-complex-gives-resolution}
as the $q$th row of the double complex
$K_\bullet(A, f_1, \ldots, f_n) \otimes_A K_\bullet(A, g_1, \ldots, g_m)$
is a resolution of $A/I \otimes_A K_q(A, g_1, \ldots, g_m)$.
The second equality is clear. The last quasi-isomorphism by assumption.
Hence we win.
\end{proof}

\noindent
To conclude in the following lemma it is necessary to assume that both
$f_1, \ldots, f_n$ and $f_1, \ldots, f_n, g_1, \ldots, g_m$
are Koszul-regular. A counter example to dropping the assumption
that $f_1, \ldots, f_n$ is Koszul-regular is
Examples, Lemma \ref{examples-lemma-strange-regular-sequence}.

\begin{lemma}
\label{lemma-truncate-koszul-regular}
Let $A$ be a ring. Let $f_1, \ldots, f_n, g_1, \ldots, g_m \in A$.
If both $f_1, \ldots, f_n$ and $f_1, \ldots, f_n, g_1, \ldots, g_m$
are Koszul-regular sequences in $A$, then
$\overline{g}_1, \ldots, \overline{g}_m$ in $A/(f_1, \ldots, f_n)$
form a Koszul-regular sequence.
\end{lemma}

\begin{proof}
Set $I = (f_1, \ldots, f_n)$.
Our assumptions say that $K_\bullet(A, f_1, \ldots, f_n)$ is a finite free
resolution of $A/I$ and
$K_\bullet(A, f_1, \ldots, f_n, g_1, \ldots, g_m)$ is a
finite free resolution of $A/(f_i, g_j)$ over $A$. Then
\begin{align*}
A/(f_i, g_j) & \cong K_\bullet(A, f_1, \ldots, f_n, g_1, \ldots, g_m) \\
& = \text{Tot}(K_\bullet(A, f_1, \ldots, f_n) \otimes_A
K_\bullet(A, g_1, \ldots, g_m)) \\
& \cong A/I \otimes_A K_\bullet(A, g_1, \ldots, g_m) \\
& = K_\bullet(A/I, \overline{g}_1, \ldots, \overline{g}_m)
\end{align*}
The first quasi-isomorphism $\cong$ by assumption. The first equality by
Lemma \ref{lemma-join-sequences-koszul-complex}.
The second quasi-isomorphism by (the dual of)
Homology, Lemma \ref{homology-lemma-double-complex-gives-resolution}
as the $q$th row of the double complex
$K_\bullet(A, f_1, \ldots, f_n) \otimes_A K_\bullet(A, g_1, \ldots, g_m)$
is a resolution of $A/I \otimes_A K_q(A, g_1, \ldots, g_m)$.
The second equality is clear. Hence we win.
\end{proof}

\begin{lemma}
\label{lemma-independence-of-generators}
Let $R$ be a ring. Let $I$ be an ideal generated by $f_1, \ldots, f_r \in R$.
\begin{enumerate}
\item If $I$ can be generated by a quasi-regular sequence of length $r$,
then $f_1, \ldots, f_r$ is a quasi-regular sequence.
\item If $I$ can be generated by an $H_1$-regular sequence of length $r$,
then $f_1, \ldots, f_r$ is an $H_1$-regular sequence.
\item If $I$ can be generated by a Koszul-regular sequence of length $r$,
then $f_1, \ldots, f_r$ is a Koszul-regular sequence.
\end{enumerate}
\end{lemma}

\begin{proof}
If $I$ can be generated by a quasi-regular sequence of length $r$,
then $I/I^2$ is free of rank $r$ over $R/I$. Since $f_1, \ldots, f_r$
generate by assumption we see that the images $\overline{f}_i$ form a basis of
$I/I^2$ over $R/I$. It follows that $f_1, \ldots, f_r$ is a quasi-regular
sequence as all this means, besides the freeness of $I/I^2$, is that the maps
$\text{Sym}^n_{R/I}(I/I^2) \to I^n/I^{n + 1}$ are isomorphisms.

\medskip\noindent
We continue to assume that $I$ can be generated by a
quasi-regular sequence, say
$g_1, \ldots, g_r$. Write $g_j = \sum a_{ij}f_i$. As $f_1, \ldots, f_r$
is quasi-regular according to the previous paragraph, we see that
$\det(a_{ij})$ is invertible mod $I$. The matrix
$a_{ij}$ gives a map $R^{\oplus r} \to R^{\oplus r}$ which induces
a map of Koszul complexes
$\alpha : K_\bullet(R, f_1, \ldots, f_r) \to K_\bullet(R, g_1, \ldots, g_r)$,
see
Lemma \ref{lemma-functorial}.
This map becomes an isomorphism on inverting $\det(a_{ij})$.
Since the cohomology modules of both $K_\bullet(R, f_1, \ldots, f_r)$ and
$K_\bullet(R, g_1, \ldots, g_r)$ are annihilated by $I$, see
Lemma \ref{lemma-homotopy-koszul},
we see that $\alpha$ is a quasi-isomorphism.

\medskip\noindent
Now assume that $g_1, \ldots, g_r$ is a $H_1$-regular sequence generating $I$.
Then $g_1, \ldots, g_r$ is a quasi-regular sequence by
Lemma \ref{lemma-H1-regular-quasi-regular}. By the previous paragraph
we conclude that $f_1, \ldots, f_r$ is a $H_1$-regular sequence.
Similarly for Koszul-regular sequences.
\end{proof}

\begin{lemma}
\label{lemma-make-nonzero-divisor}
\begin{reference}
This is a particular case of \cite[Corollary]{McCoy}
\end{reference}
Let $R$ be a ring. Let $a_1, \ldots, a_n \in R$ be elements such
that $R \to R^{\oplus n}$, $x \mapsto (xa_1, \ldots, xa_n)$ is injective.
Then the element $\sum a_i t_i$ of the polynomial ring $R[t_1, \ldots, t_n]$
is a nonzerodivisor.
\end{lemma}

\begin{proof}
If one of the $a_i$ is a unit this is just the statement that any
element of the form $t_1 + a_2 t_2 + \ldots + a_n t_n$ is a nonzerodivisor
in the polynomial ring over $R$.

\medskip\noindent
Case I: $R$ is Noetherian. Let $\mathfrak q_j$, $j = 1, \ldots, m$
be the associated primes of $R$. We have to show that
each of the maps
$$
\sum a_i t_i :
\text{Sym}^d(R^{\oplus n})
\longrightarrow
\text{Sym}^{d + 1}(R^{\oplus n})
$$
is injective. As $\text{Sym}^d(R^{\oplus n})$ is a free $R$-module its
associated primes are $\mathfrak q_j$, $j = 1, \ldots, m$. For each $j$
there exists an $i = i(j)$ such that $a_i \not \in \mathfrak q_j$ because
there exists an $x \in R$ with $\mathfrak q_jx = 0$ but $a_i x \not = 0$
for some $i$ by assumption. Hence $a_i$ is a unit in $R_{\mathfrak q_j}$
and the map is injective after localizing at $\mathfrak q_j$. Thus the map
is injective, see
Algebra, Lemma \ref{algebra-lemma-zero-at-ass-zero}.

\medskip\noindent
Case II: $R$ general. We can write $R$ as the union of Noetherian
rings $R_\lambda$ with $a_1, \ldots, a_n \in R_\lambda$. For each $R_\lambda$
the result holds, hence the result holds for $R$.
\end{proof}

\begin{lemma}
\label{lemma-Koszul-regular-flat-locally-regular}
Let $R$ be a ring. Let $f_1, \ldots, f_n$ be a Koszul-regular sequence
in $R$ such that $(f_1, \ldots, f_n) \not = R$.
Consider the faithfully flat, smooth ring map
$$
R \longrightarrow
S = R[\{t_{ij}\}_{i \leq j}, t_{11}^{-1}, t_{22}^{-1}, \ldots, t_{nn}^{-1}]
$$
For $1 \leq i \leq n$ set
$$
g_i = \sum\nolimits_{i \leq j} t_{ij} f_j \in S.
$$
Then $g_1, \ldots, g_n$ is a regular sequence in $S$ and
$(f_1, \ldots, f_n)S = (g_1, \ldots, g_n)$.
\end{lemma}

\begin{proof}
The equality of ideals is obvious as the matrix
$$
\left(
\begin{matrix}
t_{11} & t_{12} & t_{13} & \ldots \\
0 & t_{22} & t_{23} & \ldots \\
0 & 0 & t_{33} & \ldots \\
\ldots & \ldots & \ldots & \ldots
\end{matrix}
\right)
$$
is invertible in $S$.
Because $f_1, \ldots, f_n$ is a Koszul-regular sequence we see that
the kernel of
$R \to R^{\oplus n}$, $x \mapsto (xf_1, \ldots, xf_n)$ is zero (as it
computes the $n$the Koszul homology of $R$ w.r.t.\ $f_1, \ldots, f_n$).
Hence by
Lemma \ref{lemma-make-nonzero-divisor}
we see that $g_1 = f_1 t_{11} + \ldots + f_n t_{1n}$ is a nonzerodivisor
in $S' = R[t_{11}, t_{12}, \ldots, t_{1n}, t_{11}^{-1}]$. We see that
$g_1, f_2, \ldots, f_n$ is a Koszul-sequence in $S'$ by
Lemma \ref{lemma-koszul-regular-flat-base-change} and
\ref{lemma-independence-of-generators}.
We conclude that
$\overline{f}_2, \ldots, \overline{f}_n$ is a Koszul-regular sequence
in $S'/(g_1)$ by
Lemma \ref{lemma-truncate-koszul-regular}.
Hence by induction on $n$ we see that the images
$\overline{g}_2, \ldots, \overline{g}_n$ of $g_2, \ldots, g_n$ in
$S'/(g_1)[\{t_{ij}\}_{2 \leq i \leq j}, t_{22}^{-1}, \ldots, t_{nn}^{-1}]$
form a regular sequence. This in turn means that
$g_1, \ldots, g_n$ forms a regular sequence in $S$.
\end{proof}







\section{More on Koszul regular sequences}
\label{section-more-koszul-regular}

\noindent
We continue the discussion from Section \ref{section-koszul-regular}.

\begin{lemma}
\label{lemma-vanishing-extended-alternating-koszul}
Let $R$ be a ring. Let $f_1, \ldots, f_r \in R$ be an
Koszul-regular sequence. Then the extended alternating
{\v C}ech complex $R \to \bigoplus\nolimits_{i_0} R_{f_{i_0}} \to
\bigoplus\nolimits_{i_0 < i_1} R_{f_{i_0}f_{i_1}} \to
\ldots \to R_{f_1\ldots f_r}$ from Section \ref{section-alternating-cech}
only has cohomology in degree $r$.
\end{lemma}

\begin{proof}
By Lemma \ref{lemma-mult-koszul-regular} and induction the sequence
$f_1, \ldots, f_{r - 1}, f_r^n$ is Koszul regular for all $n \geq 1$.
By Lemma \ref{lemma-change-basis} any permutation of a Koszul regular
sequence is a Koszul regular sequence. Hence we see that we may replace
any (or all) $f_i$ by its $n$th power and still have a Koszul regular
sequence. Thus $K_\bullet(R, f_1^n, \ldots, f_r^n)$ has nonzero
cohomology only in homological degree $0$. This implies what we want by
Lemma \ref{lemma-extended-alternating-Cech-is-colimit-koszul}.
\end{proof}

\begin{lemma}
\label{lemma-blowup-regular-sequence}
Let $a, a_2, \ldots, a_r$ be an $H_1$-regular sequence in a ring $R$
(for example a Koszul regular sequence or a regular sequence, see
Lemmas \ref{lemma-regular-koszul-regular} and
\ref{lemma-koszul-regular-H1-regular}).
With $I = (a, a_2, \ldots, a_r)$ the blowup algebra $R' = R[\frac{I}{a}]$
is isomorphic to $R'' = R[y_2, \ldots, y_r]/(a y_i - a_i)$.
\end{lemma}

\begin{proof}
By Algebra, Lemma \ref{algebra-lemma-affine-blowup-quotient-description}
it suffices to show that $R''$ is $a$-torsion free.

\medskip\noindent
We claim $a, ay_2 - a_2, \ldots, ay_n - a_r$ is a $H_1$-regular
sequence in $R[y_2, \ldots, y_r]$. Namely, the map
$$
(a, ay_2 - a_2, \ldots, ay_n - a_r) :
R[y_2, \ldots, y_r]^{\oplus r}
\longrightarrow
R[y_2, \ldots, y_r]
$$
used to define the Koszul complex on $a, ay_2 - a_2, \ldots, ay_n - a_r$
is isomorphic to the map
$$
(a, a_2, \ldots, a_r) :
R[y_2, \ldots, y_r]^{\oplus r} \longrightarrow
R[y_2, \ldots, y_r]
$$
used to the define the Koszul complex on $a, a_2, \ldots, a_r$
via the isomorphism
$$
R[y_2, \ldots, y_r]^{\oplus r}
\longrightarrow
R[y_2, \ldots, y_r]^{\oplus r}
$$
sending $(b_1, \ldots, b_r)$ to
$(b_1 - b_2y_2 \ldots - b_ry_r, -b_2, \ldots, - b_r)$.
By Lemma \ref{lemma-functorial} these Koszul complexes are isomorphic.
By Lemma \ref{lemma-koszul-regular-flat-base-change}
applied to the flat ring map $R \to R[y_2, \ldots, y_r]$
we conclude our claim is true. By Lemma \ref{lemma-cone-koszul}
we see that the Koszul complex $K$ on $a, ay_2 - a_2, \ldots, ay_n - a_r$
is the cone on $a : L \to L$ where $L$ is the Koszul complex
on $ay_2 - a_2, \ldots, ay_n - a_r$. Since $H_1(K) = 0$ by the claim,
we conclude that $a : H_0(L) \to H_0(L)$ is injective, in other words that
$R'' = R[y_2, \ldots, y_r]/(a y_i - a_i)$
has no nonzero $a$-torsion elements as desired.
\end{proof}

\begin{lemma}
\label{lemma-base-change-H1-regular}
Let $A \to B$ be a ring map.
Let $f_1, \ldots, f_r$ be a sequence in $B$ such that $B/(f_1, \ldots, f_r)$
is $A$-flat. Let $A \to A'$ be a ring map. Then the canonical map
$$
H_1(K_\bullet(B, f_1, \ldots, f_r)) \otimes_A A'
\longrightarrow
H_1(K_\bullet(B', f'_1, \ldots, f'_r))
$$
is surjective. Here $B' = B \otimes_A A'$ and $f_i' \in B'$ is the image
of $f_i$.
\end{lemma}

\begin{proof}
The sequence
$$
\wedge^2(B^{\oplus r}) \to B^{\oplus r} \to B \to B/J \to 0
$$
is a complex of $A$-modules with $B/J$ flat over $A$ and
cohomology group $H_1 = H_1(K_\bullet(B, f_1, \ldots, f_r))$ in the spot
$B^{\oplus r}$. If we tensor this with $A'$ we obtain a complex
$$
\wedge^2((B')^{\oplus r}) \to (B')^{\oplus r} \to B' \to B'/J' \to 0
$$
which is exact at $B'$ and $B'/J'$. In order to compute its
cohomology group $H'_1 = H_1(K_\bullet(B', f'_1, \ldots, f'_r))$
at $(B')^{\oplus r}$ we split the first sequence above into
the exact sequences $0 \to J \to B \to B/J \to 0$,
$0 \to K \to B^{\oplus r} \to J \to 0$, and
$\wedge^2(B^{\oplus r}) \to K \to H_1 \to 0$.
Tensoring over $A$ with $A'$ we obtain the exact sequences
$$
\begin{matrix}
0 \to J \otimes_A A' \to B \otimes_A A' \to (B/J) \otimes_A A' \to 0 \\
K \otimes_A A' \to B^{\oplus r} \otimes_A A' \to J \otimes_A A' \to 0 \\
\wedge^2(B^{\oplus r}) \otimes_A A' \to K \otimes_A A' \to H_1 \otimes_A A'
\to 0
\end{matrix}
$$
where the first one is exact as $B/J$ is flat over $A$, see
Algebra, Lemma \ref{algebra-lemma-flat-tor-zero}. We conclude
that $J' = J \otimes_A A' \subset B'$ and that
$K \otimes_A A' \to \Ker((B')^{\oplus r} \to B')$ is
surjective. Thus
\begin{align*}
H_1 \otimes_A A'
& =
\Coker\left(\wedge^2(B^{\oplus r}) \otimes_A A' \to K \otimes_A A'\right) \\
& \to
\Coker\left(
\wedge^2((B')^{\oplus r})  \to \Ker((B')^{\oplus r} \to B')
\right) = H'_1
\end{align*}
is surjective too.
\end{proof}

\begin{lemma}
\label{lemma-relative-regular-immersion-algebra}
Let $A \to B$ and $A \to A'$ be ring maps. Set $B' = B \otimes_A A'$.
Let $f_1, \ldots, f_r \in B$. Assume $B/(f_1, \ldots, f_r)B$ is flat over $A$
\begin{enumerate}
\item If $f_1, \ldots, f_r$ is a quasi-regular sequence, then
the image in $B'$ is a quasi-regular sequence.
\item If $f_1, \ldots, f_r$ is a $H_1$-regular sequence, then
the image in $B'$ is a $H_1$-regular sequence.
\end{enumerate}
\end{lemma}

\begin{proof}
Assume $f_1, \ldots, f_r$ is quasi-regular. Set $J = (f_1, \ldots, f_r)$.
By assumption $J^n/J^{n + 1}$ is isomorphic to a direct sum of copies of
$B/J$ hence flat over $A$. By induction and
Algebra, Lemma \ref{algebra-lemma-flat-ses}
we conclude that $B/J^n$ is flat over $A$. The ideal $(J')^n$ is equal to
$J^n \otimes_A A'$, see
Algebra, Lemma \ref{algebra-lemma-flat-tor-zero}.
Hence $(J')^n/(J')^{n + 1} = J^n/J^{n + 1} \otimes_A A'$ which clearly
implies that $f_1, \ldots, f_r$ is a quasi-regular sequence in $B'$.

\medskip\noindent
Assume $f_1, \ldots, f_r$ is $H_1$-regular. By
Lemma \ref{lemma-base-change-H1-regular}
the vanishing of the Koszul homology group
$H_1(K_\bullet(B, f_1, \ldots, f_r))$
implies the vanishing of $H_1(K_\bullet(B', f'_1, \ldots, f'_r))$
and we win.
\end{proof}

\begin{lemma}
\label{lemma-cut-by-koszul-flat}
Let $A' \to B'$ be a ring map. Let $I \subset A'$ be an ideal.
Set $A = A'/I$ and $B = B'/IB'$. Let $f'_1, \ldots, f'_r \in B'$. Assume
\begin{enumerate}
\item $A' \to B'$ is flat and of finite presentation,
\item $I$ is locally nilpotent,
\item the images $f_1, \ldots, f_r \in B$ form a quasi-regular sequence,
\item $B/(f_1, \ldots, f_r)$ is flat over $A$.
\end{enumerate}
Then $B'/(f'_1, \ldots, f'_r)$ is flat over $A'$.
\end{lemma}

\begin{proof}
Set $C' = B'/(f'_1, \ldots, f'_r)$. We have to show $A' \to C'$ is flat.
Let $\mathfrak r' \subset C'$ be a prime ideal lying over
$\mathfrak p' \subset A'$. We let $\mathfrak q' \subset B'$
be the inverse image of $\mathfrak r'$.
By Algebra, Lemma \ref{algebra-lemma-flat-localization}
it suffices to show that $A'_{\mathfrak p'} \to C'_{\mathfrak q'}$ is flat.
Algebra, Lemma \ref{algebra-lemma-grothendieck-regular-sequence-general}
tells us it suffices to show that $f'_1, \ldots, f'_r$ map to
a regular sequence in
$$
B'_{\mathfrak q'}/\mathfrak p'B'_{\mathfrak q'} =
B_\mathfrak q/\mathfrak p B_\mathfrak q =
(B \otimes_A \kappa(\mathfrak p))_\mathfrak q
$$
with obvious notation. What we know is that $f_1, \ldots, f_r$
is a quasi-regular sequence in $B$ and that $B/(f_1, \ldots, f_r)$
is flat over $A$. By Lemma \ref{lemma-relative-regular-immersion-algebra}
the images $\overline{f}_1, \ldots, \overline{f}_r$
of $f'_1, \ldots, f'_r$ in $B \otimes_A \kappa(\mathfrak p)$ form a
quasi-regular sequence. Since $(B \otimes_A \kappa(\mathfrak p))_\mathfrak q$
is a Noetherian local ring, we conclude by Lemma
\ref{lemma-noetherian-finite-all-equivalent}.
\end{proof}

\begin{lemma}
\label{lemma-cut-by-koszul}
Let $A' \to B'$ be a ring map. Let $I \subset A'$ be an ideal.
Set $A = A'/I$ and $B = B'/IB'$. Let $f'_1, \ldots, f'_r \in B'$. Assume
\begin{enumerate}
\item $A' \to B'$ is flat and of finite presentation (for example smooth),
\item $I$ is locally nilpotent,
\item the images $f_1, \ldots, f_r \in B$ form a quasi-regular sequence,
\item $B/(f_1, \ldots, f_r)$ is smooth over $A$.
\end{enumerate}
Then $B'/(f'_1, \ldots, f'_r)$ is smooth over $A'$.
\end{lemma}

\begin{proof}
Set $C' = B'/(f'_1, \ldots, f'_r)$ and $C = B/(f_1, \ldots, f_r)$.
Then $A' \to C'$ is of finite presentation.
By Lemma \ref{lemma-cut-by-koszul-flat} we see that $A' \to C'$ is flat.
The fibre rings of $A' \to C'$ are equal to the fibre rings of $A \to C$
and hence smooth by assumption (4). It follows that
$A' \to C'$ is smooth by
Algebra, Lemma \ref{algebra-lemma-flat-fibre-smooth}.
\end{proof}








\section{Regular ideals}
\label{section-ideals}

\noindent
We will discuss the notion of a regular ideal sheaf in great generality in
Divisors, Section \ref{divisors-section-regular-ideal-sheaves}.
Here we define the corresponding notion in the affine case, i.e., in
the case of an ideal in a ring.

\begin{definition}
\label{definition-regular-ideal}
Let $R$ be a ring and let $I \subset R$ be an ideal.
\begin{enumerate}
\item We say $I$ is a {\it regular ideal} if for every
$\mathfrak p \in V(I)$ there exists a $g \in R$, $g \not \in \mathfrak p$
and a regular sequence $f_1, \ldots, f_r \in R_g$ such that $I_g$
is generated by $f_1, \ldots, f_r$.
\item We say $I$ is a {\it Koszul-regular ideal} if for every
$\mathfrak p \in V(I)$ there exists a $g \in R$, $g \not \in \mathfrak p$
and a Koszul-regular sequence $f_1, \ldots, f_r \in R_g$ such that $I_g$
is generated by $f_1, \ldots, f_r$.
\item We say $I$ is a {\it $H_1$-regular ideal} if for every
$\mathfrak p \in V(I)$ there exists a $g \in R$, $g \not \in \mathfrak p$
and an $H_1$-regular sequence $f_1, \ldots, f_r \in R_g$ such that $I_g$
is generated by $f_1, \ldots, f_r$.
\item We say $I$ is a {\it quasi-regular ideal} if for every
$\mathfrak p \in V(I)$ there exists a $g \in R$, $g \not \in \mathfrak p$
and a quasi-regular sequence $f_1, \ldots, f_r \in R_g$ such that $I_g$
is generated by $f_1, \ldots, f_r$.
\end{enumerate}
\end{definition}

\noindent
It is clear that given $I \subset R$ we have the implications
\begin{align*}
I\text{ is a regular ideal}
& \Rightarrow
I\text{ is a Koszul-regular ideal} \\
& \Rightarrow
I\text{ is a }H_1\text{-regular ideal} \\
& \Rightarrow
I\text{ is a quasi-regular ideal}
\end{align*}
see Lemmas \ref{lemma-regular-koszul-regular},
\ref{lemma-koszul-regular-H1-regular}, and
\ref{lemma-H1-regular-quasi-regular}. Such an ideal is always finitely
generated.

\begin{lemma}
\label{lemma-quasi-regular-ideal-finite}
A quasi-regular ideal is finitely generated.
\end{lemma}

\begin{proof}
Let $I \subset R$ be a quasi-regular ideal. Since $V(I)$ is quasi-compact,
there exist $g_1, \ldots, g_m \in R$ such that
$V(I) \subset D(g_1) \cup \ldots \cup D(g_m)$
and such that $I_{g_j}$ is generated by a quasi-regular sequence
$g_{j1}, \ldots, g_{jr_j} \in R_{g_j}$. Write $g_{ji} = g'_{ji}/g_j^{e_{ij}}$
for some $g'_{ij} \in I$. Write $1 + x = \sum g_j h_j$
for some $x \in I$ which is possible as
$V(I) \subset D(g_1) \cup \ldots \cup D(g_m)$.
Note that $\Spec(R) = D(g_1) \cup \ldots \cup D(g_m) \bigcup D(x)$
Then $I$ is generated by the elements $g'_{ij}$ and $x$ as
these generate on each of the pieces of the cover, see
Algebra, Lemma \ref{algebra-lemma-cover}.
\end{proof}

\begin{lemma}
\label{lemma-quasi-regular-ideal-finite-projective}
Let $I \subset R$ be a quasi-regular ideal of a ring.
Then $I/I^2$ is a finite projective $R/I$-module.
\end{lemma}

\begin{proof}
This follows from Algebra, Lemma \ref{algebra-lemma-finite-projective}
and the definitions.
\end{proof}

\noindent
We prove flat descent for Koszul-regular, $H_1$-regular, quasi-regular
ideals.

\begin{lemma}
\label{lemma-flat-descent-regular-ideal}
Let $A \to B$ be a faithfully flat ring map. Let $I \subset A$ be an ideal.
If $IB$ is a Koszul-regular
(resp.\ $H_1$-regular, resp.\ quasi-regular) ideal in $B$, then
$I$ is a Koszul-regular (resp.\ $H_1$-regular, resp.\ quasi-regular)
ideal in $A$.
\end{lemma}

\begin{proof}
We fix the prime $\mathfrak p \supset I$ throughout the proof.
Assume $IB$ is quasi-regular. By
Lemma \ref{lemma-quasi-regular-ideal-finite}
$IB$ is a finite module, hence $I$ is a finite $A$-module by
Algebra, Lemma \ref{algebra-lemma-descend-properties-modules}.
As $A \to B$ is flat we see that
$$
I/I^2 \otimes_{A/I} B/IB = I/I^2 \otimes_A B = IB/(IB)^2.
$$
As $IB$ is quasi-regular, the $B/IB$-module $IB/(IB)^2$ is finite
locally free. Hence $I/I^2$ is finite projective, see
Algebra, Proposition \ref{algebra-proposition-ffdescent-finite-projectivity}.
In particular, after replacing $A$ by $A_f$ for some
$f \in A$, $f \not \in \mathfrak p$ we may assume that $I/I^2$ is free of
rank $r$. Pick $f_1, \ldots, f_r \in I$ which
give a basis of $I/I^2$. By Nakayama's lemma (see
Algebra, Lemma \ref{algebra-lemma-NAK})
we see that, after another replacement $A \leadsto A_f$ as above,
$I$ is generated by $f_1, \ldots, f_r$.

\medskip\noindent
Proof of the ``quasi-regular'' case. Above we have seen that
$I/I^2$ is free on the $r$-generators $f_1, \ldots, f_r$.
To finish the proof in this case we have to show that the maps
$\text{Sym}^d(I/I^2) \to I^d/I^{d + 1}$ are isomorphisms
for each $d \geq 2$. This is clear as the faithfully flat
base changes $\text{Sym}^d(IB/(IB)^2) \to (IB)^d/(IB)^{d + 1}$
are isomorphisms locally on $B$ by assumption.
Details omitted.

\medskip\noindent
Proof of the ``$H_1$-regular'' and ``Koszul-regular'' case.
Consider the sequence of elements $f_1, \ldots, f_r$ generating
$I$ we constructed above. By
Lemma \ref{lemma-independence-of-generators}
we see that $f_1, \ldots, f_r$ map to a $H_1$-regular or Koszul-regular
sequence in $B_g$ for any $g \in B$ such that $IB$ is generated by
an $H_1$-regular or Koszul-regular sequence. Hence
$K_\bullet(A, f_1, \ldots, f_r) \otimes_A B_g$ has vanishing
$H_1$ or $H_i$, $i > 0$. Since the homology of
$K_\bullet(B, f_1, \ldots, f_r) = K_\bullet(A, f_1, \ldots, f_r) \otimes_A B$
is annihilated by $IB$ (see
Lemma \ref{lemma-homotopy-koszul})
and since $V(IB) \subset \bigcup_{g\text{ as above}} D(g)$ we conclude that
$K_\bullet(A, f_1, \ldots, f_r) \otimes_A B$ has vanishing homology in
degree $1$ or all positive degrees. Using that $A \to B$ is faithfully
flat we conclude that the same is true for
$K_\bullet(A, f_1, \ldots, f_r)$.
\end{proof}

\begin{lemma}
\label{lemma-conormal-sequence-H1-regular-ideal}
Let $A$ be a ring. Let $I \subset J \subset A$ be ideals.
Assume that $J/I \subset A/I$ is a $H_1$-regular ideal.
Then $I \cap J^2 = IJ$.
\end{lemma}

\begin{proof}
Follows immediately from Lemma \ref{lemma-conormal-sequence-H1-regular}
by localizing.
\end{proof}







\section{Local complete intersection maps}
\label{section-lci}

\noindent
We can use the material above to define a local complete intersection
map between rings using presentations by (finite) polynomial algebras.

\begin{lemma}
\label{lemma-koszul-independence-presentation}
Let $A \to B$ be a finite type ring map. If for some presentation
$\alpha : A[x_1, \ldots, x_n] \to B$ the kernel $I$ is a Koszul-regular ideal
then for any presentation $\beta : A[y_1, \ldots, y_m] \to B$ the kernel
$J$ is a Koszul-regular ideal.
\end{lemma}

\begin{proof}
Choose $f_j \in A[x_1, \ldots, x_n]$ with $\alpha(f_j) = \beta(y_j)$
and $g_i \in A[y_1, \ldots, y_m]$ with $\beta(g_i) = \alpha(x_i)$.
Then we get a commutative diagram
$$
\xymatrix{
A[x_1, \ldots, x_n, y_1, \ldots, y_m]
\ar[d]^{x_i \mapsto g_i} \ar[rr]_-{y_j \mapsto f_j} & &
A[x_1, \ldots, x_n] \ar[d] \\
A[y_1, \ldots, y_m] \ar[rr] & & B
}
$$
Note that the kernel $K$ of $A[x_i, y_j] \to B$ is equal to
$K = (I, y_j - f_j) = (J, x_i - g_i)$. In particular, as
$I$ is finitely generated by
Lemma \ref{lemma-quasi-regular-ideal-finite}
we see that $J = K/(x_i - g_i)$ is finitely generated too.

\medskip\noindent
Pick a prime $\mathfrak q \subset B$. Since
$I/I^2 \oplus B^{\oplus m} = J/J^2 \oplus B^{\oplus n}$
(Algebra, Lemma \ref{algebra-lemma-conormal-module})
we see that
$$
\dim J/J^2 \otimes_B \kappa(\mathfrak q) + n =
\dim I/I^2 \otimes_B \kappa(\mathfrak q) + m.
$$
Pick $p_1, \ldots, p_t \in I$ which map to a basis of
$I/I^2 \otimes \kappa(\mathfrak q) = I \otimes_{A[x_i]} \kappa(\mathfrak q)$.
Pick $q_1, \ldots, q_s \in J$ which map to a basis of
$J/J^2 \otimes \kappa(\mathfrak q) = J \otimes_{A[y_j]} \kappa(\mathfrak q)$.
So $s + n = t + m$. By Nakayama's lemma there exist $h \in A[x_i]$ and
$h' \in A[y_j]$ both mapping to a nonzero element of $\kappa(\mathfrak q)$
such that $I_h = (p_1, \ldots, p_t)$ in $A[x_i, 1/h]$ and
$J_{h'} = (q_1, \ldots, q_s)$ in $A[y_j, 1/h']$.
As $I$ is Koszul-regular we may also assume that $I_h$ is generated
by a Koszul regular sequence. This sequence must necessarily have length
$t = \dim I/I^2 \otimes_B \kappa(\mathfrak q)$, hence we see that
$p_1, \ldots, p_t$ is a Koszul-regular sequence by
Lemma \ref{lemma-independence-of-generators}.
As also $y_1 - f_1, \ldots, y_m - f_m$ is a regular sequence we
conclude
$$
y_1 - f_1, \ldots, y_m - f_m, p_1, \ldots, p_t
$$
is a Koszul-regular sequence in $A[x_i, y_j, 1/h]$
(see Lemma \ref{lemma-join-koszul-regular-sequences}).
This sequence generates the ideal $K_h$. Hence the
ideal $K_{hh'}$ is generated by a Koszul-regular sequence
of length $m + t = n + s$. But it is also generated by the sequence
$$
x_1 - g_1, \ldots, x_n - g_n, q_1, \ldots, q_s
$$
of the same length which is thus a Koszul-regular sequence by
Lemma \ref{lemma-independence-of-generators}.
Finally, by
Lemma \ref{lemma-truncate-koszul-regular}
we conclude that the images of $q_1, \ldots, q_s$ in
$$
A[x_i, y_j, 1/hh']/(x_1 - g_1, \ldots, x_n - g_n)
\cong A[y_j, 1/h'']
$$
form a Koszul-regular sequence generating $J_{h''}$. Since $h''$
is the image of $hh'$ it doesn't map to zero in $\kappa(\mathfrak q)$
and we win.
\end{proof}

\noindent
This lemma allows us to make the following definition.

\begin{definition}
\label{definition-local-complete-intersection}
A ring map $A \to B$ is called a {\it local complete intersection}
if it is of finite type and for some (equivalently any) presentation
$B = A[x_1, \ldots, x_n]/I$ the ideal $I$ is Koszul-regular.
\end{definition}

\noindent
This notion is local.

\begin{lemma}
\label{lemma-lci-local}
Let $R \to S$ be a ring map. Let $g_1, \ldots, g_m \in S$
generate the unit ideal. If each $R \to S_{g_j}$ is a local
complete intersection so is $R \to S$.
\end{lemma}

\begin{proof}
Let $S = R[x_1, \ldots, x_n]/I$ be a presentation. Pick
$h_j \in R[x_1, \ldots, x_n]$ mapping to $g_j$ in $S$.
Then $R[x_1, \ldots, x_n, x_{n + 1}]/(I, x_{n + 1}h_j - 1)$
is a presentation of $S_{g_j}$. Hence
$I_j = (I, x_{n + 1}h_j - 1)$ is a Koszul-regular ideal in
$R[x_1, \ldots, x_n, x_{n + 1}]$. Pick a prime
$I \subset \mathfrak q \subset R[x_1, \ldots, x_n]$.
Then $h_j \not \in \mathfrak q$ for some $j$ and
$\mathfrak q_j = (\mathfrak q, x_{n + 1}h_j - 1)$ is a prime
ideal of $V(I_j)$ lying over $\mathfrak q$.
Pick $f_1, \ldots, f_r \in I$ which map to a basis of
$I/I^2 \otimes \kappa(\mathfrak q)$. Then
$x_{n + 1}h_j - 1, f_1, \ldots, f_r$ is a sequence of elements of $I_j$
which map to a basis of $I_j \otimes \kappa(\mathfrak q_j)$, see
Algebra, Lemma \ref{algebra-lemma-principal-localization-NL}.
By Nakayama's lemma there exists an $h \in R[x_1, \ldots, x_n, x_{n + 1}]$
such that $(I_j)_h$ is generated by $x_{n + 1}h_j - 1, f_1, \ldots, f_r$.
We may also assume that $(I_j)_h$ is generated by a Koszul regular
sequence of some length $e$. Looking at the dimension of
$I_j \otimes \kappa(\mathfrak q_j)$ we see that $e = r + 1$.
Hence by
Lemma \ref{lemma-independence-of-generators}
we see that $x_{n + 1}h_j - 1, f_1, \ldots, f_r$ is a
Koszul-regular sequence generating $(I_j)_h$ for some
$h \in R[x_1, \ldots, x_n, x_{n + 1}]$, $h \not \in \mathfrak q_j$. By
Lemma \ref{lemma-truncate-koszul-regular}
we see that $I_{h'}$ is generated by a Koszul-regular sequence for some
$h' \in R[x_1, \ldots, x_n]$, $h' \not \in \mathfrak q$ as desired.
\end{proof}

\begin{lemma}
\label{lemma-relative-global-complete-intersection-koszul}
Let $R$ be a ring. If $R[x_1, \ldots, x_n]/(f_1, \ldots, f_c)$
is a relative global complete intersection, then $f_1, \ldots, f_c$
is a Koszul regular sequence.
\end{lemma}

\begin{proof}
Recall that the homology groups $H_i(K_\bullet(f_\bullet))$ are
annihilated by the ideal $(f_1, \ldots, f_c)$. Hence it suffices
to show that $H_i(K_\bullet(f_\bullet))_\mathfrak q$ is zero for
all primes $\mathfrak q \subset R[x_1, \ldots, x_n]$
containing $(f_1, \ldots, f_c)$. This follows from
Algebra, Lemma
\ref{algebra-lemma-relative-global-complete-intersection-conormal}
and the fact that a regular sequence is Koszul regular
(Lemma \ref{lemma-regular-koszul-regular}).
\end{proof}

\begin{lemma}
\label{lemma-syntomic-lci}
Let $R \to S$ be a ring map. The following are equivalent
\begin{enumerate}
\item $R \to S$ is syntomic
(Algebra, Definition \ref{algebra-definition-lci}), and
\item $R \to S$ is flat and a local complete intersection.
\end{enumerate}
\end{lemma}

\begin{proof}
Assume (1). Then $R \to S$ is flat by definition.
By Algebra, Lemma \ref{algebra-lemma-syntomic} and
Lemma \ref{lemma-lci-local} we see that it suffices to
show a relative global complete intersection is
a local complete intersection homomorphism which is
Lemma \ref{lemma-relative-global-complete-intersection-koszul}.

\medskip\noindent
Assume (2). A local complete intersection is of finite presentation
because a Koszul-regular ideal is finitely generated.
Let $R \to k$ be a map to a field. It suffices to show that
$S' = S \otimes_R k$ is a local complete intersection over $k$, see
Algebra, Definition \ref{algebra-definition-lci-field}.
Choose a prime $\mathfrak q' \subset S'$.
Write $S = R[x_1, \ldots, x_n]/I$.
Then $S' = k[x_1, \ldots, x_n]/I'$ where
$I' \subset k[x_1, \ldots, x_n]$ is the image of $I$.
Let $\mathfrak p' \subset k[x_1, \ldots, x_n]$,
$\mathfrak q \subset S$,
and $\mathfrak p \subset R[x_1, \ldots, x_n]$
be the corresponding primes.
By Definition \ref{definition-regular-ideal}
exists an $g \in R[x_1, \ldots, x_n]$, $g \not \in \mathfrak p$
and $f_1, \ldots, f_r \in R[x_1, \ldots, x_n]_g$ which form
a Koszul-regular sequence generating $I_g$.
Since $S$ and hence $S_g$ is flat over $R$
we see that the images $f'_1, \ldots, f'_r$ in
$k[x_1, \ldots, x_n]_g$ form a $H_1$-regular sequence
generating $I'_g$, see
Lemma \ref{lemma-relative-regular-immersion-algebra}.
Thus $f'_1, \ldots, f'_r$ map to a regular sequence in
$k[x_1, \ldots, x_n]_{\mathfrak p'}$ generating
$I'_{\mathfrak p'}$ by Lemma \ref{lemma-noetherian-finite-all-equivalent}.
Applying Algebra, Lemma \ref{algebra-lemma-lci}
we conclude $S'_{gg'}$ for some $g' \in S$, $g' \not \in \mathfrak q'$
is a global complete intersection over $k$ as desired.
\end{proof}

\noindent
For a local complete intersection $R \to S$ we have $H_n(L_{S/R}) = 0$ for
$n \geq 2$. Since we haven't (yet) defined the full cotangent complex
we can't state and prove this, but we can deduce one of the consequences.

\begin{lemma}
\label{lemma-transitive-lci-at-end}
Let $A \to B \to C$ be ring maps. Assume $B \to C$ is a local complete
intersection homomorphism. Choose a presentation
$\alpha : A[x_s, s \in S] \to B$ with kernel $I$. Choose a presentation
$\beta : B[y_1, \ldots, y_m] \to C$ with kernel $J$. Let
$\gamma : A[x_s, y_t] \to C$ be the induced presentation of $C$ with kernel
$K$. Then we get a canonical commutative diagram
$$
\xymatrix{
0 \ar[r] &
\Omega_{A[x_s]/A} \otimes C \ar[r] &
\Omega_{A[x_s, y_t]/A} \otimes C \ar[r] &
\Omega_{B[y_t]/B} \otimes C \ar[r] &
0 \\
0 \ar[r] &
I/I^2 \otimes C \ar[r] \ar[u] &
K/K^2 \ar[r] \ar[u] &
J/J^2 \ar[r] \ar[u] &
0
}
$$
with exact rows. In particular, the six term exact sequence of
Algebra, Lemma \ref{algebra-lemma-exact-sequence-NL}
can be completed with a zero on the left, i.e., the sequence
$$
0 \to H_1(\NL_{B/A} \otimes_B C) \to
H_1(L_{C/A}) \to
H_1(L_{C/B}) \to
\Omega_{B/A} \otimes_B C \to
\Omega_{C/A} \to
\Omega_{C/B} \to 0
$$
is exact.
\end{lemma}

\begin{proof}
The only thing to prove is the injectivity of the map
$I/I^2 \otimes C \to K/K^2$.
By assumption the ideal $J$ is Koszul-regular.
Hence we have $IA[x_s, y_j] \cap K^2 = IK$ by
Lemma \ref{lemma-conormal-sequence-H1-regular-ideal}.
This means that the kernel of $K/K^2 \to J/J^2$ is
isomorphic to $IA[x_s, y_j]/IK$. Since
$I/I^2 \otimes_A C = IA[x_s, y_j]/IK$ by right exactness
of tensor product, this provides us with the desired injectivity of
$I/I^2 \otimes_A C \to K/K^2$.
\end{proof}

\begin{lemma}
\label{lemma-transitive-colimit-lci-at-end}
Let $A \to B \to C$ be ring maps.
If $B \to C$ is a filtered colimit of local complete intersection
homomorphisms then the conclusion of
Lemma \ref{lemma-transitive-lci-at-end}
remains valid.
\end{lemma}

\begin{proof}
Follows from
Lemma \ref{lemma-transitive-lci-at-end}
and
Algebra, Lemma \ref{algebra-lemma-colimits-NL}.
\end{proof}

\begin{lemma}
\label{lemma-henselization-NL}
Let $A \to B$ be a local homomorphism of local rings.
Let $A^h \to B^h$, resp.\ $A^{sh} \to B^{sh}$ be the induced
map on henselizations, resp.\ strict henselizations
(Algebra, Lemma \ref{algebra-lemma-henselian-functorial},
resp.\ Lemma \ref{algebra-lemma-strictly-henselian-functorial}).
Then $\NL_{B/A} \otimes_B B^h \to \NL_{B^h/A^h}$ and
$\NL_{B/A} \otimes_B B^{sh} \to \NL_{B^{sh}/A^{sh}}$
induce isomorphisms on cohomology groups.
\end{lemma}

\begin{proof}
Since $A^h$ is a filtered colimit of \'etale algebras over $A$
we see that $\NL_{A^h/A}$ is an acyclic complex by
Algebra, Lemma \ref{algebra-lemma-colimits-NL} and
Algebra, Definition \ref{algebra-definition-etale}.
The same is true for $B^h/B$.
Using the Jacobi-Zariski sequence
(Algebra, Lemma \ref{algebra-lemma-exact-sequence-NL})
for $A \to A^h \to B^h$ we find that
$\NL_{B^h/A} \to \NL_{B^h/A^h}$ induces isomorphisms
on cohomology groups. Moreover, an \'etale ring map is
a local complete intersection as it is even a global complete
intersection, see Algebra, Lemma \ref{algebra-lemma-etale-standard-smooth}.
By Lemma \ref{lemma-transitive-colimit-lci-at-end}
we get a six term exact Jacobi-Zariski sequence
associated to $A \to B \to B^h$ which proves that
$\NL_{B/A} \otimes_B B^h \to \NL_{B^h/A}$ induces isomorphisms
on cohomology groups. This finishes the proof in the
case of the map on henselizations. The case of strict henselization
is proved in exactly the same manner.
\end{proof}







\section{Cartier's equality and geometric regularity}
\label{section-cartier-equality}

\noindent
A reference for this section and the next is \cite[Section 39]{MatCA}.
In order to comfortably read this section the reader should be
familiar with the naive cotangent complex and its properties, see
Algebra, Section \ref{algebra-section-netherlander}.

\begin{lemma}[Cartier equality]
\label{lemma-cartier-equality}
Let $K/k$ be a finitely generated field extension.
Then $\Omega_{K/k}$ and $H_1(L_{K/k})$ are finite dimensional and
$\text{trdeg}_k(K) = \dim_K \Omega_{K/k} - \dim_K H_1(L_{K/k})$.
\end{lemma}

\begin{proof}
We can find a global complete intersection
$A = k[x_1, \ldots, x_n]/(f_1, \ldots, f_c)$
over $k$ such that
$K$ is isomorphic to the fraction field of $A$, see
Algebra, Lemma \ref{algebra-lemma-colimit-syntomic}
and its proof. In this case we see that $\NL_{K/k}$ is homotopy
equivalent to the complex
$$
\bigoplus\nolimits_{j = 1, \ldots, c} K \longrightarrow
\bigoplus\nolimits_{i = 1, \ldots, n} K\text{d}x_i
$$
by Algebra, Lemmas \ref{algebra-lemma-NL-homotopy} and
\ref{algebra-lemma-localize-NL}.
The transcendence degree of $K$ over $k$ is the dimension of $A$
(by Algebra, Lemma \ref{algebra-lemma-dimension-prime-polynomial-ring})
which is $n - c$ and we win.
\end{proof}

\begin{lemma}
\label{lemma-transitivity-gamma}
Let $M/L/K$ be field extensions. Then the Jacobi-Zariski
sequence
$$
0 \to H_1(L_{L/K}) \otimes_L M \to
H_1(L_{M/K}) \to
H_1(L_{M/L}) \to
\Omega_{L/K} \otimes_L M \to
\Omega_{M/K} \to
\Omega_{M/L} \to 0
$$
is exact.
\end{lemma}

\begin{proof}
Combine
Lemma \ref{lemma-transitive-colimit-lci-at-end}
with
Algebra, Lemma \ref{algebra-lemma-colimit-syntomic}.
\end{proof}

\begin{lemma}
\label{lemma-gamma-commutative-diagram}
Given a commutative diagram of fields
$$
\xymatrix{
K \ar[r] & K' \\
k \ar[u] \ar[r] & k' \ar[u]
}
$$
with $k'/k$ and $K'/K$ finitely generated field extensions
the kernel and cokernel of the maps
$$
\alpha : \Omega_{K/k} \otimes_K K' \to \Omega_{K'/k'}
\quad\text{and}\quad
\beta : H_1(L_{K/k}) \otimes_K K' \to H_1(L_{K'/k'})
$$
are finite dimensional and
$$
\dim \Ker(\alpha) - \dim \Coker(\alpha)
-\dim \Ker(\beta) + \dim \Coker(\beta)
=
\text{trdeg}_k(k') - \text{trdeg}_K(K')
$$
\end{lemma}

\begin{proof}
The Jacobi-Zariski sequences for $k \subset k' \subset K'$ and
$k \subset K \subset K'$ are
$$
0 \to H_1(L_{k'/k}) \otimes K'  \to
H_1(L_{K'/k}) \to
H_1(L_{K'/k'}) \to
\Omega_{k'/k} \otimes K' \to
\Omega_{K'/k} \to
\Omega_{K'/k'} \to 0
$$
and
$$
0 \to H_1(L_{K/k}) \otimes K' \to
H_1(L_{K'/k}) \to
H_1(L_{K'/K}) \to
\Omega_{K/k} \otimes K' \to
\Omega_{K'/k} \to
\Omega_{K'/K} \to 0
$$
By
Lemma \ref{lemma-cartier-equality}
the vector spaces $\Omega_{k'/k}$, $\Omega_{K'/K}$, $H_1(L_{K'/K})$, and
$H_1(L_{k'/k})$ are finite dimensional and the alternating sum of their
dimensions is $\text{trdeg}_k(k') - \text{trdeg}_K(K')$.
The lemma follows.
\end{proof}





\section{Geometric regularity}
\label{section-geometrically-regular}

\noindent
Let $k$ be a field. Let $(A, \mathfrak m, K)$ be a Noetherian local
$k$-algebra. The Jacobi-Zariski sequence
(Algebra, Lemma \ref{algebra-lemma-exact-sequence-NL})
is a canonical exact sequence
$$
H_1(L_{K/k}) \to \mathfrak m/\mathfrak m^2 \to \Omega_{A/k} \otimes_A K \to
\Omega_{K/k} \to 0
$$
because $H_1(L_{K/A}) = \mathfrak m/\mathfrak m^2$ by
Algebra, Lemma \ref{algebra-lemma-NL-surjection}.
We will show that exactness on the left of this sequence characterizes
whether or not a regular local ring $A$ is geometrically regular over $k$.
We will link this to the notion of formal smoothness in
Section \ref{section-regular-fs}.

\begin{proposition}
\label{proposition-characterization-geometrically-regular}
Let $k$ be a field of characteristic $p > 0$.
Let $(A, \mathfrak m, K)$ be a Noetherian local
$k$-algebra. The following are equivalent
\begin{enumerate}
\item $A$ is geometrically regular over $k$,
\item for all $k \subset k' \subset k^{1/p}$
finite over $k$ the ring $A \otimes_k k'$ is regular,
\item $A$ is regular and the canonical map
$H_1(L_{K/k}) \to \mathfrak m/\mathfrak m^2$ is injective, and
\item $A$ is regular and the map
$\Omega_{k/\mathbf{F}_p} \otimes_k K \to \Omega_{A/\mathbf{F}_p} \otimes_A K$
is injective.
\end{enumerate}
\end{proposition}

\begin{proof}
Proof of (3) $\Rightarrow$ (1).
Assume (3). Let $k'/k$ be a finite purely inseparable extension.
Set $A' = A \otimes_k k'$. This is a local ring with maximal ideal
$\mathfrak m'$. Set $K' = A'/\mathfrak m'$. We get a commutative
diagram
$$
\xymatrix{
0 \ar[r] &
H_1(L_{K/k}) \otimes K' \ar[r] \ar[d]_\beta &
\mathfrak m/\mathfrak m^2 \otimes K' \ar[r] \ar[d] &
\Omega_{A/k} \otimes_A K' \ar[r] \ar[d]_{\cong} &
\Omega_{K/k} \otimes K' \ar[r] \ar[d]_\alpha &
0
\\
 &
H_1(L_{K'/k'}) \ar[r] &
\mathfrak m'/(\mathfrak m')^2 \ar[r] &
\Omega_{A'/k'} \otimes_{A'} K' \ar[r] &
\Omega_{K'/k'} \ar[r] &
0
}
$$
with exact rows. The third vertical arrow is an isomorphism by base
change for modules of differentials
(Algebra, Lemma \ref{algebra-lemma-differentials-base-change}).
Thus $\alpha$ is surjective. By
Lemma \ref{lemma-gamma-commutative-diagram} we have
$$
\dim \Ker(\alpha) - \dim \Ker(\beta) + \dim \Coker(\beta) = 0
$$
(and these dimensions are all finite). A diagram chase shows that
$\dim \mathfrak m'/(\mathfrak m')^2 \leq \dim \mathfrak m/\mathfrak m^2$.
However, since $A \to A'$ is finite flat we see that
$\dim(A) = \dim(A')$, see
Algebra, Lemma \ref{algebra-lemma-dimension-base-fibre-total}.
Hence $A'$ is regular by definition.

\medskip\noindent
Equivalence of (3) and (4). Consider the Jacobi-Zariski sequences
for rows of the commutative diagram
$$
\xymatrix{
\mathbf{F}_p \ar[r] & A \ar[r] & K \\
\mathbf{F}_p \ar[r] \ar[u] & k \ar[r] \ar[u] & K \ar[u]
}
$$
to get a commutative diagram
$$
\xymatrix{
0 \ar[r] &
\mathfrak m/\mathfrak m^2 \ar[r] &
\Omega_{A/\mathbf{F}_p} \otimes_A K \ar[r] &
\Omega_{K/\mathbf{F}_p} \ar[r] & 0 & \\
0 \ar[r] &
H_1(L_{K/k}) \ar[r] \ar[u] &
\Omega_{k/\mathbf{F}_p} \otimes_k K \ar[r] \ar[u] &
\Omega_{K/\mathbf{F}_p} \ar[r] \ar[u] &
\Omega_{K/k} \ar[r] \ar[u] &
0
}
$$
with exact rows. We have used that $H_1(L_{K/A}) = \mathfrak m/\mathfrak m^2$
and that $H_1(L_{K/\mathbf{F}_p}) = 0$ as $K/\mathbf{F}_p$ is separable, see
Algebra, Proposition
\ref{algebra-proposition-characterize-separable-field-extensions}.
Thus it is clear that the kernels of
$H_1(L_{K/k}) \to \mathfrak m/\mathfrak m^2$
and
$\Omega_{k/\mathbf{F}_p} \otimes_k K \to \Omega_{A/\mathbf{F}_p} \otimes_A K$
have the same dimension.

\medskip\noindent
Proof of (2) $\Rightarrow$ (4) following Faltings, see
\cite{Faltings-einfacher}. Let $a_1, \ldots, a_n \in k$ be elements
such that $\text{d}a_1, \ldots, \text{d}a_n$ are linearly independent
in $\Omega_{k/\mathbf{F}_p}$. Consider the field extension
$k' = k(a_1^{1/p}, \ldots, a_n^{1/p})$. By
Algebra, Lemma \ref{algebra-lemma-size-extension-pth-roots}
we see that $k' = k[x_1, \ldots, x_n]/(x_1^p - a_1, \ldots, x_n^p - a_n)$.
In particular we see that the naive cotangent complex of $k'/k$
is homotopic to the complex
$\bigoplus_{j = 1, \ldots, n} k' \rightarrow \bigoplus_{i = 1, \ldots, n} k'$
with the zero differential as
$\text{d}(x_j^p - a_j) = 0$ in $\Omega_{k[x_1, \ldots, x_n]/k}$.
Set $A' = A \otimes_k k'$ and $K' = A'/\mathfrak m'$ as above.
By Algebra, Lemma \ref{algebra-lemma-change-base-NL}
we see that $\NL_{A'/A}$ is homotopy equivalent to the complex
$\bigoplus_{j = 1, \ldots, n} A' \rightarrow \bigoplus_{i = 1, \ldots, n} A'$
with the zero differential, i.e., $H_1(L_{A'/A})$ and
$\Omega_{A'/A}$ are free of rank $n$. The Jacobi-Zariski sequence for
$\mathbf{F}_p \to A \to A'$ is
$$
H_1(L_{A'/A}) \to \Omega_{A/\mathbf{F}_p} \otimes_A A'
\to \Omega_{A'/\mathbf{F}_p} \to \Omega_{A'/A} \to 0
$$
Using the presentation $A[x_1, \ldots, x_n] \to A'$ with
kernel $(x_j^p - a_j)$ we see, unwinding the maps in
Algebra, Lemma \ref{algebra-lemma-exact-sequence-NL},
that the $j$th basis vector of $H_1(L_{A'/A})$ maps to
$\text{d}a_j \otimes 1$ in $\Omega_{A/\mathbf{F}_p} \otimes A'$.
As $\Omega_{A'/A}$ is free (hence flat) we get on tensoring with $K'$
an exact sequence
$$
K'^{\oplus n} \to \Omega_{A/\mathbf{F}_p} \otimes_A K'
\xrightarrow{\beta} \Omega_{A'/\mathbf{F}_p} \otimes_{A'} K' \to
K'^{\oplus n} \to 0
$$
We conclude that the elements $\text{d}a_j \otimes 1$ generate
$\Ker(\beta)$ and we have to show that are linearly independent, i.e.,
we have to show $\dim(\Ker(\beta)) = n$.
Consider the following big diagram
$$
\xymatrix{
0 \ar[r] &
\mathfrak m'/(\mathfrak m')^2 \ar[r] &
\Omega_{A'/\mathbf{F}_p} \otimes K' \ar[r] &
\Omega_{K'/\mathbf{F}_p} \ar[r] & 0 \\
0 \ar[r] &
\mathfrak m/\mathfrak m^2 \otimes K' \ar[r] \ar[u]^\alpha &
\Omega_{A/\mathbf{F}_p} \otimes K' \ar[r] \ar[u]^\beta &
\Omega_{K/\mathbf{F}_p} \otimes K' \ar[r] \ar[u]^\gamma & 0
}
$$
By Lemma \ref{lemma-cartier-equality} and the Jacobi-Zariski sequence for
$\mathbf{F}_p \to K \to K'$ we see that the kernel and cokernel of
$\gamma$ have the same finite dimension.
By assumption $A'$ is regular (and of the same dimension as $A$, see
above) hence the kernel and cokernel of $\alpha$ have the same dimension.
It follows that the kernel and cokernel of $\beta$ have the same
dimension which is what we wanted to show.

\medskip\noindent
The implication (1) $\Rightarrow$ (2) is trivial. This finishes
the proof of the proposition.
\end{proof}

\begin{lemma}
\label{lemma-geometrically-regular-over-field}
Let $k$ be a field of characteristic $p > 0$. Let $(A, \mathfrak m, K)$
be a Noetherian local $k$-algebra. Assume $A$ is geometrically regular
over $k$. Let $K/F/k$ be a finitely generated subextension.
Let $\varphi : k[y_1, \ldots, y_m] \to A$ be a $k$-algebra map
such that $y_i$ maps to an element of $F$ in $K$ and such that
$\text{d}y_1, \ldots, \text{d}y_m$ map to a basis of $\Omega_{F/k}$.
Set $\mathfrak p = \varphi^{-1}(\mathfrak m)$. Then
$$
k[y_1, \ldots, y_m]_\mathfrak p \to A
$$
is flat and $A/\mathfrak pA$ is regular.
\end{lemma}

\begin{proof}
Set $A_0 = k[y_1, \ldots, y_m]_\mathfrak p$ with maximal ideal
$\mathfrak m_0$ and residue field $K_0$. Note that
$\Omega_{A_0/k}$ is free of rank $m$ and
$\Omega_{A_0/k} \otimes K_0 \to \Omega_{K_0/k}$ is an isomorphism.
It is clear that $A_0$ is geometrically regular over $k$. Hence
$H_1(L_{K_0/k}) \to \mathfrak m_0/\mathfrak m_0^2$ is an isomorphism, see
Proposition \ref{proposition-characterization-geometrically-regular}.
Now consider
$$
\xymatrix{
H_1(L_{K_0/k}) \otimes K \ar[d] \ar[r] &
\mathfrak m_0/\mathfrak m_0^2 \otimes K \ar[d] \\
H_1(L_{K/k}) \ar[r] & \mathfrak m/\mathfrak m^2
}
$$
Since the left vertical arrow is injective by
Lemma \ref{lemma-transitivity-gamma}
and the lower horizontal by
Proposition \ref{proposition-characterization-geometrically-regular}
we conclude that the right vertical one is too.
Hence a regular system of parameters in $A_0$ maps to
part of a regular system of parameters in $A$.
We win by
Algebra, Lemmas \ref{algebra-lemma-flat-over-regular} and
\ref{algebra-lemma-regular-ring-CM}.
\end{proof}












\section{Topological rings and modules}
\label{section-topological-ring}

\noindent
Let's quickly discuss some properties of topological abelian groups.
An abelian group $M$ is a {\it topological abelian group} if $M$ is
endowed with a topology such that addition $M \times M \to M$,
$(x, y) \mapsto x + y$ and inverse $M \to M$, $x \mapsto -x$ are
continuous. A {\it homomorphism of topological abelian groups}
is just a homomorphism of abelian groups which is continuous.
The category of commutative topological groups is additive and
has kernels and cokernels, but is not abelian (as the axiom
$\Im = \Coim$ doesn't hold). If $N \subset M$ is a
subgroup, then we think of $N$ and $M/N$ as topological groups
also, namely using the induced topology on $N$ and the quotient
topology on $M/N$ (i.e., such that $M \to M/N$ is submersive).
Note that if $N \subset M$ is an open subgroup, then the topology
on $M/N$ is discrete.

\medskip\noindent
We say the topology on $M$ is {\it linear} if there exists a fundamental
system of neighbourhoods of $0$ consisting of subgroups. If so then these
subgroups are also open. An example is the following. Let $I$ be a directed set
and let $G_i$ be an inverse system of (discrete)
abelian groups over $I$. Then
$$
G = \lim_{i \in I} G_i
$$
with the inverse limit topology is linearly topologized with a fundamental
system of neighbourhoods of $0$ given by $\Ker(G \to G_i)$.
Conversely, let $M$ be a linearly topologized abelian group.
Choose any fundamental system of open subgroups $U_i \subset M$, $i \in I$
(i.e., the $U_i$ form a fundamental system of open neighbourhoods and
each $U_i$ is a subgroup of $M$). Setting
$i \geq i' \Leftrightarrow U_i \subset U_{i'}$ we see that $I$
is a directed set. We obtain a homomorphism of
linearly topologized abelian groups
$$
c : M \longrightarrow \lim_{i \in I} M/U_i.
$$
It is clear that $M$ is {\it separated} (as a topological space)
if and only if $c$ is injective. We say that $M$ is {\it complete}
if $c$ is an isomorphism\footnote{We include being separated as part
of being complete as we'd like to have a unique limits in complete
groups. There is a definition of completeness for any topological group,
agreeing, modulo the separation issue, with this one in our special case.}.
We leave it to the reader to check that this condition is independent of
the choice of fundamental system of open subgroups $\{U_i\}_{i \in I}$
chosen above. In fact the topological abelian group
$M^\wedge = \lim_{i \in I} M/U_i$ is independent of this
choice and is sometimes called the {\it completion} of $M$.
Any $G = \lim G_i$ as above is complete, in particular,
the completion $M^\wedge$ is always complete.

\begin{definition}[Topological rings]
\label{definition-topological-ring}
\begin{reference}
\cite[Chapter 0, Sections 7.1 and 7.2]{EGA1}
\end{reference}
Let $R$ be a ring and let $M$ be an $R$-module.
\begin{enumerate}
\item We say $R$ is a {\it topological ring} if $R$ is endowed with a topology
such that both addition and multiplication are continuous as maps
$R \times R \to R$ where $R \times R$ has the product topology.
In this case we say $M$ is a {\it topological module} if $M$ is endowed
with a topology such that addition $M \times M \to M$ and
scalar multiplication $R \times M \to M$ are continuous.
\item A {\it homomorphism of topological modules} is just a continuous
$R$-module map. A {\it homomorphism of topological rings} is a
ring homomorphism which is continuous for the given topologies.
\item We say $M$ is {\it linearly topologized} if $0$ has a fundamental
system of neighbourhoods consisting of submodules. We say $R$ is
{\it linearly topologized} if $0$ has a fundamental system of neighbourhoods
consisting of ideals.
\item If $R$ is linearly topologized, we say that $I \subset R$ is an
{\it ideal of definition} if $I$ is open and if every neighbourhood
of $0$ contains $I^n$ for some $n$.
\item If $R$ is linearly topologized, we say that $R$ is {\it pre-admissible}
if $R$ has an ideal of definition.
\item If $R$ is linearly topologized, we say that $R$ is {\it admissible} if
it is pre-admissible and
complete\footnote{By our conventions this includes separated.}.
\item If $R$ is linearly topologized, we say that $R$ is {\it pre-adic} if
there exists an ideal of definition $I$ such that $\{I^n\}_{n \geq 0}$
forms a fundamental system of neighbourhoods of $0$.
\item If $R$ is linearly topologized, we say that $R$ is {\it adic} if
$R$ is pre-adic and complete.
\end{enumerate}
Note that a (pre)adic topological ring is the same thing as a (pre)admissible
topological ring which has an ideal of definition $I$ such that $I^n$ is
open for all $n \geq 1$.
\end{definition}

\noindent
Let $R$ be a ring and let $M$ be an $R$-module. Let $I \subset R$ be
an ideal. Then we can consider the linear topology on $R$ which has
$\{I^n\}_{n \geq 0}$ as a fundamental system of neighbourhoods of $0$.
This topology is called the $I$-adic topology; $R$ is a pre-adic
topological ring in the $I$-adic topology\footnote{Thus the $I$-adic
topology is sometimes called the $I$-pre-adic topology.}. Moreover, the
linear topology
on $M$ which has $\{I^nM\}_{n \geq 0}$ as a fundamental system of open
neighbourhoods of $0$ turns $M$ into a topological $R$-module.
This is called the {\it $I$-adic topology} on $M$. We see that
$M$ is $I$-adically complete (as defined in
Algebra, Definition \ref{algebra-definition-complete})
if and only if $M$ is complete in the $I$-adic topology\footnote{
It may happen that the $I$-adic completion $M^\wedge$
is not $I$-adically complete, even though $M^\wedge$
is always complete with respect to the limit topology. If $I$ is finitely
generated then the $I$-adic topology and the limit topology on $M^\wedge$
agree, see Algebra, Lemma \ref{algebra-lemma-hathat-finitely-generated} and
its proof.}.
In particular, we see that $R$ is $I$-adically complete if and
only if $R$ is an adic topological ring in the $I$-adic topology.

\medskip\noindent
As a special case, note that the discrete topology is the $0$-adic
topology and that any ring in the discrete topology is adic.

\begin{lemma}
\label{lemma-continuous}
Let $\varphi : R \to S$ be a ring map.
Let $I \subset R$ and $J \subset S$ be ideals
and endow $R$ with the $I$-adic topology and $S$ with the $J$-adic
topology. Then $\varphi$ is a homomorphism of topological rings
if and only if $\varphi(I^n) \subset J$ for some $n \geq 1$.
\end{lemma}

\begin{proof}
Omitted.
\end{proof}

\begin{lemma}[Baire category theorem]
\label{lemma-baire-category-complete-module}
Let $M$ be a topological abelian group. Assume $M$ is linearly
topologized, complete, and has a countable fundamental system of
neighbourhoods of $0$. If $U_n \subset M$, $n \geq 1$
are open dense subsets, then $\bigcap_{n \geq 1} U_n$ is dense.
\end{lemma}

\begin{proof}
Let $U_n$ be as in the statement of the lemma. After replacing $U_n$ by
$U_1 \cap \ldots \cap U_n$, we may assume that
$U_1 \supset U_2 \supset \ldots$.
Let $M_n$, $n \in \mathbf{N}$ be a fundamental
system of neighbourhoods of $0$. We may assume that
$M_{n + 1} \subset M_n$. Pick $x \in M$. We will show that
for every $k \geq 1$ there exists a $y \in \bigcap_{n \geq 1} U_n$
with $x - y \in M_k$.

\medskip\noindent
To construct $y$ we argue as follows.
First, we pick a $y_1 \in U_1$ with $y_1 \in x + M_k$.
This is possible because $U_1$ is dense and $x + M_k$ is open.
Then we pick a $k_1 > k$ such that $y_1 + M_{k_1} \subset U_1$.
This is possible because $U_1$ is open.
Next, we pick a $y_2 \in U_2$ with $y_2 \in y_1 + M_{k_1}$.
This is possible because $U_2$ is dense and $y_2 + M_{k_1}$ is open.
Then we pick a $k_2 > k_1$ such that $y_2 + M_{k_2} \subset U_2$.
This is possible because $U_2$ is open.

\medskip\noindent
Continuing in this fashion we get a converging sequence
$y_i$ of elements of $M$ with limit $y$. By construction
$x - y \in M_k$. Since
$$
y - y_i = (y_{i + 1} - y_i) + (y_{i + 2} - y_{i + 1}) + \ldots
$$
is in $M_{k_i}$ we see that $y \in y_i + M_{k_i} \subset U_i$
for all $i$ as desired.
\end{proof}

\begin{lemma}
\label{lemma-consequence-baire-complete-module}
With same assumptions as Lemma \ref{lemma-baire-category-complete-module}
if $M = \bigcup_{n \geq 1} N_n$ for some closed subgroups $N_n$,
then $N_n$ is open for some $n$.
\end{lemma}

\begin{proof}
If not, then $U_n = M \setminus N_n$ is dense for all $n$ and
we get a contradiction with Lemma \ref{lemma-baire-category-complete-module}.
\end{proof}

\begin{lemma}[Open mapping lemma]
\label{lemma-open-mapping}
Let $u : N \to M$ be a continuous map of topologized
abelian groups. Assume that $M$ separated and that $N$ is complete, linearly
topologized, and has a countable fundamental system of neighbourhoods of $0$.
Then exactly one of the following holds
\begin{enumerate}
\item $u$ is open, or
\item for some open subgroup $N' \subset N$ the image
$u(N')$ is nowhere dense in $M$.
\end{enumerate}
\end{lemma}

\begin{proof}
Let $N_n$, $n \in \mathbf{N}$ be a fundamental system of neighbourhoods
of $0$. We may assume that $N_n$ is a subgroup and that
$N_{n + 1} \subset N_n$. If (2) does not hold, then
the closure $M_n$ of $u(N_n)$ is an open subgroup for $n = 1, 2, 3, \ldots$.
Since $u$ is continuous, we see that $M_n$, $n \in \mathbf{N}$
must be a fundamental system of open neighbourhoods of $0$ in $M$.
Also, since $M_n$ is the closure of $u(N_n)$ we see that
$$
u(N_n) + M_{n + 1} = M_n
$$
for all $n \geq 1$. Pick $x_1 \in M_1$. Then we can inductively choose
$y_i \in N_i$ and $x_{i + 1} \in M_{i + 1}$ such that
$$
u(y_i) + x_{i + 1} = x_i
$$
The element $y = y_1 + y_2 + y_3 + \ldots$ of $N$ exists because
$N$ is complete. Whereupon we see that $x_1 = u(y)$ because $M$ is separated.
Thus $M_1 = u(N_1)$. In exactly the same way the reader shows that
$M_i = u(N_i)$ for all $i \geq 2$ and we see that $u$ is open.
\end{proof}









\section{Formally smooth maps of topological rings}
\label{section-formally-smooth}

\noindent
There is a version of formal smoothness which applies to
homomorphisms of topological rings.

\begin{definition}
\label{definition-formally-smooth}
Let $R \to S$ be a homomorphism of topological rings with $R$ and $S$
linearly topologized. We say $S$ is {\it formally smooth over $R$} if
for every commutative solid diagram
$$
\xymatrix{
S \ar[r] \ar@{-->}[rd] & A/J \\
R \ar[r] \ar[u] & A \ar[u]
}
$$
of homomorphisms of topological rings where $A$ is a discrete ring and
$J \subset A$ is an ideal of square zero, a dotted arrow exists which
makes the diagram commute.
\end{definition}

\noindent
We will mostly use this notion when given ideals $\mathfrak m \subset R$
and $\mathfrak n \subset S$ and we endow $R$ with the $\mathfrak m$-adic
topology and $S$ with the $\mathfrak n$-adic topology. Continuity of
$\varphi : R \to S$ holds if and only if
$\varphi(\mathfrak m^m) \subset  \mathfrak n$ for some $m \geq 1$, see
Lemma \ref{lemma-continuous}. It turns out that
in this case only the topology on $S$ is relevant.

\begin{lemma}
\label{lemma-formally-smooth}
Let $\varphi : R \to S$ be a ring map.
\begin{enumerate}
\item If $R \to S$ is formally smooth in
the sense of Algebra, Definition \ref{algebra-definition-formally-smooth},
then $R \to S$ is formally smooth for any linear topology on $R$ and
any pre-adic topology on $S$ such that $R \to S$ is continuous.
\item Let $\mathfrak n \subset S$ and $\mathfrak m \subset R$
ideals such that $\varphi$ is continuous for the $\mathfrak m$-adic
topology on $R$ and the $\mathfrak n$-adic topology
on $S$. Then the following are equivalent
\begin{enumerate}
\item $\varphi$ is formally smooth for the $\mathfrak m$-adic topology on
$R$ and the $\mathfrak n$-adic topology on $S$, and
\item $\varphi$ is formally smooth for the discrete topology
on $R$ and the $\mathfrak n$-adic topology on $S$.
\end{enumerate}
\end{enumerate}
\end{lemma}

\begin{proof}
Assume $R \to S$ is formally smooth in
the sense of Algebra, Definition \ref{algebra-definition-formally-smooth}.
If $S$ has a pre-adic topology, then
there exists an ideal $\mathfrak n \subset S$ such that $S$ has the
$\mathfrak n$-adic topology. Suppose given a solid commutative diagram as in
Definition \ref{definition-formally-smooth}.
Continuity of $S \to A/J$ means that $\mathfrak n^k$ maps to zero
in $A/J$ for some $k \geq 1$, see Lemma \ref{lemma-continuous}.
We obtain a ring map $\psi : S \to A$ from the assumed formal smoothness of
$S$ over $R$. Then $\psi(\mathfrak n^k) \subset J$ hence
$\psi(\mathfrak n^{2k}) = 0$ as $J^2 = 0$. Hence $\psi$ is continuous by
Lemma \ref{lemma-continuous}. This proves (1).

\medskip\noindent
The proof of (2)(b) $\Rightarrow$ (2)(a) is the same as the proof of (1).
Assume (2)(a). Suppose given a solid commutative diagram as in
Definition \ref{definition-formally-smooth} where we use the discrete
topology on $R$. Since $\varphi$ is continuous we see that
$\varphi(\mathfrak m^n) \subset \mathfrak n$ for some $n \geq 1$.
As $S \to A/J$ is continuous we see that $\mathfrak n^k$ maps to
zero in $A/J$ for some $k \geq 1$. Hence $\mathfrak m^{nk}$ maps
into $J$ under the map $R \to A$. Thus $\mathfrak m^{2nk}$ maps to zero
in $A$ and we see that $R \to A$ is continuous in the $\mathfrak m$-adic
topology. Thus (2)(a) gives a dotted arrow as desired.
\end{proof}

\begin{definition}
\label{definition-formally-smooth-adic}
Let $R \to S$ be a ring map. Let $\mathfrak n \subset S$ be an
ideal. If the equivalent conditions (2)(a) and (2)(b) of
Lemma \ref{lemma-formally-smooth} hold, then we say
$R \to S$ is {\it formally smooth for the $\mathfrak n$-adic topology}.
\end{definition}

\noindent
This property is inherited by the completions.

\begin{lemma}
\label{lemma-formally-smooth-completion}
Let $(R, \mathfrak m)$ and $(S, \mathfrak n)$ be rings endowed
with finitely generated ideals. Endow $R$ and $S$ with the
$\mathfrak m$-adic and $\mathfrak n$-adic topologies.
Let $R \to S$ be a homomorphism of topological rings.
The following are equivalent
\begin{enumerate}
\item $R \to S$ is formally smooth for the $\mathfrak n$-adic topology,
\item $R \to S^\wedge$ is formally smooth for the $\mathfrak n^\wedge$-adic
topology,
\item $R^\wedge \to S^\wedge$ is formally smooth for the
$\mathfrak n^\wedge$-adic topology.
\end{enumerate}
Here $R^\wedge$ and $S^\wedge$ are the $\mathfrak m$-adic and
$\mathfrak n$-adic completions of $R$ and $S$.
\end{lemma}

\begin{proof}
The assumption that $\mathfrak m$ is finitely generated implies that
$R^\wedge$ is $\mathfrak mR^\wedge$-adically complete, that
$\mathfrak mR^\wedge = \mathfrak m^\wedge$ and that
$R^\wedge/\mathfrak m^nR^\wedge = R/\mathfrak m^n$,
see Algebra, Lemma \ref{algebra-lemma-hathat-finitely-generated}
and its proof. Similarly for $(S, \mathfrak n)$. Thus it is clear that
diagrams as in Definition \ref{definition-formally-smooth}
for the cases (1), (2), and (3) are in 1-to-1 correspondence.
\end{proof}

\noindent
The advantage of working with adic rings is that one gets a
stronger lifting property.

\begin{lemma}
\label{lemma-lift-continuous}
Let $R \to S$ be a ring map. Let $\mathfrak n$ be an ideal of $S$.
Assume that $R \to S$ is formally smooth in the $\mathfrak n$-adic
topology. Consider a solid commutative diagram
$$
\xymatrix{
S \ar[r]_\psi \ar@{-->}[rd] & A/J \\
R \ar[r] \ar[u] & A \ar[u]
}
$$
of homomorphisms of topological rings where $A$ is adic
and $A/J$ is the quotient (as topological ring) of $A$ by a closed ideal
$J \subset A$ such that $J^t$ is contained in an ideal of definition
of $A$ for some $t \geq 1$. Then there exists a dotted arrow in the category of
topological rings which makes the diagram commute.
\end{lemma}

\begin{proof}
Let $I \subset A$ be an ideal of definition so that $I \supset J^t$
for some $t$. Then $A = \lim A/I^n$ and $A/J = \lim A/J + I^n$
because $J$ is assumed closed. Consider the following diagram of
discrete $R$ algebras $A_{n, m} = A/J^n + I^m$:
$$
\xymatrix{
A/J^3 + I^3 \ar[r] \ar[d] &
A/J^2 + I^3 \ar[r] \ar[d] &
A/J + I^3 \ar[d] \\
A/J^3 + I^2 \ar[r] \ar[d] &
A/J^2 + I^2 \ar[r] \ar[d] &
A/J + I^2 \ar[d] \\
A/J^3 + I \ar[r] &
A/J^2 + I \ar[r] &
A/J + I
}
$$
Note that each of the commutative squares defines a surjection
$$
A_{n + 1, m + 1} \longrightarrow A_{n + 1, m} \times_{A_{n, m}} A_{n, m + 1}
$$
of $R$-algebras whose kernel has square zero.
We will inductively construct $R$-algebra maps
$\varphi_{n, m} : S \to A_{n, m}$.
Namely, we have the maps $\varphi_{1, m} = \psi \bmod J + I^m$.
Note that each of these maps is continuous as $\psi$ is.
We can inductively choose the maps $\varphi_{n, 1}$ by starting
with our choice of $\varphi_{1, 1}$ and lifting up, using the
formal smoothness of $S$ over $R$, along the bottom row of the
diagram above. We construct the remaining maps $\varphi_{n, m}$
by induction on $n + m$. Namely, we choose $\varphi_{n + 1, m + 1}$
by lifting the pair $(\varphi_{n + 1, m}, \varphi_{n, m + 1})$
along the displayed surjection above (again using the formal smoothness
of $S$ over $R$). In this way all of the maps $\varphi_{n, m}$ are
compatible with the transition maps of the system.
As $J^t \subset I$ we see that for example
$\varphi_n = \varphi_{nt, n} \bmod I^n$ induces a map $S \to A/I^n$.
Taking the limit $\varphi = \lim \varphi_n$ we obtain a map
$S \to A = \lim A/I^n$. The composition into $A/J$ agrees
with $\psi$ as we have seen that $A/J = \lim A/J + I^n$.
Finally we show that $\varphi$ is continuous. Namely, we know that
$\psi(\mathfrak n^r) \subset J + I/J$ for some $r \geq 1$ by our
assumption that $\psi$ is a morphism of topological rings, see
Lemma \ref{lemma-continuous}. Hence $\varphi(\mathfrak n^r) \subset J + I$
hence $\varphi(\mathfrak n^{rt}) \subset I$ as desired.
\end{proof}

\begin{lemma}
\label{lemma-increase-ideal}
Let $R \to S$ be a ring map. Let $\mathfrak n \subset \mathfrak n' \subset S$
be ideals. If $R \to S$ is formally smooth for the $\mathfrak n$-adic
topology, then $R \to S$ is formally smooth for the $\mathfrak n'$-adic
topology.
\end{lemma}

\begin{proof}
Omitted.
\end{proof}

\begin{lemma}
\label{lemma-compose-formally-smooth}
A composition of formally smooth continuous homomorphisms of linearly
topologized rings is formally smooth.
\end{lemma}

\begin{proof}
Omitted. (Hint: This is completely formal, and follows from considering
a suitable diagram.)
\end{proof}

\begin{lemma}
\label{lemma-base-change-fs}
Let $R$, $S$ be rings. Let $\mathfrak n \subset S$ be an ideal.
Let $R \to S$ be formally smooth for the $\mathfrak n$-adic topology.
Let $R \to R'$ be any ring map. Then $R' \to S' = S \otimes_R R'$
is formally smooth in the $\mathfrak n' = \mathfrak nS'$-adic
topology.
\end{lemma}

\begin{proof}
Let a solid diagram
$$
\xymatrix{
S \ar[r] \ar@{-->}[rrd] & S' \ar[r] \ar@{-->}[rd] & A/J \\
R  \ar[u] \ar[r] & R' \ar[r] \ar[u] & A \ar[u]
}
$$
as in Definition \ref{definition-formally-smooth} be given.
Then the composition $S \to S' \to A/J$ is continuous.
By assumption the longer dotted arrow exists. By the universal
property of tensor product we obtain the shorter dotted arrow.
\end{proof}

\noindent
We have seen descent for formal smoothness along faithfully flat ring
maps in Algebra, Lemma \ref{algebra-lemma-descent-formally-smooth}.
Something similar holds in the current setting of topological rings.
However, here we just prove the following very simple and easy to prove
version which is already quite useful.

\begin{lemma}
\label{lemma-descent-fs}
Let $R$, $S$ be rings. Let $\mathfrak n \subset S$ be an ideal.
Let $R \to R'$ be a ring map. Set $S' = S \otimes_R R'$ and
$\mathfrak n' = \mathfrak nS$. If
\begin{enumerate}
\item the map $R \to R'$ embeds $R$ as a direct summand of $R'$
as an $R$-module, and
\item $R' \to S'$ is formally smooth for the $\mathfrak n'$-adic topology,
\end{enumerate}
then $R \to S$ is formally smooth in the $\mathfrak n$-adic topology.
\end{lemma}

\begin{proof}
Let a solid diagram
$$
\xymatrix{
S \ar[r] & A/J \\
R \ar[u] \ar[r] & A \ar[u]
}
$$
as in Definition \ref{definition-formally-smooth} be given.
Set $A' = A \otimes_R R'$ and $J' = \Im(J \otimes_R R' \to A')$.
The base change of the diagram above is the diagram
$$
\xymatrix{
S' \ar[r] \ar@{-->}[rd]^{\psi'} & A'/J' \\
R' \ar[u] \ar[r] & A' \ar[u]
}
$$
with continuous arrows. By condition (2) we obtain the dotted arrow
$\psi' : S' \to A'$. Using condition (1) choose a direct summand decomposition
$R' = R \oplus C$ as $R$-modules. (Warning: $C$ isn't an ideal in $R'$.)
Then $A' = A \oplus A \otimes_R C$. Set
$$
J'' = \Im(J \otimes_R C \to A \otimes_R C) \subset J' \subset A'.
$$
Then $J' = J \oplus J''$ as $A$-modules. The image of the composition
$\psi : S \to A'$ of $\psi'$ with $S \to S'$ is contained in
$A + J' = A \oplus J''$. However, in the ring $A + J' = A \oplus J''$
the $A$-submodule $J''$ is an ideal! (Use that $J^2 = 0$.) Hence the
composition $S \to A + J' \to (A + J')/J'' = A$ is the arrow we were
looking for.
\end{proof}








\section{Formally smooth maps of local rings}
\label{section-formally-smooth-local}

\noindent
In the case of a local homomorphism of local rings one can
limit the diagrams for which the lifting property has to be checked.
Please compare with Algebra, Lemma \ref{algebra-lemma-smooth-test-artinian}.

\begin{lemma}
\label{lemma-fs-local}
Let $(R, \mathfrak m) \to (S, \mathfrak n)$ be a local homomorphism
of local rings. The following are equivalent
\begin{enumerate}
\item $R \to S$ is formally smooth in the $\mathfrak n$-adic topology,
\item for every solid commutative diagram
$$
\xymatrix{
S \ar[r] \ar@{-->}[rd] & A/J \\
R \ar[r] \ar[u] & A \ar[u]
}
$$
of local homomorphisms of local rings where $J \subset A$ is
an ideal of square zero, $\mathfrak m_A^n = 0$ for some $n > 0$, and
$S \to A/J$ induces an isomorphism on residue fields, a dotted
arrow exists which makes the diagram commute.
\end{enumerate}
If $S$ is Noetherian these conditions are also equivalent to
\begin{enumerate}
\item[(3)] same as in (2) but only for diagrams where in addition
$A \to A/J$ is a small extension
(Algebra, Definition \ref{algebra-definition-small-extension}).
\end{enumerate}
\end{lemma}

\begin{proof}
The implication (1) $\Rightarrow$ (2) follows from the definitions.
Consider a diagram
$$
\xymatrix{
S \ar[r] \ar@{-->}[rd] & A/J \\
R \ar[r] \ar[u] & A \ar[u]
}
$$
as in Definition \ref{definition-formally-smooth} for the
$\mathfrak m$-adic topology on $R$ and the $\mathfrak n$-adic topology on $S$.
Pick $m > 0$ with $\mathfrak n^m(A/J) = 0$
(possible by continuity of maps in diagram).
Consider the subring $A'$ of $A$ which is the inverse image
of the image of $S$ in $A/J$. Set $J' = J$ viewed as an ideal in $A'$.
Then $J'$ is an ideal of square zero in $A'$ and $A'/J'$
is a quotient of $S/\mathfrak n^m$. Hence $A'$ is local
and $\mathfrak m_{A'}^{2m} = 0$. Thus we get a diagram
$$
\xymatrix{
S \ar[r] \ar@{-->}[rd] & A'/J' \\
R \ar[r] \ar[u] & A' \ar[u]
}
$$
as in (2). If we can construct the dotted arrow in this diagram,
then we obtain the dotted arrow in the original one by composing
with $A' \to A$. In this way we see that (2) implies (1).

\medskip\noindent
Assume $S$ Noetherian. The implication (1) $\Rightarrow$ (3) is immediate.
Assume (3) and suppose a diagram as in (2) is given.
Then $\mathfrak m_A^n J = 0$ for some $n > 0$.
Considering the maps
$$
A \to A/\mathfrak m_A^{n - 1}J \to \ldots \to A/\mathfrak mJ \to A/J
$$
we see that it suffices to produce the lifting if $\mathfrak m_A J = 0$.
Assume $\mathfrak m_A J = 0$ and let $A' \subset A$ be the ring
constructed above. Then $A'/J'$ is Artinian
as a quotient of the Artinian local ring $S/\mathfrak n^m$.
Thus it suffices to show that given property (3) we can find the dotted
arrow in diagrams as in (2) with $A/J$ Artinian and
$\mathfrak m_A J = 0$.
Let $\kappa$ be the common residue field of $A$,
$A/J$, and $S$. By (3), if $J_0 \subset J$ is an ideal with
$\dim_\kappa(J/J_0) = 1$, then we can produce a dotted arrow
$S \to A/J_0$. Taking the product we obtain
$$
S \longrightarrow \prod\nolimits_{J_0 \text{ as above}} A/J_0
$$
Clearly the image of this arrow is contained in the sub $R$-algebra
$A'$ of elements which map into the small diagonal
$A/J \subset \prod_{J_0} A/J$.
Let $J' \subset A'$ be the elements mapping to zero in $A/J$.
Then $J'$ is an ideal of square zero and as $\kappa$-vector space
equal to
$$
J' = \prod\nolimits_{J_0 \text{ as above}} J/J_0
$$
Thus the map $J \to J'$ is injective. By the theory of vector spaces
we can choose a splitting $J' = J \oplus M$. It follows that
$$
A' = A \oplus M
$$
as an $R$-algebra. Hence the map $S \to A'$ can be composed with
the projection $A' \to A$ to give the desired dotted arrow thereby
finishing the proof of the lemma.
\end{proof}

\noindent
The following lemma will be improved on in
Section \ref{section-regular-fs}.

\begin{lemma}
\label{lemma-fs-implies-regular}
Let $k$ be a field and let $(A, \mathfrak m, K)$ be a Noetherian
local $k$-algebra. If $k \to A$ is formally smooth for the
$\mathfrak m$-adic topology, then $A$ is a regular local ring.
\end{lemma}

\begin{proof}
Let $k_0 \subset k$ be the prime field. Then $k_0$ is perfect, hence
$k / k_0$ is separable, hence formally smooth by
Algebra, Lemma \ref{algebra-lemma-formally-smooth-extensions-easy}. By
Lemmas \ref{lemma-formally-smooth} and \ref{lemma-compose-formally-smooth}
we see that $k_0 \to A$ is formally smooth for the $\mathfrak m$-adic
topology on $A$. Hence we may assume $k = \mathbf{Q}$ or $k = \mathbf{F}_p$.

\medskip\noindent
By Algebra, Lemmas \ref{algebra-lemma-completion-faithfully-flat} and
\ref{algebra-lemma-flat-under-regular} it
suffices to prove the completion $A^\wedge$ is regular.
By Lemma \ref{lemma-formally-smooth-completion} we may replace
$A$ by $A^\wedge$. Thus we may assume that $A$ is a Noetherian
complete local ring. By the Cohen structure theorem
(Algebra, Theorem \ref{algebra-theorem-cohen-structure-theorem})
there exist a map $K \to A$. As $k$ is the prime field we see that
$K \to A$ is a $k$-algebra map.

\medskip\noindent
Let $x_1, \ldots, x_n \in \mathfrak m$ be elements whose images
form a basis of $\mathfrak m/\mathfrak m^2$.
Set $T = K[[X_1, \ldots, X_n]]$. Note that
$$
A/\mathfrak m^2 \cong K[x_1, \ldots, x_n]/(x_ix_j)
$$
and
$$
T/\mathfrak m_T^2 \cong K[X_1, \ldots, X_n]/(X_iX_j).
$$
Let $A/\mathfrak m^2 \to T/m_T^2$ be the local $K$-algebra isomorphism
given by mapping the class of $x_i$ to the class of $X_i$.
Denote $f_1 : A \to T/\mathfrak m_T^2$ the composition of this
isomorphism with the quotient map $A \to A/\mathfrak m^2$.
The assumption that $k \to A$ is formally smooth in the $\mathfrak m$-adic
topology means we can lift $f_1$ to a map
$f_2 : A \to T/\mathfrak{m}_T^3$, then to a map
$f_3 : A \to T/\mathfrak{m}_T^4$, and so on, for all $n \geq 1$.
Warning: the maps $f_n$ are continuous $k$-algebra maps and may not
be $K$-algebra maps. We get an induced map
$f : A \to T = \lim T/\mathfrak m_T^n$ of local $k$-algebras.
By our choice of $f_1$, the map $f$ induces an
isomorphism $\mathfrak m/\mathfrak m^2 \to \mathfrak m_T/\mathfrak m_T^2$
hence each $f_n$ is surjective and we conclude $f$ is surjective as $A$ is
complete. This implies $\dim(A) \geq \dim(T) = n$. Hence $A$ is regular
by definition. (It also follows that $f$ is an isomorphism.)
\end{proof}

\begin{lemma}
\label{lemma-lift-residue-field}
Let $k$ be a field. Let $(A, \mathfrak m, \kappa)$ be a complete
local $k$-algebra. If $\kappa/k$ is separable, then there exists
a $k$-algebra map $\kappa \to A$ such that $\kappa \to A \to \kappa$
is $\text{id}_\kappa$.
\end{lemma}

\begin{proof}
By Algebra, Proposition
\ref{algebra-proposition-characterize-separable-field-extensions}
the extension $\kappa/k$ is formally smooth. By
Lemma \ref{lemma-formally-smooth}
$k \to \kappa$ is formally smooth in the sense of
Definition \ref{definition-formally-smooth}.
Then we get $\kappa \to A$ from Lemma \ref{lemma-lift-continuous}.
\end{proof}

\begin{lemma}
\label{lemma-power-series-over-residue-field}
Let $k$ be a field. Let $(A, \mathfrak m, \kappa)$ be a complete
local $k$-algebra. If $\kappa/k$ is separable and $A$ regular, then
there exists an isomorphism of $A \cong \kappa[[t_1, \ldots, t_d]]$
as $k$-algebras.
\end{lemma}

\begin{proof}
Choose $\kappa \to A$ as in Lemma \ref{lemma-lift-residue-field}
and apply Algebra, Lemma
\ref{algebra-lemma-regular-complete-containing-coefficient-field}.
\end{proof}

\noindent
The following result will be improved on in
Section \ref{section-regular-fs}

\begin{lemma}
\label{lemma-regular-implies-fs}
Let $k$ be a field. Let $(A, \mathfrak m, K)$ be a regular local
$k$-algebra such that $K/k$ is separable. Then $k \to A$
is formally smooth in the $\mathfrak m$-adic topology.
\end{lemma}

\begin{proof}
It suffices to prove that the completion of $A$ is formally
smooth over $k$, see Lemma \ref{lemma-formally-smooth-completion}.
Hence we may assume that $A$ is a complete local regular $k$-algebra
with residue field $K$ separable over $k$.
By Lemma \ref{lemma-power-series-over-residue-field} we see that
$A = K[[x_1, \ldots, x_n]]$.

\medskip\noindent
The power series ring $K[[x_1, \ldots, x_n]]$ is formally
smooth over $k$. Namely, $K$ is formally smooth over $k$ and
$K[x_1, \ldots, x_n]$ is formally smooth over $K$ as a polynomial algebra.
Hence $K[x_1, \ldots, x_n]$ is formally smooth over $k$ by
Algebra, Lemma \ref{algebra-lemma-compose-formally-smooth}.
It follows that $k \to K[x_1, \ldots, x_n]$ is formally smooth
for the $(x_1, \ldots, x_n)$-adic topology by
Lemma \ref{lemma-formally-smooth}.
Finally, it follows that $k \to K[[x_1, \ldots, x_n]]$ is formally
smooth for the $(x_1, \ldots, x_n)$-adic topology by
Lemma \ref{lemma-formally-smooth-completion}.
\end{proof}

\begin{lemma}
\label{lemma-formally-smooth-finite-type}
Let $A \to B$ be a finite type ring map with $A$ Noetherian.
Let $\mathfrak q \subset B$ be a prime ideal lying over
$\mathfrak p \subset A$. The following are equivalent
\begin{enumerate}
\item $A \to B$ is smooth at $\mathfrak q$, and
\item $A_\mathfrak p \to B_\mathfrak q$ is formally smooth in
the $\mathfrak q$-adic topology.
\end{enumerate}
\end{lemma}

\begin{proof}
The implication (2) $\Rightarrow$ (1) follows from
Algebra, Lemma \ref{algebra-lemma-smooth-test-artinian}.
Conversely, if $A \to B$ is smooth at $\mathfrak q$, then
$A \to B_g$ is smooth for some $g \in B$, $g \not \in \mathfrak q$.
Then $A \to B_g$ is formally smooth by
Algebra, Proposition \ref{algebra-proposition-smooth-formally-smooth}.
Hence $A_\mathfrak p \to B_\mathfrak q$ is formally smooth
as localization preserves formal smoothness (for example by
the criterion of Algebra, Proposition
\ref{algebra-proposition-characterize-formally-smooth}
and the fact that the cotangent complex behaves well with respect to
localization, see
Algebra, Lemmas \ref{algebra-lemma-NL-localize-bottom} and
\ref{algebra-lemma-localize-NL}).
Finally, Lemma \ref{lemma-formally-smooth} implies that
$A_\mathfrak p \to B_\mathfrak q$ is formally smooth in the
$\mathfrak q$-adic topology.
\end{proof}









\section{Some results on power series rings}
\label{section-power-series}

\noindent
Questions on formally smooth maps between Noetherian local rings
can often be reduced to questions on maps between power series
rings. In this section we prove some helper lemmas to facilitate
this kind of argument.

\begin{lemma}
\label{lemma-power-series-ring-over-Cohen-fs}
Let $K$ be a field of characteristic $0$ and $A = K[[x_1, \ldots, x_n]]$.
Let $L$ be a field of characteristic $p > 0$ and $B = L[[x_1, \ldots, x_n]]$.
Let $\Lambda$ be a Cohen ring. Let $C = \Lambda[[x_1, \ldots, x_n]]$.
\begin{enumerate}
\item $\mathbf{Q} \to A$ is formally smooth in the 
$\mathfrak m_A$-adic topology.
\item $\mathbf{F}_p \to B$ is formally smooth in the 
$\mathfrak m_B$-adic topology.
\item $\mathbf{Z} \to C$ is formally smooth in the
$\mathfrak m_C$-adic topology.
\end{enumerate}
\end{lemma}

\begin{proof}
By the universal property of power series rings it suffices to prove:
\begin{enumerate}
\item $\mathbf{Q} \to K$ is formally smooth.
\item $\mathbf{F}_p \to L$ is formally smooth.
\item $\mathbf{Z} \to \Lambda$ is formally smooth in the
$\mathfrak m_\Lambda$-adic topology.
\end{enumerate}
The first two are
Algebra, Proposition
\ref{algebra-proposition-characterize-separable-field-extensions}.
The third follows from
Algebra, Lemma \ref{algebra-lemma-cohen-ring-formally-smooth}
since for any test diagram as in Definition \ref{definition-formally-smooth}
some power of $p$ will be zero in $A/J$ and hence some power of $p$ will
be zero in $A$.
\end{proof}

\begin{lemma}
\label{lemma-quotient-power-series-ring-over-Cohen}
Let $K$ be a field and $A = K[[x_1, \ldots, x_n]]$.
Let $\Lambda$ be a Cohen ring and let $B = \Lambda[[x_1, \ldots, x_n]]$.
\begin{enumerate}
\item If $y_1, \ldots, y_n \in A$ is a regular system of parameters
then $K[[y_1, \ldots, y_n]] \to A$ is an isomorphism.
\item If $z_1, \ldots, z_r \in A$ form part of a regular system of
parameters for $A$, then $r \leq n$ and
$A/(z_1, \ldots, z_r) \cong K[[y_1, \ldots, y_{n - r}]]$.
\item If $p, y_1, \ldots, y_n \in B$ is a regular system of parameters
then $\Lambda[[y_1, \ldots, y_n]] \to B$ is an isomorphism.
\item If $p, z_1, \ldots, z_r \in B$ form part of a regular system of
parameters for $B$, then $r \leq n$ and
$B/(z_1, \ldots, z_r) \cong \Lambda[[y_1, \ldots, y_{n - r}]]$.
\end{enumerate}
\end{lemma}

\begin{proof}
Proof of (1). Set $A' = K[[y_1, \ldots, y_n]]$. It is clear that
the map $A' \to A$ induces an isomorphism
$A'/\mathfrak m_{A'}^n \to A/\mathfrak m_A^n$ for all $n \geq 1$.
Since $A$ and $A'$ are both complete we deduce that $A' \to A$ is an
isomorphism.
Proof of (2). Extend $z_1, \ldots, z_r$ to a regular system of parameters
$z_1, \ldots, z_r, y_1, \ldots, y_{n - r}$ of $A$. Consider the map
$A' = K[[z_1, \ldots, z_r, y_1, \ldots, y_{n - r}]] \to A$.
This is an isomorphism by (1). Hence (2) follows as it is clear that
$A'/(z_1, \ldots, z_r) \cong K[[y_1, \ldots, y_{n - r}]]$.
The proofs of (3) and (4) are exactly the same as the proofs of (1) and (2).
\end{proof}

\begin{lemma}
\label{lemma-embed-map-Noetherian-complete-local-rings}
Let $A \to B$ be a local homomorphism of Noetherian complete local rings.
Then there exists a commutative diagram
$$
\xymatrix{
S \ar[r] & B \\
R \ar[u] \ar[r] & A \ar[u]
}
$$
with the following properties:
\begin{enumerate}
\item the horizontal arrows are surjective,
\item if the characteristic of $A/\mathfrak m_A$ is zero, then $S$ and $R$
are power series rings over fields,
\item if the characteristic of $A/\mathfrak m_A$ is $p > 0$, then $S$ and $R$
are power series rings over Cohen rings, and
\item $R \to S$ maps a regular system of parameters of $R$ to part of a
regular system of parameters of $S$.
\end{enumerate}
In particular $R \to S$ is flat (see Algebra,
Lemma \ref{algebra-lemma-flat-over-regular}) with regular fibre
$S/\mathfrak m_R S$ (see Algebra, Lemma \ref{algebra-lemma-regular-ring-CM}).
\end{lemma}

\begin{proof}
Use the Cohen structure theorem
(Algebra, Theorem \ref{algebra-theorem-cohen-structure-theorem})
to choose a surjection $S \to B$ as in the statement of the lemma
where we choose $S$ to be a power series over a Cohen ring if the
residue characteristic is $p > 0$ and a power series over a field else.
Let $J \subset S$ be the kernel of $S \to B$.
Next, choose a surjection $R = \Lambda[[x_1, \ldots, x_n]] \to A$ where
we choose $\Lambda$ to be a Cohen ring if the residue characteristic of
$A$ is $p > 0$ and $\Lambda$ equal to the residue field of $A$ otherwise.
We lift the composition $\Lambda[[x_1, \ldots, x_n]] \to A \to B$
to a map $\varphi : R \to S$. This is possible because
$\Lambda[[x_1, \ldots, x_n]]$ is formally smooth over $\mathbf{Z}$
in the $\mathfrak m$-adic topology (see
Lemma \ref{lemma-power-series-ring-over-Cohen-fs})
by an application of Lemma \ref{lemma-lift-continuous}.
Finally, we replace $\varphi$ by the map
$\varphi' : R = \Lambda[[x_1, \ldots, x_n]] \to S' = S[[y_1, \ldots, y_n]]$
with $\varphi'|_\Lambda = \varphi|_\Lambda$ and
$\varphi'(x_i) = \varphi(x_i) + y_i$. We also replace $S \to B$
by the map $S' \to B$ which maps $y_i$ to zero. After this replace